% =====================================================================
%  couverture.tex — couverture dos carré collé pour les livres 50ohm-fr
% =====================================================================
%
%  Produit une couverture d'un seul tenant : 4e de couverture | dos |
%  1re de couverture, aux dimensions exactes du livre relié, fond perdu
%  compris. L'épaisseur du dos est CALCULÉE à partir du nombre de pages,
%  du grammage et du bouffant du papier.
%
%  ORIGINE ET ÉTAT. Écrit le 20/08/2026 dans Fichierstravail/ (hors git),
%  versé au dépôt le 16/09/2026 pour que fenetre_compilation.py puisse
%  l'appeler. Ce qui a changé à cette occasion :
%
%    - le FORMAT est un paramètre (\Format = a4, 20x24 ou 20x24-marge,
%      mêmes noms que --format de build_book.py) ;
%    - la pagination et la version peuvent être IMPOSÉES de l'extérieur
%      (\PagesForcees, \VersionForcee) : la fenêtre lit la pagination dans
%      le journal du livre qu'elle vient de compiler, au lieu de se fier à
%      une table qui vieillit ;
%    - un dos imposé par l'imprimeur n'est appliqué QUE pour la pagination
%      et le format pour lesquels il a été donné ;
%    - un CONTRÔLE DE SUPERPOSITION est calculé par LaTeX lui-même, à partir
%      des boîtes réelles des nœuds, et écrit dans le journal ;
%    - le tome SWL est refusé explicitement : sa 4e de couverture n'est pas
%      rédigée (le texte commun parle des classes N, E et A).
%
%    Neutralité vérifiée : en A4, pour N (258 p., a.2) et NEA (816 p., a.3),
%    en épreuve comme en impression, le rendu est identique au pixel près à
%    celui du gabarit d'origine.
%
%  COMPILATION (LuaLaTeX, depuis la racine du dépôt) :
%
%      lualatex couverture.tex
%      lualatex couverture.tex        <-- DEUX FOIS, obligatoirement
%
%  DEUX PASSES SONT NÉCESSAIRES. Le dessin s'ancre sur « current page »,
%  dont la position n'est connue qu'après un premier passage : à la
%  première compilation, tout se tasse dans le coin supérieur et la
%  couverture paraît ratée alors qu'elle ne l'est pas. Aucune erreur n'est
%  émise — seul le rendu le montre. Le contrôle de superposition, lui
%  aussi, ne vaut qu'à la seconde passe.
%
%  Le plus simple est de passer par fenetre_compilation.py (case
%  « couverture pour l'imprimeur »), qui enchaîne les deux passes, produit
%  l'épreuve ET le fichier d'impression, et relit le contrôle.
%
%  À la main, pour une autre classe, un autre format, une pagination mesurée :
%
%      lualatex "\def\Classe{A}\def\Format{20x24}\def\PagesForcees{482}% =====================================================================
%  couverture.tex — couverture dos carré collé pour les livres 50ohm-fr
% =====================================================================
%
%  Produit une couverture d'un seul tenant : 4e de couverture | dos |
%  1re de couverture, aux dimensions exactes du livre relié, fond perdu
%  compris. L'épaisseur du dos est CALCULÉE à partir du nombre de pages,
%  du grammage et du bouffant du papier.
%
%  ORIGINE ET ÉTAT. Écrit le 20/08/2026 dans Fichierstravail/ (hors git),
%  versé au dépôt le 16/09/2026 pour que fenetre_compilation.py puisse
%  l'appeler. Ce qui a changé à cette occasion :
%
%    - le FORMAT est un paramètre (\Format = a4, 20x24 ou 20x24-marge,
%      mêmes noms que --format de build_book.py) ;
%    - la pagination et la version peuvent être IMPOSÉES de l'extérieur
%      (\PagesForcees, \VersionForcee) : la fenêtre lit la pagination dans
%      le journal du livre qu'elle vient de compiler, au lieu de se fier à
%      une table qui vieillit ;
%    - un dos imposé par l'imprimeur n'est appliqué QUE pour la pagination
%      et le format pour lesquels il a été donné ;
%    - un CONTRÔLE DE SUPERPOSITION est calculé par LaTeX lui-même, à partir
%      des boîtes réelles des nœuds, et écrit dans le journal ;
%    - le tome SWL est refusé explicitement : sa 4e de couverture n'est pas
%      rédigée (le texte commun parle des classes N, E et A).
%
%    Neutralité vérifiée : en A4, pour N (258 p., a.2) et NEA (816 p., a.3),
%    en épreuve comme en impression, le rendu est identique au pixel près à
%    celui du gabarit d'origine.
%
%  COMPILATION (LuaLaTeX, depuis la racine du dépôt) :
%
%      lualatex couverture.tex
%      lualatex couverture.tex        <-- DEUX FOIS, obligatoirement
%
%  DEUX PASSES SONT NÉCESSAIRES. Le dessin s'ancre sur « current page »,
%  dont la position n'est connue qu'après un premier passage : à la
%  première compilation, tout se tasse dans le coin supérieur et la
%  couverture paraît ratée alors qu'elle ne l'est pas. Aucune erreur n'est
%  émise — seul le rendu le montre. Le contrôle de superposition, lui
%  aussi, ne vaut qu'à la seconde passe.
%
%  Le plus simple est de passer par fenetre_compilation.py (case
%  « couverture pour l'imprimeur »), qui enchaîne les deux passes, produit
%  l'épreuve ET le fichier d'impression, et relit le contrôle.
%
%  À la main, pour une autre classe, un autre format, une pagination mesurée :
%
%      lualatex "\def\Classe{A}\def\Format{20x24}\def\PagesForcees{482}% =====================================================================
%  couverture.tex — couverture dos carré collé pour les livres 50ohm-fr
% =====================================================================
%
%  Produit une couverture d'un seul tenant : 4e de couverture | dos |
%  1re de couverture, aux dimensions exactes du livre relié, fond perdu
%  compris. L'épaisseur du dos est CALCULÉE à partir du nombre de pages,
%  du grammage et du bouffant du papier.
%
%  ORIGINE ET ÉTAT. Écrit le 20/08/2026 dans Fichierstravail/ (hors git),
%  versé au dépôt le 16/09/2026 pour que fenetre_compilation.py puisse
%  l'appeler. Ce qui a changé à cette occasion :
%
%    - le FORMAT est un paramètre (\Format = a4, 20x24 ou 20x24-marge,
%      mêmes noms que --format de build_book.py) ;
%    - la pagination et la version peuvent être IMPOSÉES de l'extérieur
%      (\PagesForcees, \VersionForcee) : la fenêtre lit la pagination dans
%      le journal du livre qu'elle vient de compiler, au lieu de se fier à
%      une table qui vieillit ;
%    - un dos imposé par l'imprimeur n'est appliqué QUE pour la pagination
%      et le format pour lesquels il a été donné ;
%    - un CONTRÔLE DE SUPERPOSITION est calculé par LaTeX lui-même, à partir
%      des boîtes réelles des nœuds, et écrit dans le journal ;
%    - le tome SWL est refusé explicitement : sa 4e de couverture n'est pas
%      rédigée (le texte commun parle des classes N, E et A).
%
%    Neutralité vérifiée : en A4, pour N (258 p., a.2) et NEA (816 p., a.3),
%    en épreuve comme en impression, le rendu est identique au pixel près à
%    celui du gabarit d'origine.
%
%  COMPILATION (LuaLaTeX, depuis la racine du dépôt) :
%
%      lualatex couverture.tex
%      lualatex couverture.tex        <-- DEUX FOIS, obligatoirement
%
%  DEUX PASSES SONT NÉCESSAIRES. Le dessin s'ancre sur « current page »,
%  dont la position n'est connue qu'après un premier passage : à la
%  première compilation, tout se tasse dans le coin supérieur et la
%  couverture paraît ratée alors qu'elle ne l'est pas. Aucune erreur n'est
%  émise — seul le rendu le montre. Le contrôle de superposition, lui
%  aussi, ne vaut qu'à la seconde passe.
%
%  Le plus simple est de passer par fenetre_compilation.py (case
%  « couverture pour l'imprimeur »), qui enchaîne les deux passes, produit
%  l'épreuve ET le fichier d'impression, et relit le contrôle.
%
%  À la main, pour une autre classe, un autre format, une pagination mesurée :
%
%      lualatex "\def\Classe{A}\def\Format{20x24}\def\PagesForcees{482}% =====================================================================
%  couverture.tex — couverture dos carré collé pour les livres 50ohm-fr
% =====================================================================
%
%  Produit une couverture d'un seul tenant : 4e de couverture | dos |
%  1re de couverture, aux dimensions exactes du livre relié, fond perdu
%  compris. L'épaisseur du dos est CALCULÉE à partir du nombre de pages,
%  du grammage et du bouffant du papier.
%
%  ORIGINE ET ÉTAT. Écrit le 20/08/2026 dans Fichierstravail/ (hors git),
%  versé au dépôt le 16/09/2026 pour que fenetre_compilation.py puisse
%  l'appeler. Ce qui a changé à cette occasion :
%
%    - le FORMAT est un paramètre (\Format = a4, 20x24 ou 20x24-marge,
%      mêmes noms que --format de build_book.py) ;
%    - la pagination et la version peuvent être IMPOSÉES de l'extérieur
%      (\PagesForcees, \VersionForcee) : la fenêtre lit la pagination dans
%      le journal du livre qu'elle vient de compiler, au lieu de se fier à
%      une table qui vieillit ;
%    - un dos imposé par l'imprimeur n'est appliqué QUE pour la pagination
%      et le format pour lesquels il a été donné ;
%    - un CONTRÔLE DE SUPERPOSITION est calculé par LaTeX lui-même, à partir
%      des boîtes réelles des nœuds, et écrit dans le journal ;
%    - le tome SWL est refusé explicitement : sa 4e de couverture n'est pas
%      rédigée (le texte commun parle des classes N, E et A).
%
%    Neutralité vérifiée : en A4, pour N (258 p., a.2) et NEA (816 p., a.3),
%    en épreuve comme en impression, le rendu est identique au pixel près à
%    celui du gabarit d'origine.
%
%  COMPILATION (LuaLaTeX, depuis la racine du dépôt) :
%
%      lualatex couverture.tex
%      lualatex couverture.tex        <-- DEUX FOIS, obligatoirement
%
%  DEUX PASSES SONT NÉCESSAIRES. Le dessin s'ancre sur « current page »,
%  dont la position n'est connue qu'après un premier passage : à la
%  première compilation, tout se tasse dans le coin supérieur et la
%  couverture paraît ratée alors qu'elle ne l'est pas. Aucune erreur n'est
%  émise — seul le rendu le montre. Le contrôle de superposition, lui
%  aussi, ne vaut qu'à la seconde passe.
%
%  Le plus simple est de passer par fenetre_compilation.py (case
%  « couverture pour l'imprimeur »), qui enchaîne les deux passes, produit
%  l'épreuve ET le fichier d'impression, et relit le contrôle.
%
%  À la main, pour une autre classe, un autre format, une pagination mesurée :
%
%      lualatex "\def\Classe{A}\def\Format{20x24}\def\PagesForcees{482}\input{couverture.tex}"
%
%  Pour le fichier destiné à l'imprimeur (sans repères) :
%
%      lualatex "\def\Reperes{non}\input{couverture.tex}"
%
%  ---------------------------------------------------------------------
%  DEUX MODES, un seul paramètre : \Reperes
%
%    oui  — mode ÉPREUVE. Ajoute une marge technique, les traits de coupe,
%           les limites du dos, la zone de sécurité et les cotes. C'est ce
%           qu'il faut pour vérifier à l'écran.
%    non  — mode PRODUCTION. Papier = format final + fond perdu, rien
%           d'autre. C'est ce fichier-là qu'on envoie à l'imprimeur.
%
%  ---------------------------------------------------------------------
%  CONTRÔLE DE SUPERPOSITION
%
%  À la fin du dessin, chaque bloc de texte est mesuré par ses ancres TikZ
%  (quatre coins, ce qui vaut aussi pour le texte pivoté du dos) et comparé
%  à la zone où il doit tenir, puis aux blocs voisins. Le journal porte une
%  ligne par vérification :
%
%      COUVERTURE-CONTROLE OK      <quoi> : ...
%      COUVERTURE-CONTROLE ALERTE  <quoi> : ...
%
%  et un bilan : COUVERTURE-BILAN alertes=<n> ; ...
%
%  Ce que le contrôle VOIT : un bloc qui franchit le pli du dos, la coupe,
%  la zone de sécurité ou le bord du bandeau ; deux blocs qui se chevauchent ;
%  un bloc qui empiète sur l'emplacement réservé au code-barres ISBN.
%  Ce qu'il NE VOIT PAS : une ligne trop longue à l'INTÉRIEUR d'un bloc à
%  largeur fixe (la boîte garde sa largeur, le texte déborde). Celle-là se
%  signale par un « Overfull \hbox » dans le même journal, que la fenêtre
%  compte aussi.
%
%  ---------------------------------------------------------------------
%  AVERTISSEMENT SUR L'ÉPAISSEUR DU DOS
%
%  Le dos se calcule ainsi :
%
%      épaisseur d'une feuille (mm) = grammage x bouffant / 1000
%      dos (mm)                     = (nombre de pages / 2) x épaisseur
%
%  On divise par deux parce qu'une feuille porte DEUX pages.
%
%  Le « bouffant » (main, volume, bulk) dépend du papier et n'est pas
%  devinable : un offset 80 g courant est autour de 1,30, un bouffant
%  d'édition monte à 1,75 et plus. UN ÉCART DE BOUFFANT DÉPLACE LE DOS DE
%  PLUSIEURS MILLIMÈTRES. Demander à l'imprimeur soit son bouffant, soit
%  directement l'épaisseur du dos, et reporter ci-dessous. Tant que cette
%  valeur n'est pas confirmée, la couverture n'est pas bonne à tirer.
%
%  Si l'imprimeur donne le dos en clair, forcer la valeur avec \DosImpose.
% =====================================================================

\documentclass[11pt]{article}

% ---------------------------------------------------------------------
% 1. PARAMÈTRES — c'est la seule zone à modifier
% ---------------------------------------------------------------------

% Classe de l'ouvrage : N, E, A, NE, EA ou NEA.
% (\providecommand : une valeur passée en ligne de commande l'emporte.)
\providecommand{\Classe}{NEA}

% Mode épreuve (oui) ou production (non).
\providecommand{\Reperes}{oui}

% Format du livre FINI : a4, 20x24 ou 20x24-marge — les noms de --format
% dans build_book.py. Les deux maquettes 20 x 24 ne diffèrent que par leur
% colonne de marge intérieure : le livre fini, donc la couverture, est le
% même (200 x 240 mm).
\providecommand{\Format}{a4}

% Reliure : « souple » (dos carré collé, fond perdu 3 mm) ou « rigide »
% (couverture rigide, dite hardcover). En rigide, consignes de l'imprimeur
% du 17/09/2026 : chaque plat déborde du livre de 5 mm en largeur et en
% hauteur (le carton dépasse du bloc), et le fichier porte 15 mm de rabat
% autour du carton au lieu de 3 mm de fond perdu. Les aplats vont déjà au
% bord du papier : ils remplissent le rabat d'eux-mêmes.
\providecommand{\Reliure}{souple}

% Pagination et version imposées de l'extérieur. Vides = tables du § 2.
% La fenêtre les renseigne toujours : la pagination vient du journal du
% livre réellement compilé, dans le format demandé.
\providecommand{\PagesForcees}{}
\providecommand{\VersionForcee}{}

% Papier de l'INTÉRIEUR du livre (pas de la couverture).
% CoolLibri, papier « Standard 90 g blanc ».
\def\Grammage{90}        % g/m²
\def\Bouffant{1.206}     % main DÉDUITE, voir ci-dessous

% Épaisseur du dos IMPOSÉE par l'imprimeur, en mm, par classe.
% Une valeur > 0 prime sur le calcul — MAIS SEULEMENT pour la pagination et
% le format pour lesquels l'imprimeur l'a donnée (\DosImposePages<classe>,
% \DosImposeFormat<classe>). Hors de ce cas, le dos est calculé et le
% journal le dit : un dos de 14 mm donné pour 258 pages A4 ne vaut rien pour
% 262 pages, ni pour le même livre recomposé en 20 x 24.
%
% ATTENTION : cette valeur est propre à CHAQUE ÉDITION. Un dos unique
% appliqué aux quatre livres serait faux pour trois d'entre eux.
%
%  N   — 14 mm, VALEUR DONNÉE PAR COOLIBRI pour 258 pages en 90 g, A4.
%        Fichier attendu : 440 x 303 mm, ce que ce gabarit produit en mode
%        production (2 x 210 + 14 + 2 x 3 = 440 ; 297 + 2 x 3 = 303).
%        Le N de la a.3 compte 262 pages : cette valeur ne s'y applique plus.
%  E, A, NEA — non demandés à l'imprimeur : ils sont CALCULÉS avec la main
%        déduite ci-dessous. À faire confirmer avant tout tirage.
%
% La main de 1,206 n'est pas publiée par CoolLibri : elle est déduite de
% leur propre chiffre, 14 mm / (0,090 x 258/2) = 1,206. Elle reproduit donc
% exactement le dos de N, et n'est qu'une ESTIMATION pour les autres.
% Écarts que cela recouvre, pour mémoire : une main de 1,0 donnerait
% 11,61 mm pour N, une main de 1,3 en donnerait 16,77. Il ne fallait pas
% deviner.
%  E   — 17 mm, VALEUR DONNÉE PAR L'IMPRIMEUR le 17/09/2026 pour 276 pages,
%        20 x 24, couverture RIGIDE, reliure collée (Klebebindung), papier
%        115 g. En reliure cousue (Fadenheftung), il annonce 18 mm.
%
% \DosImposeReliure<classe> : la reliure pour laquelle la valeur a été
% donnée. Le dos d'une couverture rigide n'est pas celui d'un dos carré
% collé souple, même à pagination égale.
\def\DosImposeN{14}   \def\DosImposePagesN{258}  \def\DosImposeFormatN{a4}
\def\DosImposeE{17}   \def\DosImposePagesE{276}  \def\DosImposeFormatE{20x24-marge}
\def\DosImposeA{0}    \def\DosImposePagesA{0}    \def\DosImposeFormatA{}
\def\DosImposeNE{0}   \def\DosImposePagesNE{0}   \def\DosImposeFormatNE{}
\def\DosImposeEA{0}   \def\DosImposePagesEA{0}   \def\DosImposeFormatEA{}
\def\DosImposeNEA{0}  \def\DosImposePagesNEA{0}  \def\DosImposeFormatNEA{}
\def\DosImposeReliureN{souple}  \def\DosImposeReliureE{rigide}
\def\DosImposeReliureA{souple}  \def\DosImposeReliureNE{souple}
\def\DosImposeReliureEA{souple} \def\DosImposeReliureNEA{souple}

% Fond perdu, en mm. 3 mm est la valeur courante ; certains imprimeurs
% demandent 5 mm.
\def\FondPerdu{3}

% Marge technique autour du fond perdu, en mm. Sert à loger les traits de
% coupe en mode épreuve. Ignorée en mode production.
\def\MargeTech{12}

% Zone de sécurité, en mm : aucun texte ne doit s'en approcher davantage
% du bord de coupe.
\def\Securite{8}

% ---------------------------------------------------------------------
% 2. DONNÉES PAR CLASSE
%    pages : nombre de pages du PDF intérieur, COUVERTURE NON COMPRISE.
%    Un dos carré collé impose un nombre de pages PAIR ; le fichier
%    refuse de compiler si le compte est impair.
%
%    Ces tables ne servent que si \PagesForcees et \VersionForcee sont
%    vides, c'est-à-dire à la main. Valeurs A4 mesurées sur les PDF de la
%    a.3 (CLAUDE.md § 1, 05/09/2026). En 20 x 24, la pagination change
%    entièrement : NEA passe de 816 à 1036 pages.
% ---------------------------------------------------------------------

\def\DonneesN{262}
\def\DonneesE{214}
\def\DonneesA{386}
\def\DonneesNE{0}      % édition jamais produite
\def\DonneesEA{0}      % édition jamais produite
\def\DonneesNEA{816}

% Version imprimée sur le dos et la 1re de couverture, PAR CLASSE.
% « Les versions sont suivies par classe, chacune évoluant à son rythme »
% (CHANGELOG.md).
\def\VersionN{a.3}
\def\VersionE{a.3}
\def\VersionA{a.3}
\def\VersionNE{a.3}
\def\VersionEA{a.3}
\def\VersionNEA{a.3}

% ---------------------------------------------------------------------
% 3. TEXTES
% ---------------------------------------------------------------------

\def\TitreOuvrage{50\,Ohm}
\def\SousTitre{Préparation à l'examen radioamateur}

% Sous-titre propre à chaque classe (repris de FR_TITLES, build_book.py).
\def\SousTitreN{Cours complet — classe N}
\def\SousTitreE{Cours complet — classe E}
\def\SousTitreA{Cours complet — classe A}
\def\SousTitreNE{Cours complet — classes N et E}
\def\SousTitreEA{Cours complet — classes E et A}
\def\SousTitreNEA{Cours complet — classes N, E et A}

% Lettres du filigrane, empilées de haut en bas (cf. build_book.py v0.20).
\def\LettresN{N}
\def\LettresE{E}
\def\LettresA{A}
\def\LettresNE{N,E}
\def\LettresEA{E,A}
\def\LettresNEA{N,E,A}

% Texte de 4e de couverture.
% NB : il parle des classes N, E et A quelle que soit l'édition. Juste pour
% N, E, A et NEA ; approximatif pour NE et EA ; faux pour SWL, d'où son refus.
\def\QuatriemeTexte{%
Ce guide d'étude prépare aux examens radioamateur allemands des classes
N, E et A. Il traduit en français les contenus du projet 50ohm.de, publiés
par le référat AJW du DARC e.\,V., et les complète d'encarts « En France »
qui signalent, sans rien retirer au texte allemand, ce que la
réglementation française prévoit sur le même sujet.\par\smallskip
Les questions officielles du catalogue allemand accompagnent chaque
notion, avec leurs réponses.%
}

\def\Auteur{Traduit avec l'aide d'une IA par Pierre F4JWI}
\def\Depot{github.com/DrPLG/50ohm-fr}

% =====================================================================
%  À PARTIR D'ICI, PLUS RIEN À MODIFIER
% =====================================================================

\usepackage{iftex}
\RequireLuaTeX
\usepackage{xfp}
\usepackage{tikz}
\usetikzlibrary{calc}
\usepackage{fontspec}
\usepackage[french]{babel}
\usepackage{geometry}

% ---------------------------------------------------------------------
% 4. RÉSOLUTION DES DONNÉES DE LA CLASSE
% ---------------------------------------------------------------------

\makeatletter
% Tome SWL : refus explicite, plutôt qu'un « classe inconnue » qui laisserait
% croire à un oubli. Le texte de 4e parle des classes N, E et A.
\def\@swlmot{SWL}
\ifx\Classe\@swlmot
  \PackageError{couverture}{Pas de couverture pour le tome SWL}%
    {Sa 4e de couverture n'est pas redigee : le texte commun parle des
     classes N, E et A. C'est une decision editoriale a prendre d'abord.}%
\fi

\def\@resoudre#1#2{%
  \expandafter\let\expandafter#1\csname #2\Classe\endcsname
  \ifx#1\relax
    \PackageError{couverture}{Classe \Classe\space inconnue : pas de #2\Classe}%
      {Valeurs admises : N, E, A, NE, EA, NEA.}%
  \fi}
\@resoudre\NbPages{Donnees}
\@resoudre\SousTitreClasse{SousTitre}
\@resoudre\Lettres{Lettres}
\@resoudre\DosImpose{DosImpose}
\@resoudre\DosImposePages{DosImposePages}
\@resoudre\DosImposeFormat{DosImposeFormat}
\@resoudre\DosImposeReliure{DosImposeReliure}
\@resoudre\Version{Version}

% Valeurs imposées de l'extérieur. Test de vacuité par \detokenize : il ne
% dépend pas de la façon dont la macro a été définie.
\if\relax\detokenize\expandafter{\PagesForcees}\relax\else
  \let\NbPages\PagesForcees
\fi
\if\relax\detokenize\expandafter{\VersionForcee}\relax\else
  \let\Version\VersionForcee
\fi

% Format du livre fini.
\def\@formatAquatre{a4}
\def\@formatVingt{20x24}
\def\@formatVingtMarge{20x24-marge}
\ifx\Format\@formatAquatre
  \def\PageLargeur{210}\def\PageHauteur{297}
\else\ifx\Format\@formatVingt
  \def\PageLargeur{200}\def\PageHauteur{240}
\else\ifx\Format\@formatVingtMarge
  \def\PageLargeur{200}\def\PageHauteur{240}
\else
  \PackageError{couverture}{Format \Format\space inconnu}%
    {Valeurs admises : a4, 20x24, 20x24-marge.}%
  \def\PageLargeur{210}\def\PageHauteur{297}
\fi\fi\fi
% Reliure. Hors de « souple » et « rigide » : erreur, pas de supposition.
\def\@reliuresouple{souple}
\def\@reliurerigide{rigide}
\ifx\Reliure\@reliurerigide
  \edef\PageLargeur{\fpeval{\PageLargeur + 5}}
  \edef\PageHauteur{\fpeval{\PageHauteur + 5}}
  \def\FondPerdu{15}
\else\ifx\Reliure\@reliuresouple\else
  \PackageError{couverture}{Reliure \Reliure\space inconnue}%
    {Valeurs admises : souple, rigide.}%
\fi\fi
\makeatother

% Garde-fou : une édition jamais produite n'a pas de couverture.
\ifnum\NbPages=0
  \PackageError{couverture}{L'edition \Classe\space n'a pas de pagination}%
    {Cette edition n'a jamais ete compilee. Renseigner \string\Donnees\Classe
     ou passer \string\PagesForcees.}%
\fi

% Avertissement : l'imprimeur demande un nombre de pages multiple de 4
% (17/09/2026). Pas une erreur : c'est son exigence, pas une loi du papier.
\ifnum\fpeval{\NbPages - 4*floor(\NbPages/4)}=0 \else
  \typeout{COUVERTURE-AVERTISSEMENT \NbPages\space pages : pas un multiple
    de 4, que l'imprimeur exige -- ajouter des pages a l'interieur}%
\fi

% Garde-fou : un dos carré collé impose un nombre de pages pair.
\ifodd\NbPages
  \PackageError{couverture}{Nombre de pages impair (\NbPages)}%
    {Un dos carre colle exige un nombre PAIR de pages : une feuille en
     porte deux. Ajouter une page blanche a l'interieur.}%
\fi

% ---------------------------------------------------------------------
% 5. GÉOMÉTRIE — c'est ici que se joue l'exactitude du dos
% ---------------------------------------------------------------------

% Épaisseur du dos : valeur imposée si elle est donnée ET qu'elle vaut pour
% cette pagination et ce format ; calcul sinon.
\edef\DosCalcule{\fpeval{round(\NbPages / 2 * \Grammage * \Bouffant / 1000, 2)}}
\def\DosOrigine{calculee}
\ifnum\fpeval{\DosImpose > 0}=1
  \ifnum\fpeval{\NbPages = \DosImposePages}=1
    \ifx\Format\DosImposeFormat
      \ifx\Reliure\DosImposeReliure
        \def\DosOrigine{imposee}
      \fi
    \fi
  \fi
  \def\MotCalculee{calculee}
  \ifx\DosOrigine\MotCalculee
    \typeout{COUVERTURE-AVERTISSEMENT dos impose de \DosImpose\space mm
      donne pour \DosImposePages\space pages en \DosImposeFormat\space
      (reliure \DosImposeReliure) -- inapplicable a \NbPages\space pages en
      \Format\space (reliure \Reliure), dos calcule}%
  \fi
\fi
% NB : pas de \@ dans ces noms. Cette section est hors \makeatletter, et
% « \def\@imposee{...} » y redéfinit \@ (suivi du texte « imposee ») : le
% test échouait toujours, en silence. Constaté le 16/09/2026 — le dos de N
% tombait juste par coïncidence, 14,0 mm calculés pour 14 imposés.
\def\MotRigide{rigide}
\def\MotImposee{imposee}
\ifx\DosOrigine\MotImposee
  \edef\Dos{\DosImpose}
\else
  \edef\Dos{\DosCalcule}
\fi

% Marge technique : uniquement en mode épreuve.
\def\ouimot{oui}
\ifx\Reperes\ouimot \edef\Tech{\MargeTech}\else \def\Tech{0}\fi

% Dimensions du papier.
%   largeur = marge + fond perdu + 4e + dos + 1re + fond perdu + marge
%   hauteur = marge + fond perdu + hauteur + fond perdu + marge
\edef\PapierL{\fpeval{2*\PageLargeur + \Dos + 2*\FondPerdu + 2*\Tech}}
\edef\PapierH{\fpeval{\PageHauteur + 2*\FondPerdu + 2*\Tech}}

\geometry{paperwidth=\PapierL mm, paperheight=\PapierH mm,
          margin=0pt, nohead, nofoot}

% ---------------------------------------------------------------------
%  TOUTES LES COTES SONT FIGÉES ICI, ET NULLE PART AILLEURS.
%
%  Raison, payée par trois compilations ratées : un \fpeval laissé dans
%  une coordonnée ou dans le texte d'un nœud TikZ n'est pas développé au
%  moment où TikZ analyse le chemin. On obtient « Giving up on this path.
%  Did you forget a semicolon? », qui accuse une ponctuation alors que la
%  faute est ailleurs. Le corps du document ci-dessous n'emploie donc que
%  des macros déjà calculées.
% ---------------------------------------------------------------------

% Repères horizontaux, mesurés depuis le bord GAUCHE du papier.
\edef\CoupeG{\fpeval{\Tech + \FondPerdu}}          % coupe gauche (4e de couv.)
\edef\DosG{\fpeval{\CoupeG + \PageLargeur}}         % dos, début
\edef\DosD{\fpeval{\DosG + \Dos}}                   % dos, fin
\edef\CoupeD{\fpeval{\DosD + \PageLargeur}}         % coupe droite (1re de couv.)
% Repères verticaux, depuis le bord BAS.
\edef\CoupeB{\fpeval{\Tech + \FondPerdu}}
\edef\CoupeH{\fpeval{\CoupeB + \PageHauteur}}
% Axes.
\edef\DosAxe{\fpeval{(\DosG + \DosD)/2}}
\edef\MilieuH{\fpeval{(\CoupeB + \CoupeH)/2}}
\edef\MilieuL{\fpeval{\PapierL/2}}

% Format fini, pour les cotes affichées.
\edef\FiniL{\fpeval{2*\PageLargeur + \Dos}}

% Bandeau de la 1re de couverture : même proportion que la page de titre.
\edef\BandeauL{\fpeval{round(0.34*\PageLargeur, 2)}}
\edef\BandeauG{\fpeval{\CoupeD - \BandeauL}}
\edef\BandeauAccent{\fpeval{\BandeauG - 0.7}}
\edef\DosAccentG{\fpeval{\DosG - 0.7}}
\edef\DosAccentD{\fpeval{\DosD + 0.7}}
% Couverture rigide : le bleu du dos déborde de 2 mm sur chaque plat, pour
% absorber la tolérance de l'imprimeur sur l'épaisseur du dos (±10 %).
\edef\DosEtenduG{\fpeval{\DosG - 2}}
\edef\DosEtenduD{\fpeval{\DosD + 2}}

% Filigrane de classe.
\edef\FiligraneX{\fpeval{(\BandeauG + \CoupeD)/2}}
\edef\FiligraneY{\fpeval{\CoupeH - 6}}

% Bloc-titre de la 1re de couverture.
\edef\TitreX{\fpeval{\DosD + 12}}
\edef\TitreY{\fpeval{\CoupeH - 72}}
\edef\TitreW{\fpeval{\BandeauG - \DosD - 16}}
\edef\VersionY{\fpeval{\CoupeB + 14}}

% 4e de couverture.
\edef\QuatreX{\fpeval{\CoupeG + \Securite + 6}}
\edef\QuatreY{\fpeval{\CoupeH - \Securite - 14}}
\edef\QuatreBasY{\fpeval{\CoupeB + \Securite + 22}}
\edef\QuatreW{\fpeval{\PageLargeur - 2*\Securite - 12}}

% Zone de sécurité.
\edef\SecG{\fpeval{\CoupeG + \Securite}}
\edef\SecD{\fpeval{\CoupeD - \Securite}}
\edef\SecB{\fpeval{\CoupeB + \Securite}}
\edef\SecH{\fpeval{\CoupeH - \Securite}}

% Zones de contrôle : ce qui doit rester de part et d'autre du pli du dos.
\edef\PliQuatre{\fpeval{\DosG - \Securite}}         % 4e : limite côté dos
\edef\PliPremiere{\fpeval{\DosD + \Securite}}       % 1re : limite côté dos

% Traits de coupe et cotes, dans la marge technique.
\edef\TechExt{\fpeval{\Tech - 8}}
\edef\TechInt{\fpeval{\Tech - 1}}
\edef\TechCote{\fpeval{\Tech - 9}}
\edef\HautInt{\fpeval{\PapierH - \Tech + 1}}
\edef\HautExt{\fpeval{\PapierH - \Tech + 8}}
\edef\HautCote{\fpeval{\PapierH - \Tech + 9}}
\edef\DroitInt{\fpeval{\PapierL - \Tech + 1}}
\edef\DroitExt{\fpeval{\PapierL - \Tech + 8}}

% Emplacement du code-barres ISBN : en bas à DROITE de la 4e de couverture,
% selon l'usage éditorial, et non à gauche où il chevaucherait les mentions
% de licence.
\edef\IsbnD{\fpeval{\DosG - \Securite}}
\edef\IsbnG{\fpeval{\IsbnD - 40}}
\edef\IsbnH{\fpeval{\SecB + 25}}
\edef\IsbnCX{\fpeval{\IsbnG + 20}}
\edef\IsbnCY{\fpeval{\SecB + 12.5}}

% Cartouches de cotes, dans la marge technique. Ancrés vers l'INTÉRIEUR du
% papier : ancrés vers l'extérieur, ils débordaient et étaient rognés.
\edef\CoteBasY{\fpeval{1}}
\edef\CoteHautY{\fpeval{\PapierH - 1}}

% ---------------------------------------------------------------------
% 6. COULEURS ET POLICES — reprises de la page de titre (build_book.py)
% ---------------------------------------------------------------------

\definecolor{TitleBand}{cmyk}{.9,.55,.1,.35}   % bleu profond
\definecolor{TitleAccent}{cmyk}{.05,.8,1,0}    % rouge DARC (accent fin)
\definecolor{Repere}{cmyk}{0,1,1,0}            % magenta : repères de montage

\setmainfont{Libertinus Serif}[
  ItalicFont     = Libertinus Serif Italic,
  BoldFont       = Libertinus Serif Bold,
  BoldItalicFont = Libertinus Serif Bold Italic]

\pagestyle{empty}

% ---------------------------------------------------------------------
% 7. FILIGRANE DE CLASSE
%    Le texte est construit ici, et non dans le nœud : un \\ produit à
%    l'intérieur d'un \foreach referme le groupe d'alignement de TikZ.
%    Une lettre seule est grande ; plusieurs s'empilent et se réduisent
%    (build_book.py v0.20 : « NEA » à 105 pt débordait encore).
% ---------------------------------------------------------------------

\makeatletter
\def\FiligraneTexte{}
\newif\ifpremierelettre \premierelettretrue
\count255=0
\@for\@tempa:=\Lettres\do{%
  \advance\count255 by 1\relax
  \ifpremierelettre \premierelettrefalse
  \else \g@addto@macro\FiligraneTexte{\\}\fi
  \expandafter\g@addto@macro\expandafter\FiligraneTexte\expandafter{\@tempa}}%
\xdef\NbLettres{\the\count255}
\makeatother

\ifnum\NbLettres=1 \def\CorpsF{150}\else\def\CorpsF{86}\fi
\edef\CorpsFI{\fpeval{round(\CorpsF*1.02, 0)}}

% Texte du cartouche de cotes, lui aussi construit hors du nœud TikZ.
%
% NE PAS y employer \textsuperscript : dans un nœud « align=center », il
% casse le chemin TikZ (« Giving up on this path. Did you forget a
% semicolon? »), et l'erreur est signalée sur le nœud SUIVANT — elle
% accuse donc une ponctuation là où la faute est ici. Isolé par essais :
% $^2$ et le caractère ² passent, \textsuperscript non.
\def\CotesTexte{%
  \Classe\ --- \NbPages\ pages --- papier \Grammage\ g/m$^2$,
  bouffant \Bouffant\ --- fond perdu \FondPerdu\ mm\\
  papier \PapierL\ mm sur \PapierH\ mm --- fini \FiniL\ mm sur
  \PageHauteur\ mm\\
  EPREUVE, ne pas envoyer a l'impression%
}

% ---------------------------------------------------------------------
% 8. CONTRÔLE DE SUPERPOSITION
%
%  \MesureNoeud{nœud}{préfixe}{retrait} mesure la boîte d'un nœud par ses
%  quatre coins (un nœud pivoté a ses ancres pivotées : le min et le max des
%  quatre suffisent) et la ramène au repère du papier, en mm. Le point (0,0)
%  est lu par le même « let » que les ancres : les deux passent donc par la
%  même transformation, et la comparaison ne dépend pas de l'endroit où
%  TeX a posé l'image sur la page.
%
%  Le retrait sert au seul filigrane, pour mesurer les lettres et non leur
%  cadre : c'est le bord des lettres que le détourage rognerait.
%
%  PIÈGE MESURÉ le 16/09/2026 : la marge interne TikZ (.3333em) n'est PAS
%  évaluée dans la police du nœud, mais dans celle qui entoure le dessin.
%  Instrumenté : dans le nœud, em = 150 pt ; boîte brute du « N » :
%  41,7 x 37,5 mm, soit la lettre plus 2 x 1,3 mm — 1,3 mm étant .3333em
%  à 11 pt. Un retrait calculé à 150 pt (17,6 mm) ramenait la lettre à
%  6,6 mm de large, et le contrôle aurait dit OK sur une mesure absurde.
% ---------------------------------------------------------------------

\newcount\NbAlertes
\newdimen\FiligraneSep
\newsavebox\CreditsBoite

% Crédits de 4e : lignes coupées à la main. Leur largeur RÉELLE est celle
% de la plus longue ligne, mesurée dans une boîte à part (voir § 8, ISBN).
\def\CreditsTexte{%
    Réalisé à partir des contenus de 50ohm.de (en allemand)\\
    50ohm.de-Autorenteam, coordonné par le référat AJW du DARC e.\,V.\\
    \Auteur\\[3pt]
    Licence CC BY 4.0 --- \texttt{\Depot}}

\def\MesureNoeud#1#2#3{%
  \path let \p1=(#1.south west), \p2=(#1.north east),
            \p3=(#1.north west), \p4=(#1.south east), \p5=(0,0) in
    \pgfextra{%
      \expandafter\xdef\csname #2xmin\endcsname{\fpeval{round((min(\x1,\x2,\x3,\x4) - \x5 + #3)/1mm, 2)}}%
      \expandafter\xdef\csname #2xmax\endcsname{\fpeval{round((max(\x1,\x2,\x3,\x4) - \x5 - #3)/1mm, 2)}}%
      \expandafter\xdef\csname #2ymin\endcsname{\fpeval{round((min(\y1,\y2,\y3,\y4) - \y5 + #3)/1mm, 2)}}%
      \expandafter\xdef\csname #2ymax\endcsname{\fpeval{round((max(\y1,\y2,\y3,\y4) - \y5 - #3)/1mm, 2)}}%
    };}

% \ControleDans{libellé}{préfixe}{xmin}{xmax}{ymin}{ymax}
\def\ControleDans#1#2#3#4#5#6{%
  \ifnum\fpeval{\csname #2xmin\endcsname >= #3 && \csname #2xmax\endcsname <= #4
             && \csname #2ymin\endcsname >= #5 && \csname #2ymax\endcsname <= #6}=1
    \typeout{COUVERTURE-CONTROLE OK #1 : x \csname #2xmin\endcsname-\csname #2xmax\endcsname,
      y \csname #2ymin\endcsname-\csname #2ymax\endcsname\space dans x #3-#4, y #5-#6}%
  \else
    \global\advance\NbAlertes by 1
    \typeout{COUVERTURE-CONTROLE ALERTE #1 : x \csname #2xmin\endcsname-\csname #2xmax\endcsname,
      y \csname #2ymin\endcsname-\csname #2ymax\endcsname\space HORS de x #3-#4, y #5-#6}%
  \fi}

% \ControleDisjoints{libellé}{préfixe}{xmin}{xmax}{ymin}{ymax}
% Deux rectangles se chevauchent si et seulement s'ils se recouvrent sur
% les DEUX axes à la fois.
\def\ControleDisjoints#1#2#3#4#5#6{%
  \ifnum\fpeval{\csname #2xmin\endcsname < #4 && #3 < \csname #2xmax\endcsname
             && \csname #2ymin\endcsname < #6 && #5 < \csname #2ymax\endcsname}=1
    \global\advance\NbAlertes by 1
    \typeout{COUVERTURE-CONTROLE ALERTE #1 : chevauchement}%
  \else
    \typeout{COUVERTURE-CONTROLE OK #1 : disjoints}%
  \fi}

\begin{document}
% Marge interne du filigrane, évaluée dans la police qui ENTOURE le dessin :
% c'est là que TikZ évalue le .3333em de « inner sep » (cf. § 8).
\FiligraneSep=.3333em
% Largeur réelle des crédits : la plus longue ligne, dans leur police.
\sbox\CreditsBoite{\small\begin{tabular}{@{}l@{}}\CreditsTexte\end{tabular}}%
\begin{tikzpicture}[remember picture, overlay,
                    x=1mm, y=1mm, shift={(current page.south west)}]

% ---- Fonds ----------------------------------------------------------
% La 4e de couverture reste sur le blanc du papier.
% Le bandeau bleu couvre le dos, puis reprend au tiers extérieur de la
% 1re de couverture — même partition 2/3-1/3 que la page de titre.
% Couverture rigide : pas de filets au dos (ils tomberaient à côté des
% plis, vu la tolérance), et un dos bleu élargi de 2 mm de chaque côté.
\ifx\Reliure\MotRigide
  \fill[TitleBand] (\DosEtenduG, 0) rectangle (\DosEtenduD, \PapierH);
\else
  \fill[TitleBand] (\DosG, 0) rectangle (\DosD, \PapierH);
\fi
\fill[TitleBand] (\BandeauG, 0) rectangle (\PapierL, \PapierH);
% Filets d'accent : à la jonction sur la 1re, de part et d'autre du dos.
\fill[TitleAccent] (\BandeauAccent, 0) rectangle (\BandeauG, \PapierH);
\ifx\Reliure\MotRigide\else
  \fill[TitleAccent] (\DosAccentG, 0) rectangle (\DosG, \PapierH);
  \fill[TitleAccent] (\DosD, 0) rectangle (\DosAccentD, \PapierH);
\fi

% ---- 1re de couverture ----------------------------------------------
% Filigrane de classe, en réserve claire, détouré au bandeau.
\begin{scope}
\clip (\BandeauG, 0) rectangle (\PapierL, \PapierH);
\node[name=filigrane, anchor=north, align=center, text=white, opacity=0.16,
      font=\fontsize{\CorpsF}{\CorpsFI}\selectfont\bfseries]
  at (\FiligraneX, \FiligraneY) {\FiligraneTexte};
\end{scope}

\node[name=titre, anchor=west, align=left, text width=\TitreW mm]
  at (\TitreX, \TitreY) {%
    {\fontsize{54}{56}\selectfont\bfseries \TitreOuvrage}\\[10pt]
    {\Large\color{TitleBand}\bfseries \SousTitre}\\[16pt]
    {\large\color{black!75} \SousTitreClasse}};

\node[name=version, anchor=south west, align=left, font=\small, text=black!70]
  at (\TitreX, \VersionY) {Version \Version\\ Licence CC BY 4.0};

% ---- Dos -------------------------------------------------------------
% Un dos étroit ne reçoit pas de texte : en deçà de 6 mm, le débord au
% massicot ferait mordre le texte vertical sur les plats.
\ifnum\fpeval{\Dos >= 6}=1
  \node[name=dos, rotate=90, anchor=center, text=white, font=\bfseries,
        inner sep=0pt]
    at (\DosAxe, \MilieuH)
    {\TitreOuvrage\ \ --- \ \SousTitreClasse\ \ --- \ \Version};
\fi

% ---- 4e de couverture ------------------------------------------------
\node[name=quatre, anchor=north west, align=justify, text width=\QuatreW mm]
  at (\QuatreX, \QuatreY) {\QuatriemeTexte};

% Le texte vit dans \CreditsTexte (§ 8) pour être mesuré à part : la boîte
% du nœud a la largeur fixe \QuatreW et couvre l'emplacement ISBN sans
% qu'aucun caractère n'y soit. Retirer « text width » pour l'éviter change
% l'interlignage — essayé le 16/09/2026, 6 934 pixels différents.
\node[name=credits, anchor=south west, align=left, font=\small, text=black!70,
      text width=\QuatreW mm]
  at (\QuatreX, \QuatreBasY) {\CreditsTexte};

% ---- Repères de fabrication (mode épreuve uniquement) ----------------
\ifx\Reperes\ouimot
  % Traits de coupe, hors du fond perdu.
  \foreach \x in {\CoupeG, \CoupeD} {
    \draw[Repere, line width=0.25pt] (\x, \TechExt) -- (\x, \TechInt);
    \draw[Repere, line width=0.25pt] (\x, \HautInt) -- (\x, \HautExt);}
  \foreach \y in {\CoupeB, \CoupeH} {
    \draw[Repere, line width=0.25pt] (\TechExt, \y) -- (\TechInt, \y);
    \draw[Repere, line width=0.25pt] (\DroitInt, \y) -- (\DroitExt, \y);}
  % Limites du dos, en pointillés.
  \foreach \x in {\DosG, \DosD} {
    \draw[Repere, line width=0.25pt, dashed] (\x, \TechExt) -- (\x, \TechInt);
    \draw[Repere, line width=0.25pt, dashed] (\x, \HautInt) -- (\x, \HautExt);}
  % Zone de sécurité.
  \draw[Repere, line width=0.2pt, dotted]
    (\SecG, \SecB) rectangle (\SecD, \SecH);
  % Emplacement réservé au code-barres ISBN.
  \draw[Repere, line width=0.2pt] (\IsbnG, \SecB) rectangle (\IsbnD, \IsbnH);
  \node[font=\tiny\sffamily, text=Repere, align=center]
    at (\IsbnCX, \IsbnCY) {emplacement\\ISBN};
  % Cotes.
  \node[anchor=north, font=\tiny\sffamily, text=Repere]
    at (\DosAxe, \CoteHautY) {dos \Dos\ mm (\DosOrigine)};
  \node[anchor=south, font=\tiny\sffamily, text=Repere, align=center]
    at (\MilieuL, \CoteBasY) {\CotesTexte};
\fi

% ---- Contrôle de superposition (§ 8) ----------------------------------
% Rien n'est dessiné : des mesures et des lignes de journal.
\MesureNoeud{filigrane}{Fil}{\the\FiligraneSep}
\MesureNoeud{titre}{Tit}{0pt}
\MesureNoeud{version}{Ver}{0pt}
\MesureNoeud{quatre}{Qua}{0pt}
\MesureNoeud{credits}{Cre}{0pt}
% Bord droit des crédits = bord gauche + marge interne + plus longue ligne.
\xdef\Crexmax{\fpeval{round(\Crexmin + (\the\FiligraneSep + \the\wd\CreditsBoite)/1mm, 2)}}

% Le filigrane peut entrer dans la zone de sécurité (il est décoratif),
% pas franchir la coupe ni sortir du bandeau — le détourage le rognerait.
\ControleDans{filigrane dans le bandeau}{Fil}{\BandeauG}{\CoupeD}{\CoupeB}{\CoupeH}
\ControleDans{titre sur la 1re, hors bandeau}{Tit}{\PliPremiere}{\BandeauAccent}{\SecB}{\SecH}
\ControleDans{version sur la 1re, hors bandeau}{Ver}{\PliPremiere}{\BandeauAccent}{\SecB}{\SecH}
\ControleDans{texte de 4e dans la zone de securite}{Qua}{\SecG}{\PliQuatre}{\SecB}{\SecH}
\ControleDans{credits de 4e dans la zone de securite}{Cre}{\SecG}{\PliQuatre}{\SecB}{\SecH}
\ControleDisjoints{titre / version}{Tit}{\Verxmin}{\Verxmax}{\Verymin}{\Verymax}
\ControleDisjoints{texte de 4e / credits}{Qua}{\Crexmin}{\Crexmax}{\Creymin}{\Creymax}
\ControleDisjoints{texte de 4e / emplacement ISBN}{Qua}{\IsbnG}{\IsbnD}{\SecB}{\IsbnH}
\ControleDisjoints{credits / emplacement ISBN}{Cre}{\IsbnG}{\IsbnD}{\SecB}{\IsbnH}
\ifnum\fpeval{\Dos >= 6}=1
  \MesureNoeud{dos}{Dos}{0pt}
  \ControleDans{texte du dos entre les plis}{Dos}{\DosG}{\DosD}{\SecB}{\SecH}
\else
  \typeout{COUVERTURE-CONTROLE OK dos de \Dos\space mm : trop etroit, sans texte}%
\fi

\end{tikzpicture}
\typeout{COUVERTURE-BILAN alertes=\the\NbAlertes\space; classe=\Classe\space;
  format=\Format\space; reliure=\Reliure\space; pages=\NbPages\space;
  version=\Version\space;
  dos=\Dos\space mm (\DosOrigine); fini=\FiniL x \PageHauteur\space mm;
  papier=\PapierL x \PapierH\space mm}
\end{document}
"
%
%  Pour le fichier destiné à l'imprimeur (sans repères) :
%
%      lualatex "\def\Reperes{non}% =====================================================================
%  couverture.tex — couverture dos carré collé pour les livres 50ohm-fr
% =====================================================================
%
%  Produit une couverture d'un seul tenant : 4e de couverture | dos |
%  1re de couverture, aux dimensions exactes du livre relié, fond perdu
%  compris. L'épaisseur du dos est CALCULÉE à partir du nombre de pages,
%  du grammage et du bouffant du papier.
%
%  ORIGINE ET ÉTAT. Écrit le 20/08/2026 dans Fichierstravail/ (hors git),
%  versé au dépôt le 16/09/2026 pour que fenetre_compilation.py puisse
%  l'appeler. Ce qui a changé à cette occasion :
%
%    - le FORMAT est un paramètre (\Format = a4, 20x24 ou 20x24-marge,
%      mêmes noms que --format de build_book.py) ;
%    - la pagination et la version peuvent être IMPOSÉES de l'extérieur
%      (\PagesForcees, \VersionForcee) : la fenêtre lit la pagination dans
%      le journal du livre qu'elle vient de compiler, au lieu de se fier à
%      une table qui vieillit ;
%    - un dos imposé par l'imprimeur n'est appliqué QUE pour la pagination
%      et le format pour lesquels il a été donné ;
%    - un CONTRÔLE DE SUPERPOSITION est calculé par LaTeX lui-même, à partir
%      des boîtes réelles des nœuds, et écrit dans le journal ;
%    - le tome SWL est refusé explicitement : sa 4e de couverture n'est pas
%      rédigée (le texte commun parle des classes N, E et A).
%
%    Neutralité vérifiée : en A4, pour N (258 p., a.2) et NEA (816 p., a.3),
%    en épreuve comme en impression, le rendu est identique au pixel près à
%    celui du gabarit d'origine.
%
%  COMPILATION (LuaLaTeX, depuis la racine du dépôt) :
%
%      lualatex couverture.tex
%      lualatex couverture.tex        <-- DEUX FOIS, obligatoirement
%
%  DEUX PASSES SONT NÉCESSAIRES. Le dessin s'ancre sur « current page »,
%  dont la position n'est connue qu'après un premier passage : à la
%  première compilation, tout se tasse dans le coin supérieur et la
%  couverture paraît ratée alors qu'elle ne l'est pas. Aucune erreur n'est
%  émise — seul le rendu le montre. Le contrôle de superposition, lui
%  aussi, ne vaut qu'à la seconde passe.
%
%  Le plus simple est de passer par fenetre_compilation.py (case
%  « couverture pour l'imprimeur »), qui enchaîne les deux passes, produit
%  l'épreuve ET le fichier d'impression, et relit le contrôle.
%
%  À la main, pour une autre classe, un autre format, une pagination mesurée :
%
%      lualatex "\def\Classe{A}\def\Format{20x24}\def\PagesForcees{482}\input{couverture.tex}"
%
%  Pour le fichier destiné à l'imprimeur (sans repères) :
%
%      lualatex "\def\Reperes{non}\input{couverture.tex}"
%
%  ---------------------------------------------------------------------
%  DEUX MODES, un seul paramètre : \Reperes
%
%    oui  — mode ÉPREUVE. Ajoute une marge technique, les traits de coupe,
%           les limites du dos, la zone de sécurité et les cotes. C'est ce
%           qu'il faut pour vérifier à l'écran.
%    non  — mode PRODUCTION. Papier = format final + fond perdu, rien
%           d'autre. C'est ce fichier-là qu'on envoie à l'imprimeur.
%
%  ---------------------------------------------------------------------
%  CONTRÔLE DE SUPERPOSITION
%
%  À la fin du dessin, chaque bloc de texte est mesuré par ses ancres TikZ
%  (quatre coins, ce qui vaut aussi pour le texte pivoté du dos) et comparé
%  à la zone où il doit tenir, puis aux blocs voisins. Le journal porte une
%  ligne par vérification :
%
%      COUVERTURE-CONTROLE OK      <quoi> : ...
%      COUVERTURE-CONTROLE ALERTE  <quoi> : ...
%
%  et un bilan : COUVERTURE-BILAN alertes=<n> ; ...
%
%  Ce que le contrôle VOIT : un bloc qui franchit le pli du dos, la coupe,
%  la zone de sécurité ou le bord du bandeau ; deux blocs qui se chevauchent ;
%  un bloc qui empiète sur l'emplacement réservé au code-barres ISBN.
%  Ce qu'il NE VOIT PAS : une ligne trop longue à l'INTÉRIEUR d'un bloc à
%  largeur fixe (la boîte garde sa largeur, le texte déborde). Celle-là se
%  signale par un « Overfull \hbox » dans le même journal, que la fenêtre
%  compte aussi.
%
%  ---------------------------------------------------------------------
%  AVERTISSEMENT SUR L'ÉPAISSEUR DU DOS
%
%  Le dos se calcule ainsi :
%
%      épaisseur d'une feuille (mm) = grammage x bouffant / 1000
%      dos (mm)                     = (nombre de pages / 2) x épaisseur
%
%  On divise par deux parce qu'une feuille porte DEUX pages.
%
%  Le « bouffant » (main, volume, bulk) dépend du papier et n'est pas
%  devinable : un offset 80 g courant est autour de 1,30, un bouffant
%  d'édition monte à 1,75 et plus. UN ÉCART DE BOUFFANT DÉPLACE LE DOS DE
%  PLUSIEURS MILLIMÈTRES. Demander à l'imprimeur soit son bouffant, soit
%  directement l'épaisseur du dos, et reporter ci-dessous. Tant que cette
%  valeur n'est pas confirmée, la couverture n'est pas bonne à tirer.
%
%  Si l'imprimeur donne le dos en clair, forcer la valeur avec \DosImpose.
% =====================================================================

\documentclass[11pt]{article}

% ---------------------------------------------------------------------
% 1. PARAMÈTRES — c'est la seule zone à modifier
% ---------------------------------------------------------------------

% Classe de l'ouvrage : N, E, A, NE, EA ou NEA.
% (\providecommand : une valeur passée en ligne de commande l'emporte.)
\providecommand{\Classe}{NEA}

% Mode épreuve (oui) ou production (non).
\providecommand{\Reperes}{oui}

% Format du livre FINI : a4, 20x24 ou 20x24-marge — les noms de --format
% dans build_book.py. Les deux maquettes 20 x 24 ne diffèrent que par leur
% colonne de marge intérieure : le livre fini, donc la couverture, est le
% même (200 x 240 mm).
\providecommand{\Format}{a4}

% Reliure : « souple » (dos carré collé, fond perdu 3 mm) ou « rigide »
% (couverture rigide, dite hardcover). En rigide, consignes de l'imprimeur
% du 17/09/2026 : chaque plat déborde du livre de 5 mm en largeur et en
% hauteur (le carton dépasse du bloc), et le fichier porte 15 mm de rabat
% autour du carton au lieu de 3 mm de fond perdu. Les aplats vont déjà au
% bord du papier : ils remplissent le rabat d'eux-mêmes.
\providecommand{\Reliure}{souple}

% Pagination et version imposées de l'extérieur. Vides = tables du § 2.
% La fenêtre les renseigne toujours : la pagination vient du journal du
% livre réellement compilé, dans le format demandé.
\providecommand{\PagesForcees}{}
\providecommand{\VersionForcee}{}

% Papier de l'INTÉRIEUR du livre (pas de la couverture).
% CoolLibri, papier « Standard 90 g blanc ».
\def\Grammage{90}        % g/m²
\def\Bouffant{1.206}     % main DÉDUITE, voir ci-dessous

% Épaisseur du dos IMPOSÉE par l'imprimeur, en mm, par classe.
% Une valeur > 0 prime sur le calcul — MAIS SEULEMENT pour la pagination et
% le format pour lesquels l'imprimeur l'a donnée (\DosImposePages<classe>,
% \DosImposeFormat<classe>). Hors de ce cas, le dos est calculé et le
% journal le dit : un dos de 14 mm donné pour 258 pages A4 ne vaut rien pour
% 262 pages, ni pour le même livre recomposé en 20 x 24.
%
% ATTENTION : cette valeur est propre à CHAQUE ÉDITION. Un dos unique
% appliqué aux quatre livres serait faux pour trois d'entre eux.
%
%  N   — 14 mm, VALEUR DONNÉE PAR COOLIBRI pour 258 pages en 90 g, A4.
%        Fichier attendu : 440 x 303 mm, ce que ce gabarit produit en mode
%        production (2 x 210 + 14 + 2 x 3 = 440 ; 297 + 2 x 3 = 303).
%        Le N de la a.3 compte 262 pages : cette valeur ne s'y applique plus.
%  E, A, NEA — non demandés à l'imprimeur : ils sont CALCULÉS avec la main
%        déduite ci-dessous. À faire confirmer avant tout tirage.
%
% La main de 1,206 n'est pas publiée par CoolLibri : elle est déduite de
% leur propre chiffre, 14 mm / (0,090 x 258/2) = 1,206. Elle reproduit donc
% exactement le dos de N, et n'est qu'une ESTIMATION pour les autres.
% Écarts que cela recouvre, pour mémoire : une main de 1,0 donnerait
% 11,61 mm pour N, une main de 1,3 en donnerait 16,77. Il ne fallait pas
% deviner.
%  E   — 17 mm, VALEUR DONNÉE PAR L'IMPRIMEUR le 17/09/2026 pour 276 pages,
%        20 x 24, couverture RIGIDE, reliure collée (Klebebindung), papier
%        115 g. En reliure cousue (Fadenheftung), il annonce 18 mm.
%
% \DosImposeReliure<classe> : la reliure pour laquelle la valeur a été
% donnée. Le dos d'une couverture rigide n'est pas celui d'un dos carré
% collé souple, même à pagination égale.
\def\DosImposeN{14}   \def\DosImposePagesN{258}  \def\DosImposeFormatN{a4}
\def\DosImposeE{17}   \def\DosImposePagesE{276}  \def\DosImposeFormatE{20x24-marge}
\def\DosImposeA{0}    \def\DosImposePagesA{0}    \def\DosImposeFormatA{}
\def\DosImposeNE{0}   \def\DosImposePagesNE{0}   \def\DosImposeFormatNE{}
\def\DosImposeEA{0}   \def\DosImposePagesEA{0}   \def\DosImposeFormatEA{}
\def\DosImposeNEA{0}  \def\DosImposePagesNEA{0}  \def\DosImposeFormatNEA{}
\def\DosImposeReliureN{souple}  \def\DosImposeReliureE{rigide}
\def\DosImposeReliureA{souple}  \def\DosImposeReliureNE{souple}
\def\DosImposeReliureEA{souple} \def\DosImposeReliureNEA{souple}

% Fond perdu, en mm. 3 mm est la valeur courante ; certains imprimeurs
% demandent 5 mm.
\def\FondPerdu{3}

% Marge technique autour du fond perdu, en mm. Sert à loger les traits de
% coupe en mode épreuve. Ignorée en mode production.
\def\MargeTech{12}

% Zone de sécurité, en mm : aucun texte ne doit s'en approcher davantage
% du bord de coupe.
\def\Securite{8}

% ---------------------------------------------------------------------
% 2. DONNÉES PAR CLASSE
%    pages : nombre de pages du PDF intérieur, COUVERTURE NON COMPRISE.
%    Un dos carré collé impose un nombre de pages PAIR ; le fichier
%    refuse de compiler si le compte est impair.
%
%    Ces tables ne servent que si \PagesForcees et \VersionForcee sont
%    vides, c'est-à-dire à la main. Valeurs A4 mesurées sur les PDF de la
%    a.3 (CLAUDE.md § 1, 05/09/2026). En 20 x 24, la pagination change
%    entièrement : NEA passe de 816 à 1036 pages.
% ---------------------------------------------------------------------

\def\DonneesN{262}
\def\DonneesE{214}
\def\DonneesA{386}
\def\DonneesNE{0}      % édition jamais produite
\def\DonneesEA{0}      % édition jamais produite
\def\DonneesNEA{816}

% Version imprimée sur le dos et la 1re de couverture, PAR CLASSE.
% « Les versions sont suivies par classe, chacune évoluant à son rythme »
% (CHANGELOG.md).
\def\VersionN{a.3}
\def\VersionE{a.3}
\def\VersionA{a.3}
\def\VersionNE{a.3}
\def\VersionEA{a.3}
\def\VersionNEA{a.3}

% ---------------------------------------------------------------------
% 3. TEXTES
% ---------------------------------------------------------------------

\def\TitreOuvrage{50\,Ohm}
\def\SousTitre{Préparation à l'examen radioamateur}

% Sous-titre propre à chaque classe (repris de FR_TITLES, build_book.py).
\def\SousTitreN{Cours complet — classe N}
\def\SousTitreE{Cours complet — classe E}
\def\SousTitreA{Cours complet — classe A}
\def\SousTitreNE{Cours complet — classes N et E}
\def\SousTitreEA{Cours complet — classes E et A}
\def\SousTitreNEA{Cours complet — classes N, E et A}

% Lettres du filigrane, empilées de haut en bas (cf. build_book.py v0.20).
\def\LettresN{N}
\def\LettresE{E}
\def\LettresA{A}
\def\LettresNE{N,E}
\def\LettresEA{E,A}
\def\LettresNEA{N,E,A}

% Texte de 4e de couverture.
% NB : il parle des classes N, E et A quelle que soit l'édition. Juste pour
% N, E, A et NEA ; approximatif pour NE et EA ; faux pour SWL, d'où son refus.
\def\QuatriemeTexte{%
Ce guide d'étude prépare aux examens radioamateur allemands des classes
N, E et A. Il traduit en français les contenus du projet 50ohm.de, publiés
par le référat AJW du DARC e.\,V., et les complète d'encarts « En France »
qui signalent, sans rien retirer au texte allemand, ce que la
réglementation française prévoit sur le même sujet.\par\smallskip
Les questions officielles du catalogue allemand accompagnent chaque
notion, avec leurs réponses.%
}

\def\Auteur{Traduit avec l'aide d'une IA par Pierre F4JWI}
\def\Depot{github.com/DrPLG/50ohm-fr}

% =====================================================================
%  À PARTIR D'ICI, PLUS RIEN À MODIFIER
% =====================================================================

\usepackage{iftex}
\RequireLuaTeX
\usepackage{xfp}
\usepackage{tikz}
\usetikzlibrary{calc}
\usepackage{fontspec}
\usepackage[french]{babel}
\usepackage{geometry}

% ---------------------------------------------------------------------
% 4. RÉSOLUTION DES DONNÉES DE LA CLASSE
% ---------------------------------------------------------------------

\makeatletter
% Tome SWL : refus explicite, plutôt qu'un « classe inconnue » qui laisserait
% croire à un oubli. Le texte de 4e parle des classes N, E et A.
\def\@swlmot{SWL}
\ifx\Classe\@swlmot
  \PackageError{couverture}{Pas de couverture pour le tome SWL}%
    {Sa 4e de couverture n'est pas redigee : le texte commun parle des
     classes N, E et A. C'est une decision editoriale a prendre d'abord.}%
\fi

\def\@resoudre#1#2{%
  \expandafter\let\expandafter#1\csname #2\Classe\endcsname
  \ifx#1\relax
    \PackageError{couverture}{Classe \Classe\space inconnue : pas de #2\Classe}%
      {Valeurs admises : N, E, A, NE, EA, NEA.}%
  \fi}
\@resoudre\NbPages{Donnees}
\@resoudre\SousTitreClasse{SousTitre}
\@resoudre\Lettres{Lettres}
\@resoudre\DosImpose{DosImpose}
\@resoudre\DosImposePages{DosImposePages}
\@resoudre\DosImposeFormat{DosImposeFormat}
\@resoudre\DosImposeReliure{DosImposeReliure}
\@resoudre\Version{Version}

% Valeurs imposées de l'extérieur. Test de vacuité par \detokenize : il ne
% dépend pas de la façon dont la macro a été définie.
\if\relax\detokenize\expandafter{\PagesForcees}\relax\else
  \let\NbPages\PagesForcees
\fi
\if\relax\detokenize\expandafter{\VersionForcee}\relax\else
  \let\Version\VersionForcee
\fi

% Format du livre fini.
\def\@formatAquatre{a4}
\def\@formatVingt{20x24}
\def\@formatVingtMarge{20x24-marge}
\ifx\Format\@formatAquatre
  \def\PageLargeur{210}\def\PageHauteur{297}
\else\ifx\Format\@formatVingt
  \def\PageLargeur{200}\def\PageHauteur{240}
\else\ifx\Format\@formatVingtMarge
  \def\PageLargeur{200}\def\PageHauteur{240}
\else
  \PackageError{couverture}{Format \Format\space inconnu}%
    {Valeurs admises : a4, 20x24, 20x24-marge.}%
  \def\PageLargeur{210}\def\PageHauteur{297}
\fi\fi\fi
% Reliure. Hors de « souple » et « rigide » : erreur, pas de supposition.
\def\@reliuresouple{souple}
\def\@reliurerigide{rigide}
\ifx\Reliure\@reliurerigide
  \edef\PageLargeur{\fpeval{\PageLargeur + 5}}
  \edef\PageHauteur{\fpeval{\PageHauteur + 5}}
  \def\FondPerdu{15}
\else\ifx\Reliure\@reliuresouple\else
  \PackageError{couverture}{Reliure \Reliure\space inconnue}%
    {Valeurs admises : souple, rigide.}%
\fi\fi
\makeatother

% Garde-fou : une édition jamais produite n'a pas de couverture.
\ifnum\NbPages=0
  \PackageError{couverture}{L'edition \Classe\space n'a pas de pagination}%
    {Cette edition n'a jamais ete compilee. Renseigner \string\Donnees\Classe
     ou passer \string\PagesForcees.}%
\fi

% Avertissement : l'imprimeur demande un nombre de pages multiple de 4
% (17/09/2026). Pas une erreur : c'est son exigence, pas une loi du papier.
\ifnum\fpeval{\NbPages - 4*floor(\NbPages/4)}=0 \else
  \typeout{COUVERTURE-AVERTISSEMENT \NbPages\space pages : pas un multiple
    de 4, que l'imprimeur exige -- ajouter des pages a l'interieur}%
\fi

% Garde-fou : un dos carré collé impose un nombre de pages pair.
\ifodd\NbPages
  \PackageError{couverture}{Nombre de pages impair (\NbPages)}%
    {Un dos carre colle exige un nombre PAIR de pages : une feuille en
     porte deux. Ajouter une page blanche a l'interieur.}%
\fi

% ---------------------------------------------------------------------
% 5. GÉOMÉTRIE — c'est ici que se joue l'exactitude du dos
% ---------------------------------------------------------------------

% Épaisseur du dos : valeur imposée si elle est donnée ET qu'elle vaut pour
% cette pagination et ce format ; calcul sinon.
\edef\DosCalcule{\fpeval{round(\NbPages / 2 * \Grammage * \Bouffant / 1000, 2)}}
\def\DosOrigine{calculee}
\ifnum\fpeval{\DosImpose > 0}=1
  \ifnum\fpeval{\NbPages = \DosImposePages}=1
    \ifx\Format\DosImposeFormat
      \ifx\Reliure\DosImposeReliure
        \def\DosOrigine{imposee}
      \fi
    \fi
  \fi
  \def\MotCalculee{calculee}
  \ifx\DosOrigine\MotCalculee
    \typeout{COUVERTURE-AVERTISSEMENT dos impose de \DosImpose\space mm
      donne pour \DosImposePages\space pages en \DosImposeFormat\space
      (reliure \DosImposeReliure) -- inapplicable a \NbPages\space pages en
      \Format\space (reliure \Reliure), dos calcule}%
  \fi
\fi
% NB : pas de \@ dans ces noms. Cette section est hors \makeatletter, et
% « \def\@imposee{...} » y redéfinit \@ (suivi du texte « imposee ») : le
% test échouait toujours, en silence. Constaté le 16/09/2026 — le dos de N
% tombait juste par coïncidence, 14,0 mm calculés pour 14 imposés.
\def\MotRigide{rigide}
\def\MotImposee{imposee}
\ifx\DosOrigine\MotImposee
  \edef\Dos{\DosImpose}
\else
  \edef\Dos{\DosCalcule}
\fi

% Marge technique : uniquement en mode épreuve.
\def\ouimot{oui}
\ifx\Reperes\ouimot \edef\Tech{\MargeTech}\else \def\Tech{0}\fi

% Dimensions du papier.
%   largeur = marge + fond perdu + 4e + dos + 1re + fond perdu + marge
%   hauteur = marge + fond perdu + hauteur + fond perdu + marge
\edef\PapierL{\fpeval{2*\PageLargeur + \Dos + 2*\FondPerdu + 2*\Tech}}
\edef\PapierH{\fpeval{\PageHauteur + 2*\FondPerdu + 2*\Tech}}

\geometry{paperwidth=\PapierL mm, paperheight=\PapierH mm,
          margin=0pt, nohead, nofoot}

% ---------------------------------------------------------------------
%  TOUTES LES COTES SONT FIGÉES ICI, ET NULLE PART AILLEURS.
%
%  Raison, payée par trois compilations ratées : un \fpeval laissé dans
%  une coordonnée ou dans le texte d'un nœud TikZ n'est pas développé au
%  moment où TikZ analyse le chemin. On obtient « Giving up on this path.
%  Did you forget a semicolon? », qui accuse une ponctuation alors que la
%  faute est ailleurs. Le corps du document ci-dessous n'emploie donc que
%  des macros déjà calculées.
% ---------------------------------------------------------------------

% Repères horizontaux, mesurés depuis le bord GAUCHE du papier.
\edef\CoupeG{\fpeval{\Tech + \FondPerdu}}          % coupe gauche (4e de couv.)
\edef\DosG{\fpeval{\CoupeG + \PageLargeur}}         % dos, début
\edef\DosD{\fpeval{\DosG + \Dos}}                   % dos, fin
\edef\CoupeD{\fpeval{\DosD + \PageLargeur}}         % coupe droite (1re de couv.)
% Repères verticaux, depuis le bord BAS.
\edef\CoupeB{\fpeval{\Tech + \FondPerdu}}
\edef\CoupeH{\fpeval{\CoupeB + \PageHauteur}}
% Axes.
\edef\DosAxe{\fpeval{(\DosG + \DosD)/2}}
\edef\MilieuH{\fpeval{(\CoupeB + \CoupeH)/2}}
\edef\MilieuL{\fpeval{\PapierL/2}}

% Format fini, pour les cotes affichées.
\edef\FiniL{\fpeval{2*\PageLargeur + \Dos}}

% Bandeau de la 1re de couverture : même proportion que la page de titre.
\edef\BandeauL{\fpeval{round(0.34*\PageLargeur, 2)}}
\edef\BandeauG{\fpeval{\CoupeD - \BandeauL}}
\edef\BandeauAccent{\fpeval{\BandeauG - 0.7}}
\edef\DosAccentG{\fpeval{\DosG - 0.7}}
\edef\DosAccentD{\fpeval{\DosD + 0.7}}
% Couverture rigide : le bleu du dos déborde de 2 mm sur chaque plat, pour
% absorber la tolérance de l'imprimeur sur l'épaisseur du dos (±10 %).
\edef\DosEtenduG{\fpeval{\DosG - 2}}
\edef\DosEtenduD{\fpeval{\DosD + 2}}

% Filigrane de classe.
\edef\FiligraneX{\fpeval{(\BandeauG + \CoupeD)/2}}
\edef\FiligraneY{\fpeval{\CoupeH - 6}}

% Bloc-titre de la 1re de couverture.
\edef\TitreX{\fpeval{\DosD + 12}}
\edef\TitreY{\fpeval{\CoupeH - 72}}
\edef\TitreW{\fpeval{\BandeauG - \DosD - 16}}
\edef\VersionY{\fpeval{\CoupeB + 14}}

% 4e de couverture.
\edef\QuatreX{\fpeval{\CoupeG + \Securite + 6}}
\edef\QuatreY{\fpeval{\CoupeH - \Securite - 14}}
\edef\QuatreBasY{\fpeval{\CoupeB + \Securite + 22}}
\edef\QuatreW{\fpeval{\PageLargeur - 2*\Securite - 12}}

% Zone de sécurité.
\edef\SecG{\fpeval{\CoupeG + \Securite}}
\edef\SecD{\fpeval{\CoupeD - \Securite}}
\edef\SecB{\fpeval{\CoupeB + \Securite}}
\edef\SecH{\fpeval{\CoupeH - \Securite}}

% Zones de contrôle : ce qui doit rester de part et d'autre du pli du dos.
\edef\PliQuatre{\fpeval{\DosG - \Securite}}         % 4e : limite côté dos
\edef\PliPremiere{\fpeval{\DosD + \Securite}}       % 1re : limite côté dos

% Traits de coupe et cotes, dans la marge technique.
\edef\TechExt{\fpeval{\Tech - 8}}
\edef\TechInt{\fpeval{\Tech - 1}}
\edef\TechCote{\fpeval{\Tech - 9}}
\edef\HautInt{\fpeval{\PapierH - \Tech + 1}}
\edef\HautExt{\fpeval{\PapierH - \Tech + 8}}
\edef\HautCote{\fpeval{\PapierH - \Tech + 9}}
\edef\DroitInt{\fpeval{\PapierL - \Tech + 1}}
\edef\DroitExt{\fpeval{\PapierL - \Tech + 8}}

% Emplacement du code-barres ISBN : en bas à DROITE de la 4e de couverture,
% selon l'usage éditorial, et non à gauche où il chevaucherait les mentions
% de licence.
\edef\IsbnD{\fpeval{\DosG - \Securite}}
\edef\IsbnG{\fpeval{\IsbnD - 40}}
\edef\IsbnH{\fpeval{\SecB + 25}}
\edef\IsbnCX{\fpeval{\IsbnG + 20}}
\edef\IsbnCY{\fpeval{\SecB + 12.5}}

% Cartouches de cotes, dans la marge technique. Ancrés vers l'INTÉRIEUR du
% papier : ancrés vers l'extérieur, ils débordaient et étaient rognés.
\edef\CoteBasY{\fpeval{1}}
\edef\CoteHautY{\fpeval{\PapierH - 1}}

% ---------------------------------------------------------------------
% 6. COULEURS ET POLICES — reprises de la page de titre (build_book.py)
% ---------------------------------------------------------------------

\definecolor{TitleBand}{cmyk}{.9,.55,.1,.35}   % bleu profond
\definecolor{TitleAccent}{cmyk}{.05,.8,1,0}    % rouge DARC (accent fin)
\definecolor{Repere}{cmyk}{0,1,1,0}            % magenta : repères de montage

\setmainfont{Libertinus Serif}[
  ItalicFont     = Libertinus Serif Italic,
  BoldFont       = Libertinus Serif Bold,
  BoldItalicFont = Libertinus Serif Bold Italic]

\pagestyle{empty}

% ---------------------------------------------------------------------
% 7. FILIGRANE DE CLASSE
%    Le texte est construit ici, et non dans le nœud : un \\ produit à
%    l'intérieur d'un \foreach referme le groupe d'alignement de TikZ.
%    Une lettre seule est grande ; plusieurs s'empilent et se réduisent
%    (build_book.py v0.20 : « NEA » à 105 pt débordait encore).
% ---------------------------------------------------------------------

\makeatletter
\def\FiligraneTexte{}
\newif\ifpremierelettre \premierelettretrue
\count255=0
\@for\@tempa:=\Lettres\do{%
  \advance\count255 by 1\relax
  \ifpremierelettre \premierelettrefalse
  \else \g@addto@macro\FiligraneTexte{\\}\fi
  \expandafter\g@addto@macro\expandafter\FiligraneTexte\expandafter{\@tempa}}%
\xdef\NbLettres{\the\count255}
\makeatother

\ifnum\NbLettres=1 \def\CorpsF{150}\else\def\CorpsF{86}\fi
\edef\CorpsFI{\fpeval{round(\CorpsF*1.02, 0)}}

% Texte du cartouche de cotes, lui aussi construit hors du nœud TikZ.
%
% NE PAS y employer \textsuperscript : dans un nœud « align=center », il
% casse le chemin TikZ (« Giving up on this path. Did you forget a
% semicolon? »), et l'erreur est signalée sur le nœud SUIVANT — elle
% accuse donc une ponctuation là où la faute est ici. Isolé par essais :
% $^2$ et le caractère ² passent, \textsuperscript non.
\def\CotesTexte{%
  \Classe\ --- \NbPages\ pages --- papier \Grammage\ g/m$^2$,
  bouffant \Bouffant\ --- fond perdu \FondPerdu\ mm\\
  papier \PapierL\ mm sur \PapierH\ mm --- fini \FiniL\ mm sur
  \PageHauteur\ mm\\
  EPREUVE, ne pas envoyer a l'impression%
}

% ---------------------------------------------------------------------
% 8. CONTRÔLE DE SUPERPOSITION
%
%  \MesureNoeud{nœud}{préfixe}{retrait} mesure la boîte d'un nœud par ses
%  quatre coins (un nœud pivoté a ses ancres pivotées : le min et le max des
%  quatre suffisent) et la ramène au repère du papier, en mm. Le point (0,0)
%  est lu par le même « let » que les ancres : les deux passent donc par la
%  même transformation, et la comparaison ne dépend pas de l'endroit où
%  TeX a posé l'image sur la page.
%
%  Le retrait sert au seul filigrane, pour mesurer les lettres et non leur
%  cadre : c'est le bord des lettres que le détourage rognerait.
%
%  PIÈGE MESURÉ le 16/09/2026 : la marge interne TikZ (.3333em) n'est PAS
%  évaluée dans la police du nœud, mais dans celle qui entoure le dessin.
%  Instrumenté : dans le nœud, em = 150 pt ; boîte brute du « N » :
%  41,7 x 37,5 mm, soit la lettre plus 2 x 1,3 mm — 1,3 mm étant .3333em
%  à 11 pt. Un retrait calculé à 150 pt (17,6 mm) ramenait la lettre à
%  6,6 mm de large, et le contrôle aurait dit OK sur une mesure absurde.
% ---------------------------------------------------------------------

\newcount\NbAlertes
\newdimen\FiligraneSep
\newsavebox\CreditsBoite

% Crédits de 4e : lignes coupées à la main. Leur largeur RÉELLE est celle
% de la plus longue ligne, mesurée dans une boîte à part (voir § 8, ISBN).
\def\CreditsTexte{%
    Réalisé à partir des contenus de 50ohm.de (en allemand)\\
    50ohm.de-Autorenteam, coordonné par le référat AJW du DARC e.\,V.\\
    \Auteur\\[3pt]
    Licence CC BY 4.0 --- \texttt{\Depot}}

\def\MesureNoeud#1#2#3{%
  \path let \p1=(#1.south west), \p2=(#1.north east),
            \p3=(#1.north west), \p4=(#1.south east), \p5=(0,0) in
    \pgfextra{%
      \expandafter\xdef\csname #2xmin\endcsname{\fpeval{round((min(\x1,\x2,\x3,\x4) - \x5 + #3)/1mm, 2)}}%
      \expandafter\xdef\csname #2xmax\endcsname{\fpeval{round((max(\x1,\x2,\x3,\x4) - \x5 - #3)/1mm, 2)}}%
      \expandafter\xdef\csname #2ymin\endcsname{\fpeval{round((min(\y1,\y2,\y3,\y4) - \y5 + #3)/1mm, 2)}}%
      \expandafter\xdef\csname #2ymax\endcsname{\fpeval{round((max(\y1,\y2,\y3,\y4) - \y5 - #3)/1mm, 2)}}%
    };}

% \ControleDans{libellé}{préfixe}{xmin}{xmax}{ymin}{ymax}
\def\ControleDans#1#2#3#4#5#6{%
  \ifnum\fpeval{\csname #2xmin\endcsname >= #3 && \csname #2xmax\endcsname <= #4
             && \csname #2ymin\endcsname >= #5 && \csname #2ymax\endcsname <= #6}=1
    \typeout{COUVERTURE-CONTROLE OK #1 : x \csname #2xmin\endcsname-\csname #2xmax\endcsname,
      y \csname #2ymin\endcsname-\csname #2ymax\endcsname\space dans x #3-#4, y #5-#6}%
  \else
    \global\advance\NbAlertes by 1
    \typeout{COUVERTURE-CONTROLE ALERTE #1 : x \csname #2xmin\endcsname-\csname #2xmax\endcsname,
      y \csname #2ymin\endcsname-\csname #2ymax\endcsname\space HORS de x #3-#4, y #5-#6}%
  \fi}

% \ControleDisjoints{libellé}{préfixe}{xmin}{xmax}{ymin}{ymax}
% Deux rectangles se chevauchent si et seulement s'ils se recouvrent sur
% les DEUX axes à la fois.
\def\ControleDisjoints#1#2#3#4#5#6{%
  \ifnum\fpeval{\csname #2xmin\endcsname < #4 && #3 < \csname #2xmax\endcsname
             && \csname #2ymin\endcsname < #6 && #5 < \csname #2ymax\endcsname}=1
    \global\advance\NbAlertes by 1
    \typeout{COUVERTURE-CONTROLE ALERTE #1 : chevauchement}%
  \else
    \typeout{COUVERTURE-CONTROLE OK #1 : disjoints}%
  \fi}

\begin{document}
% Marge interne du filigrane, évaluée dans la police qui ENTOURE le dessin :
% c'est là que TikZ évalue le .3333em de « inner sep » (cf. § 8).
\FiligraneSep=.3333em
% Largeur réelle des crédits : la plus longue ligne, dans leur police.
\sbox\CreditsBoite{\small\begin{tabular}{@{}l@{}}\CreditsTexte\end{tabular}}%
\begin{tikzpicture}[remember picture, overlay,
                    x=1mm, y=1mm, shift={(current page.south west)}]

% ---- Fonds ----------------------------------------------------------
% La 4e de couverture reste sur le blanc du papier.
% Le bandeau bleu couvre le dos, puis reprend au tiers extérieur de la
% 1re de couverture — même partition 2/3-1/3 que la page de titre.
% Couverture rigide : pas de filets au dos (ils tomberaient à côté des
% plis, vu la tolérance), et un dos bleu élargi de 2 mm de chaque côté.
\ifx\Reliure\MotRigide
  \fill[TitleBand] (\DosEtenduG, 0) rectangle (\DosEtenduD, \PapierH);
\else
  \fill[TitleBand] (\DosG, 0) rectangle (\DosD, \PapierH);
\fi
\fill[TitleBand] (\BandeauG, 0) rectangle (\PapierL, \PapierH);
% Filets d'accent : à la jonction sur la 1re, de part et d'autre du dos.
\fill[TitleAccent] (\BandeauAccent, 0) rectangle (\BandeauG, \PapierH);
\ifx\Reliure\MotRigide\else
  \fill[TitleAccent] (\DosAccentG, 0) rectangle (\DosG, \PapierH);
  \fill[TitleAccent] (\DosD, 0) rectangle (\DosAccentD, \PapierH);
\fi

% ---- 1re de couverture ----------------------------------------------
% Filigrane de classe, en réserve claire, détouré au bandeau.
\begin{scope}
\clip (\BandeauG, 0) rectangle (\PapierL, \PapierH);
\node[name=filigrane, anchor=north, align=center, text=white, opacity=0.16,
      font=\fontsize{\CorpsF}{\CorpsFI}\selectfont\bfseries]
  at (\FiligraneX, \FiligraneY) {\FiligraneTexte};
\end{scope}

\node[name=titre, anchor=west, align=left, text width=\TitreW mm]
  at (\TitreX, \TitreY) {%
    {\fontsize{54}{56}\selectfont\bfseries \TitreOuvrage}\\[10pt]
    {\Large\color{TitleBand}\bfseries \SousTitre}\\[16pt]
    {\large\color{black!75} \SousTitreClasse}};

\node[name=version, anchor=south west, align=left, font=\small, text=black!70]
  at (\TitreX, \VersionY) {Version \Version\\ Licence CC BY 4.0};

% ---- Dos -------------------------------------------------------------
% Un dos étroit ne reçoit pas de texte : en deçà de 6 mm, le débord au
% massicot ferait mordre le texte vertical sur les plats.
\ifnum\fpeval{\Dos >= 6}=1
  \node[name=dos, rotate=90, anchor=center, text=white, font=\bfseries,
        inner sep=0pt]
    at (\DosAxe, \MilieuH)
    {\TitreOuvrage\ \ --- \ \SousTitreClasse\ \ --- \ \Version};
\fi

% ---- 4e de couverture ------------------------------------------------
\node[name=quatre, anchor=north west, align=justify, text width=\QuatreW mm]
  at (\QuatreX, \QuatreY) {\QuatriemeTexte};

% Le texte vit dans \CreditsTexte (§ 8) pour être mesuré à part : la boîte
% du nœud a la largeur fixe \QuatreW et couvre l'emplacement ISBN sans
% qu'aucun caractère n'y soit. Retirer « text width » pour l'éviter change
% l'interlignage — essayé le 16/09/2026, 6 934 pixels différents.
\node[name=credits, anchor=south west, align=left, font=\small, text=black!70,
      text width=\QuatreW mm]
  at (\QuatreX, \QuatreBasY) {\CreditsTexte};

% ---- Repères de fabrication (mode épreuve uniquement) ----------------
\ifx\Reperes\ouimot
  % Traits de coupe, hors du fond perdu.
  \foreach \x in {\CoupeG, \CoupeD} {
    \draw[Repere, line width=0.25pt] (\x, \TechExt) -- (\x, \TechInt);
    \draw[Repere, line width=0.25pt] (\x, \HautInt) -- (\x, \HautExt);}
  \foreach \y in {\CoupeB, \CoupeH} {
    \draw[Repere, line width=0.25pt] (\TechExt, \y) -- (\TechInt, \y);
    \draw[Repere, line width=0.25pt] (\DroitInt, \y) -- (\DroitExt, \y);}
  % Limites du dos, en pointillés.
  \foreach \x in {\DosG, \DosD} {
    \draw[Repere, line width=0.25pt, dashed] (\x, \TechExt) -- (\x, \TechInt);
    \draw[Repere, line width=0.25pt, dashed] (\x, \HautInt) -- (\x, \HautExt);}
  % Zone de sécurité.
  \draw[Repere, line width=0.2pt, dotted]
    (\SecG, \SecB) rectangle (\SecD, \SecH);
  % Emplacement réservé au code-barres ISBN.
  \draw[Repere, line width=0.2pt] (\IsbnG, \SecB) rectangle (\IsbnD, \IsbnH);
  \node[font=\tiny\sffamily, text=Repere, align=center]
    at (\IsbnCX, \IsbnCY) {emplacement\\ISBN};
  % Cotes.
  \node[anchor=north, font=\tiny\sffamily, text=Repere]
    at (\DosAxe, \CoteHautY) {dos \Dos\ mm (\DosOrigine)};
  \node[anchor=south, font=\tiny\sffamily, text=Repere, align=center]
    at (\MilieuL, \CoteBasY) {\CotesTexte};
\fi

% ---- Contrôle de superposition (§ 8) ----------------------------------
% Rien n'est dessiné : des mesures et des lignes de journal.
\MesureNoeud{filigrane}{Fil}{\the\FiligraneSep}
\MesureNoeud{titre}{Tit}{0pt}
\MesureNoeud{version}{Ver}{0pt}
\MesureNoeud{quatre}{Qua}{0pt}
\MesureNoeud{credits}{Cre}{0pt}
% Bord droit des crédits = bord gauche + marge interne + plus longue ligne.
\xdef\Crexmax{\fpeval{round(\Crexmin + (\the\FiligraneSep + \the\wd\CreditsBoite)/1mm, 2)}}

% Le filigrane peut entrer dans la zone de sécurité (il est décoratif),
% pas franchir la coupe ni sortir du bandeau — le détourage le rognerait.
\ControleDans{filigrane dans le bandeau}{Fil}{\BandeauG}{\CoupeD}{\CoupeB}{\CoupeH}
\ControleDans{titre sur la 1re, hors bandeau}{Tit}{\PliPremiere}{\BandeauAccent}{\SecB}{\SecH}
\ControleDans{version sur la 1re, hors bandeau}{Ver}{\PliPremiere}{\BandeauAccent}{\SecB}{\SecH}
\ControleDans{texte de 4e dans la zone de securite}{Qua}{\SecG}{\PliQuatre}{\SecB}{\SecH}
\ControleDans{credits de 4e dans la zone de securite}{Cre}{\SecG}{\PliQuatre}{\SecB}{\SecH}
\ControleDisjoints{titre / version}{Tit}{\Verxmin}{\Verxmax}{\Verymin}{\Verymax}
\ControleDisjoints{texte de 4e / credits}{Qua}{\Crexmin}{\Crexmax}{\Creymin}{\Creymax}
\ControleDisjoints{texte de 4e / emplacement ISBN}{Qua}{\IsbnG}{\IsbnD}{\SecB}{\IsbnH}
\ControleDisjoints{credits / emplacement ISBN}{Cre}{\IsbnG}{\IsbnD}{\SecB}{\IsbnH}
\ifnum\fpeval{\Dos >= 6}=1
  \MesureNoeud{dos}{Dos}{0pt}
  \ControleDans{texte du dos entre les plis}{Dos}{\DosG}{\DosD}{\SecB}{\SecH}
\else
  \typeout{COUVERTURE-CONTROLE OK dos de \Dos\space mm : trop etroit, sans texte}%
\fi

\end{tikzpicture}
\typeout{COUVERTURE-BILAN alertes=\the\NbAlertes\space; classe=\Classe\space;
  format=\Format\space; reliure=\Reliure\space; pages=\NbPages\space;
  version=\Version\space;
  dos=\Dos\space mm (\DosOrigine); fini=\FiniL x \PageHauteur\space mm;
  papier=\PapierL x \PapierH\space mm}
\end{document}
"
%
%  ---------------------------------------------------------------------
%  DEUX MODES, un seul paramètre : \Reperes
%
%    oui  — mode ÉPREUVE. Ajoute une marge technique, les traits de coupe,
%           les limites du dos, la zone de sécurité et les cotes. C'est ce
%           qu'il faut pour vérifier à l'écran.
%    non  — mode PRODUCTION. Papier = format final + fond perdu, rien
%           d'autre. C'est ce fichier-là qu'on envoie à l'imprimeur.
%
%  ---------------------------------------------------------------------
%  CONTRÔLE DE SUPERPOSITION
%
%  À la fin du dessin, chaque bloc de texte est mesuré par ses ancres TikZ
%  (quatre coins, ce qui vaut aussi pour le texte pivoté du dos) et comparé
%  à la zone où il doit tenir, puis aux blocs voisins. Le journal porte une
%  ligne par vérification :
%
%      COUVERTURE-CONTROLE OK      <quoi> : ...
%      COUVERTURE-CONTROLE ALERTE  <quoi> : ...
%
%  et un bilan : COUVERTURE-BILAN alertes=<n> ; ...
%
%  Ce que le contrôle VOIT : un bloc qui franchit le pli du dos, la coupe,
%  la zone de sécurité ou le bord du bandeau ; deux blocs qui se chevauchent ;
%  un bloc qui empiète sur l'emplacement réservé au code-barres ISBN.
%  Ce qu'il NE VOIT PAS : une ligne trop longue à l'INTÉRIEUR d'un bloc à
%  largeur fixe (la boîte garde sa largeur, le texte déborde). Celle-là se
%  signale par un « Overfull \hbox » dans le même journal, que la fenêtre
%  compte aussi.
%
%  ---------------------------------------------------------------------
%  AVERTISSEMENT SUR L'ÉPAISSEUR DU DOS
%
%  Le dos se calcule ainsi :
%
%      épaisseur d'une feuille (mm) = grammage x bouffant / 1000
%      dos (mm)                     = (nombre de pages / 2) x épaisseur
%
%  On divise par deux parce qu'une feuille porte DEUX pages.
%
%  Le « bouffant » (main, volume, bulk) dépend du papier et n'est pas
%  devinable : un offset 80 g courant est autour de 1,30, un bouffant
%  d'édition monte à 1,75 et plus. UN ÉCART DE BOUFFANT DÉPLACE LE DOS DE
%  PLUSIEURS MILLIMÈTRES. Demander à l'imprimeur soit son bouffant, soit
%  directement l'épaisseur du dos, et reporter ci-dessous. Tant que cette
%  valeur n'est pas confirmée, la couverture n'est pas bonne à tirer.
%
%  Si l'imprimeur donne le dos en clair, forcer la valeur avec \DosImpose.
% =====================================================================

\documentclass[11pt]{article}

% ---------------------------------------------------------------------
% 1. PARAMÈTRES — c'est la seule zone à modifier
% ---------------------------------------------------------------------

% Classe de l'ouvrage : N, E, A, NE, EA ou NEA.
% (\providecommand : une valeur passée en ligne de commande l'emporte.)
\providecommand{\Classe}{NEA}

% Mode épreuve (oui) ou production (non).
\providecommand{\Reperes}{oui}

% Format du livre FINI : a4, 20x24 ou 20x24-marge — les noms de --format
% dans build_book.py. Les deux maquettes 20 x 24 ne diffèrent que par leur
% colonne de marge intérieure : le livre fini, donc la couverture, est le
% même (200 x 240 mm).
\providecommand{\Format}{a4}

% Reliure : « souple » (dos carré collé, fond perdu 3 mm) ou « rigide »
% (couverture rigide, dite hardcover). En rigide, consignes de l'imprimeur
% du 17/09/2026 : chaque plat déborde du livre de 5 mm en largeur et en
% hauteur (le carton dépasse du bloc), et le fichier porte 15 mm de rabat
% autour du carton au lieu de 3 mm de fond perdu. Les aplats vont déjà au
% bord du papier : ils remplissent le rabat d'eux-mêmes.
\providecommand{\Reliure}{souple}

% Pagination et version imposées de l'extérieur. Vides = tables du § 2.
% La fenêtre les renseigne toujours : la pagination vient du journal du
% livre réellement compilé, dans le format demandé.
\providecommand{\PagesForcees}{}
\providecommand{\VersionForcee}{}

% Papier de l'INTÉRIEUR du livre (pas de la couverture).
% CoolLibri, papier « Standard 90 g blanc ».
\def\Grammage{90}        % g/m²
\def\Bouffant{1.206}     % main DÉDUITE, voir ci-dessous

% Épaisseur du dos IMPOSÉE par l'imprimeur, en mm, par classe.
% Une valeur > 0 prime sur le calcul — MAIS SEULEMENT pour la pagination et
% le format pour lesquels l'imprimeur l'a donnée (\DosImposePages<classe>,
% \DosImposeFormat<classe>). Hors de ce cas, le dos est calculé et le
% journal le dit : un dos de 14 mm donné pour 258 pages A4 ne vaut rien pour
% 262 pages, ni pour le même livre recomposé en 20 x 24.
%
% ATTENTION : cette valeur est propre à CHAQUE ÉDITION. Un dos unique
% appliqué aux quatre livres serait faux pour trois d'entre eux.
%
%  N   — 14 mm, VALEUR DONNÉE PAR COOLIBRI pour 258 pages en 90 g, A4.
%        Fichier attendu : 440 x 303 mm, ce que ce gabarit produit en mode
%        production (2 x 210 + 14 + 2 x 3 = 440 ; 297 + 2 x 3 = 303).
%        Le N de la a.3 compte 262 pages : cette valeur ne s'y applique plus.
%  E, A, NEA — non demandés à l'imprimeur : ils sont CALCULÉS avec la main
%        déduite ci-dessous. À faire confirmer avant tout tirage.
%
% La main de 1,206 n'est pas publiée par CoolLibri : elle est déduite de
% leur propre chiffre, 14 mm / (0,090 x 258/2) = 1,206. Elle reproduit donc
% exactement le dos de N, et n'est qu'une ESTIMATION pour les autres.
% Écarts que cela recouvre, pour mémoire : une main de 1,0 donnerait
% 11,61 mm pour N, une main de 1,3 en donnerait 16,77. Il ne fallait pas
% deviner.
%  E   — 17 mm, VALEUR DONNÉE PAR L'IMPRIMEUR le 17/09/2026 pour 276 pages,
%        20 x 24, couverture RIGIDE, reliure collée (Klebebindung), papier
%        115 g. En reliure cousue (Fadenheftung), il annonce 18 mm.
%
% \DosImposeReliure<classe> : la reliure pour laquelle la valeur a été
% donnée. Le dos d'une couverture rigide n'est pas celui d'un dos carré
% collé souple, même à pagination égale.
\def\DosImposeN{14}   \def\DosImposePagesN{258}  \def\DosImposeFormatN{a4}
\def\DosImposeE{17}   \def\DosImposePagesE{276}  \def\DosImposeFormatE{20x24-marge}
\def\DosImposeA{0}    \def\DosImposePagesA{0}    \def\DosImposeFormatA{}
\def\DosImposeNE{0}   \def\DosImposePagesNE{0}   \def\DosImposeFormatNE{}
\def\DosImposeEA{0}   \def\DosImposePagesEA{0}   \def\DosImposeFormatEA{}
\def\DosImposeNEA{0}  \def\DosImposePagesNEA{0}  \def\DosImposeFormatNEA{}
\def\DosImposeReliureN{souple}  \def\DosImposeReliureE{rigide}
\def\DosImposeReliureA{souple}  \def\DosImposeReliureNE{souple}
\def\DosImposeReliureEA{souple} \def\DosImposeReliureNEA{souple}

% Fond perdu, en mm. 3 mm est la valeur courante ; certains imprimeurs
% demandent 5 mm.
\def\FondPerdu{3}

% Marge technique autour du fond perdu, en mm. Sert à loger les traits de
% coupe en mode épreuve. Ignorée en mode production.
\def\MargeTech{12}

% Zone de sécurité, en mm : aucun texte ne doit s'en approcher davantage
% du bord de coupe.
\def\Securite{8}

% ---------------------------------------------------------------------
% 2. DONNÉES PAR CLASSE
%    pages : nombre de pages du PDF intérieur, COUVERTURE NON COMPRISE.
%    Un dos carré collé impose un nombre de pages PAIR ; le fichier
%    refuse de compiler si le compte est impair.
%
%    Ces tables ne servent que si \PagesForcees et \VersionForcee sont
%    vides, c'est-à-dire à la main. Valeurs A4 mesurées sur les PDF de la
%    a.3 (CLAUDE.md § 1, 05/09/2026). En 20 x 24, la pagination change
%    entièrement : NEA passe de 816 à 1036 pages.
% ---------------------------------------------------------------------

\def\DonneesN{262}
\def\DonneesE{214}
\def\DonneesA{386}
\def\DonneesNE{0}      % édition jamais produite
\def\DonneesEA{0}      % édition jamais produite
\def\DonneesNEA{816}

% Version imprimée sur le dos et la 1re de couverture, PAR CLASSE.
% « Les versions sont suivies par classe, chacune évoluant à son rythme »
% (CHANGELOG.md).
\def\VersionN{a.3}
\def\VersionE{a.3}
\def\VersionA{a.3}
\def\VersionNE{a.3}
\def\VersionEA{a.3}
\def\VersionNEA{a.3}

% ---------------------------------------------------------------------
% 3. TEXTES
% ---------------------------------------------------------------------

\def\TitreOuvrage{50\,Ohm}
\def\SousTitre{Préparation à l'examen radioamateur}

% Sous-titre propre à chaque classe (repris de FR_TITLES, build_book.py).
\def\SousTitreN{Cours complet — classe N}
\def\SousTitreE{Cours complet — classe E}
\def\SousTitreA{Cours complet — classe A}
\def\SousTitreNE{Cours complet — classes N et E}
\def\SousTitreEA{Cours complet — classes E et A}
\def\SousTitreNEA{Cours complet — classes N, E et A}

% Lettres du filigrane, empilées de haut en bas (cf. build_book.py v0.20).
\def\LettresN{N}
\def\LettresE{E}
\def\LettresA{A}
\def\LettresNE{N,E}
\def\LettresEA{E,A}
\def\LettresNEA{N,E,A}

% Texte de 4e de couverture.
% NB : il parle des classes N, E et A quelle que soit l'édition. Juste pour
% N, E, A et NEA ; approximatif pour NE et EA ; faux pour SWL, d'où son refus.
\def\QuatriemeTexte{%
Ce guide d'étude prépare aux examens radioamateur allemands des classes
N, E et A. Il traduit en français les contenus du projet 50ohm.de, publiés
par le référat AJW du DARC e.\,V., et les complète d'encarts « En France »
qui signalent, sans rien retirer au texte allemand, ce que la
réglementation française prévoit sur le même sujet.\par\smallskip
Les questions officielles du catalogue allemand accompagnent chaque
notion, avec leurs réponses.%
}

\def\Auteur{Traduit avec l'aide d'une IA par Pierre F4JWI}
\def\Depot{github.com/DrPLG/50ohm-fr}

% =====================================================================
%  À PARTIR D'ICI, PLUS RIEN À MODIFIER
% =====================================================================

\usepackage{iftex}
\RequireLuaTeX
\usepackage{xfp}
\usepackage{tikz}
\usetikzlibrary{calc}
\usepackage{fontspec}
\usepackage[french]{babel}
\usepackage{geometry}

% ---------------------------------------------------------------------
% 4. RÉSOLUTION DES DONNÉES DE LA CLASSE
% ---------------------------------------------------------------------

\makeatletter
% Tome SWL : refus explicite, plutôt qu'un « classe inconnue » qui laisserait
% croire à un oubli. Le texte de 4e parle des classes N, E et A.
\def\@swlmot{SWL}
\ifx\Classe\@swlmot
  \PackageError{couverture}{Pas de couverture pour le tome SWL}%
    {Sa 4e de couverture n'est pas redigee : le texte commun parle des
     classes N, E et A. C'est une decision editoriale a prendre d'abord.}%
\fi

\def\@resoudre#1#2{%
  \expandafter\let\expandafter#1\csname #2\Classe\endcsname
  \ifx#1\relax
    \PackageError{couverture}{Classe \Classe\space inconnue : pas de #2\Classe}%
      {Valeurs admises : N, E, A, NE, EA, NEA.}%
  \fi}
\@resoudre\NbPages{Donnees}
\@resoudre\SousTitreClasse{SousTitre}
\@resoudre\Lettres{Lettres}
\@resoudre\DosImpose{DosImpose}
\@resoudre\DosImposePages{DosImposePages}
\@resoudre\DosImposeFormat{DosImposeFormat}
\@resoudre\DosImposeReliure{DosImposeReliure}
\@resoudre\Version{Version}

% Valeurs imposées de l'extérieur. Test de vacuité par \detokenize : il ne
% dépend pas de la façon dont la macro a été définie.
\if\relax\detokenize\expandafter{\PagesForcees}\relax\else
  \let\NbPages\PagesForcees
\fi
\if\relax\detokenize\expandafter{\VersionForcee}\relax\else
  \let\Version\VersionForcee
\fi

% Format du livre fini.
\def\@formatAquatre{a4}
\def\@formatVingt{20x24}
\def\@formatVingtMarge{20x24-marge}
\ifx\Format\@formatAquatre
  \def\PageLargeur{210}\def\PageHauteur{297}
\else\ifx\Format\@formatVingt
  \def\PageLargeur{200}\def\PageHauteur{240}
\else\ifx\Format\@formatVingtMarge
  \def\PageLargeur{200}\def\PageHauteur{240}
\else
  \PackageError{couverture}{Format \Format\space inconnu}%
    {Valeurs admises : a4, 20x24, 20x24-marge.}%
  \def\PageLargeur{210}\def\PageHauteur{297}
\fi\fi\fi
% Reliure. Hors de « souple » et « rigide » : erreur, pas de supposition.
\def\@reliuresouple{souple}
\def\@reliurerigide{rigide}
\ifx\Reliure\@reliurerigide
  \edef\PageLargeur{\fpeval{\PageLargeur + 5}}
  \edef\PageHauteur{\fpeval{\PageHauteur + 5}}
  \def\FondPerdu{15}
\else\ifx\Reliure\@reliuresouple\else
  \PackageError{couverture}{Reliure \Reliure\space inconnue}%
    {Valeurs admises : souple, rigide.}%
\fi\fi
\makeatother

% Garde-fou : une édition jamais produite n'a pas de couverture.
\ifnum\NbPages=0
  \PackageError{couverture}{L'edition \Classe\space n'a pas de pagination}%
    {Cette edition n'a jamais ete compilee. Renseigner \string\Donnees\Classe
     ou passer \string\PagesForcees.}%
\fi

% Avertissement : l'imprimeur demande un nombre de pages multiple de 4
% (17/09/2026). Pas une erreur : c'est son exigence, pas une loi du papier.
\ifnum\fpeval{\NbPages - 4*floor(\NbPages/4)}=0 \else
  \typeout{COUVERTURE-AVERTISSEMENT \NbPages\space pages : pas un multiple
    de 4, que l'imprimeur exige -- ajouter des pages a l'interieur}%
\fi

% Garde-fou : un dos carré collé impose un nombre de pages pair.
\ifodd\NbPages
  \PackageError{couverture}{Nombre de pages impair (\NbPages)}%
    {Un dos carre colle exige un nombre PAIR de pages : une feuille en
     porte deux. Ajouter une page blanche a l'interieur.}%
\fi

% ---------------------------------------------------------------------
% 5. GÉOMÉTRIE — c'est ici que se joue l'exactitude du dos
% ---------------------------------------------------------------------

% Épaisseur du dos : valeur imposée si elle est donnée ET qu'elle vaut pour
% cette pagination et ce format ; calcul sinon.
\edef\DosCalcule{\fpeval{round(\NbPages / 2 * \Grammage * \Bouffant / 1000, 2)}}
\def\DosOrigine{calculee}
\ifnum\fpeval{\DosImpose > 0}=1
  \ifnum\fpeval{\NbPages = \DosImposePages}=1
    \ifx\Format\DosImposeFormat
      \ifx\Reliure\DosImposeReliure
        \def\DosOrigine{imposee}
      \fi
    \fi
  \fi
  \def\MotCalculee{calculee}
  \ifx\DosOrigine\MotCalculee
    \typeout{COUVERTURE-AVERTISSEMENT dos impose de \DosImpose\space mm
      donne pour \DosImposePages\space pages en \DosImposeFormat\space
      (reliure \DosImposeReliure) -- inapplicable a \NbPages\space pages en
      \Format\space (reliure \Reliure), dos calcule}%
  \fi
\fi
% NB : pas de \@ dans ces noms. Cette section est hors \makeatletter, et
% « \def\@imposee{...} » y redéfinit \@ (suivi du texte « imposee ») : le
% test échouait toujours, en silence. Constaté le 16/09/2026 — le dos de N
% tombait juste par coïncidence, 14,0 mm calculés pour 14 imposés.
\def\MotRigide{rigide}
\def\MotImposee{imposee}
\ifx\DosOrigine\MotImposee
  \edef\Dos{\DosImpose}
\else
  \edef\Dos{\DosCalcule}
\fi

% Marge technique : uniquement en mode épreuve.
\def\ouimot{oui}
\ifx\Reperes\ouimot \edef\Tech{\MargeTech}\else \def\Tech{0}\fi

% Dimensions du papier.
%   largeur = marge + fond perdu + 4e + dos + 1re + fond perdu + marge
%   hauteur = marge + fond perdu + hauteur + fond perdu + marge
\edef\PapierL{\fpeval{2*\PageLargeur + \Dos + 2*\FondPerdu + 2*\Tech}}
\edef\PapierH{\fpeval{\PageHauteur + 2*\FondPerdu + 2*\Tech}}

\geometry{paperwidth=\PapierL mm, paperheight=\PapierH mm,
          margin=0pt, nohead, nofoot}

% ---------------------------------------------------------------------
%  TOUTES LES COTES SONT FIGÉES ICI, ET NULLE PART AILLEURS.
%
%  Raison, payée par trois compilations ratées : un \fpeval laissé dans
%  une coordonnée ou dans le texte d'un nœud TikZ n'est pas développé au
%  moment où TikZ analyse le chemin. On obtient « Giving up on this path.
%  Did you forget a semicolon? », qui accuse une ponctuation alors que la
%  faute est ailleurs. Le corps du document ci-dessous n'emploie donc que
%  des macros déjà calculées.
% ---------------------------------------------------------------------

% Repères horizontaux, mesurés depuis le bord GAUCHE du papier.
\edef\CoupeG{\fpeval{\Tech + \FondPerdu}}          % coupe gauche (4e de couv.)
\edef\DosG{\fpeval{\CoupeG + \PageLargeur}}         % dos, début
\edef\DosD{\fpeval{\DosG + \Dos}}                   % dos, fin
\edef\CoupeD{\fpeval{\DosD + \PageLargeur}}         % coupe droite (1re de couv.)
% Repères verticaux, depuis le bord BAS.
\edef\CoupeB{\fpeval{\Tech + \FondPerdu}}
\edef\CoupeH{\fpeval{\CoupeB + \PageHauteur}}
% Axes.
\edef\DosAxe{\fpeval{(\DosG + \DosD)/2}}
\edef\MilieuH{\fpeval{(\CoupeB + \CoupeH)/2}}
\edef\MilieuL{\fpeval{\PapierL/2}}

% Format fini, pour les cotes affichées.
\edef\FiniL{\fpeval{2*\PageLargeur + \Dos}}

% Bandeau de la 1re de couverture : même proportion que la page de titre.
\edef\BandeauL{\fpeval{round(0.34*\PageLargeur, 2)}}
\edef\BandeauG{\fpeval{\CoupeD - \BandeauL}}
\edef\BandeauAccent{\fpeval{\BandeauG - 0.7}}
\edef\DosAccentG{\fpeval{\DosG - 0.7}}
\edef\DosAccentD{\fpeval{\DosD + 0.7}}
% Couverture rigide : le bleu du dos déborde de 2 mm sur chaque plat, pour
% absorber la tolérance de l'imprimeur sur l'épaisseur du dos (±10 %).
\edef\DosEtenduG{\fpeval{\DosG - 2}}
\edef\DosEtenduD{\fpeval{\DosD + 2}}

% Filigrane de classe.
\edef\FiligraneX{\fpeval{(\BandeauG + \CoupeD)/2}}
\edef\FiligraneY{\fpeval{\CoupeH - 6}}

% Bloc-titre de la 1re de couverture.
\edef\TitreX{\fpeval{\DosD + 12}}
\edef\TitreY{\fpeval{\CoupeH - 72}}
\edef\TitreW{\fpeval{\BandeauG - \DosD - 16}}
\edef\VersionY{\fpeval{\CoupeB + 14}}

% 4e de couverture.
\edef\QuatreX{\fpeval{\CoupeG + \Securite + 6}}
\edef\QuatreY{\fpeval{\CoupeH - \Securite - 14}}
\edef\QuatreBasY{\fpeval{\CoupeB + \Securite + 22}}
\edef\QuatreW{\fpeval{\PageLargeur - 2*\Securite - 12}}

% Zone de sécurité.
\edef\SecG{\fpeval{\CoupeG + \Securite}}
\edef\SecD{\fpeval{\CoupeD - \Securite}}
\edef\SecB{\fpeval{\CoupeB + \Securite}}
\edef\SecH{\fpeval{\CoupeH - \Securite}}

% Zones de contrôle : ce qui doit rester de part et d'autre du pli du dos.
\edef\PliQuatre{\fpeval{\DosG - \Securite}}         % 4e : limite côté dos
\edef\PliPremiere{\fpeval{\DosD + \Securite}}       % 1re : limite côté dos

% Traits de coupe et cotes, dans la marge technique.
\edef\TechExt{\fpeval{\Tech - 8}}
\edef\TechInt{\fpeval{\Tech - 1}}
\edef\TechCote{\fpeval{\Tech - 9}}
\edef\HautInt{\fpeval{\PapierH - \Tech + 1}}
\edef\HautExt{\fpeval{\PapierH - \Tech + 8}}
\edef\HautCote{\fpeval{\PapierH - \Tech + 9}}
\edef\DroitInt{\fpeval{\PapierL - \Tech + 1}}
\edef\DroitExt{\fpeval{\PapierL - \Tech + 8}}

% Emplacement du code-barres ISBN : en bas à DROITE de la 4e de couverture,
% selon l'usage éditorial, et non à gauche où il chevaucherait les mentions
% de licence.
\edef\IsbnD{\fpeval{\DosG - \Securite}}
\edef\IsbnG{\fpeval{\IsbnD - 40}}
\edef\IsbnH{\fpeval{\SecB + 25}}
\edef\IsbnCX{\fpeval{\IsbnG + 20}}
\edef\IsbnCY{\fpeval{\SecB + 12.5}}

% Cartouches de cotes, dans la marge technique. Ancrés vers l'INTÉRIEUR du
% papier : ancrés vers l'extérieur, ils débordaient et étaient rognés.
\edef\CoteBasY{\fpeval{1}}
\edef\CoteHautY{\fpeval{\PapierH - 1}}

% ---------------------------------------------------------------------
% 6. COULEURS ET POLICES — reprises de la page de titre (build_book.py)
% ---------------------------------------------------------------------

\definecolor{TitleBand}{cmyk}{.9,.55,.1,.35}   % bleu profond
\definecolor{TitleAccent}{cmyk}{.05,.8,1,0}    % rouge DARC (accent fin)
\definecolor{Repere}{cmyk}{0,1,1,0}            % magenta : repères de montage

\setmainfont{Libertinus Serif}[
  ItalicFont     = Libertinus Serif Italic,
  BoldFont       = Libertinus Serif Bold,
  BoldItalicFont = Libertinus Serif Bold Italic]

\pagestyle{empty}

% ---------------------------------------------------------------------
% 7. FILIGRANE DE CLASSE
%    Le texte est construit ici, et non dans le nœud : un \\ produit à
%    l'intérieur d'un \foreach referme le groupe d'alignement de TikZ.
%    Une lettre seule est grande ; plusieurs s'empilent et se réduisent
%    (build_book.py v0.20 : « NEA » à 105 pt débordait encore).
% ---------------------------------------------------------------------

\makeatletter
\def\FiligraneTexte{}
\newif\ifpremierelettre \premierelettretrue
\count255=0
\@for\@tempa:=\Lettres\do{%
  \advance\count255 by 1\relax
  \ifpremierelettre \premierelettrefalse
  \else \g@addto@macro\FiligraneTexte{\\}\fi
  \expandafter\g@addto@macro\expandafter\FiligraneTexte\expandafter{\@tempa}}%
\xdef\NbLettres{\the\count255}
\makeatother

\ifnum\NbLettres=1 \def\CorpsF{150}\else\def\CorpsF{86}\fi
\edef\CorpsFI{\fpeval{round(\CorpsF*1.02, 0)}}

% Texte du cartouche de cotes, lui aussi construit hors du nœud TikZ.
%
% NE PAS y employer \textsuperscript : dans un nœud « align=center », il
% casse le chemin TikZ (« Giving up on this path. Did you forget a
% semicolon? »), et l'erreur est signalée sur le nœud SUIVANT — elle
% accuse donc une ponctuation là où la faute est ici. Isolé par essais :
% $^2$ et le caractère ² passent, \textsuperscript non.
\def\CotesTexte{%
  \Classe\ --- \NbPages\ pages --- papier \Grammage\ g/m$^2$,
  bouffant \Bouffant\ --- fond perdu \FondPerdu\ mm\\
  papier \PapierL\ mm sur \PapierH\ mm --- fini \FiniL\ mm sur
  \PageHauteur\ mm\\
  EPREUVE, ne pas envoyer a l'impression%
}

% ---------------------------------------------------------------------
% 8. CONTRÔLE DE SUPERPOSITION
%
%  \MesureNoeud{nœud}{préfixe}{retrait} mesure la boîte d'un nœud par ses
%  quatre coins (un nœud pivoté a ses ancres pivotées : le min et le max des
%  quatre suffisent) et la ramène au repère du papier, en mm. Le point (0,0)
%  est lu par le même « let » que les ancres : les deux passent donc par la
%  même transformation, et la comparaison ne dépend pas de l'endroit où
%  TeX a posé l'image sur la page.
%
%  Le retrait sert au seul filigrane, pour mesurer les lettres et non leur
%  cadre : c'est le bord des lettres que le détourage rognerait.
%
%  PIÈGE MESURÉ le 16/09/2026 : la marge interne TikZ (.3333em) n'est PAS
%  évaluée dans la police du nœud, mais dans celle qui entoure le dessin.
%  Instrumenté : dans le nœud, em = 150 pt ; boîte brute du « N » :
%  41,7 x 37,5 mm, soit la lettre plus 2 x 1,3 mm — 1,3 mm étant .3333em
%  à 11 pt. Un retrait calculé à 150 pt (17,6 mm) ramenait la lettre à
%  6,6 mm de large, et le contrôle aurait dit OK sur une mesure absurde.
% ---------------------------------------------------------------------

\newcount\NbAlertes
\newdimen\FiligraneSep
\newsavebox\CreditsBoite

% Crédits de 4e : lignes coupées à la main. Leur largeur RÉELLE est celle
% de la plus longue ligne, mesurée dans une boîte à part (voir § 8, ISBN).
\def\CreditsTexte{%
    Réalisé à partir des contenus de 50ohm.de (en allemand)\\
    50ohm.de-Autorenteam, coordonné par le référat AJW du DARC e.\,V.\\
    \Auteur\\[3pt]
    Licence CC BY 4.0 --- \texttt{\Depot}}

\def\MesureNoeud#1#2#3{%
  \path let \p1=(#1.south west), \p2=(#1.north east),
            \p3=(#1.north west), \p4=(#1.south east), \p5=(0,0) in
    \pgfextra{%
      \expandafter\xdef\csname #2xmin\endcsname{\fpeval{round((min(\x1,\x2,\x3,\x4) - \x5 + #3)/1mm, 2)}}%
      \expandafter\xdef\csname #2xmax\endcsname{\fpeval{round((max(\x1,\x2,\x3,\x4) - \x5 - #3)/1mm, 2)}}%
      \expandafter\xdef\csname #2ymin\endcsname{\fpeval{round((min(\y1,\y2,\y3,\y4) - \y5 + #3)/1mm, 2)}}%
      \expandafter\xdef\csname #2ymax\endcsname{\fpeval{round((max(\y1,\y2,\y3,\y4) - \y5 - #3)/1mm, 2)}}%
    };}

% \ControleDans{libellé}{préfixe}{xmin}{xmax}{ymin}{ymax}
\def\ControleDans#1#2#3#4#5#6{%
  \ifnum\fpeval{\csname #2xmin\endcsname >= #3 && \csname #2xmax\endcsname <= #4
             && \csname #2ymin\endcsname >= #5 && \csname #2ymax\endcsname <= #6}=1
    \typeout{COUVERTURE-CONTROLE OK #1 : x \csname #2xmin\endcsname-\csname #2xmax\endcsname,
      y \csname #2ymin\endcsname-\csname #2ymax\endcsname\space dans x #3-#4, y #5-#6}%
  \else
    \global\advance\NbAlertes by 1
    \typeout{COUVERTURE-CONTROLE ALERTE #1 : x \csname #2xmin\endcsname-\csname #2xmax\endcsname,
      y \csname #2ymin\endcsname-\csname #2ymax\endcsname\space HORS de x #3-#4, y #5-#6}%
  \fi}

% \ControleDisjoints{libellé}{préfixe}{xmin}{xmax}{ymin}{ymax}
% Deux rectangles se chevauchent si et seulement s'ils se recouvrent sur
% les DEUX axes à la fois.
\def\ControleDisjoints#1#2#3#4#5#6{%
  \ifnum\fpeval{\csname #2xmin\endcsname < #4 && #3 < \csname #2xmax\endcsname
             && \csname #2ymin\endcsname < #6 && #5 < \csname #2ymax\endcsname}=1
    \global\advance\NbAlertes by 1
    \typeout{COUVERTURE-CONTROLE ALERTE #1 : chevauchement}%
  \else
    \typeout{COUVERTURE-CONTROLE OK #1 : disjoints}%
  \fi}

\begin{document}
% Marge interne du filigrane, évaluée dans la police qui ENTOURE le dessin :
% c'est là que TikZ évalue le .3333em de « inner sep » (cf. § 8).
\FiligraneSep=.3333em
% Largeur réelle des crédits : la plus longue ligne, dans leur police.
\sbox\CreditsBoite{\small\begin{tabular}{@{}l@{}}\CreditsTexte\end{tabular}}%
\begin{tikzpicture}[remember picture, overlay,
                    x=1mm, y=1mm, shift={(current page.south west)}]

% ---- Fonds ----------------------------------------------------------
% La 4e de couverture reste sur le blanc du papier.
% Le bandeau bleu couvre le dos, puis reprend au tiers extérieur de la
% 1re de couverture — même partition 2/3-1/3 que la page de titre.
% Couverture rigide : pas de filets au dos (ils tomberaient à côté des
% plis, vu la tolérance), et un dos bleu élargi de 2 mm de chaque côté.
\ifx\Reliure\MotRigide
  \fill[TitleBand] (\DosEtenduG, 0) rectangle (\DosEtenduD, \PapierH);
\else
  \fill[TitleBand] (\DosG, 0) rectangle (\DosD, \PapierH);
\fi
\fill[TitleBand] (\BandeauG, 0) rectangle (\PapierL, \PapierH);
% Filets d'accent : à la jonction sur la 1re, de part et d'autre du dos.
\fill[TitleAccent] (\BandeauAccent, 0) rectangle (\BandeauG, \PapierH);
\ifx\Reliure\MotRigide\else
  \fill[TitleAccent] (\DosAccentG, 0) rectangle (\DosG, \PapierH);
  \fill[TitleAccent] (\DosD, 0) rectangle (\DosAccentD, \PapierH);
\fi

% ---- 1re de couverture ----------------------------------------------
% Filigrane de classe, en réserve claire, détouré au bandeau.
\begin{scope}
\clip (\BandeauG, 0) rectangle (\PapierL, \PapierH);
\node[name=filigrane, anchor=north, align=center, text=white, opacity=0.16,
      font=\fontsize{\CorpsF}{\CorpsFI}\selectfont\bfseries]
  at (\FiligraneX, \FiligraneY) {\FiligraneTexte};
\end{scope}

\node[name=titre, anchor=west, align=left, text width=\TitreW mm]
  at (\TitreX, \TitreY) {%
    {\fontsize{54}{56}\selectfont\bfseries \TitreOuvrage}\\[10pt]
    {\Large\color{TitleBand}\bfseries \SousTitre}\\[16pt]
    {\large\color{black!75} \SousTitreClasse}};

\node[name=version, anchor=south west, align=left, font=\small, text=black!70]
  at (\TitreX, \VersionY) {Version \Version\\ Licence CC BY 4.0};

% ---- Dos -------------------------------------------------------------
% Un dos étroit ne reçoit pas de texte : en deçà de 6 mm, le débord au
% massicot ferait mordre le texte vertical sur les plats.
\ifnum\fpeval{\Dos >= 6}=1
  \node[name=dos, rotate=90, anchor=center, text=white, font=\bfseries,
        inner sep=0pt]
    at (\DosAxe, \MilieuH)
    {\TitreOuvrage\ \ --- \ \SousTitreClasse\ \ --- \ \Version};
\fi

% ---- 4e de couverture ------------------------------------------------
\node[name=quatre, anchor=north west, align=justify, text width=\QuatreW mm]
  at (\QuatreX, \QuatreY) {\QuatriemeTexte};

% Le texte vit dans \CreditsTexte (§ 8) pour être mesuré à part : la boîte
% du nœud a la largeur fixe \QuatreW et couvre l'emplacement ISBN sans
% qu'aucun caractère n'y soit. Retirer « text width » pour l'éviter change
% l'interlignage — essayé le 16/09/2026, 6 934 pixels différents.
\node[name=credits, anchor=south west, align=left, font=\small, text=black!70,
      text width=\QuatreW mm]
  at (\QuatreX, \QuatreBasY) {\CreditsTexte};

% ---- Repères de fabrication (mode épreuve uniquement) ----------------
\ifx\Reperes\ouimot
  % Traits de coupe, hors du fond perdu.
  \foreach \x in {\CoupeG, \CoupeD} {
    \draw[Repere, line width=0.25pt] (\x, \TechExt) -- (\x, \TechInt);
    \draw[Repere, line width=0.25pt] (\x, \HautInt) -- (\x, \HautExt);}
  \foreach \y in {\CoupeB, \CoupeH} {
    \draw[Repere, line width=0.25pt] (\TechExt, \y) -- (\TechInt, \y);
    \draw[Repere, line width=0.25pt] (\DroitInt, \y) -- (\DroitExt, \y);}
  % Limites du dos, en pointillés.
  \foreach \x in {\DosG, \DosD} {
    \draw[Repere, line width=0.25pt, dashed] (\x, \TechExt) -- (\x, \TechInt);
    \draw[Repere, line width=0.25pt, dashed] (\x, \HautInt) -- (\x, \HautExt);}
  % Zone de sécurité.
  \draw[Repere, line width=0.2pt, dotted]
    (\SecG, \SecB) rectangle (\SecD, \SecH);
  % Emplacement réservé au code-barres ISBN.
  \draw[Repere, line width=0.2pt] (\IsbnG, \SecB) rectangle (\IsbnD, \IsbnH);
  \node[font=\tiny\sffamily, text=Repere, align=center]
    at (\IsbnCX, \IsbnCY) {emplacement\\ISBN};
  % Cotes.
  \node[anchor=north, font=\tiny\sffamily, text=Repere]
    at (\DosAxe, \CoteHautY) {dos \Dos\ mm (\DosOrigine)};
  \node[anchor=south, font=\tiny\sffamily, text=Repere, align=center]
    at (\MilieuL, \CoteBasY) {\CotesTexte};
\fi

% ---- Contrôle de superposition (§ 8) ----------------------------------
% Rien n'est dessiné : des mesures et des lignes de journal.
\MesureNoeud{filigrane}{Fil}{\the\FiligraneSep}
\MesureNoeud{titre}{Tit}{0pt}
\MesureNoeud{version}{Ver}{0pt}
\MesureNoeud{quatre}{Qua}{0pt}
\MesureNoeud{credits}{Cre}{0pt}
% Bord droit des crédits = bord gauche + marge interne + plus longue ligne.
\xdef\Crexmax{\fpeval{round(\Crexmin + (\the\FiligraneSep + \the\wd\CreditsBoite)/1mm, 2)}}

% Le filigrane peut entrer dans la zone de sécurité (il est décoratif),
% pas franchir la coupe ni sortir du bandeau — le détourage le rognerait.
\ControleDans{filigrane dans le bandeau}{Fil}{\BandeauG}{\CoupeD}{\CoupeB}{\CoupeH}
\ControleDans{titre sur la 1re, hors bandeau}{Tit}{\PliPremiere}{\BandeauAccent}{\SecB}{\SecH}
\ControleDans{version sur la 1re, hors bandeau}{Ver}{\PliPremiere}{\BandeauAccent}{\SecB}{\SecH}
\ControleDans{texte de 4e dans la zone de securite}{Qua}{\SecG}{\PliQuatre}{\SecB}{\SecH}
\ControleDans{credits de 4e dans la zone de securite}{Cre}{\SecG}{\PliQuatre}{\SecB}{\SecH}
\ControleDisjoints{titre / version}{Tit}{\Verxmin}{\Verxmax}{\Verymin}{\Verymax}
\ControleDisjoints{texte de 4e / credits}{Qua}{\Crexmin}{\Crexmax}{\Creymin}{\Creymax}
\ControleDisjoints{texte de 4e / emplacement ISBN}{Qua}{\IsbnG}{\IsbnD}{\SecB}{\IsbnH}
\ControleDisjoints{credits / emplacement ISBN}{Cre}{\IsbnG}{\IsbnD}{\SecB}{\IsbnH}
\ifnum\fpeval{\Dos >= 6}=1
  \MesureNoeud{dos}{Dos}{0pt}
  \ControleDans{texte du dos entre les plis}{Dos}{\DosG}{\DosD}{\SecB}{\SecH}
\else
  \typeout{COUVERTURE-CONTROLE OK dos de \Dos\space mm : trop etroit, sans texte}%
\fi

\end{tikzpicture}
\typeout{COUVERTURE-BILAN alertes=\the\NbAlertes\space; classe=\Classe\space;
  format=\Format\space; reliure=\Reliure\space; pages=\NbPages\space;
  version=\Version\space;
  dos=\Dos\space mm (\DosOrigine); fini=\FiniL x \PageHauteur\space mm;
  papier=\PapierL x \PapierH\space mm}
\end{document}
"
%
%  Pour le fichier destiné à l'imprimeur (sans repères) :
%
%      lualatex "\def\Reperes{non}% =====================================================================
%  couverture.tex — couverture dos carré collé pour les livres 50ohm-fr
% =====================================================================
%
%  Produit une couverture d'un seul tenant : 4e de couverture | dos |
%  1re de couverture, aux dimensions exactes du livre relié, fond perdu
%  compris. L'épaisseur du dos est CALCULÉE à partir du nombre de pages,
%  du grammage et du bouffant du papier.
%
%  ORIGINE ET ÉTAT. Écrit le 20/08/2026 dans Fichierstravail/ (hors git),
%  versé au dépôt le 16/09/2026 pour que fenetre_compilation.py puisse
%  l'appeler. Ce qui a changé à cette occasion :
%
%    - le FORMAT est un paramètre (\Format = a4, 20x24 ou 20x24-marge,
%      mêmes noms que --format de build_book.py) ;
%    - la pagination et la version peuvent être IMPOSÉES de l'extérieur
%      (\PagesForcees, \VersionForcee) : la fenêtre lit la pagination dans
%      le journal du livre qu'elle vient de compiler, au lieu de se fier à
%      une table qui vieillit ;
%    - un dos imposé par l'imprimeur n'est appliqué QUE pour la pagination
%      et le format pour lesquels il a été donné ;
%    - un CONTRÔLE DE SUPERPOSITION est calculé par LaTeX lui-même, à partir
%      des boîtes réelles des nœuds, et écrit dans le journal ;
%    - le tome SWL est refusé explicitement : sa 4e de couverture n'est pas
%      rédigée (le texte commun parle des classes N, E et A).
%
%    Neutralité vérifiée : en A4, pour N (258 p., a.2) et NEA (816 p., a.3),
%    en épreuve comme en impression, le rendu est identique au pixel près à
%    celui du gabarit d'origine.
%
%  COMPILATION (LuaLaTeX, depuis la racine du dépôt) :
%
%      lualatex couverture.tex
%      lualatex couverture.tex        <-- DEUX FOIS, obligatoirement
%
%  DEUX PASSES SONT NÉCESSAIRES. Le dessin s'ancre sur « current page »,
%  dont la position n'est connue qu'après un premier passage : à la
%  première compilation, tout se tasse dans le coin supérieur et la
%  couverture paraît ratée alors qu'elle ne l'est pas. Aucune erreur n'est
%  émise — seul le rendu le montre. Le contrôle de superposition, lui
%  aussi, ne vaut qu'à la seconde passe.
%
%  Le plus simple est de passer par fenetre_compilation.py (case
%  « couverture pour l'imprimeur »), qui enchaîne les deux passes, produit
%  l'épreuve ET le fichier d'impression, et relit le contrôle.
%
%  À la main, pour une autre classe, un autre format, une pagination mesurée :
%
%      lualatex "\def\Classe{A}\def\Format{20x24}\def\PagesForcees{482}% =====================================================================
%  couverture.tex — couverture dos carré collé pour les livres 50ohm-fr
% =====================================================================
%
%  Produit une couverture d'un seul tenant : 4e de couverture | dos |
%  1re de couverture, aux dimensions exactes du livre relié, fond perdu
%  compris. L'épaisseur du dos est CALCULÉE à partir du nombre de pages,
%  du grammage et du bouffant du papier.
%
%  ORIGINE ET ÉTAT. Écrit le 20/08/2026 dans Fichierstravail/ (hors git),
%  versé au dépôt le 16/09/2026 pour que fenetre_compilation.py puisse
%  l'appeler. Ce qui a changé à cette occasion :
%
%    - le FORMAT est un paramètre (\Format = a4, 20x24 ou 20x24-marge,
%      mêmes noms que --format de build_book.py) ;
%    - la pagination et la version peuvent être IMPOSÉES de l'extérieur
%      (\PagesForcees, \VersionForcee) : la fenêtre lit la pagination dans
%      le journal du livre qu'elle vient de compiler, au lieu de se fier à
%      une table qui vieillit ;
%    - un dos imposé par l'imprimeur n'est appliqué QUE pour la pagination
%      et le format pour lesquels il a été donné ;
%    - un CONTRÔLE DE SUPERPOSITION est calculé par LaTeX lui-même, à partir
%      des boîtes réelles des nœuds, et écrit dans le journal ;
%    - le tome SWL est refusé explicitement : sa 4e de couverture n'est pas
%      rédigée (le texte commun parle des classes N, E et A).
%
%    Neutralité vérifiée : en A4, pour N (258 p., a.2) et NEA (816 p., a.3),
%    en épreuve comme en impression, le rendu est identique au pixel près à
%    celui du gabarit d'origine.
%
%  COMPILATION (LuaLaTeX, depuis la racine du dépôt) :
%
%      lualatex couverture.tex
%      lualatex couverture.tex        <-- DEUX FOIS, obligatoirement
%
%  DEUX PASSES SONT NÉCESSAIRES. Le dessin s'ancre sur « current page »,
%  dont la position n'est connue qu'après un premier passage : à la
%  première compilation, tout se tasse dans le coin supérieur et la
%  couverture paraît ratée alors qu'elle ne l'est pas. Aucune erreur n'est
%  émise — seul le rendu le montre. Le contrôle de superposition, lui
%  aussi, ne vaut qu'à la seconde passe.
%
%  Le plus simple est de passer par fenetre_compilation.py (case
%  « couverture pour l'imprimeur »), qui enchaîne les deux passes, produit
%  l'épreuve ET le fichier d'impression, et relit le contrôle.
%
%  À la main, pour une autre classe, un autre format, une pagination mesurée :
%
%      lualatex "\def\Classe{A}\def\Format{20x24}\def\PagesForcees{482}\input{couverture.tex}"
%
%  Pour le fichier destiné à l'imprimeur (sans repères) :
%
%      lualatex "\def\Reperes{non}\input{couverture.tex}"
%
%  ---------------------------------------------------------------------
%  DEUX MODES, un seul paramètre : \Reperes
%
%    oui  — mode ÉPREUVE. Ajoute une marge technique, les traits de coupe,
%           les limites du dos, la zone de sécurité et les cotes. C'est ce
%           qu'il faut pour vérifier à l'écran.
%    non  — mode PRODUCTION. Papier = format final + fond perdu, rien
%           d'autre. C'est ce fichier-là qu'on envoie à l'imprimeur.
%
%  ---------------------------------------------------------------------
%  CONTRÔLE DE SUPERPOSITION
%
%  À la fin du dessin, chaque bloc de texte est mesuré par ses ancres TikZ
%  (quatre coins, ce qui vaut aussi pour le texte pivoté du dos) et comparé
%  à la zone où il doit tenir, puis aux blocs voisins. Le journal porte une
%  ligne par vérification :
%
%      COUVERTURE-CONTROLE OK      <quoi> : ...
%      COUVERTURE-CONTROLE ALERTE  <quoi> : ...
%
%  et un bilan : COUVERTURE-BILAN alertes=<n> ; ...
%
%  Ce que le contrôle VOIT : un bloc qui franchit le pli du dos, la coupe,
%  la zone de sécurité ou le bord du bandeau ; deux blocs qui se chevauchent ;
%  un bloc qui empiète sur l'emplacement réservé au code-barres ISBN.
%  Ce qu'il NE VOIT PAS : une ligne trop longue à l'INTÉRIEUR d'un bloc à
%  largeur fixe (la boîte garde sa largeur, le texte déborde). Celle-là se
%  signale par un « Overfull \hbox » dans le même journal, que la fenêtre
%  compte aussi.
%
%  ---------------------------------------------------------------------
%  AVERTISSEMENT SUR L'ÉPAISSEUR DU DOS
%
%  Le dos se calcule ainsi :
%
%      épaisseur d'une feuille (mm) = grammage x bouffant / 1000
%      dos (mm)                     = (nombre de pages / 2) x épaisseur
%
%  On divise par deux parce qu'une feuille porte DEUX pages.
%
%  Le « bouffant » (main, volume, bulk) dépend du papier et n'est pas
%  devinable : un offset 80 g courant est autour de 1,30, un bouffant
%  d'édition monte à 1,75 et plus. UN ÉCART DE BOUFFANT DÉPLACE LE DOS DE
%  PLUSIEURS MILLIMÈTRES. Demander à l'imprimeur soit son bouffant, soit
%  directement l'épaisseur du dos, et reporter ci-dessous. Tant que cette
%  valeur n'est pas confirmée, la couverture n'est pas bonne à tirer.
%
%  Si l'imprimeur donne le dos en clair, forcer la valeur avec \DosImpose.
% =====================================================================

\documentclass[11pt]{article}

% ---------------------------------------------------------------------
% 1. PARAMÈTRES — c'est la seule zone à modifier
% ---------------------------------------------------------------------

% Classe de l'ouvrage : N, E, A, NE, EA ou NEA.
% (\providecommand : une valeur passée en ligne de commande l'emporte.)
\providecommand{\Classe}{NEA}

% Mode épreuve (oui) ou production (non).
\providecommand{\Reperes}{oui}

% Format du livre FINI : a4, 20x24 ou 20x24-marge — les noms de --format
% dans build_book.py. Les deux maquettes 20 x 24 ne diffèrent que par leur
% colonne de marge intérieure : le livre fini, donc la couverture, est le
% même (200 x 240 mm).
\providecommand{\Format}{a4}

% Reliure : « souple » (dos carré collé, fond perdu 3 mm) ou « rigide »
% (couverture rigide, dite hardcover). En rigide, consignes de l'imprimeur
% du 17/09/2026 : chaque plat déborde du livre de 5 mm en largeur et en
% hauteur (le carton dépasse du bloc), et le fichier porte 15 mm de rabat
% autour du carton au lieu de 3 mm de fond perdu. Les aplats vont déjà au
% bord du papier : ils remplissent le rabat d'eux-mêmes.
\providecommand{\Reliure}{souple}

% Pagination et version imposées de l'extérieur. Vides = tables du § 2.
% La fenêtre les renseigne toujours : la pagination vient du journal du
% livre réellement compilé, dans le format demandé.
\providecommand{\PagesForcees}{}
\providecommand{\VersionForcee}{}

% Papier de l'INTÉRIEUR du livre (pas de la couverture).
% CoolLibri, papier « Standard 90 g blanc ».
\def\Grammage{90}        % g/m²
\def\Bouffant{1.206}     % main DÉDUITE, voir ci-dessous

% Épaisseur du dos IMPOSÉE par l'imprimeur, en mm, par classe.
% Une valeur > 0 prime sur le calcul — MAIS SEULEMENT pour la pagination et
% le format pour lesquels l'imprimeur l'a donnée (\DosImposePages<classe>,
% \DosImposeFormat<classe>). Hors de ce cas, le dos est calculé et le
% journal le dit : un dos de 14 mm donné pour 258 pages A4 ne vaut rien pour
% 262 pages, ni pour le même livre recomposé en 20 x 24.
%
% ATTENTION : cette valeur est propre à CHAQUE ÉDITION. Un dos unique
% appliqué aux quatre livres serait faux pour trois d'entre eux.
%
%  N   — 14 mm, VALEUR DONNÉE PAR COOLIBRI pour 258 pages en 90 g, A4.
%        Fichier attendu : 440 x 303 mm, ce que ce gabarit produit en mode
%        production (2 x 210 + 14 + 2 x 3 = 440 ; 297 + 2 x 3 = 303).
%        Le N de la a.3 compte 262 pages : cette valeur ne s'y applique plus.
%  E, A, NEA — non demandés à l'imprimeur : ils sont CALCULÉS avec la main
%        déduite ci-dessous. À faire confirmer avant tout tirage.
%
% La main de 1,206 n'est pas publiée par CoolLibri : elle est déduite de
% leur propre chiffre, 14 mm / (0,090 x 258/2) = 1,206. Elle reproduit donc
% exactement le dos de N, et n'est qu'une ESTIMATION pour les autres.
% Écarts que cela recouvre, pour mémoire : une main de 1,0 donnerait
% 11,61 mm pour N, une main de 1,3 en donnerait 16,77. Il ne fallait pas
% deviner.
%  E   — 17 mm, VALEUR DONNÉE PAR L'IMPRIMEUR le 17/09/2026 pour 276 pages,
%        20 x 24, couverture RIGIDE, reliure collée (Klebebindung), papier
%        115 g. En reliure cousue (Fadenheftung), il annonce 18 mm.
%
% \DosImposeReliure<classe> : la reliure pour laquelle la valeur a été
% donnée. Le dos d'une couverture rigide n'est pas celui d'un dos carré
% collé souple, même à pagination égale.
\def\DosImposeN{14}   \def\DosImposePagesN{258}  \def\DosImposeFormatN{a4}
\def\DosImposeE{17}   \def\DosImposePagesE{276}  \def\DosImposeFormatE{20x24-marge}
\def\DosImposeA{0}    \def\DosImposePagesA{0}    \def\DosImposeFormatA{}
\def\DosImposeNE{0}   \def\DosImposePagesNE{0}   \def\DosImposeFormatNE{}
\def\DosImposeEA{0}   \def\DosImposePagesEA{0}   \def\DosImposeFormatEA{}
\def\DosImposeNEA{0}  \def\DosImposePagesNEA{0}  \def\DosImposeFormatNEA{}
\def\DosImposeReliureN{souple}  \def\DosImposeReliureE{rigide}
\def\DosImposeReliureA{souple}  \def\DosImposeReliureNE{souple}
\def\DosImposeReliureEA{souple} \def\DosImposeReliureNEA{souple}

% Fond perdu, en mm. 3 mm est la valeur courante ; certains imprimeurs
% demandent 5 mm.
\def\FondPerdu{3}

% Marge technique autour du fond perdu, en mm. Sert à loger les traits de
% coupe en mode épreuve. Ignorée en mode production.
\def\MargeTech{12}

% Zone de sécurité, en mm : aucun texte ne doit s'en approcher davantage
% du bord de coupe.
\def\Securite{8}

% ---------------------------------------------------------------------
% 2. DONNÉES PAR CLASSE
%    pages : nombre de pages du PDF intérieur, COUVERTURE NON COMPRISE.
%    Un dos carré collé impose un nombre de pages PAIR ; le fichier
%    refuse de compiler si le compte est impair.
%
%    Ces tables ne servent que si \PagesForcees et \VersionForcee sont
%    vides, c'est-à-dire à la main. Valeurs A4 mesurées sur les PDF de la
%    a.3 (CLAUDE.md § 1, 05/09/2026). En 20 x 24, la pagination change
%    entièrement : NEA passe de 816 à 1036 pages.
% ---------------------------------------------------------------------

\def\DonneesN{262}
\def\DonneesE{214}
\def\DonneesA{386}
\def\DonneesNE{0}      % édition jamais produite
\def\DonneesEA{0}      % édition jamais produite
\def\DonneesNEA{816}

% Version imprimée sur le dos et la 1re de couverture, PAR CLASSE.
% « Les versions sont suivies par classe, chacune évoluant à son rythme »
% (CHANGELOG.md).
\def\VersionN{a.3}
\def\VersionE{a.3}
\def\VersionA{a.3}
\def\VersionNE{a.3}
\def\VersionEA{a.3}
\def\VersionNEA{a.3}

% ---------------------------------------------------------------------
% 3. TEXTES
% ---------------------------------------------------------------------

\def\TitreOuvrage{50\,Ohm}
\def\SousTitre{Préparation à l'examen radioamateur}

% Sous-titre propre à chaque classe (repris de FR_TITLES, build_book.py).
\def\SousTitreN{Cours complet — classe N}
\def\SousTitreE{Cours complet — classe E}
\def\SousTitreA{Cours complet — classe A}
\def\SousTitreNE{Cours complet — classes N et E}
\def\SousTitreEA{Cours complet — classes E et A}
\def\SousTitreNEA{Cours complet — classes N, E et A}

% Lettres du filigrane, empilées de haut en bas (cf. build_book.py v0.20).
\def\LettresN{N}
\def\LettresE{E}
\def\LettresA{A}
\def\LettresNE{N,E}
\def\LettresEA{E,A}
\def\LettresNEA{N,E,A}

% Texte de 4e de couverture.
% NB : il parle des classes N, E et A quelle que soit l'édition. Juste pour
% N, E, A et NEA ; approximatif pour NE et EA ; faux pour SWL, d'où son refus.
\def\QuatriemeTexte{%
Ce guide d'étude prépare aux examens radioamateur allemands des classes
N, E et A. Il traduit en français les contenus du projet 50ohm.de, publiés
par le référat AJW du DARC e.\,V., et les complète d'encarts « En France »
qui signalent, sans rien retirer au texte allemand, ce que la
réglementation française prévoit sur le même sujet.\par\smallskip
Les questions officielles du catalogue allemand accompagnent chaque
notion, avec leurs réponses.%
}

\def\Auteur{Traduit avec l'aide d'une IA par Pierre F4JWI}
\def\Depot{github.com/DrPLG/50ohm-fr}

% =====================================================================
%  À PARTIR D'ICI, PLUS RIEN À MODIFIER
% =====================================================================

\usepackage{iftex}
\RequireLuaTeX
\usepackage{xfp}
\usepackage{tikz}
\usetikzlibrary{calc}
\usepackage{fontspec}
\usepackage[french]{babel}
\usepackage{geometry}

% ---------------------------------------------------------------------
% 4. RÉSOLUTION DES DONNÉES DE LA CLASSE
% ---------------------------------------------------------------------

\makeatletter
% Tome SWL : refus explicite, plutôt qu'un « classe inconnue » qui laisserait
% croire à un oubli. Le texte de 4e parle des classes N, E et A.
\def\@swlmot{SWL}
\ifx\Classe\@swlmot
  \PackageError{couverture}{Pas de couverture pour le tome SWL}%
    {Sa 4e de couverture n'est pas redigee : le texte commun parle des
     classes N, E et A. C'est une decision editoriale a prendre d'abord.}%
\fi

\def\@resoudre#1#2{%
  \expandafter\let\expandafter#1\csname #2\Classe\endcsname
  \ifx#1\relax
    \PackageError{couverture}{Classe \Classe\space inconnue : pas de #2\Classe}%
      {Valeurs admises : N, E, A, NE, EA, NEA.}%
  \fi}
\@resoudre\NbPages{Donnees}
\@resoudre\SousTitreClasse{SousTitre}
\@resoudre\Lettres{Lettres}
\@resoudre\DosImpose{DosImpose}
\@resoudre\DosImposePages{DosImposePages}
\@resoudre\DosImposeFormat{DosImposeFormat}
\@resoudre\DosImposeReliure{DosImposeReliure}
\@resoudre\Version{Version}

% Valeurs imposées de l'extérieur. Test de vacuité par \detokenize : il ne
% dépend pas de la façon dont la macro a été définie.
\if\relax\detokenize\expandafter{\PagesForcees}\relax\else
  \let\NbPages\PagesForcees
\fi
\if\relax\detokenize\expandafter{\VersionForcee}\relax\else
  \let\Version\VersionForcee
\fi

% Format du livre fini.
\def\@formatAquatre{a4}
\def\@formatVingt{20x24}
\def\@formatVingtMarge{20x24-marge}
\ifx\Format\@formatAquatre
  \def\PageLargeur{210}\def\PageHauteur{297}
\else\ifx\Format\@formatVingt
  \def\PageLargeur{200}\def\PageHauteur{240}
\else\ifx\Format\@formatVingtMarge
  \def\PageLargeur{200}\def\PageHauteur{240}
\else
  \PackageError{couverture}{Format \Format\space inconnu}%
    {Valeurs admises : a4, 20x24, 20x24-marge.}%
  \def\PageLargeur{210}\def\PageHauteur{297}
\fi\fi\fi
% Reliure. Hors de « souple » et « rigide » : erreur, pas de supposition.
\def\@reliuresouple{souple}
\def\@reliurerigide{rigide}
\ifx\Reliure\@reliurerigide
  \edef\PageLargeur{\fpeval{\PageLargeur + 5}}
  \edef\PageHauteur{\fpeval{\PageHauteur + 5}}
  \def\FondPerdu{15}
\else\ifx\Reliure\@reliuresouple\else
  \PackageError{couverture}{Reliure \Reliure\space inconnue}%
    {Valeurs admises : souple, rigide.}%
\fi\fi
\makeatother

% Garde-fou : une édition jamais produite n'a pas de couverture.
\ifnum\NbPages=0
  \PackageError{couverture}{L'edition \Classe\space n'a pas de pagination}%
    {Cette edition n'a jamais ete compilee. Renseigner \string\Donnees\Classe
     ou passer \string\PagesForcees.}%
\fi

% Avertissement : l'imprimeur demande un nombre de pages multiple de 4
% (17/09/2026). Pas une erreur : c'est son exigence, pas une loi du papier.
\ifnum\fpeval{\NbPages - 4*floor(\NbPages/4)}=0 \else
  \typeout{COUVERTURE-AVERTISSEMENT \NbPages\space pages : pas un multiple
    de 4, que l'imprimeur exige -- ajouter des pages a l'interieur}%
\fi

% Garde-fou : un dos carré collé impose un nombre de pages pair.
\ifodd\NbPages
  \PackageError{couverture}{Nombre de pages impair (\NbPages)}%
    {Un dos carre colle exige un nombre PAIR de pages : une feuille en
     porte deux. Ajouter une page blanche a l'interieur.}%
\fi

% ---------------------------------------------------------------------
% 5. GÉOMÉTRIE — c'est ici que se joue l'exactitude du dos
% ---------------------------------------------------------------------

% Épaisseur du dos : valeur imposée si elle est donnée ET qu'elle vaut pour
% cette pagination et ce format ; calcul sinon.
\edef\DosCalcule{\fpeval{round(\NbPages / 2 * \Grammage * \Bouffant / 1000, 2)}}
\def\DosOrigine{calculee}
\ifnum\fpeval{\DosImpose > 0}=1
  \ifnum\fpeval{\NbPages = \DosImposePages}=1
    \ifx\Format\DosImposeFormat
      \ifx\Reliure\DosImposeReliure
        \def\DosOrigine{imposee}
      \fi
    \fi
  \fi
  \def\MotCalculee{calculee}
  \ifx\DosOrigine\MotCalculee
    \typeout{COUVERTURE-AVERTISSEMENT dos impose de \DosImpose\space mm
      donne pour \DosImposePages\space pages en \DosImposeFormat\space
      (reliure \DosImposeReliure) -- inapplicable a \NbPages\space pages en
      \Format\space (reliure \Reliure), dos calcule}%
  \fi
\fi
% NB : pas de \@ dans ces noms. Cette section est hors \makeatletter, et
% « \def\@imposee{...} » y redéfinit \@ (suivi du texte « imposee ») : le
% test échouait toujours, en silence. Constaté le 16/09/2026 — le dos de N
% tombait juste par coïncidence, 14,0 mm calculés pour 14 imposés.
\def\MotRigide{rigide}
\def\MotImposee{imposee}
\ifx\DosOrigine\MotImposee
  \edef\Dos{\DosImpose}
\else
  \edef\Dos{\DosCalcule}
\fi

% Marge technique : uniquement en mode épreuve.
\def\ouimot{oui}
\ifx\Reperes\ouimot \edef\Tech{\MargeTech}\else \def\Tech{0}\fi

% Dimensions du papier.
%   largeur = marge + fond perdu + 4e + dos + 1re + fond perdu + marge
%   hauteur = marge + fond perdu + hauteur + fond perdu + marge
\edef\PapierL{\fpeval{2*\PageLargeur + \Dos + 2*\FondPerdu + 2*\Tech}}
\edef\PapierH{\fpeval{\PageHauteur + 2*\FondPerdu + 2*\Tech}}

\geometry{paperwidth=\PapierL mm, paperheight=\PapierH mm,
          margin=0pt, nohead, nofoot}

% ---------------------------------------------------------------------
%  TOUTES LES COTES SONT FIGÉES ICI, ET NULLE PART AILLEURS.
%
%  Raison, payée par trois compilations ratées : un \fpeval laissé dans
%  une coordonnée ou dans le texte d'un nœud TikZ n'est pas développé au
%  moment où TikZ analyse le chemin. On obtient « Giving up on this path.
%  Did you forget a semicolon? », qui accuse une ponctuation alors que la
%  faute est ailleurs. Le corps du document ci-dessous n'emploie donc que
%  des macros déjà calculées.
% ---------------------------------------------------------------------

% Repères horizontaux, mesurés depuis le bord GAUCHE du papier.
\edef\CoupeG{\fpeval{\Tech + \FondPerdu}}          % coupe gauche (4e de couv.)
\edef\DosG{\fpeval{\CoupeG + \PageLargeur}}         % dos, début
\edef\DosD{\fpeval{\DosG + \Dos}}                   % dos, fin
\edef\CoupeD{\fpeval{\DosD + \PageLargeur}}         % coupe droite (1re de couv.)
% Repères verticaux, depuis le bord BAS.
\edef\CoupeB{\fpeval{\Tech + \FondPerdu}}
\edef\CoupeH{\fpeval{\CoupeB + \PageHauteur}}
% Axes.
\edef\DosAxe{\fpeval{(\DosG + \DosD)/2}}
\edef\MilieuH{\fpeval{(\CoupeB + \CoupeH)/2}}
\edef\MilieuL{\fpeval{\PapierL/2}}

% Format fini, pour les cotes affichées.
\edef\FiniL{\fpeval{2*\PageLargeur + \Dos}}

% Bandeau de la 1re de couverture : même proportion que la page de titre.
\edef\BandeauL{\fpeval{round(0.34*\PageLargeur, 2)}}
\edef\BandeauG{\fpeval{\CoupeD - \BandeauL}}
\edef\BandeauAccent{\fpeval{\BandeauG - 0.7}}
\edef\DosAccentG{\fpeval{\DosG - 0.7}}
\edef\DosAccentD{\fpeval{\DosD + 0.7}}
% Couverture rigide : le bleu du dos déborde de 2 mm sur chaque plat, pour
% absorber la tolérance de l'imprimeur sur l'épaisseur du dos (±10 %).
\edef\DosEtenduG{\fpeval{\DosG - 2}}
\edef\DosEtenduD{\fpeval{\DosD + 2}}

% Filigrane de classe.
\edef\FiligraneX{\fpeval{(\BandeauG + \CoupeD)/2}}
\edef\FiligraneY{\fpeval{\CoupeH - 6}}

% Bloc-titre de la 1re de couverture.
\edef\TitreX{\fpeval{\DosD + 12}}
\edef\TitreY{\fpeval{\CoupeH - 72}}
\edef\TitreW{\fpeval{\BandeauG - \DosD - 16}}
\edef\VersionY{\fpeval{\CoupeB + 14}}

% 4e de couverture.
\edef\QuatreX{\fpeval{\CoupeG + \Securite + 6}}
\edef\QuatreY{\fpeval{\CoupeH - \Securite - 14}}
\edef\QuatreBasY{\fpeval{\CoupeB + \Securite + 22}}
\edef\QuatreW{\fpeval{\PageLargeur - 2*\Securite - 12}}

% Zone de sécurité.
\edef\SecG{\fpeval{\CoupeG + \Securite}}
\edef\SecD{\fpeval{\CoupeD - \Securite}}
\edef\SecB{\fpeval{\CoupeB + \Securite}}
\edef\SecH{\fpeval{\CoupeH - \Securite}}

% Zones de contrôle : ce qui doit rester de part et d'autre du pli du dos.
\edef\PliQuatre{\fpeval{\DosG - \Securite}}         % 4e : limite côté dos
\edef\PliPremiere{\fpeval{\DosD + \Securite}}       % 1re : limite côté dos

% Traits de coupe et cotes, dans la marge technique.
\edef\TechExt{\fpeval{\Tech - 8}}
\edef\TechInt{\fpeval{\Tech - 1}}
\edef\TechCote{\fpeval{\Tech - 9}}
\edef\HautInt{\fpeval{\PapierH - \Tech + 1}}
\edef\HautExt{\fpeval{\PapierH - \Tech + 8}}
\edef\HautCote{\fpeval{\PapierH - \Tech + 9}}
\edef\DroitInt{\fpeval{\PapierL - \Tech + 1}}
\edef\DroitExt{\fpeval{\PapierL - \Tech + 8}}

% Emplacement du code-barres ISBN : en bas à DROITE de la 4e de couverture,
% selon l'usage éditorial, et non à gauche où il chevaucherait les mentions
% de licence.
\edef\IsbnD{\fpeval{\DosG - \Securite}}
\edef\IsbnG{\fpeval{\IsbnD - 40}}
\edef\IsbnH{\fpeval{\SecB + 25}}
\edef\IsbnCX{\fpeval{\IsbnG + 20}}
\edef\IsbnCY{\fpeval{\SecB + 12.5}}

% Cartouches de cotes, dans la marge technique. Ancrés vers l'INTÉRIEUR du
% papier : ancrés vers l'extérieur, ils débordaient et étaient rognés.
\edef\CoteBasY{\fpeval{1}}
\edef\CoteHautY{\fpeval{\PapierH - 1}}

% ---------------------------------------------------------------------
% 6. COULEURS ET POLICES — reprises de la page de titre (build_book.py)
% ---------------------------------------------------------------------

\definecolor{TitleBand}{cmyk}{.9,.55,.1,.35}   % bleu profond
\definecolor{TitleAccent}{cmyk}{.05,.8,1,0}    % rouge DARC (accent fin)
\definecolor{Repere}{cmyk}{0,1,1,0}            % magenta : repères de montage

\setmainfont{Libertinus Serif}[
  ItalicFont     = Libertinus Serif Italic,
  BoldFont       = Libertinus Serif Bold,
  BoldItalicFont = Libertinus Serif Bold Italic]

\pagestyle{empty}

% ---------------------------------------------------------------------
% 7. FILIGRANE DE CLASSE
%    Le texte est construit ici, et non dans le nœud : un \\ produit à
%    l'intérieur d'un \foreach referme le groupe d'alignement de TikZ.
%    Une lettre seule est grande ; plusieurs s'empilent et se réduisent
%    (build_book.py v0.20 : « NEA » à 105 pt débordait encore).
% ---------------------------------------------------------------------

\makeatletter
\def\FiligraneTexte{}
\newif\ifpremierelettre \premierelettretrue
\count255=0
\@for\@tempa:=\Lettres\do{%
  \advance\count255 by 1\relax
  \ifpremierelettre \premierelettrefalse
  \else \g@addto@macro\FiligraneTexte{\\}\fi
  \expandafter\g@addto@macro\expandafter\FiligraneTexte\expandafter{\@tempa}}%
\xdef\NbLettres{\the\count255}
\makeatother

\ifnum\NbLettres=1 \def\CorpsF{150}\else\def\CorpsF{86}\fi
\edef\CorpsFI{\fpeval{round(\CorpsF*1.02, 0)}}

% Texte du cartouche de cotes, lui aussi construit hors du nœud TikZ.
%
% NE PAS y employer \textsuperscript : dans un nœud « align=center », il
% casse le chemin TikZ (« Giving up on this path. Did you forget a
% semicolon? »), et l'erreur est signalée sur le nœud SUIVANT — elle
% accuse donc une ponctuation là où la faute est ici. Isolé par essais :
% $^2$ et le caractère ² passent, \textsuperscript non.
\def\CotesTexte{%
  \Classe\ --- \NbPages\ pages --- papier \Grammage\ g/m$^2$,
  bouffant \Bouffant\ --- fond perdu \FondPerdu\ mm\\
  papier \PapierL\ mm sur \PapierH\ mm --- fini \FiniL\ mm sur
  \PageHauteur\ mm\\
  EPREUVE, ne pas envoyer a l'impression%
}

% ---------------------------------------------------------------------
% 8. CONTRÔLE DE SUPERPOSITION
%
%  \MesureNoeud{nœud}{préfixe}{retrait} mesure la boîte d'un nœud par ses
%  quatre coins (un nœud pivoté a ses ancres pivotées : le min et le max des
%  quatre suffisent) et la ramène au repère du papier, en mm. Le point (0,0)
%  est lu par le même « let » que les ancres : les deux passent donc par la
%  même transformation, et la comparaison ne dépend pas de l'endroit où
%  TeX a posé l'image sur la page.
%
%  Le retrait sert au seul filigrane, pour mesurer les lettres et non leur
%  cadre : c'est le bord des lettres que le détourage rognerait.
%
%  PIÈGE MESURÉ le 16/09/2026 : la marge interne TikZ (.3333em) n'est PAS
%  évaluée dans la police du nœud, mais dans celle qui entoure le dessin.
%  Instrumenté : dans le nœud, em = 150 pt ; boîte brute du « N » :
%  41,7 x 37,5 mm, soit la lettre plus 2 x 1,3 mm — 1,3 mm étant .3333em
%  à 11 pt. Un retrait calculé à 150 pt (17,6 mm) ramenait la lettre à
%  6,6 mm de large, et le contrôle aurait dit OK sur une mesure absurde.
% ---------------------------------------------------------------------

\newcount\NbAlertes
\newdimen\FiligraneSep
\newsavebox\CreditsBoite

% Crédits de 4e : lignes coupées à la main. Leur largeur RÉELLE est celle
% de la plus longue ligne, mesurée dans une boîte à part (voir § 8, ISBN).
\def\CreditsTexte{%
    Réalisé à partir des contenus de 50ohm.de (en allemand)\\
    50ohm.de-Autorenteam, coordonné par le référat AJW du DARC e.\,V.\\
    \Auteur\\[3pt]
    Licence CC BY 4.0 --- \texttt{\Depot}}

\def\MesureNoeud#1#2#3{%
  \path let \p1=(#1.south west), \p2=(#1.north east),
            \p3=(#1.north west), \p4=(#1.south east), \p5=(0,0) in
    \pgfextra{%
      \expandafter\xdef\csname #2xmin\endcsname{\fpeval{round((min(\x1,\x2,\x3,\x4) - \x5 + #3)/1mm, 2)}}%
      \expandafter\xdef\csname #2xmax\endcsname{\fpeval{round((max(\x1,\x2,\x3,\x4) - \x5 - #3)/1mm, 2)}}%
      \expandafter\xdef\csname #2ymin\endcsname{\fpeval{round((min(\y1,\y2,\y3,\y4) - \y5 + #3)/1mm, 2)}}%
      \expandafter\xdef\csname #2ymax\endcsname{\fpeval{round((max(\y1,\y2,\y3,\y4) - \y5 - #3)/1mm, 2)}}%
    };}

% \ControleDans{libellé}{préfixe}{xmin}{xmax}{ymin}{ymax}
\def\ControleDans#1#2#3#4#5#6{%
  \ifnum\fpeval{\csname #2xmin\endcsname >= #3 && \csname #2xmax\endcsname <= #4
             && \csname #2ymin\endcsname >= #5 && \csname #2ymax\endcsname <= #6}=1
    \typeout{COUVERTURE-CONTROLE OK #1 : x \csname #2xmin\endcsname-\csname #2xmax\endcsname,
      y \csname #2ymin\endcsname-\csname #2ymax\endcsname\space dans x #3-#4, y #5-#6}%
  \else
    \global\advance\NbAlertes by 1
    \typeout{COUVERTURE-CONTROLE ALERTE #1 : x \csname #2xmin\endcsname-\csname #2xmax\endcsname,
      y \csname #2ymin\endcsname-\csname #2ymax\endcsname\space HORS de x #3-#4, y #5-#6}%
  \fi}

% \ControleDisjoints{libellé}{préfixe}{xmin}{xmax}{ymin}{ymax}
% Deux rectangles se chevauchent si et seulement s'ils se recouvrent sur
% les DEUX axes à la fois.
\def\ControleDisjoints#1#2#3#4#5#6{%
  \ifnum\fpeval{\csname #2xmin\endcsname < #4 && #3 < \csname #2xmax\endcsname
             && \csname #2ymin\endcsname < #6 && #5 < \csname #2ymax\endcsname}=1
    \global\advance\NbAlertes by 1
    \typeout{COUVERTURE-CONTROLE ALERTE #1 : chevauchement}%
  \else
    \typeout{COUVERTURE-CONTROLE OK #1 : disjoints}%
  \fi}

\begin{document}
% Marge interne du filigrane, évaluée dans la police qui ENTOURE le dessin :
% c'est là que TikZ évalue le .3333em de « inner sep » (cf. § 8).
\FiligraneSep=.3333em
% Largeur réelle des crédits : la plus longue ligne, dans leur police.
\sbox\CreditsBoite{\small\begin{tabular}{@{}l@{}}\CreditsTexte\end{tabular}}%
\begin{tikzpicture}[remember picture, overlay,
                    x=1mm, y=1mm, shift={(current page.south west)}]

% ---- Fonds ----------------------------------------------------------
% La 4e de couverture reste sur le blanc du papier.
% Le bandeau bleu couvre le dos, puis reprend au tiers extérieur de la
% 1re de couverture — même partition 2/3-1/3 que la page de titre.
% Couverture rigide : pas de filets au dos (ils tomberaient à côté des
% plis, vu la tolérance), et un dos bleu élargi de 2 mm de chaque côté.
\ifx\Reliure\MotRigide
  \fill[TitleBand] (\DosEtenduG, 0) rectangle (\DosEtenduD, \PapierH);
\else
  \fill[TitleBand] (\DosG, 0) rectangle (\DosD, \PapierH);
\fi
\fill[TitleBand] (\BandeauG, 0) rectangle (\PapierL, \PapierH);
% Filets d'accent : à la jonction sur la 1re, de part et d'autre du dos.
\fill[TitleAccent] (\BandeauAccent, 0) rectangle (\BandeauG, \PapierH);
\ifx\Reliure\MotRigide\else
  \fill[TitleAccent] (\DosAccentG, 0) rectangle (\DosG, \PapierH);
  \fill[TitleAccent] (\DosD, 0) rectangle (\DosAccentD, \PapierH);
\fi

% ---- 1re de couverture ----------------------------------------------
% Filigrane de classe, en réserve claire, détouré au bandeau.
\begin{scope}
\clip (\BandeauG, 0) rectangle (\PapierL, \PapierH);
\node[name=filigrane, anchor=north, align=center, text=white, opacity=0.16,
      font=\fontsize{\CorpsF}{\CorpsFI}\selectfont\bfseries]
  at (\FiligraneX, \FiligraneY) {\FiligraneTexte};
\end{scope}

\node[name=titre, anchor=west, align=left, text width=\TitreW mm]
  at (\TitreX, \TitreY) {%
    {\fontsize{54}{56}\selectfont\bfseries \TitreOuvrage}\\[10pt]
    {\Large\color{TitleBand}\bfseries \SousTitre}\\[16pt]
    {\large\color{black!75} \SousTitreClasse}};

\node[name=version, anchor=south west, align=left, font=\small, text=black!70]
  at (\TitreX, \VersionY) {Version \Version\\ Licence CC BY 4.0};

% ---- Dos -------------------------------------------------------------
% Un dos étroit ne reçoit pas de texte : en deçà de 6 mm, le débord au
% massicot ferait mordre le texte vertical sur les plats.
\ifnum\fpeval{\Dos >= 6}=1
  \node[name=dos, rotate=90, anchor=center, text=white, font=\bfseries,
        inner sep=0pt]
    at (\DosAxe, \MilieuH)
    {\TitreOuvrage\ \ --- \ \SousTitreClasse\ \ --- \ \Version};
\fi

% ---- 4e de couverture ------------------------------------------------
\node[name=quatre, anchor=north west, align=justify, text width=\QuatreW mm]
  at (\QuatreX, \QuatreY) {\QuatriemeTexte};

% Le texte vit dans \CreditsTexte (§ 8) pour être mesuré à part : la boîte
% du nœud a la largeur fixe \QuatreW et couvre l'emplacement ISBN sans
% qu'aucun caractère n'y soit. Retirer « text width » pour l'éviter change
% l'interlignage — essayé le 16/09/2026, 6 934 pixels différents.
\node[name=credits, anchor=south west, align=left, font=\small, text=black!70,
      text width=\QuatreW mm]
  at (\QuatreX, \QuatreBasY) {\CreditsTexte};

% ---- Repères de fabrication (mode épreuve uniquement) ----------------
\ifx\Reperes\ouimot
  % Traits de coupe, hors du fond perdu.
  \foreach \x in {\CoupeG, \CoupeD} {
    \draw[Repere, line width=0.25pt] (\x, \TechExt) -- (\x, \TechInt);
    \draw[Repere, line width=0.25pt] (\x, \HautInt) -- (\x, \HautExt);}
  \foreach \y in {\CoupeB, \CoupeH} {
    \draw[Repere, line width=0.25pt] (\TechExt, \y) -- (\TechInt, \y);
    \draw[Repere, line width=0.25pt] (\DroitInt, \y) -- (\DroitExt, \y);}
  % Limites du dos, en pointillés.
  \foreach \x in {\DosG, \DosD} {
    \draw[Repere, line width=0.25pt, dashed] (\x, \TechExt) -- (\x, \TechInt);
    \draw[Repere, line width=0.25pt, dashed] (\x, \HautInt) -- (\x, \HautExt);}
  % Zone de sécurité.
  \draw[Repere, line width=0.2pt, dotted]
    (\SecG, \SecB) rectangle (\SecD, \SecH);
  % Emplacement réservé au code-barres ISBN.
  \draw[Repere, line width=0.2pt] (\IsbnG, \SecB) rectangle (\IsbnD, \IsbnH);
  \node[font=\tiny\sffamily, text=Repere, align=center]
    at (\IsbnCX, \IsbnCY) {emplacement\\ISBN};
  % Cotes.
  \node[anchor=north, font=\tiny\sffamily, text=Repere]
    at (\DosAxe, \CoteHautY) {dos \Dos\ mm (\DosOrigine)};
  \node[anchor=south, font=\tiny\sffamily, text=Repere, align=center]
    at (\MilieuL, \CoteBasY) {\CotesTexte};
\fi

% ---- Contrôle de superposition (§ 8) ----------------------------------
% Rien n'est dessiné : des mesures et des lignes de journal.
\MesureNoeud{filigrane}{Fil}{\the\FiligraneSep}
\MesureNoeud{titre}{Tit}{0pt}
\MesureNoeud{version}{Ver}{0pt}
\MesureNoeud{quatre}{Qua}{0pt}
\MesureNoeud{credits}{Cre}{0pt}
% Bord droit des crédits = bord gauche + marge interne + plus longue ligne.
\xdef\Crexmax{\fpeval{round(\Crexmin + (\the\FiligraneSep + \the\wd\CreditsBoite)/1mm, 2)}}

% Le filigrane peut entrer dans la zone de sécurité (il est décoratif),
% pas franchir la coupe ni sortir du bandeau — le détourage le rognerait.
\ControleDans{filigrane dans le bandeau}{Fil}{\BandeauG}{\CoupeD}{\CoupeB}{\CoupeH}
\ControleDans{titre sur la 1re, hors bandeau}{Tit}{\PliPremiere}{\BandeauAccent}{\SecB}{\SecH}
\ControleDans{version sur la 1re, hors bandeau}{Ver}{\PliPremiere}{\BandeauAccent}{\SecB}{\SecH}
\ControleDans{texte de 4e dans la zone de securite}{Qua}{\SecG}{\PliQuatre}{\SecB}{\SecH}
\ControleDans{credits de 4e dans la zone de securite}{Cre}{\SecG}{\PliQuatre}{\SecB}{\SecH}
\ControleDisjoints{titre / version}{Tit}{\Verxmin}{\Verxmax}{\Verymin}{\Verymax}
\ControleDisjoints{texte de 4e / credits}{Qua}{\Crexmin}{\Crexmax}{\Creymin}{\Creymax}
\ControleDisjoints{texte de 4e / emplacement ISBN}{Qua}{\IsbnG}{\IsbnD}{\SecB}{\IsbnH}
\ControleDisjoints{credits / emplacement ISBN}{Cre}{\IsbnG}{\IsbnD}{\SecB}{\IsbnH}
\ifnum\fpeval{\Dos >= 6}=1
  \MesureNoeud{dos}{Dos}{0pt}
  \ControleDans{texte du dos entre les plis}{Dos}{\DosG}{\DosD}{\SecB}{\SecH}
\else
  \typeout{COUVERTURE-CONTROLE OK dos de \Dos\space mm : trop etroit, sans texte}%
\fi

\end{tikzpicture}
\typeout{COUVERTURE-BILAN alertes=\the\NbAlertes\space; classe=\Classe\space;
  format=\Format\space; reliure=\Reliure\space; pages=\NbPages\space;
  version=\Version\space;
  dos=\Dos\space mm (\DosOrigine); fini=\FiniL x \PageHauteur\space mm;
  papier=\PapierL x \PapierH\space mm}
\end{document}
"
%
%  Pour le fichier destiné à l'imprimeur (sans repères) :
%
%      lualatex "\def\Reperes{non}% =====================================================================
%  couverture.tex — couverture dos carré collé pour les livres 50ohm-fr
% =====================================================================
%
%  Produit une couverture d'un seul tenant : 4e de couverture | dos |
%  1re de couverture, aux dimensions exactes du livre relié, fond perdu
%  compris. L'épaisseur du dos est CALCULÉE à partir du nombre de pages,
%  du grammage et du bouffant du papier.
%
%  ORIGINE ET ÉTAT. Écrit le 20/08/2026 dans Fichierstravail/ (hors git),
%  versé au dépôt le 16/09/2026 pour que fenetre_compilation.py puisse
%  l'appeler. Ce qui a changé à cette occasion :
%
%    - le FORMAT est un paramètre (\Format = a4, 20x24 ou 20x24-marge,
%      mêmes noms que --format de build_book.py) ;
%    - la pagination et la version peuvent être IMPOSÉES de l'extérieur
%      (\PagesForcees, \VersionForcee) : la fenêtre lit la pagination dans
%      le journal du livre qu'elle vient de compiler, au lieu de se fier à
%      une table qui vieillit ;
%    - un dos imposé par l'imprimeur n'est appliqué QUE pour la pagination
%      et le format pour lesquels il a été donné ;
%    - un CONTRÔLE DE SUPERPOSITION est calculé par LaTeX lui-même, à partir
%      des boîtes réelles des nœuds, et écrit dans le journal ;
%    - le tome SWL est refusé explicitement : sa 4e de couverture n'est pas
%      rédigée (le texte commun parle des classes N, E et A).
%
%    Neutralité vérifiée : en A4, pour N (258 p., a.2) et NEA (816 p., a.3),
%    en épreuve comme en impression, le rendu est identique au pixel près à
%    celui du gabarit d'origine.
%
%  COMPILATION (LuaLaTeX, depuis la racine du dépôt) :
%
%      lualatex couverture.tex
%      lualatex couverture.tex        <-- DEUX FOIS, obligatoirement
%
%  DEUX PASSES SONT NÉCESSAIRES. Le dessin s'ancre sur « current page »,
%  dont la position n'est connue qu'après un premier passage : à la
%  première compilation, tout se tasse dans le coin supérieur et la
%  couverture paraît ratée alors qu'elle ne l'est pas. Aucune erreur n'est
%  émise — seul le rendu le montre. Le contrôle de superposition, lui
%  aussi, ne vaut qu'à la seconde passe.
%
%  Le plus simple est de passer par fenetre_compilation.py (case
%  « couverture pour l'imprimeur »), qui enchaîne les deux passes, produit
%  l'épreuve ET le fichier d'impression, et relit le contrôle.
%
%  À la main, pour une autre classe, un autre format, une pagination mesurée :
%
%      lualatex "\def\Classe{A}\def\Format{20x24}\def\PagesForcees{482}\input{couverture.tex}"
%
%  Pour le fichier destiné à l'imprimeur (sans repères) :
%
%      lualatex "\def\Reperes{non}\input{couverture.tex}"
%
%  ---------------------------------------------------------------------
%  DEUX MODES, un seul paramètre : \Reperes
%
%    oui  — mode ÉPREUVE. Ajoute une marge technique, les traits de coupe,
%           les limites du dos, la zone de sécurité et les cotes. C'est ce
%           qu'il faut pour vérifier à l'écran.
%    non  — mode PRODUCTION. Papier = format final + fond perdu, rien
%           d'autre. C'est ce fichier-là qu'on envoie à l'imprimeur.
%
%  ---------------------------------------------------------------------
%  CONTRÔLE DE SUPERPOSITION
%
%  À la fin du dessin, chaque bloc de texte est mesuré par ses ancres TikZ
%  (quatre coins, ce qui vaut aussi pour le texte pivoté du dos) et comparé
%  à la zone où il doit tenir, puis aux blocs voisins. Le journal porte une
%  ligne par vérification :
%
%      COUVERTURE-CONTROLE OK      <quoi> : ...
%      COUVERTURE-CONTROLE ALERTE  <quoi> : ...
%
%  et un bilan : COUVERTURE-BILAN alertes=<n> ; ...
%
%  Ce que le contrôle VOIT : un bloc qui franchit le pli du dos, la coupe,
%  la zone de sécurité ou le bord du bandeau ; deux blocs qui se chevauchent ;
%  un bloc qui empiète sur l'emplacement réservé au code-barres ISBN.
%  Ce qu'il NE VOIT PAS : une ligne trop longue à l'INTÉRIEUR d'un bloc à
%  largeur fixe (la boîte garde sa largeur, le texte déborde). Celle-là se
%  signale par un « Overfull \hbox » dans le même journal, que la fenêtre
%  compte aussi.
%
%  ---------------------------------------------------------------------
%  AVERTISSEMENT SUR L'ÉPAISSEUR DU DOS
%
%  Le dos se calcule ainsi :
%
%      épaisseur d'une feuille (mm) = grammage x bouffant / 1000
%      dos (mm)                     = (nombre de pages / 2) x épaisseur
%
%  On divise par deux parce qu'une feuille porte DEUX pages.
%
%  Le « bouffant » (main, volume, bulk) dépend du papier et n'est pas
%  devinable : un offset 80 g courant est autour de 1,30, un bouffant
%  d'édition monte à 1,75 et plus. UN ÉCART DE BOUFFANT DÉPLACE LE DOS DE
%  PLUSIEURS MILLIMÈTRES. Demander à l'imprimeur soit son bouffant, soit
%  directement l'épaisseur du dos, et reporter ci-dessous. Tant que cette
%  valeur n'est pas confirmée, la couverture n'est pas bonne à tirer.
%
%  Si l'imprimeur donne le dos en clair, forcer la valeur avec \DosImpose.
% =====================================================================

\documentclass[11pt]{article}

% ---------------------------------------------------------------------
% 1. PARAMÈTRES — c'est la seule zone à modifier
% ---------------------------------------------------------------------

% Classe de l'ouvrage : N, E, A, NE, EA ou NEA.
% (\providecommand : une valeur passée en ligne de commande l'emporte.)
\providecommand{\Classe}{NEA}

% Mode épreuve (oui) ou production (non).
\providecommand{\Reperes}{oui}

% Format du livre FINI : a4, 20x24 ou 20x24-marge — les noms de --format
% dans build_book.py. Les deux maquettes 20 x 24 ne diffèrent que par leur
% colonne de marge intérieure : le livre fini, donc la couverture, est le
% même (200 x 240 mm).
\providecommand{\Format}{a4}

% Reliure : « souple » (dos carré collé, fond perdu 3 mm) ou « rigide »
% (couverture rigide, dite hardcover). En rigide, consignes de l'imprimeur
% du 17/09/2026 : chaque plat déborde du livre de 5 mm en largeur et en
% hauteur (le carton dépasse du bloc), et le fichier porte 15 mm de rabat
% autour du carton au lieu de 3 mm de fond perdu. Les aplats vont déjà au
% bord du papier : ils remplissent le rabat d'eux-mêmes.
\providecommand{\Reliure}{souple}

% Pagination et version imposées de l'extérieur. Vides = tables du § 2.
% La fenêtre les renseigne toujours : la pagination vient du journal du
% livre réellement compilé, dans le format demandé.
\providecommand{\PagesForcees}{}
\providecommand{\VersionForcee}{}

% Papier de l'INTÉRIEUR du livre (pas de la couverture).
% CoolLibri, papier « Standard 90 g blanc ».
\def\Grammage{90}        % g/m²
\def\Bouffant{1.206}     % main DÉDUITE, voir ci-dessous

% Épaisseur du dos IMPOSÉE par l'imprimeur, en mm, par classe.
% Une valeur > 0 prime sur le calcul — MAIS SEULEMENT pour la pagination et
% le format pour lesquels l'imprimeur l'a donnée (\DosImposePages<classe>,
% \DosImposeFormat<classe>). Hors de ce cas, le dos est calculé et le
% journal le dit : un dos de 14 mm donné pour 258 pages A4 ne vaut rien pour
% 262 pages, ni pour le même livre recomposé en 20 x 24.
%
% ATTENTION : cette valeur est propre à CHAQUE ÉDITION. Un dos unique
% appliqué aux quatre livres serait faux pour trois d'entre eux.
%
%  N   — 14 mm, VALEUR DONNÉE PAR COOLIBRI pour 258 pages en 90 g, A4.
%        Fichier attendu : 440 x 303 mm, ce que ce gabarit produit en mode
%        production (2 x 210 + 14 + 2 x 3 = 440 ; 297 + 2 x 3 = 303).
%        Le N de la a.3 compte 262 pages : cette valeur ne s'y applique plus.
%  E, A, NEA — non demandés à l'imprimeur : ils sont CALCULÉS avec la main
%        déduite ci-dessous. À faire confirmer avant tout tirage.
%
% La main de 1,206 n'est pas publiée par CoolLibri : elle est déduite de
% leur propre chiffre, 14 mm / (0,090 x 258/2) = 1,206. Elle reproduit donc
% exactement le dos de N, et n'est qu'une ESTIMATION pour les autres.
% Écarts que cela recouvre, pour mémoire : une main de 1,0 donnerait
% 11,61 mm pour N, une main de 1,3 en donnerait 16,77. Il ne fallait pas
% deviner.
%  E   — 17 mm, VALEUR DONNÉE PAR L'IMPRIMEUR le 17/09/2026 pour 276 pages,
%        20 x 24, couverture RIGIDE, reliure collée (Klebebindung), papier
%        115 g. En reliure cousue (Fadenheftung), il annonce 18 mm.
%
% \DosImposeReliure<classe> : la reliure pour laquelle la valeur a été
% donnée. Le dos d'une couverture rigide n'est pas celui d'un dos carré
% collé souple, même à pagination égale.
\def\DosImposeN{14}   \def\DosImposePagesN{258}  \def\DosImposeFormatN{a4}
\def\DosImposeE{17}   \def\DosImposePagesE{276}  \def\DosImposeFormatE{20x24-marge}
\def\DosImposeA{0}    \def\DosImposePagesA{0}    \def\DosImposeFormatA{}
\def\DosImposeNE{0}   \def\DosImposePagesNE{0}   \def\DosImposeFormatNE{}
\def\DosImposeEA{0}   \def\DosImposePagesEA{0}   \def\DosImposeFormatEA{}
\def\DosImposeNEA{0}  \def\DosImposePagesNEA{0}  \def\DosImposeFormatNEA{}
\def\DosImposeReliureN{souple}  \def\DosImposeReliureE{rigide}
\def\DosImposeReliureA{souple}  \def\DosImposeReliureNE{souple}
\def\DosImposeReliureEA{souple} \def\DosImposeReliureNEA{souple}

% Fond perdu, en mm. 3 mm est la valeur courante ; certains imprimeurs
% demandent 5 mm.
\def\FondPerdu{3}

% Marge technique autour du fond perdu, en mm. Sert à loger les traits de
% coupe en mode épreuve. Ignorée en mode production.
\def\MargeTech{12}

% Zone de sécurité, en mm : aucun texte ne doit s'en approcher davantage
% du bord de coupe.
\def\Securite{8}

% ---------------------------------------------------------------------
% 2. DONNÉES PAR CLASSE
%    pages : nombre de pages du PDF intérieur, COUVERTURE NON COMPRISE.
%    Un dos carré collé impose un nombre de pages PAIR ; le fichier
%    refuse de compiler si le compte est impair.
%
%    Ces tables ne servent que si \PagesForcees et \VersionForcee sont
%    vides, c'est-à-dire à la main. Valeurs A4 mesurées sur les PDF de la
%    a.3 (CLAUDE.md § 1, 05/09/2026). En 20 x 24, la pagination change
%    entièrement : NEA passe de 816 à 1036 pages.
% ---------------------------------------------------------------------

\def\DonneesN{262}
\def\DonneesE{214}
\def\DonneesA{386}
\def\DonneesNE{0}      % édition jamais produite
\def\DonneesEA{0}      % édition jamais produite
\def\DonneesNEA{816}

% Version imprimée sur le dos et la 1re de couverture, PAR CLASSE.
% « Les versions sont suivies par classe, chacune évoluant à son rythme »
% (CHANGELOG.md).
\def\VersionN{a.3}
\def\VersionE{a.3}
\def\VersionA{a.3}
\def\VersionNE{a.3}
\def\VersionEA{a.3}
\def\VersionNEA{a.3}

% ---------------------------------------------------------------------
% 3. TEXTES
% ---------------------------------------------------------------------

\def\TitreOuvrage{50\,Ohm}
\def\SousTitre{Préparation à l'examen radioamateur}

% Sous-titre propre à chaque classe (repris de FR_TITLES, build_book.py).
\def\SousTitreN{Cours complet — classe N}
\def\SousTitreE{Cours complet — classe E}
\def\SousTitreA{Cours complet — classe A}
\def\SousTitreNE{Cours complet — classes N et E}
\def\SousTitreEA{Cours complet — classes E et A}
\def\SousTitreNEA{Cours complet — classes N, E et A}

% Lettres du filigrane, empilées de haut en bas (cf. build_book.py v0.20).
\def\LettresN{N}
\def\LettresE{E}
\def\LettresA{A}
\def\LettresNE{N,E}
\def\LettresEA{E,A}
\def\LettresNEA{N,E,A}

% Texte de 4e de couverture.
% NB : il parle des classes N, E et A quelle que soit l'édition. Juste pour
% N, E, A et NEA ; approximatif pour NE et EA ; faux pour SWL, d'où son refus.
\def\QuatriemeTexte{%
Ce guide d'étude prépare aux examens radioamateur allemands des classes
N, E et A. Il traduit en français les contenus du projet 50ohm.de, publiés
par le référat AJW du DARC e.\,V., et les complète d'encarts « En France »
qui signalent, sans rien retirer au texte allemand, ce que la
réglementation française prévoit sur le même sujet.\par\smallskip
Les questions officielles du catalogue allemand accompagnent chaque
notion, avec leurs réponses.%
}

\def\Auteur{Traduit avec l'aide d'une IA par Pierre F4JWI}
\def\Depot{github.com/DrPLG/50ohm-fr}

% =====================================================================
%  À PARTIR D'ICI, PLUS RIEN À MODIFIER
% =====================================================================

\usepackage{iftex}
\RequireLuaTeX
\usepackage{xfp}
\usepackage{tikz}
\usetikzlibrary{calc}
\usepackage{fontspec}
\usepackage[french]{babel}
\usepackage{geometry}

% ---------------------------------------------------------------------
% 4. RÉSOLUTION DES DONNÉES DE LA CLASSE
% ---------------------------------------------------------------------

\makeatletter
% Tome SWL : refus explicite, plutôt qu'un « classe inconnue » qui laisserait
% croire à un oubli. Le texte de 4e parle des classes N, E et A.
\def\@swlmot{SWL}
\ifx\Classe\@swlmot
  \PackageError{couverture}{Pas de couverture pour le tome SWL}%
    {Sa 4e de couverture n'est pas redigee : le texte commun parle des
     classes N, E et A. C'est une decision editoriale a prendre d'abord.}%
\fi

\def\@resoudre#1#2{%
  \expandafter\let\expandafter#1\csname #2\Classe\endcsname
  \ifx#1\relax
    \PackageError{couverture}{Classe \Classe\space inconnue : pas de #2\Classe}%
      {Valeurs admises : N, E, A, NE, EA, NEA.}%
  \fi}
\@resoudre\NbPages{Donnees}
\@resoudre\SousTitreClasse{SousTitre}
\@resoudre\Lettres{Lettres}
\@resoudre\DosImpose{DosImpose}
\@resoudre\DosImposePages{DosImposePages}
\@resoudre\DosImposeFormat{DosImposeFormat}
\@resoudre\DosImposeReliure{DosImposeReliure}
\@resoudre\Version{Version}

% Valeurs imposées de l'extérieur. Test de vacuité par \detokenize : il ne
% dépend pas de la façon dont la macro a été définie.
\if\relax\detokenize\expandafter{\PagesForcees}\relax\else
  \let\NbPages\PagesForcees
\fi
\if\relax\detokenize\expandafter{\VersionForcee}\relax\else
  \let\Version\VersionForcee
\fi

% Format du livre fini.
\def\@formatAquatre{a4}
\def\@formatVingt{20x24}
\def\@formatVingtMarge{20x24-marge}
\ifx\Format\@formatAquatre
  \def\PageLargeur{210}\def\PageHauteur{297}
\else\ifx\Format\@formatVingt
  \def\PageLargeur{200}\def\PageHauteur{240}
\else\ifx\Format\@formatVingtMarge
  \def\PageLargeur{200}\def\PageHauteur{240}
\else
  \PackageError{couverture}{Format \Format\space inconnu}%
    {Valeurs admises : a4, 20x24, 20x24-marge.}%
  \def\PageLargeur{210}\def\PageHauteur{297}
\fi\fi\fi
% Reliure. Hors de « souple » et « rigide » : erreur, pas de supposition.
\def\@reliuresouple{souple}
\def\@reliurerigide{rigide}
\ifx\Reliure\@reliurerigide
  \edef\PageLargeur{\fpeval{\PageLargeur + 5}}
  \edef\PageHauteur{\fpeval{\PageHauteur + 5}}
  \def\FondPerdu{15}
\else\ifx\Reliure\@reliuresouple\else
  \PackageError{couverture}{Reliure \Reliure\space inconnue}%
    {Valeurs admises : souple, rigide.}%
\fi\fi
\makeatother

% Garde-fou : une édition jamais produite n'a pas de couverture.
\ifnum\NbPages=0
  \PackageError{couverture}{L'edition \Classe\space n'a pas de pagination}%
    {Cette edition n'a jamais ete compilee. Renseigner \string\Donnees\Classe
     ou passer \string\PagesForcees.}%
\fi

% Avertissement : l'imprimeur demande un nombre de pages multiple de 4
% (17/09/2026). Pas une erreur : c'est son exigence, pas une loi du papier.
\ifnum\fpeval{\NbPages - 4*floor(\NbPages/4)}=0 \else
  \typeout{COUVERTURE-AVERTISSEMENT \NbPages\space pages : pas un multiple
    de 4, que l'imprimeur exige -- ajouter des pages a l'interieur}%
\fi

% Garde-fou : un dos carré collé impose un nombre de pages pair.
\ifodd\NbPages
  \PackageError{couverture}{Nombre de pages impair (\NbPages)}%
    {Un dos carre colle exige un nombre PAIR de pages : une feuille en
     porte deux. Ajouter une page blanche a l'interieur.}%
\fi

% ---------------------------------------------------------------------
% 5. GÉOMÉTRIE — c'est ici que se joue l'exactitude du dos
% ---------------------------------------------------------------------

% Épaisseur du dos : valeur imposée si elle est donnée ET qu'elle vaut pour
% cette pagination et ce format ; calcul sinon.
\edef\DosCalcule{\fpeval{round(\NbPages / 2 * \Grammage * \Bouffant / 1000, 2)}}
\def\DosOrigine{calculee}
\ifnum\fpeval{\DosImpose > 0}=1
  \ifnum\fpeval{\NbPages = \DosImposePages}=1
    \ifx\Format\DosImposeFormat
      \ifx\Reliure\DosImposeReliure
        \def\DosOrigine{imposee}
      \fi
    \fi
  \fi
  \def\MotCalculee{calculee}
  \ifx\DosOrigine\MotCalculee
    \typeout{COUVERTURE-AVERTISSEMENT dos impose de \DosImpose\space mm
      donne pour \DosImposePages\space pages en \DosImposeFormat\space
      (reliure \DosImposeReliure) -- inapplicable a \NbPages\space pages en
      \Format\space (reliure \Reliure), dos calcule}%
  \fi
\fi
% NB : pas de \@ dans ces noms. Cette section est hors \makeatletter, et
% « \def\@imposee{...} » y redéfinit \@ (suivi du texte « imposee ») : le
% test échouait toujours, en silence. Constaté le 16/09/2026 — le dos de N
% tombait juste par coïncidence, 14,0 mm calculés pour 14 imposés.
\def\MotRigide{rigide}
\def\MotImposee{imposee}
\ifx\DosOrigine\MotImposee
  \edef\Dos{\DosImpose}
\else
  \edef\Dos{\DosCalcule}
\fi

% Marge technique : uniquement en mode épreuve.
\def\ouimot{oui}
\ifx\Reperes\ouimot \edef\Tech{\MargeTech}\else \def\Tech{0}\fi

% Dimensions du papier.
%   largeur = marge + fond perdu + 4e + dos + 1re + fond perdu + marge
%   hauteur = marge + fond perdu + hauteur + fond perdu + marge
\edef\PapierL{\fpeval{2*\PageLargeur + \Dos + 2*\FondPerdu + 2*\Tech}}
\edef\PapierH{\fpeval{\PageHauteur + 2*\FondPerdu + 2*\Tech}}

\geometry{paperwidth=\PapierL mm, paperheight=\PapierH mm,
          margin=0pt, nohead, nofoot}

% ---------------------------------------------------------------------
%  TOUTES LES COTES SONT FIGÉES ICI, ET NULLE PART AILLEURS.
%
%  Raison, payée par trois compilations ratées : un \fpeval laissé dans
%  une coordonnée ou dans le texte d'un nœud TikZ n'est pas développé au
%  moment où TikZ analyse le chemin. On obtient « Giving up on this path.
%  Did you forget a semicolon? », qui accuse une ponctuation alors que la
%  faute est ailleurs. Le corps du document ci-dessous n'emploie donc que
%  des macros déjà calculées.
% ---------------------------------------------------------------------

% Repères horizontaux, mesurés depuis le bord GAUCHE du papier.
\edef\CoupeG{\fpeval{\Tech + \FondPerdu}}          % coupe gauche (4e de couv.)
\edef\DosG{\fpeval{\CoupeG + \PageLargeur}}         % dos, début
\edef\DosD{\fpeval{\DosG + \Dos}}                   % dos, fin
\edef\CoupeD{\fpeval{\DosD + \PageLargeur}}         % coupe droite (1re de couv.)
% Repères verticaux, depuis le bord BAS.
\edef\CoupeB{\fpeval{\Tech + \FondPerdu}}
\edef\CoupeH{\fpeval{\CoupeB + \PageHauteur}}
% Axes.
\edef\DosAxe{\fpeval{(\DosG + \DosD)/2}}
\edef\MilieuH{\fpeval{(\CoupeB + \CoupeH)/2}}
\edef\MilieuL{\fpeval{\PapierL/2}}

% Format fini, pour les cotes affichées.
\edef\FiniL{\fpeval{2*\PageLargeur + \Dos}}

% Bandeau de la 1re de couverture : même proportion que la page de titre.
\edef\BandeauL{\fpeval{round(0.34*\PageLargeur, 2)}}
\edef\BandeauG{\fpeval{\CoupeD - \BandeauL}}
\edef\BandeauAccent{\fpeval{\BandeauG - 0.7}}
\edef\DosAccentG{\fpeval{\DosG - 0.7}}
\edef\DosAccentD{\fpeval{\DosD + 0.7}}
% Couverture rigide : le bleu du dos déborde de 2 mm sur chaque plat, pour
% absorber la tolérance de l'imprimeur sur l'épaisseur du dos (±10 %).
\edef\DosEtenduG{\fpeval{\DosG - 2}}
\edef\DosEtenduD{\fpeval{\DosD + 2}}

% Filigrane de classe.
\edef\FiligraneX{\fpeval{(\BandeauG + \CoupeD)/2}}
\edef\FiligraneY{\fpeval{\CoupeH - 6}}

% Bloc-titre de la 1re de couverture.
\edef\TitreX{\fpeval{\DosD + 12}}
\edef\TitreY{\fpeval{\CoupeH - 72}}
\edef\TitreW{\fpeval{\BandeauG - \DosD - 16}}
\edef\VersionY{\fpeval{\CoupeB + 14}}

% 4e de couverture.
\edef\QuatreX{\fpeval{\CoupeG + \Securite + 6}}
\edef\QuatreY{\fpeval{\CoupeH - \Securite - 14}}
\edef\QuatreBasY{\fpeval{\CoupeB + \Securite + 22}}
\edef\QuatreW{\fpeval{\PageLargeur - 2*\Securite - 12}}

% Zone de sécurité.
\edef\SecG{\fpeval{\CoupeG + \Securite}}
\edef\SecD{\fpeval{\CoupeD - \Securite}}
\edef\SecB{\fpeval{\CoupeB + \Securite}}
\edef\SecH{\fpeval{\CoupeH - \Securite}}

% Zones de contrôle : ce qui doit rester de part et d'autre du pli du dos.
\edef\PliQuatre{\fpeval{\DosG - \Securite}}         % 4e : limite côté dos
\edef\PliPremiere{\fpeval{\DosD + \Securite}}       % 1re : limite côté dos

% Traits de coupe et cotes, dans la marge technique.
\edef\TechExt{\fpeval{\Tech - 8}}
\edef\TechInt{\fpeval{\Tech - 1}}
\edef\TechCote{\fpeval{\Tech - 9}}
\edef\HautInt{\fpeval{\PapierH - \Tech + 1}}
\edef\HautExt{\fpeval{\PapierH - \Tech + 8}}
\edef\HautCote{\fpeval{\PapierH - \Tech + 9}}
\edef\DroitInt{\fpeval{\PapierL - \Tech + 1}}
\edef\DroitExt{\fpeval{\PapierL - \Tech + 8}}

% Emplacement du code-barres ISBN : en bas à DROITE de la 4e de couverture,
% selon l'usage éditorial, et non à gauche où il chevaucherait les mentions
% de licence.
\edef\IsbnD{\fpeval{\DosG - \Securite}}
\edef\IsbnG{\fpeval{\IsbnD - 40}}
\edef\IsbnH{\fpeval{\SecB + 25}}
\edef\IsbnCX{\fpeval{\IsbnG + 20}}
\edef\IsbnCY{\fpeval{\SecB + 12.5}}

% Cartouches de cotes, dans la marge technique. Ancrés vers l'INTÉRIEUR du
% papier : ancrés vers l'extérieur, ils débordaient et étaient rognés.
\edef\CoteBasY{\fpeval{1}}
\edef\CoteHautY{\fpeval{\PapierH - 1}}

% ---------------------------------------------------------------------
% 6. COULEURS ET POLICES — reprises de la page de titre (build_book.py)
% ---------------------------------------------------------------------

\definecolor{TitleBand}{cmyk}{.9,.55,.1,.35}   % bleu profond
\definecolor{TitleAccent}{cmyk}{.05,.8,1,0}    % rouge DARC (accent fin)
\definecolor{Repere}{cmyk}{0,1,1,0}            % magenta : repères de montage

\setmainfont{Libertinus Serif}[
  ItalicFont     = Libertinus Serif Italic,
  BoldFont       = Libertinus Serif Bold,
  BoldItalicFont = Libertinus Serif Bold Italic]

\pagestyle{empty}

% ---------------------------------------------------------------------
% 7. FILIGRANE DE CLASSE
%    Le texte est construit ici, et non dans le nœud : un \\ produit à
%    l'intérieur d'un \foreach referme le groupe d'alignement de TikZ.
%    Une lettre seule est grande ; plusieurs s'empilent et se réduisent
%    (build_book.py v0.20 : « NEA » à 105 pt débordait encore).
% ---------------------------------------------------------------------

\makeatletter
\def\FiligraneTexte{}
\newif\ifpremierelettre \premierelettretrue
\count255=0
\@for\@tempa:=\Lettres\do{%
  \advance\count255 by 1\relax
  \ifpremierelettre \premierelettrefalse
  \else \g@addto@macro\FiligraneTexte{\\}\fi
  \expandafter\g@addto@macro\expandafter\FiligraneTexte\expandafter{\@tempa}}%
\xdef\NbLettres{\the\count255}
\makeatother

\ifnum\NbLettres=1 \def\CorpsF{150}\else\def\CorpsF{86}\fi
\edef\CorpsFI{\fpeval{round(\CorpsF*1.02, 0)}}

% Texte du cartouche de cotes, lui aussi construit hors du nœud TikZ.
%
% NE PAS y employer \textsuperscript : dans un nœud « align=center », il
% casse le chemin TikZ (« Giving up on this path. Did you forget a
% semicolon? »), et l'erreur est signalée sur le nœud SUIVANT — elle
% accuse donc une ponctuation là où la faute est ici. Isolé par essais :
% $^2$ et le caractère ² passent, \textsuperscript non.
\def\CotesTexte{%
  \Classe\ --- \NbPages\ pages --- papier \Grammage\ g/m$^2$,
  bouffant \Bouffant\ --- fond perdu \FondPerdu\ mm\\
  papier \PapierL\ mm sur \PapierH\ mm --- fini \FiniL\ mm sur
  \PageHauteur\ mm\\
  EPREUVE, ne pas envoyer a l'impression%
}

% ---------------------------------------------------------------------
% 8. CONTRÔLE DE SUPERPOSITION
%
%  \MesureNoeud{nœud}{préfixe}{retrait} mesure la boîte d'un nœud par ses
%  quatre coins (un nœud pivoté a ses ancres pivotées : le min et le max des
%  quatre suffisent) et la ramène au repère du papier, en mm. Le point (0,0)
%  est lu par le même « let » que les ancres : les deux passent donc par la
%  même transformation, et la comparaison ne dépend pas de l'endroit où
%  TeX a posé l'image sur la page.
%
%  Le retrait sert au seul filigrane, pour mesurer les lettres et non leur
%  cadre : c'est le bord des lettres que le détourage rognerait.
%
%  PIÈGE MESURÉ le 16/09/2026 : la marge interne TikZ (.3333em) n'est PAS
%  évaluée dans la police du nœud, mais dans celle qui entoure le dessin.
%  Instrumenté : dans le nœud, em = 150 pt ; boîte brute du « N » :
%  41,7 x 37,5 mm, soit la lettre plus 2 x 1,3 mm — 1,3 mm étant .3333em
%  à 11 pt. Un retrait calculé à 150 pt (17,6 mm) ramenait la lettre à
%  6,6 mm de large, et le contrôle aurait dit OK sur une mesure absurde.
% ---------------------------------------------------------------------

\newcount\NbAlertes
\newdimen\FiligraneSep
\newsavebox\CreditsBoite

% Crédits de 4e : lignes coupées à la main. Leur largeur RÉELLE est celle
% de la plus longue ligne, mesurée dans une boîte à part (voir § 8, ISBN).
\def\CreditsTexte{%
    Réalisé à partir des contenus de 50ohm.de (en allemand)\\
    50ohm.de-Autorenteam, coordonné par le référat AJW du DARC e.\,V.\\
    \Auteur\\[3pt]
    Licence CC BY 4.0 --- \texttt{\Depot}}

\def\MesureNoeud#1#2#3{%
  \path let \p1=(#1.south west), \p2=(#1.north east),
            \p3=(#1.north west), \p4=(#1.south east), \p5=(0,0) in
    \pgfextra{%
      \expandafter\xdef\csname #2xmin\endcsname{\fpeval{round((min(\x1,\x2,\x3,\x4) - \x5 + #3)/1mm, 2)}}%
      \expandafter\xdef\csname #2xmax\endcsname{\fpeval{round((max(\x1,\x2,\x3,\x4) - \x5 - #3)/1mm, 2)}}%
      \expandafter\xdef\csname #2ymin\endcsname{\fpeval{round((min(\y1,\y2,\y3,\y4) - \y5 + #3)/1mm, 2)}}%
      \expandafter\xdef\csname #2ymax\endcsname{\fpeval{round((max(\y1,\y2,\y3,\y4) - \y5 - #3)/1mm, 2)}}%
    };}

% \ControleDans{libellé}{préfixe}{xmin}{xmax}{ymin}{ymax}
\def\ControleDans#1#2#3#4#5#6{%
  \ifnum\fpeval{\csname #2xmin\endcsname >= #3 && \csname #2xmax\endcsname <= #4
             && \csname #2ymin\endcsname >= #5 && \csname #2ymax\endcsname <= #6}=1
    \typeout{COUVERTURE-CONTROLE OK #1 : x \csname #2xmin\endcsname-\csname #2xmax\endcsname,
      y \csname #2ymin\endcsname-\csname #2ymax\endcsname\space dans x #3-#4, y #5-#6}%
  \else
    \global\advance\NbAlertes by 1
    \typeout{COUVERTURE-CONTROLE ALERTE #1 : x \csname #2xmin\endcsname-\csname #2xmax\endcsname,
      y \csname #2ymin\endcsname-\csname #2ymax\endcsname\space HORS de x #3-#4, y #5-#6}%
  \fi}

% \ControleDisjoints{libellé}{préfixe}{xmin}{xmax}{ymin}{ymax}
% Deux rectangles se chevauchent si et seulement s'ils se recouvrent sur
% les DEUX axes à la fois.
\def\ControleDisjoints#1#2#3#4#5#6{%
  \ifnum\fpeval{\csname #2xmin\endcsname < #4 && #3 < \csname #2xmax\endcsname
             && \csname #2ymin\endcsname < #6 && #5 < \csname #2ymax\endcsname}=1
    \global\advance\NbAlertes by 1
    \typeout{COUVERTURE-CONTROLE ALERTE #1 : chevauchement}%
  \else
    \typeout{COUVERTURE-CONTROLE OK #1 : disjoints}%
  \fi}

\begin{document}
% Marge interne du filigrane, évaluée dans la police qui ENTOURE le dessin :
% c'est là que TikZ évalue le .3333em de « inner sep » (cf. § 8).
\FiligraneSep=.3333em
% Largeur réelle des crédits : la plus longue ligne, dans leur police.
\sbox\CreditsBoite{\small\begin{tabular}{@{}l@{}}\CreditsTexte\end{tabular}}%
\begin{tikzpicture}[remember picture, overlay,
                    x=1mm, y=1mm, shift={(current page.south west)}]

% ---- Fonds ----------------------------------------------------------
% La 4e de couverture reste sur le blanc du papier.
% Le bandeau bleu couvre le dos, puis reprend au tiers extérieur de la
% 1re de couverture — même partition 2/3-1/3 que la page de titre.
% Couverture rigide : pas de filets au dos (ils tomberaient à côté des
% plis, vu la tolérance), et un dos bleu élargi de 2 mm de chaque côté.
\ifx\Reliure\MotRigide
  \fill[TitleBand] (\DosEtenduG, 0) rectangle (\DosEtenduD, \PapierH);
\else
  \fill[TitleBand] (\DosG, 0) rectangle (\DosD, \PapierH);
\fi
\fill[TitleBand] (\BandeauG, 0) rectangle (\PapierL, \PapierH);
% Filets d'accent : à la jonction sur la 1re, de part et d'autre du dos.
\fill[TitleAccent] (\BandeauAccent, 0) rectangle (\BandeauG, \PapierH);
\ifx\Reliure\MotRigide\else
  \fill[TitleAccent] (\DosAccentG, 0) rectangle (\DosG, \PapierH);
  \fill[TitleAccent] (\DosD, 0) rectangle (\DosAccentD, \PapierH);
\fi

% ---- 1re de couverture ----------------------------------------------
% Filigrane de classe, en réserve claire, détouré au bandeau.
\begin{scope}
\clip (\BandeauG, 0) rectangle (\PapierL, \PapierH);
\node[name=filigrane, anchor=north, align=center, text=white, opacity=0.16,
      font=\fontsize{\CorpsF}{\CorpsFI}\selectfont\bfseries]
  at (\FiligraneX, \FiligraneY) {\FiligraneTexte};
\end{scope}

\node[name=titre, anchor=west, align=left, text width=\TitreW mm]
  at (\TitreX, \TitreY) {%
    {\fontsize{54}{56}\selectfont\bfseries \TitreOuvrage}\\[10pt]
    {\Large\color{TitleBand}\bfseries \SousTitre}\\[16pt]
    {\large\color{black!75} \SousTitreClasse}};

\node[name=version, anchor=south west, align=left, font=\small, text=black!70]
  at (\TitreX, \VersionY) {Version \Version\\ Licence CC BY 4.0};

% ---- Dos -------------------------------------------------------------
% Un dos étroit ne reçoit pas de texte : en deçà de 6 mm, le débord au
% massicot ferait mordre le texte vertical sur les plats.
\ifnum\fpeval{\Dos >= 6}=1
  \node[name=dos, rotate=90, anchor=center, text=white, font=\bfseries,
        inner sep=0pt]
    at (\DosAxe, \MilieuH)
    {\TitreOuvrage\ \ --- \ \SousTitreClasse\ \ --- \ \Version};
\fi

% ---- 4e de couverture ------------------------------------------------
\node[name=quatre, anchor=north west, align=justify, text width=\QuatreW mm]
  at (\QuatreX, \QuatreY) {\QuatriemeTexte};

% Le texte vit dans \CreditsTexte (§ 8) pour être mesuré à part : la boîte
% du nœud a la largeur fixe \QuatreW et couvre l'emplacement ISBN sans
% qu'aucun caractère n'y soit. Retirer « text width » pour l'éviter change
% l'interlignage — essayé le 16/09/2026, 6 934 pixels différents.
\node[name=credits, anchor=south west, align=left, font=\small, text=black!70,
      text width=\QuatreW mm]
  at (\QuatreX, \QuatreBasY) {\CreditsTexte};

% ---- Repères de fabrication (mode épreuve uniquement) ----------------
\ifx\Reperes\ouimot
  % Traits de coupe, hors du fond perdu.
  \foreach \x in {\CoupeG, \CoupeD} {
    \draw[Repere, line width=0.25pt] (\x, \TechExt) -- (\x, \TechInt);
    \draw[Repere, line width=0.25pt] (\x, \HautInt) -- (\x, \HautExt);}
  \foreach \y in {\CoupeB, \CoupeH} {
    \draw[Repere, line width=0.25pt] (\TechExt, \y) -- (\TechInt, \y);
    \draw[Repere, line width=0.25pt] (\DroitInt, \y) -- (\DroitExt, \y);}
  % Limites du dos, en pointillés.
  \foreach \x in {\DosG, \DosD} {
    \draw[Repere, line width=0.25pt, dashed] (\x, \TechExt) -- (\x, \TechInt);
    \draw[Repere, line width=0.25pt, dashed] (\x, \HautInt) -- (\x, \HautExt);}
  % Zone de sécurité.
  \draw[Repere, line width=0.2pt, dotted]
    (\SecG, \SecB) rectangle (\SecD, \SecH);
  % Emplacement réservé au code-barres ISBN.
  \draw[Repere, line width=0.2pt] (\IsbnG, \SecB) rectangle (\IsbnD, \IsbnH);
  \node[font=\tiny\sffamily, text=Repere, align=center]
    at (\IsbnCX, \IsbnCY) {emplacement\\ISBN};
  % Cotes.
  \node[anchor=north, font=\tiny\sffamily, text=Repere]
    at (\DosAxe, \CoteHautY) {dos \Dos\ mm (\DosOrigine)};
  \node[anchor=south, font=\tiny\sffamily, text=Repere, align=center]
    at (\MilieuL, \CoteBasY) {\CotesTexte};
\fi

% ---- Contrôle de superposition (§ 8) ----------------------------------
% Rien n'est dessiné : des mesures et des lignes de journal.
\MesureNoeud{filigrane}{Fil}{\the\FiligraneSep}
\MesureNoeud{titre}{Tit}{0pt}
\MesureNoeud{version}{Ver}{0pt}
\MesureNoeud{quatre}{Qua}{0pt}
\MesureNoeud{credits}{Cre}{0pt}
% Bord droit des crédits = bord gauche + marge interne + plus longue ligne.
\xdef\Crexmax{\fpeval{round(\Crexmin + (\the\FiligraneSep + \the\wd\CreditsBoite)/1mm, 2)}}

% Le filigrane peut entrer dans la zone de sécurité (il est décoratif),
% pas franchir la coupe ni sortir du bandeau — le détourage le rognerait.
\ControleDans{filigrane dans le bandeau}{Fil}{\BandeauG}{\CoupeD}{\CoupeB}{\CoupeH}
\ControleDans{titre sur la 1re, hors bandeau}{Tit}{\PliPremiere}{\BandeauAccent}{\SecB}{\SecH}
\ControleDans{version sur la 1re, hors bandeau}{Ver}{\PliPremiere}{\BandeauAccent}{\SecB}{\SecH}
\ControleDans{texte de 4e dans la zone de securite}{Qua}{\SecG}{\PliQuatre}{\SecB}{\SecH}
\ControleDans{credits de 4e dans la zone de securite}{Cre}{\SecG}{\PliQuatre}{\SecB}{\SecH}
\ControleDisjoints{titre / version}{Tit}{\Verxmin}{\Verxmax}{\Verymin}{\Verymax}
\ControleDisjoints{texte de 4e / credits}{Qua}{\Crexmin}{\Crexmax}{\Creymin}{\Creymax}
\ControleDisjoints{texte de 4e / emplacement ISBN}{Qua}{\IsbnG}{\IsbnD}{\SecB}{\IsbnH}
\ControleDisjoints{credits / emplacement ISBN}{Cre}{\IsbnG}{\IsbnD}{\SecB}{\IsbnH}
\ifnum\fpeval{\Dos >= 6}=1
  \MesureNoeud{dos}{Dos}{0pt}
  \ControleDans{texte du dos entre les plis}{Dos}{\DosG}{\DosD}{\SecB}{\SecH}
\else
  \typeout{COUVERTURE-CONTROLE OK dos de \Dos\space mm : trop etroit, sans texte}%
\fi

\end{tikzpicture}
\typeout{COUVERTURE-BILAN alertes=\the\NbAlertes\space; classe=\Classe\space;
  format=\Format\space; reliure=\Reliure\space; pages=\NbPages\space;
  version=\Version\space;
  dos=\Dos\space mm (\DosOrigine); fini=\FiniL x \PageHauteur\space mm;
  papier=\PapierL x \PapierH\space mm}
\end{document}
"
%
%  ---------------------------------------------------------------------
%  DEUX MODES, un seul paramètre : \Reperes
%
%    oui  — mode ÉPREUVE. Ajoute une marge technique, les traits de coupe,
%           les limites du dos, la zone de sécurité et les cotes. C'est ce
%           qu'il faut pour vérifier à l'écran.
%    non  — mode PRODUCTION. Papier = format final + fond perdu, rien
%           d'autre. C'est ce fichier-là qu'on envoie à l'imprimeur.
%
%  ---------------------------------------------------------------------
%  CONTRÔLE DE SUPERPOSITION
%
%  À la fin du dessin, chaque bloc de texte est mesuré par ses ancres TikZ
%  (quatre coins, ce qui vaut aussi pour le texte pivoté du dos) et comparé
%  à la zone où il doit tenir, puis aux blocs voisins. Le journal porte une
%  ligne par vérification :
%
%      COUVERTURE-CONTROLE OK      <quoi> : ...
%      COUVERTURE-CONTROLE ALERTE  <quoi> : ...
%
%  et un bilan : COUVERTURE-BILAN alertes=<n> ; ...
%
%  Ce que le contrôle VOIT : un bloc qui franchit le pli du dos, la coupe,
%  la zone de sécurité ou le bord du bandeau ; deux blocs qui se chevauchent ;
%  un bloc qui empiète sur l'emplacement réservé au code-barres ISBN.
%  Ce qu'il NE VOIT PAS : une ligne trop longue à l'INTÉRIEUR d'un bloc à
%  largeur fixe (la boîte garde sa largeur, le texte déborde). Celle-là se
%  signale par un « Overfull \hbox » dans le même journal, que la fenêtre
%  compte aussi.
%
%  ---------------------------------------------------------------------
%  AVERTISSEMENT SUR L'ÉPAISSEUR DU DOS
%
%  Le dos se calcule ainsi :
%
%      épaisseur d'une feuille (mm) = grammage x bouffant / 1000
%      dos (mm)                     = (nombre de pages / 2) x épaisseur
%
%  On divise par deux parce qu'une feuille porte DEUX pages.
%
%  Le « bouffant » (main, volume, bulk) dépend du papier et n'est pas
%  devinable : un offset 80 g courant est autour de 1,30, un bouffant
%  d'édition monte à 1,75 et plus. UN ÉCART DE BOUFFANT DÉPLACE LE DOS DE
%  PLUSIEURS MILLIMÈTRES. Demander à l'imprimeur soit son bouffant, soit
%  directement l'épaisseur du dos, et reporter ci-dessous. Tant que cette
%  valeur n'est pas confirmée, la couverture n'est pas bonne à tirer.
%
%  Si l'imprimeur donne le dos en clair, forcer la valeur avec \DosImpose.
% =====================================================================

\documentclass[11pt]{article}

% ---------------------------------------------------------------------
% 1. PARAMÈTRES — c'est la seule zone à modifier
% ---------------------------------------------------------------------

% Classe de l'ouvrage : N, E, A, NE, EA ou NEA.
% (\providecommand : une valeur passée en ligne de commande l'emporte.)
\providecommand{\Classe}{NEA}

% Mode épreuve (oui) ou production (non).
\providecommand{\Reperes}{oui}

% Format du livre FINI : a4, 20x24 ou 20x24-marge — les noms de --format
% dans build_book.py. Les deux maquettes 20 x 24 ne diffèrent que par leur
% colonne de marge intérieure : le livre fini, donc la couverture, est le
% même (200 x 240 mm).
\providecommand{\Format}{a4}

% Reliure : « souple » (dos carré collé, fond perdu 3 mm) ou « rigide »
% (couverture rigide, dite hardcover). En rigide, consignes de l'imprimeur
% du 17/09/2026 : chaque plat déborde du livre de 5 mm en largeur et en
% hauteur (le carton dépasse du bloc), et le fichier porte 15 mm de rabat
% autour du carton au lieu de 3 mm de fond perdu. Les aplats vont déjà au
% bord du papier : ils remplissent le rabat d'eux-mêmes.
\providecommand{\Reliure}{souple}

% Pagination et version imposées de l'extérieur. Vides = tables du § 2.
% La fenêtre les renseigne toujours : la pagination vient du journal du
% livre réellement compilé, dans le format demandé.
\providecommand{\PagesForcees}{}
\providecommand{\VersionForcee}{}

% Papier de l'INTÉRIEUR du livre (pas de la couverture).
% CoolLibri, papier « Standard 90 g blanc ».
\def\Grammage{90}        % g/m²
\def\Bouffant{1.206}     % main DÉDUITE, voir ci-dessous

% Épaisseur du dos IMPOSÉE par l'imprimeur, en mm, par classe.
% Une valeur > 0 prime sur le calcul — MAIS SEULEMENT pour la pagination et
% le format pour lesquels l'imprimeur l'a donnée (\DosImposePages<classe>,
% \DosImposeFormat<classe>). Hors de ce cas, le dos est calculé et le
% journal le dit : un dos de 14 mm donné pour 258 pages A4 ne vaut rien pour
% 262 pages, ni pour le même livre recomposé en 20 x 24.
%
% ATTENTION : cette valeur est propre à CHAQUE ÉDITION. Un dos unique
% appliqué aux quatre livres serait faux pour trois d'entre eux.
%
%  N   — 14 mm, VALEUR DONNÉE PAR COOLIBRI pour 258 pages en 90 g, A4.
%        Fichier attendu : 440 x 303 mm, ce que ce gabarit produit en mode
%        production (2 x 210 + 14 + 2 x 3 = 440 ; 297 + 2 x 3 = 303).
%        Le N de la a.3 compte 262 pages : cette valeur ne s'y applique plus.
%  E, A, NEA — non demandés à l'imprimeur : ils sont CALCULÉS avec la main
%        déduite ci-dessous. À faire confirmer avant tout tirage.
%
% La main de 1,206 n'est pas publiée par CoolLibri : elle est déduite de
% leur propre chiffre, 14 mm / (0,090 x 258/2) = 1,206. Elle reproduit donc
% exactement le dos de N, et n'est qu'une ESTIMATION pour les autres.
% Écarts que cela recouvre, pour mémoire : une main de 1,0 donnerait
% 11,61 mm pour N, une main de 1,3 en donnerait 16,77. Il ne fallait pas
% deviner.
%  E   — 17 mm, VALEUR DONNÉE PAR L'IMPRIMEUR le 17/09/2026 pour 276 pages,
%        20 x 24, couverture RIGIDE, reliure collée (Klebebindung), papier
%        115 g. En reliure cousue (Fadenheftung), il annonce 18 mm.
%
% \DosImposeReliure<classe> : la reliure pour laquelle la valeur a été
% donnée. Le dos d'une couverture rigide n'est pas celui d'un dos carré
% collé souple, même à pagination égale.
\def\DosImposeN{14}   \def\DosImposePagesN{258}  \def\DosImposeFormatN{a4}
\def\DosImposeE{17}   \def\DosImposePagesE{276}  \def\DosImposeFormatE{20x24-marge}
\def\DosImposeA{0}    \def\DosImposePagesA{0}    \def\DosImposeFormatA{}
\def\DosImposeNE{0}   \def\DosImposePagesNE{0}   \def\DosImposeFormatNE{}
\def\DosImposeEA{0}   \def\DosImposePagesEA{0}   \def\DosImposeFormatEA{}
\def\DosImposeNEA{0}  \def\DosImposePagesNEA{0}  \def\DosImposeFormatNEA{}
\def\DosImposeReliureN{souple}  \def\DosImposeReliureE{rigide}
\def\DosImposeReliureA{souple}  \def\DosImposeReliureNE{souple}
\def\DosImposeReliureEA{souple} \def\DosImposeReliureNEA{souple}

% Fond perdu, en mm. 3 mm est la valeur courante ; certains imprimeurs
% demandent 5 mm.
\def\FondPerdu{3}

% Marge technique autour du fond perdu, en mm. Sert à loger les traits de
% coupe en mode épreuve. Ignorée en mode production.
\def\MargeTech{12}

% Zone de sécurité, en mm : aucun texte ne doit s'en approcher davantage
% du bord de coupe.
\def\Securite{8}

% ---------------------------------------------------------------------
% 2. DONNÉES PAR CLASSE
%    pages : nombre de pages du PDF intérieur, COUVERTURE NON COMPRISE.
%    Un dos carré collé impose un nombre de pages PAIR ; le fichier
%    refuse de compiler si le compte est impair.
%
%    Ces tables ne servent que si \PagesForcees et \VersionForcee sont
%    vides, c'est-à-dire à la main. Valeurs A4 mesurées sur les PDF de la
%    a.3 (CLAUDE.md § 1, 05/09/2026). En 20 x 24, la pagination change
%    entièrement : NEA passe de 816 à 1036 pages.
% ---------------------------------------------------------------------

\def\DonneesN{262}
\def\DonneesE{214}
\def\DonneesA{386}
\def\DonneesNE{0}      % édition jamais produite
\def\DonneesEA{0}      % édition jamais produite
\def\DonneesNEA{816}

% Version imprimée sur le dos et la 1re de couverture, PAR CLASSE.
% « Les versions sont suivies par classe, chacune évoluant à son rythme »
% (CHANGELOG.md).
\def\VersionN{a.3}
\def\VersionE{a.3}
\def\VersionA{a.3}
\def\VersionNE{a.3}
\def\VersionEA{a.3}
\def\VersionNEA{a.3}

% ---------------------------------------------------------------------
% 3. TEXTES
% ---------------------------------------------------------------------

\def\TitreOuvrage{50\,Ohm}
\def\SousTitre{Préparation à l'examen radioamateur}

% Sous-titre propre à chaque classe (repris de FR_TITLES, build_book.py).
\def\SousTitreN{Cours complet — classe N}
\def\SousTitreE{Cours complet — classe E}
\def\SousTitreA{Cours complet — classe A}
\def\SousTitreNE{Cours complet — classes N et E}
\def\SousTitreEA{Cours complet — classes E et A}
\def\SousTitreNEA{Cours complet — classes N, E et A}

% Lettres du filigrane, empilées de haut en bas (cf. build_book.py v0.20).
\def\LettresN{N}
\def\LettresE{E}
\def\LettresA{A}
\def\LettresNE{N,E}
\def\LettresEA{E,A}
\def\LettresNEA{N,E,A}

% Texte de 4e de couverture.
% NB : il parle des classes N, E et A quelle que soit l'édition. Juste pour
% N, E, A et NEA ; approximatif pour NE et EA ; faux pour SWL, d'où son refus.
\def\QuatriemeTexte{%
Ce guide d'étude prépare aux examens radioamateur allemands des classes
N, E et A. Il traduit en français les contenus du projet 50ohm.de, publiés
par le référat AJW du DARC e.\,V., et les complète d'encarts « En France »
qui signalent, sans rien retirer au texte allemand, ce que la
réglementation française prévoit sur le même sujet.\par\smallskip
Les questions officielles du catalogue allemand accompagnent chaque
notion, avec leurs réponses.%
}

\def\Auteur{Traduit avec l'aide d'une IA par Pierre F4JWI}
\def\Depot{github.com/DrPLG/50ohm-fr}

% =====================================================================
%  À PARTIR D'ICI, PLUS RIEN À MODIFIER
% =====================================================================

\usepackage{iftex}
\RequireLuaTeX
\usepackage{xfp}
\usepackage{tikz}
\usetikzlibrary{calc}
\usepackage{fontspec}
\usepackage[french]{babel}
\usepackage{geometry}

% ---------------------------------------------------------------------
% 4. RÉSOLUTION DES DONNÉES DE LA CLASSE
% ---------------------------------------------------------------------

\makeatletter
% Tome SWL : refus explicite, plutôt qu'un « classe inconnue » qui laisserait
% croire à un oubli. Le texte de 4e parle des classes N, E et A.
\def\@swlmot{SWL}
\ifx\Classe\@swlmot
  \PackageError{couverture}{Pas de couverture pour le tome SWL}%
    {Sa 4e de couverture n'est pas redigee : le texte commun parle des
     classes N, E et A. C'est une decision editoriale a prendre d'abord.}%
\fi

\def\@resoudre#1#2{%
  \expandafter\let\expandafter#1\csname #2\Classe\endcsname
  \ifx#1\relax
    \PackageError{couverture}{Classe \Classe\space inconnue : pas de #2\Classe}%
      {Valeurs admises : N, E, A, NE, EA, NEA.}%
  \fi}
\@resoudre\NbPages{Donnees}
\@resoudre\SousTitreClasse{SousTitre}
\@resoudre\Lettres{Lettres}
\@resoudre\DosImpose{DosImpose}
\@resoudre\DosImposePages{DosImposePages}
\@resoudre\DosImposeFormat{DosImposeFormat}
\@resoudre\DosImposeReliure{DosImposeReliure}
\@resoudre\Version{Version}

% Valeurs imposées de l'extérieur. Test de vacuité par \detokenize : il ne
% dépend pas de la façon dont la macro a été définie.
\if\relax\detokenize\expandafter{\PagesForcees}\relax\else
  \let\NbPages\PagesForcees
\fi
\if\relax\detokenize\expandafter{\VersionForcee}\relax\else
  \let\Version\VersionForcee
\fi

% Format du livre fini.
\def\@formatAquatre{a4}
\def\@formatVingt{20x24}
\def\@formatVingtMarge{20x24-marge}
\ifx\Format\@formatAquatre
  \def\PageLargeur{210}\def\PageHauteur{297}
\else\ifx\Format\@formatVingt
  \def\PageLargeur{200}\def\PageHauteur{240}
\else\ifx\Format\@formatVingtMarge
  \def\PageLargeur{200}\def\PageHauteur{240}
\else
  \PackageError{couverture}{Format \Format\space inconnu}%
    {Valeurs admises : a4, 20x24, 20x24-marge.}%
  \def\PageLargeur{210}\def\PageHauteur{297}
\fi\fi\fi
% Reliure. Hors de « souple » et « rigide » : erreur, pas de supposition.
\def\@reliuresouple{souple}
\def\@reliurerigide{rigide}
\ifx\Reliure\@reliurerigide
  \edef\PageLargeur{\fpeval{\PageLargeur + 5}}
  \edef\PageHauteur{\fpeval{\PageHauteur + 5}}
  \def\FondPerdu{15}
\else\ifx\Reliure\@reliuresouple\else
  \PackageError{couverture}{Reliure \Reliure\space inconnue}%
    {Valeurs admises : souple, rigide.}%
\fi\fi
\makeatother

% Garde-fou : une édition jamais produite n'a pas de couverture.
\ifnum\NbPages=0
  \PackageError{couverture}{L'edition \Classe\space n'a pas de pagination}%
    {Cette edition n'a jamais ete compilee. Renseigner \string\Donnees\Classe
     ou passer \string\PagesForcees.}%
\fi

% Avertissement : l'imprimeur demande un nombre de pages multiple de 4
% (17/09/2026). Pas une erreur : c'est son exigence, pas une loi du papier.
\ifnum\fpeval{\NbPages - 4*floor(\NbPages/4)}=0 \else
  \typeout{COUVERTURE-AVERTISSEMENT \NbPages\space pages : pas un multiple
    de 4, que l'imprimeur exige -- ajouter des pages a l'interieur}%
\fi

% Garde-fou : un dos carré collé impose un nombre de pages pair.
\ifodd\NbPages
  \PackageError{couverture}{Nombre de pages impair (\NbPages)}%
    {Un dos carre colle exige un nombre PAIR de pages : une feuille en
     porte deux. Ajouter une page blanche a l'interieur.}%
\fi

% ---------------------------------------------------------------------
% 5. GÉOMÉTRIE — c'est ici que se joue l'exactitude du dos
% ---------------------------------------------------------------------

% Épaisseur du dos : valeur imposée si elle est donnée ET qu'elle vaut pour
% cette pagination et ce format ; calcul sinon.
\edef\DosCalcule{\fpeval{round(\NbPages / 2 * \Grammage * \Bouffant / 1000, 2)}}
\def\DosOrigine{calculee}
\ifnum\fpeval{\DosImpose > 0}=1
  \ifnum\fpeval{\NbPages = \DosImposePages}=1
    \ifx\Format\DosImposeFormat
      \ifx\Reliure\DosImposeReliure
        \def\DosOrigine{imposee}
      \fi
    \fi
  \fi
  \def\MotCalculee{calculee}
  \ifx\DosOrigine\MotCalculee
    \typeout{COUVERTURE-AVERTISSEMENT dos impose de \DosImpose\space mm
      donne pour \DosImposePages\space pages en \DosImposeFormat\space
      (reliure \DosImposeReliure) -- inapplicable a \NbPages\space pages en
      \Format\space (reliure \Reliure), dos calcule}%
  \fi
\fi
% NB : pas de \@ dans ces noms. Cette section est hors \makeatletter, et
% « \def\@imposee{...} » y redéfinit \@ (suivi du texte « imposee ») : le
% test échouait toujours, en silence. Constaté le 16/09/2026 — le dos de N
% tombait juste par coïncidence, 14,0 mm calculés pour 14 imposés.
\def\MotRigide{rigide}
\def\MotImposee{imposee}
\ifx\DosOrigine\MotImposee
  \edef\Dos{\DosImpose}
\else
  \edef\Dos{\DosCalcule}
\fi

% Marge technique : uniquement en mode épreuve.
\def\ouimot{oui}
\ifx\Reperes\ouimot \edef\Tech{\MargeTech}\else \def\Tech{0}\fi

% Dimensions du papier.
%   largeur = marge + fond perdu + 4e + dos + 1re + fond perdu + marge
%   hauteur = marge + fond perdu + hauteur + fond perdu + marge
\edef\PapierL{\fpeval{2*\PageLargeur + \Dos + 2*\FondPerdu + 2*\Tech}}
\edef\PapierH{\fpeval{\PageHauteur + 2*\FondPerdu + 2*\Tech}}

\geometry{paperwidth=\PapierL mm, paperheight=\PapierH mm,
          margin=0pt, nohead, nofoot}

% ---------------------------------------------------------------------
%  TOUTES LES COTES SONT FIGÉES ICI, ET NULLE PART AILLEURS.
%
%  Raison, payée par trois compilations ratées : un \fpeval laissé dans
%  une coordonnée ou dans le texte d'un nœud TikZ n'est pas développé au
%  moment où TikZ analyse le chemin. On obtient « Giving up on this path.
%  Did you forget a semicolon? », qui accuse une ponctuation alors que la
%  faute est ailleurs. Le corps du document ci-dessous n'emploie donc que
%  des macros déjà calculées.
% ---------------------------------------------------------------------

% Repères horizontaux, mesurés depuis le bord GAUCHE du papier.
\edef\CoupeG{\fpeval{\Tech + \FondPerdu}}          % coupe gauche (4e de couv.)
\edef\DosG{\fpeval{\CoupeG + \PageLargeur}}         % dos, début
\edef\DosD{\fpeval{\DosG + \Dos}}                   % dos, fin
\edef\CoupeD{\fpeval{\DosD + \PageLargeur}}         % coupe droite (1re de couv.)
% Repères verticaux, depuis le bord BAS.
\edef\CoupeB{\fpeval{\Tech + \FondPerdu}}
\edef\CoupeH{\fpeval{\CoupeB + \PageHauteur}}
% Axes.
\edef\DosAxe{\fpeval{(\DosG + \DosD)/2}}
\edef\MilieuH{\fpeval{(\CoupeB + \CoupeH)/2}}
\edef\MilieuL{\fpeval{\PapierL/2}}

% Format fini, pour les cotes affichées.
\edef\FiniL{\fpeval{2*\PageLargeur + \Dos}}

% Bandeau de la 1re de couverture : même proportion que la page de titre.
\edef\BandeauL{\fpeval{round(0.34*\PageLargeur, 2)}}
\edef\BandeauG{\fpeval{\CoupeD - \BandeauL}}
\edef\BandeauAccent{\fpeval{\BandeauG - 0.7}}
\edef\DosAccentG{\fpeval{\DosG - 0.7}}
\edef\DosAccentD{\fpeval{\DosD + 0.7}}
% Couverture rigide : le bleu du dos déborde de 2 mm sur chaque plat, pour
% absorber la tolérance de l'imprimeur sur l'épaisseur du dos (±10 %).
\edef\DosEtenduG{\fpeval{\DosG - 2}}
\edef\DosEtenduD{\fpeval{\DosD + 2}}

% Filigrane de classe.
\edef\FiligraneX{\fpeval{(\BandeauG + \CoupeD)/2}}
\edef\FiligraneY{\fpeval{\CoupeH - 6}}

% Bloc-titre de la 1re de couverture.
\edef\TitreX{\fpeval{\DosD + 12}}
\edef\TitreY{\fpeval{\CoupeH - 72}}
\edef\TitreW{\fpeval{\BandeauG - \DosD - 16}}
\edef\VersionY{\fpeval{\CoupeB + 14}}

% 4e de couverture.
\edef\QuatreX{\fpeval{\CoupeG + \Securite + 6}}
\edef\QuatreY{\fpeval{\CoupeH - \Securite - 14}}
\edef\QuatreBasY{\fpeval{\CoupeB + \Securite + 22}}
\edef\QuatreW{\fpeval{\PageLargeur - 2*\Securite - 12}}

% Zone de sécurité.
\edef\SecG{\fpeval{\CoupeG + \Securite}}
\edef\SecD{\fpeval{\CoupeD - \Securite}}
\edef\SecB{\fpeval{\CoupeB + \Securite}}
\edef\SecH{\fpeval{\CoupeH - \Securite}}

% Zones de contrôle : ce qui doit rester de part et d'autre du pli du dos.
\edef\PliQuatre{\fpeval{\DosG - \Securite}}         % 4e : limite côté dos
\edef\PliPremiere{\fpeval{\DosD + \Securite}}       % 1re : limite côté dos

% Traits de coupe et cotes, dans la marge technique.
\edef\TechExt{\fpeval{\Tech - 8}}
\edef\TechInt{\fpeval{\Tech - 1}}
\edef\TechCote{\fpeval{\Tech - 9}}
\edef\HautInt{\fpeval{\PapierH - \Tech + 1}}
\edef\HautExt{\fpeval{\PapierH - \Tech + 8}}
\edef\HautCote{\fpeval{\PapierH - \Tech + 9}}
\edef\DroitInt{\fpeval{\PapierL - \Tech + 1}}
\edef\DroitExt{\fpeval{\PapierL - \Tech + 8}}

% Emplacement du code-barres ISBN : en bas à DROITE de la 4e de couverture,
% selon l'usage éditorial, et non à gauche où il chevaucherait les mentions
% de licence.
\edef\IsbnD{\fpeval{\DosG - \Securite}}
\edef\IsbnG{\fpeval{\IsbnD - 40}}
\edef\IsbnH{\fpeval{\SecB + 25}}
\edef\IsbnCX{\fpeval{\IsbnG + 20}}
\edef\IsbnCY{\fpeval{\SecB + 12.5}}

% Cartouches de cotes, dans la marge technique. Ancrés vers l'INTÉRIEUR du
% papier : ancrés vers l'extérieur, ils débordaient et étaient rognés.
\edef\CoteBasY{\fpeval{1}}
\edef\CoteHautY{\fpeval{\PapierH - 1}}

% ---------------------------------------------------------------------
% 6. COULEURS ET POLICES — reprises de la page de titre (build_book.py)
% ---------------------------------------------------------------------

\definecolor{TitleBand}{cmyk}{.9,.55,.1,.35}   % bleu profond
\definecolor{TitleAccent}{cmyk}{.05,.8,1,0}    % rouge DARC (accent fin)
\definecolor{Repere}{cmyk}{0,1,1,0}            % magenta : repères de montage

\setmainfont{Libertinus Serif}[
  ItalicFont     = Libertinus Serif Italic,
  BoldFont       = Libertinus Serif Bold,
  BoldItalicFont = Libertinus Serif Bold Italic]

\pagestyle{empty}

% ---------------------------------------------------------------------
% 7. FILIGRANE DE CLASSE
%    Le texte est construit ici, et non dans le nœud : un \\ produit à
%    l'intérieur d'un \foreach referme le groupe d'alignement de TikZ.
%    Une lettre seule est grande ; plusieurs s'empilent et se réduisent
%    (build_book.py v0.20 : « NEA » à 105 pt débordait encore).
% ---------------------------------------------------------------------

\makeatletter
\def\FiligraneTexte{}
\newif\ifpremierelettre \premierelettretrue
\count255=0
\@for\@tempa:=\Lettres\do{%
  \advance\count255 by 1\relax
  \ifpremierelettre \premierelettrefalse
  \else \g@addto@macro\FiligraneTexte{\\}\fi
  \expandafter\g@addto@macro\expandafter\FiligraneTexte\expandafter{\@tempa}}%
\xdef\NbLettres{\the\count255}
\makeatother

\ifnum\NbLettres=1 \def\CorpsF{150}\else\def\CorpsF{86}\fi
\edef\CorpsFI{\fpeval{round(\CorpsF*1.02, 0)}}

% Texte du cartouche de cotes, lui aussi construit hors du nœud TikZ.
%
% NE PAS y employer \textsuperscript : dans un nœud « align=center », il
% casse le chemin TikZ (« Giving up on this path. Did you forget a
% semicolon? »), et l'erreur est signalée sur le nœud SUIVANT — elle
% accuse donc une ponctuation là où la faute est ici. Isolé par essais :
% $^2$ et le caractère ² passent, \textsuperscript non.
\def\CotesTexte{%
  \Classe\ --- \NbPages\ pages --- papier \Grammage\ g/m$^2$,
  bouffant \Bouffant\ --- fond perdu \FondPerdu\ mm\\
  papier \PapierL\ mm sur \PapierH\ mm --- fini \FiniL\ mm sur
  \PageHauteur\ mm\\
  EPREUVE, ne pas envoyer a l'impression%
}

% ---------------------------------------------------------------------
% 8. CONTRÔLE DE SUPERPOSITION
%
%  \MesureNoeud{nœud}{préfixe}{retrait} mesure la boîte d'un nœud par ses
%  quatre coins (un nœud pivoté a ses ancres pivotées : le min et le max des
%  quatre suffisent) et la ramène au repère du papier, en mm. Le point (0,0)
%  est lu par le même « let » que les ancres : les deux passent donc par la
%  même transformation, et la comparaison ne dépend pas de l'endroit où
%  TeX a posé l'image sur la page.
%
%  Le retrait sert au seul filigrane, pour mesurer les lettres et non leur
%  cadre : c'est le bord des lettres que le détourage rognerait.
%
%  PIÈGE MESURÉ le 16/09/2026 : la marge interne TikZ (.3333em) n'est PAS
%  évaluée dans la police du nœud, mais dans celle qui entoure le dessin.
%  Instrumenté : dans le nœud, em = 150 pt ; boîte brute du « N » :
%  41,7 x 37,5 mm, soit la lettre plus 2 x 1,3 mm — 1,3 mm étant .3333em
%  à 11 pt. Un retrait calculé à 150 pt (17,6 mm) ramenait la lettre à
%  6,6 mm de large, et le contrôle aurait dit OK sur une mesure absurde.
% ---------------------------------------------------------------------

\newcount\NbAlertes
\newdimen\FiligraneSep
\newsavebox\CreditsBoite

% Crédits de 4e : lignes coupées à la main. Leur largeur RÉELLE est celle
% de la plus longue ligne, mesurée dans une boîte à part (voir § 8, ISBN).
\def\CreditsTexte{%
    Réalisé à partir des contenus de 50ohm.de (en allemand)\\
    50ohm.de-Autorenteam, coordonné par le référat AJW du DARC e.\,V.\\
    \Auteur\\[3pt]
    Licence CC BY 4.0 --- \texttt{\Depot}}

\def\MesureNoeud#1#2#3{%
  \path let \p1=(#1.south west), \p2=(#1.north east),
            \p3=(#1.north west), \p4=(#1.south east), \p5=(0,0) in
    \pgfextra{%
      \expandafter\xdef\csname #2xmin\endcsname{\fpeval{round((min(\x1,\x2,\x3,\x4) - \x5 + #3)/1mm, 2)}}%
      \expandafter\xdef\csname #2xmax\endcsname{\fpeval{round((max(\x1,\x2,\x3,\x4) - \x5 - #3)/1mm, 2)}}%
      \expandafter\xdef\csname #2ymin\endcsname{\fpeval{round((min(\y1,\y2,\y3,\y4) - \y5 + #3)/1mm, 2)}}%
      \expandafter\xdef\csname #2ymax\endcsname{\fpeval{round((max(\y1,\y2,\y3,\y4) - \y5 - #3)/1mm, 2)}}%
    };}

% \ControleDans{libellé}{préfixe}{xmin}{xmax}{ymin}{ymax}
\def\ControleDans#1#2#3#4#5#6{%
  \ifnum\fpeval{\csname #2xmin\endcsname >= #3 && \csname #2xmax\endcsname <= #4
             && \csname #2ymin\endcsname >= #5 && \csname #2ymax\endcsname <= #6}=1
    \typeout{COUVERTURE-CONTROLE OK #1 : x \csname #2xmin\endcsname-\csname #2xmax\endcsname,
      y \csname #2ymin\endcsname-\csname #2ymax\endcsname\space dans x #3-#4, y #5-#6}%
  \else
    \global\advance\NbAlertes by 1
    \typeout{COUVERTURE-CONTROLE ALERTE #1 : x \csname #2xmin\endcsname-\csname #2xmax\endcsname,
      y \csname #2ymin\endcsname-\csname #2ymax\endcsname\space HORS de x #3-#4, y #5-#6}%
  \fi}

% \ControleDisjoints{libellé}{préfixe}{xmin}{xmax}{ymin}{ymax}
% Deux rectangles se chevauchent si et seulement s'ils se recouvrent sur
% les DEUX axes à la fois.
\def\ControleDisjoints#1#2#3#4#5#6{%
  \ifnum\fpeval{\csname #2xmin\endcsname < #4 && #3 < \csname #2xmax\endcsname
             && \csname #2ymin\endcsname < #6 && #5 < \csname #2ymax\endcsname}=1
    \global\advance\NbAlertes by 1
    \typeout{COUVERTURE-CONTROLE ALERTE #1 : chevauchement}%
  \else
    \typeout{COUVERTURE-CONTROLE OK #1 : disjoints}%
  \fi}

\begin{document}
% Marge interne du filigrane, évaluée dans la police qui ENTOURE le dessin :
% c'est là que TikZ évalue le .3333em de « inner sep » (cf. § 8).
\FiligraneSep=.3333em
% Largeur réelle des crédits : la plus longue ligne, dans leur police.
\sbox\CreditsBoite{\small\begin{tabular}{@{}l@{}}\CreditsTexte\end{tabular}}%
\begin{tikzpicture}[remember picture, overlay,
                    x=1mm, y=1mm, shift={(current page.south west)}]

% ---- Fonds ----------------------------------------------------------
% La 4e de couverture reste sur le blanc du papier.
% Le bandeau bleu couvre le dos, puis reprend au tiers extérieur de la
% 1re de couverture — même partition 2/3-1/3 que la page de titre.
% Couverture rigide : pas de filets au dos (ils tomberaient à côté des
% plis, vu la tolérance), et un dos bleu élargi de 2 mm de chaque côté.
\ifx\Reliure\MotRigide
  \fill[TitleBand] (\DosEtenduG, 0) rectangle (\DosEtenduD, \PapierH);
\else
  \fill[TitleBand] (\DosG, 0) rectangle (\DosD, \PapierH);
\fi
\fill[TitleBand] (\BandeauG, 0) rectangle (\PapierL, \PapierH);
% Filets d'accent : à la jonction sur la 1re, de part et d'autre du dos.
\fill[TitleAccent] (\BandeauAccent, 0) rectangle (\BandeauG, \PapierH);
\ifx\Reliure\MotRigide\else
  \fill[TitleAccent] (\DosAccentG, 0) rectangle (\DosG, \PapierH);
  \fill[TitleAccent] (\DosD, 0) rectangle (\DosAccentD, \PapierH);
\fi

% ---- 1re de couverture ----------------------------------------------
% Filigrane de classe, en réserve claire, détouré au bandeau.
\begin{scope}
\clip (\BandeauG, 0) rectangle (\PapierL, \PapierH);
\node[name=filigrane, anchor=north, align=center, text=white, opacity=0.16,
      font=\fontsize{\CorpsF}{\CorpsFI}\selectfont\bfseries]
  at (\FiligraneX, \FiligraneY) {\FiligraneTexte};
\end{scope}

\node[name=titre, anchor=west, align=left, text width=\TitreW mm]
  at (\TitreX, \TitreY) {%
    {\fontsize{54}{56}\selectfont\bfseries \TitreOuvrage}\\[10pt]
    {\Large\color{TitleBand}\bfseries \SousTitre}\\[16pt]
    {\large\color{black!75} \SousTitreClasse}};

\node[name=version, anchor=south west, align=left, font=\small, text=black!70]
  at (\TitreX, \VersionY) {Version \Version\\ Licence CC BY 4.0};

% ---- Dos -------------------------------------------------------------
% Un dos étroit ne reçoit pas de texte : en deçà de 6 mm, le débord au
% massicot ferait mordre le texte vertical sur les plats.
\ifnum\fpeval{\Dos >= 6}=1
  \node[name=dos, rotate=90, anchor=center, text=white, font=\bfseries,
        inner sep=0pt]
    at (\DosAxe, \MilieuH)
    {\TitreOuvrage\ \ --- \ \SousTitreClasse\ \ --- \ \Version};
\fi

% ---- 4e de couverture ------------------------------------------------
\node[name=quatre, anchor=north west, align=justify, text width=\QuatreW mm]
  at (\QuatreX, \QuatreY) {\QuatriemeTexte};

% Le texte vit dans \CreditsTexte (§ 8) pour être mesuré à part : la boîte
% du nœud a la largeur fixe \QuatreW et couvre l'emplacement ISBN sans
% qu'aucun caractère n'y soit. Retirer « text width » pour l'éviter change
% l'interlignage — essayé le 16/09/2026, 6 934 pixels différents.
\node[name=credits, anchor=south west, align=left, font=\small, text=black!70,
      text width=\QuatreW mm]
  at (\QuatreX, \QuatreBasY) {\CreditsTexte};

% ---- Repères de fabrication (mode épreuve uniquement) ----------------
\ifx\Reperes\ouimot
  % Traits de coupe, hors du fond perdu.
  \foreach \x in {\CoupeG, \CoupeD} {
    \draw[Repere, line width=0.25pt] (\x, \TechExt) -- (\x, \TechInt);
    \draw[Repere, line width=0.25pt] (\x, \HautInt) -- (\x, \HautExt);}
  \foreach \y in {\CoupeB, \CoupeH} {
    \draw[Repere, line width=0.25pt] (\TechExt, \y) -- (\TechInt, \y);
    \draw[Repere, line width=0.25pt] (\DroitInt, \y) -- (\DroitExt, \y);}
  % Limites du dos, en pointillés.
  \foreach \x in {\DosG, \DosD} {
    \draw[Repere, line width=0.25pt, dashed] (\x, \TechExt) -- (\x, \TechInt);
    \draw[Repere, line width=0.25pt, dashed] (\x, \HautInt) -- (\x, \HautExt);}
  % Zone de sécurité.
  \draw[Repere, line width=0.2pt, dotted]
    (\SecG, \SecB) rectangle (\SecD, \SecH);
  % Emplacement réservé au code-barres ISBN.
  \draw[Repere, line width=0.2pt] (\IsbnG, \SecB) rectangle (\IsbnD, \IsbnH);
  \node[font=\tiny\sffamily, text=Repere, align=center]
    at (\IsbnCX, \IsbnCY) {emplacement\\ISBN};
  % Cotes.
  \node[anchor=north, font=\tiny\sffamily, text=Repere]
    at (\DosAxe, \CoteHautY) {dos \Dos\ mm (\DosOrigine)};
  \node[anchor=south, font=\tiny\sffamily, text=Repere, align=center]
    at (\MilieuL, \CoteBasY) {\CotesTexte};
\fi

% ---- Contrôle de superposition (§ 8) ----------------------------------
% Rien n'est dessiné : des mesures et des lignes de journal.
\MesureNoeud{filigrane}{Fil}{\the\FiligraneSep}
\MesureNoeud{titre}{Tit}{0pt}
\MesureNoeud{version}{Ver}{0pt}
\MesureNoeud{quatre}{Qua}{0pt}
\MesureNoeud{credits}{Cre}{0pt}
% Bord droit des crédits = bord gauche + marge interne + plus longue ligne.
\xdef\Crexmax{\fpeval{round(\Crexmin + (\the\FiligraneSep + \the\wd\CreditsBoite)/1mm, 2)}}

% Le filigrane peut entrer dans la zone de sécurité (il est décoratif),
% pas franchir la coupe ni sortir du bandeau — le détourage le rognerait.
\ControleDans{filigrane dans le bandeau}{Fil}{\BandeauG}{\CoupeD}{\CoupeB}{\CoupeH}
\ControleDans{titre sur la 1re, hors bandeau}{Tit}{\PliPremiere}{\BandeauAccent}{\SecB}{\SecH}
\ControleDans{version sur la 1re, hors bandeau}{Ver}{\PliPremiere}{\BandeauAccent}{\SecB}{\SecH}
\ControleDans{texte de 4e dans la zone de securite}{Qua}{\SecG}{\PliQuatre}{\SecB}{\SecH}
\ControleDans{credits de 4e dans la zone de securite}{Cre}{\SecG}{\PliQuatre}{\SecB}{\SecH}
\ControleDisjoints{titre / version}{Tit}{\Verxmin}{\Verxmax}{\Verymin}{\Verymax}
\ControleDisjoints{texte de 4e / credits}{Qua}{\Crexmin}{\Crexmax}{\Creymin}{\Creymax}
\ControleDisjoints{texte de 4e / emplacement ISBN}{Qua}{\IsbnG}{\IsbnD}{\SecB}{\IsbnH}
\ControleDisjoints{credits / emplacement ISBN}{Cre}{\IsbnG}{\IsbnD}{\SecB}{\IsbnH}
\ifnum\fpeval{\Dos >= 6}=1
  \MesureNoeud{dos}{Dos}{0pt}
  \ControleDans{texte du dos entre les plis}{Dos}{\DosG}{\DosD}{\SecB}{\SecH}
\else
  \typeout{COUVERTURE-CONTROLE OK dos de \Dos\space mm : trop etroit, sans texte}%
\fi

\end{tikzpicture}
\typeout{COUVERTURE-BILAN alertes=\the\NbAlertes\space; classe=\Classe\space;
  format=\Format\space; reliure=\Reliure\space; pages=\NbPages\space;
  version=\Version\space;
  dos=\Dos\space mm (\DosOrigine); fini=\FiniL x \PageHauteur\space mm;
  papier=\PapierL x \PapierH\space mm}
\end{document}
"
%
%  ---------------------------------------------------------------------
%  DEUX MODES, un seul paramètre : \Reperes
%
%    oui  — mode ÉPREUVE. Ajoute une marge technique, les traits de coupe,
%           les limites du dos, la zone de sécurité et les cotes. C'est ce
%           qu'il faut pour vérifier à l'écran.
%    non  — mode PRODUCTION. Papier = format final + fond perdu, rien
%           d'autre. C'est ce fichier-là qu'on envoie à l'imprimeur.
%
%  ---------------------------------------------------------------------
%  CONTRÔLE DE SUPERPOSITION
%
%  À la fin du dessin, chaque bloc de texte est mesuré par ses ancres TikZ
%  (quatre coins, ce qui vaut aussi pour le texte pivoté du dos) et comparé
%  à la zone où il doit tenir, puis aux blocs voisins. Le journal porte une
%  ligne par vérification :
%
%      COUVERTURE-CONTROLE OK      <quoi> : ...
%      COUVERTURE-CONTROLE ALERTE  <quoi> : ...
%
%  et un bilan : COUVERTURE-BILAN alertes=<n> ; ...
%
%  Ce que le contrôle VOIT : un bloc qui franchit le pli du dos, la coupe,
%  la zone de sécurité ou le bord du bandeau ; deux blocs qui se chevauchent ;
%  un bloc qui empiète sur l'emplacement réservé au code-barres ISBN.
%  Ce qu'il NE VOIT PAS : une ligne trop longue à l'INTÉRIEUR d'un bloc à
%  largeur fixe (la boîte garde sa largeur, le texte déborde). Celle-là se
%  signale par un « Overfull \hbox » dans le même journal, que la fenêtre
%  compte aussi.
%
%  ---------------------------------------------------------------------
%  AVERTISSEMENT SUR L'ÉPAISSEUR DU DOS
%
%  Le dos se calcule ainsi :
%
%      épaisseur d'une feuille (mm) = grammage x bouffant / 1000
%      dos (mm)                     = (nombre de pages / 2) x épaisseur
%
%  On divise par deux parce qu'une feuille porte DEUX pages.
%
%  Le « bouffant » (main, volume, bulk) dépend du papier et n'est pas
%  devinable : un offset 80 g courant est autour de 1,30, un bouffant
%  d'édition monte à 1,75 et plus. UN ÉCART DE BOUFFANT DÉPLACE LE DOS DE
%  PLUSIEURS MILLIMÈTRES. Demander à l'imprimeur soit son bouffant, soit
%  directement l'épaisseur du dos, et reporter ci-dessous. Tant que cette
%  valeur n'est pas confirmée, la couverture n'est pas bonne à tirer.
%
%  Si l'imprimeur donne le dos en clair, forcer la valeur avec \DosImpose.
% =====================================================================

\documentclass[11pt]{article}

% ---------------------------------------------------------------------
% 1. PARAMÈTRES — c'est la seule zone à modifier
% ---------------------------------------------------------------------

% Classe de l'ouvrage : N, E, A, NE, EA ou NEA.
% (\providecommand : une valeur passée en ligne de commande l'emporte.)
\providecommand{\Classe}{NEA}

% Mode épreuve (oui) ou production (non).
\providecommand{\Reperes}{oui}

% Format du livre FINI : a4, 20x24 ou 20x24-marge — les noms de --format
% dans build_book.py. Les deux maquettes 20 x 24 ne diffèrent que par leur
% colonne de marge intérieure : le livre fini, donc la couverture, est le
% même (200 x 240 mm).
\providecommand{\Format}{a4}

% Reliure : « souple » (dos carré collé, fond perdu 3 mm) ou « rigide »
% (couverture rigide, dite hardcover). En rigide, consignes de l'imprimeur
% du 17/09/2026 : chaque plat déborde du livre de 5 mm en largeur et en
% hauteur (le carton dépasse du bloc), et le fichier porte 15 mm de rabat
% autour du carton au lieu de 3 mm de fond perdu. Les aplats vont déjà au
% bord du papier : ils remplissent le rabat d'eux-mêmes.
\providecommand{\Reliure}{souple}

% Pagination et version imposées de l'extérieur. Vides = tables du § 2.
% La fenêtre les renseigne toujours : la pagination vient du journal du
% livre réellement compilé, dans le format demandé.
\providecommand{\PagesForcees}{}
\providecommand{\VersionForcee}{}

% Papier de l'INTÉRIEUR du livre (pas de la couverture).
% CoolLibri, papier « Standard 90 g blanc ».
\def\Grammage{90}        % g/m²
\def\Bouffant{1.206}     % main DÉDUITE, voir ci-dessous

% Épaisseur du dos IMPOSÉE par l'imprimeur, en mm, par classe.
% Une valeur > 0 prime sur le calcul — MAIS SEULEMENT pour la pagination et
% le format pour lesquels l'imprimeur l'a donnée (\DosImposePages<classe>,
% \DosImposeFormat<classe>). Hors de ce cas, le dos est calculé et le
% journal le dit : un dos de 14 mm donné pour 258 pages A4 ne vaut rien pour
% 262 pages, ni pour le même livre recomposé en 20 x 24.
%
% ATTENTION : cette valeur est propre à CHAQUE ÉDITION. Un dos unique
% appliqué aux quatre livres serait faux pour trois d'entre eux.
%
%  N   — 14 mm, VALEUR DONNÉE PAR COOLIBRI pour 258 pages en 90 g, A4.
%        Fichier attendu : 440 x 303 mm, ce que ce gabarit produit en mode
%        production (2 x 210 + 14 + 2 x 3 = 440 ; 297 + 2 x 3 = 303).
%        Le N de la a.3 compte 262 pages : cette valeur ne s'y applique plus.
%  E, A, NEA — non demandés à l'imprimeur : ils sont CALCULÉS avec la main
%        déduite ci-dessous. À faire confirmer avant tout tirage.
%
% La main de 1,206 n'est pas publiée par CoolLibri : elle est déduite de
% leur propre chiffre, 14 mm / (0,090 x 258/2) = 1,206. Elle reproduit donc
% exactement le dos de N, et n'est qu'une ESTIMATION pour les autres.
% Écarts que cela recouvre, pour mémoire : une main de 1,0 donnerait
% 11,61 mm pour N, une main de 1,3 en donnerait 16,77. Il ne fallait pas
% deviner.
%  E   — 17 mm, VALEUR DONNÉE PAR L'IMPRIMEUR le 17/09/2026 pour 276 pages,
%        20 x 24, couverture RIGIDE, reliure collée (Klebebindung), papier
%        115 g. En reliure cousue (Fadenheftung), il annonce 18 mm.
%
% \DosImposeReliure<classe> : la reliure pour laquelle la valeur a été
% donnée. Le dos d'une couverture rigide n'est pas celui d'un dos carré
% collé souple, même à pagination égale.
\def\DosImposeN{14}   \def\DosImposePagesN{258}  \def\DosImposeFormatN{a4}
\def\DosImposeE{17}   \def\DosImposePagesE{276}  \def\DosImposeFormatE{20x24-marge}
\def\DosImposeA{0}    \def\DosImposePagesA{0}    \def\DosImposeFormatA{}
\def\DosImposeNE{0}   \def\DosImposePagesNE{0}   \def\DosImposeFormatNE{}
\def\DosImposeEA{0}   \def\DosImposePagesEA{0}   \def\DosImposeFormatEA{}
\def\DosImposeNEA{0}  \def\DosImposePagesNEA{0}  \def\DosImposeFormatNEA{}
\def\DosImposeReliureN{souple}  \def\DosImposeReliureE{rigide}
\def\DosImposeReliureA{souple}  \def\DosImposeReliureNE{souple}
\def\DosImposeReliureEA{souple} \def\DosImposeReliureNEA{souple}

% Fond perdu, en mm. 3 mm est la valeur courante ; certains imprimeurs
% demandent 5 mm.
\def\FondPerdu{3}

% Marge technique autour du fond perdu, en mm. Sert à loger les traits de
% coupe en mode épreuve. Ignorée en mode production.
\def\MargeTech{12}

% Zone de sécurité, en mm : aucun texte ne doit s'en approcher davantage
% du bord de coupe.
\def\Securite{8}

% ---------------------------------------------------------------------
% 2. DONNÉES PAR CLASSE
%    pages : nombre de pages du PDF intérieur, COUVERTURE NON COMPRISE.
%    Un dos carré collé impose un nombre de pages PAIR ; le fichier
%    refuse de compiler si le compte est impair.
%
%    Ces tables ne servent que si \PagesForcees et \VersionForcee sont
%    vides, c'est-à-dire à la main. Valeurs A4 mesurées sur les PDF de la
%    a.3 (CLAUDE.md § 1, 05/09/2026). En 20 x 24, la pagination change
%    entièrement : NEA passe de 816 à 1036 pages.
% ---------------------------------------------------------------------

\def\DonneesN{262}
\def\DonneesE{214}
\def\DonneesA{386}
\def\DonneesNE{0}      % édition jamais produite
\def\DonneesEA{0}      % édition jamais produite
\def\DonneesNEA{816}

% Version imprimée sur le dos et la 1re de couverture, PAR CLASSE.
% « Les versions sont suivies par classe, chacune évoluant à son rythme »
% (CHANGELOG.md).
\def\VersionN{a.3}
\def\VersionE{a.3}
\def\VersionA{a.3}
\def\VersionNE{a.3}
\def\VersionEA{a.3}
\def\VersionNEA{a.3}

% ---------------------------------------------------------------------
% 3. TEXTES
% ---------------------------------------------------------------------

\def\TitreOuvrage{50\,Ohm}
\def\SousTitre{Préparation à l'examen radioamateur}

% Sous-titre propre à chaque classe (repris de FR_TITLES, build_book.py).
\def\SousTitreN{Cours complet — classe N}
\def\SousTitreE{Cours complet — classe E}
\def\SousTitreA{Cours complet — classe A}
\def\SousTitreNE{Cours complet — classes N et E}
\def\SousTitreEA{Cours complet — classes E et A}
\def\SousTitreNEA{Cours complet — classes N, E et A}

% Lettres du filigrane, empilées de haut en bas (cf. build_book.py v0.20).
\def\LettresN{N}
\def\LettresE{E}
\def\LettresA{A}
\def\LettresNE{N,E}
\def\LettresEA{E,A}
\def\LettresNEA{N,E,A}

% Texte de 4e de couverture.
% NB : il parle des classes N, E et A quelle que soit l'édition. Juste pour
% N, E, A et NEA ; approximatif pour NE et EA ; faux pour SWL, d'où son refus.
\def\QuatriemeTexte{%
Ce guide d'étude prépare aux examens radioamateur allemands des classes
N, E et A. Il traduit en français les contenus du projet 50ohm.de, publiés
par le référat AJW du DARC e.\,V., et les complète d'encarts « En France »
qui signalent, sans rien retirer au texte allemand, ce que la
réglementation française prévoit sur le même sujet.\par\smallskip
Les questions officielles du catalogue allemand accompagnent chaque
notion, avec leurs réponses.%
}

\def\Auteur{Traduit avec l'aide d'une IA par Pierre F4JWI}
\def\Depot{github.com/DrPLG/50ohm-fr}

% =====================================================================
%  À PARTIR D'ICI, PLUS RIEN À MODIFIER
% =====================================================================

\usepackage{iftex}
\RequireLuaTeX
\usepackage{xfp}
\usepackage{tikz}
\usetikzlibrary{calc}
\usepackage{fontspec}
\usepackage[french]{babel}
\usepackage{geometry}

% ---------------------------------------------------------------------
% 4. RÉSOLUTION DES DONNÉES DE LA CLASSE
% ---------------------------------------------------------------------

\makeatletter
% Tome SWL : refus explicite, plutôt qu'un « classe inconnue » qui laisserait
% croire à un oubli. Le texte de 4e parle des classes N, E et A.
\def\@swlmot{SWL}
\ifx\Classe\@swlmot
  \PackageError{couverture}{Pas de couverture pour le tome SWL}%
    {Sa 4e de couverture n'est pas redigee : le texte commun parle des
     classes N, E et A. C'est une decision editoriale a prendre d'abord.}%
\fi

\def\@resoudre#1#2{%
  \expandafter\let\expandafter#1\csname #2\Classe\endcsname
  \ifx#1\relax
    \PackageError{couverture}{Classe \Classe\space inconnue : pas de #2\Classe}%
      {Valeurs admises : N, E, A, NE, EA, NEA.}%
  \fi}
\@resoudre\NbPages{Donnees}
\@resoudre\SousTitreClasse{SousTitre}
\@resoudre\Lettres{Lettres}
\@resoudre\DosImpose{DosImpose}
\@resoudre\DosImposePages{DosImposePages}
\@resoudre\DosImposeFormat{DosImposeFormat}
\@resoudre\DosImposeReliure{DosImposeReliure}
\@resoudre\Version{Version}

% Valeurs imposées de l'extérieur. Test de vacuité par \detokenize : il ne
% dépend pas de la façon dont la macro a été définie.
\if\relax\detokenize\expandafter{\PagesForcees}\relax\else
  \let\NbPages\PagesForcees
\fi
\if\relax\detokenize\expandafter{\VersionForcee}\relax\else
  \let\Version\VersionForcee
\fi

% Format du livre fini.
\def\@formatAquatre{a4}
\def\@formatVingt{20x24}
\def\@formatVingtMarge{20x24-marge}
\ifx\Format\@formatAquatre
  \def\PageLargeur{210}\def\PageHauteur{297}
\else\ifx\Format\@formatVingt
  \def\PageLargeur{200}\def\PageHauteur{240}
\else\ifx\Format\@formatVingtMarge
  \def\PageLargeur{200}\def\PageHauteur{240}
\else
  \PackageError{couverture}{Format \Format\space inconnu}%
    {Valeurs admises : a4, 20x24, 20x24-marge.}%
  \def\PageLargeur{210}\def\PageHauteur{297}
\fi\fi\fi
% Reliure. Hors de « souple » et « rigide » : erreur, pas de supposition.
\def\@reliuresouple{souple}
\def\@reliurerigide{rigide}
\ifx\Reliure\@reliurerigide
  \edef\PageLargeur{\fpeval{\PageLargeur + 5}}
  \edef\PageHauteur{\fpeval{\PageHauteur + 5}}
  \def\FondPerdu{15}
\else\ifx\Reliure\@reliuresouple\else
  \PackageError{couverture}{Reliure \Reliure\space inconnue}%
    {Valeurs admises : souple, rigide.}%
\fi\fi
\makeatother

% Garde-fou : une édition jamais produite n'a pas de couverture.
\ifnum\NbPages=0
  \PackageError{couverture}{L'edition \Classe\space n'a pas de pagination}%
    {Cette edition n'a jamais ete compilee. Renseigner \string\Donnees\Classe
     ou passer \string\PagesForcees.}%
\fi

% Avertissement : l'imprimeur demande un nombre de pages multiple de 4
% (17/09/2026). Pas une erreur : c'est son exigence, pas une loi du papier.
\ifnum\fpeval{\NbPages - 4*floor(\NbPages/4)}=0 \else
  \typeout{COUVERTURE-AVERTISSEMENT \NbPages\space pages : pas un multiple
    de 4, que l'imprimeur exige -- ajouter des pages a l'interieur}%
\fi

% Garde-fou : un dos carré collé impose un nombre de pages pair.
\ifodd\NbPages
  \PackageError{couverture}{Nombre de pages impair (\NbPages)}%
    {Un dos carre colle exige un nombre PAIR de pages : une feuille en
     porte deux. Ajouter une page blanche a l'interieur.}%
\fi

% ---------------------------------------------------------------------
% 5. GÉOMÉTRIE — c'est ici que se joue l'exactitude du dos
% ---------------------------------------------------------------------

% Épaisseur du dos : valeur imposée si elle est donnée ET qu'elle vaut pour
% cette pagination et ce format ; calcul sinon.
\edef\DosCalcule{\fpeval{round(\NbPages / 2 * \Grammage * \Bouffant / 1000, 2)}}
\def\DosOrigine{calculee}
\ifnum\fpeval{\DosImpose > 0}=1
  \ifnum\fpeval{\NbPages = \DosImposePages}=1
    \ifx\Format\DosImposeFormat
      \ifx\Reliure\DosImposeReliure
        \def\DosOrigine{imposee}
      \fi
    \fi
  \fi
  \def\MotCalculee{calculee}
  \ifx\DosOrigine\MotCalculee
    \typeout{COUVERTURE-AVERTISSEMENT dos impose de \DosImpose\space mm
      donne pour \DosImposePages\space pages en \DosImposeFormat\space
      (reliure \DosImposeReliure) -- inapplicable a \NbPages\space pages en
      \Format\space (reliure \Reliure), dos calcule}%
  \fi
\fi
% NB : pas de \@ dans ces noms. Cette section est hors \makeatletter, et
% « \def\@imposee{...} » y redéfinit \@ (suivi du texte « imposee ») : le
% test échouait toujours, en silence. Constaté le 16/09/2026 — le dos de N
% tombait juste par coïncidence, 14,0 mm calculés pour 14 imposés.
\def\MotRigide{rigide}
\def\MotImposee{imposee}
\ifx\DosOrigine\MotImposee
  \edef\Dos{\DosImpose}
\else
  \edef\Dos{\DosCalcule}
\fi

% Marge technique : uniquement en mode épreuve.
\def\ouimot{oui}
\ifx\Reperes\ouimot \edef\Tech{\MargeTech}\else \def\Tech{0}\fi

% Dimensions du papier.
%   largeur = marge + fond perdu + 4e + dos + 1re + fond perdu + marge
%   hauteur = marge + fond perdu + hauteur + fond perdu + marge
\edef\PapierL{\fpeval{2*\PageLargeur + \Dos + 2*\FondPerdu + 2*\Tech}}
\edef\PapierH{\fpeval{\PageHauteur + 2*\FondPerdu + 2*\Tech}}

\geometry{paperwidth=\PapierL mm, paperheight=\PapierH mm,
          margin=0pt, nohead, nofoot}

% ---------------------------------------------------------------------
%  TOUTES LES COTES SONT FIGÉES ICI, ET NULLE PART AILLEURS.
%
%  Raison, payée par trois compilations ratées : un \fpeval laissé dans
%  une coordonnée ou dans le texte d'un nœud TikZ n'est pas développé au
%  moment où TikZ analyse le chemin. On obtient « Giving up on this path.
%  Did you forget a semicolon? », qui accuse une ponctuation alors que la
%  faute est ailleurs. Le corps du document ci-dessous n'emploie donc que
%  des macros déjà calculées.
% ---------------------------------------------------------------------

% Repères horizontaux, mesurés depuis le bord GAUCHE du papier.
\edef\CoupeG{\fpeval{\Tech + \FondPerdu}}          % coupe gauche (4e de couv.)
\edef\DosG{\fpeval{\CoupeG + \PageLargeur}}         % dos, début
\edef\DosD{\fpeval{\DosG + \Dos}}                   % dos, fin
\edef\CoupeD{\fpeval{\DosD + \PageLargeur}}         % coupe droite (1re de couv.)
% Repères verticaux, depuis le bord BAS.
\edef\CoupeB{\fpeval{\Tech + \FondPerdu}}
\edef\CoupeH{\fpeval{\CoupeB + \PageHauteur}}
% Axes.
\edef\DosAxe{\fpeval{(\DosG + \DosD)/2}}
\edef\MilieuH{\fpeval{(\CoupeB + \CoupeH)/2}}
\edef\MilieuL{\fpeval{\PapierL/2}}

% Format fini, pour les cotes affichées.
\edef\FiniL{\fpeval{2*\PageLargeur + \Dos}}

% Bandeau de la 1re de couverture : même proportion que la page de titre.
\edef\BandeauL{\fpeval{round(0.34*\PageLargeur, 2)}}
\edef\BandeauG{\fpeval{\CoupeD - \BandeauL}}
\edef\BandeauAccent{\fpeval{\BandeauG - 0.7}}
\edef\DosAccentG{\fpeval{\DosG - 0.7}}
\edef\DosAccentD{\fpeval{\DosD + 0.7}}
% Couverture rigide : le bleu du dos déborde de 2 mm sur chaque plat, pour
% absorber la tolérance de l'imprimeur sur l'épaisseur du dos (±10 %).
\edef\DosEtenduG{\fpeval{\DosG - 2}}
\edef\DosEtenduD{\fpeval{\DosD + 2}}

% Filigrane de classe.
\edef\FiligraneX{\fpeval{(\BandeauG + \CoupeD)/2}}
\edef\FiligraneY{\fpeval{\CoupeH - 6}}

% Bloc-titre de la 1re de couverture.
\edef\TitreX{\fpeval{\DosD + 12}}
\edef\TitreY{\fpeval{\CoupeH - 72}}
\edef\TitreW{\fpeval{\BandeauG - \DosD - 16}}
\edef\VersionY{\fpeval{\CoupeB + 14}}

% 4e de couverture.
\edef\QuatreX{\fpeval{\CoupeG + \Securite + 6}}
\edef\QuatreY{\fpeval{\CoupeH - \Securite - 14}}
\edef\QuatreBasY{\fpeval{\CoupeB + \Securite + 22}}
\edef\QuatreW{\fpeval{\PageLargeur - 2*\Securite - 12}}

% Zone de sécurité.
\edef\SecG{\fpeval{\CoupeG + \Securite}}
\edef\SecD{\fpeval{\CoupeD - \Securite}}
\edef\SecB{\fpeval{\CoupeB + \Securite}}
\edef\SecH{\fpeval{\CoupeH - \Securite}}

% Zones de contrôle : ce qui doit rester de part et d'autre du pli du dos.
\edef\PliQuatre{\fpeval{\DosG - \Securite}}         % 4e : limite côté dos
\edef\PliPremiere{\fpeval{\DosD + \Securite}}       % 1re : limite côté dos

% Traits de coupe et cotes, dans la marge technique.
\edef\TechExt{\fpeval{\Tech - 8}}
\edef\TechInt{\fpeval{\Tech - 1}}
\edef\TechCote{\fpeval{\Tech - 9}}
\edef\HautInt{\fpeval{\PapierH - \Tech + 1}}
\edef\HautExt{\fpeval{\PapierH - \Tech + 8}}
\edef\HautCote{\fpeval{\PapierH - \Tech + 9}}
\edef\DroitInt{\fpeval{\PapierL - \Tech + 1}}
\edef\DroitExt{\fpeval{\PapierL - \Tech + 8}}

% Emplacement du code-barres ISBN : en bas à DROITE de la 4e de couverture,
% selon l'usage éditorial, et non à gauche où il chevaucherait les mentions
% de licence.
\edef\IsbnD{\fpeval{\DosG - \Securite}}
\edef\IsbnG{\fpeval{\IsbnD - 40}}
\edef\IsbnH{\fpeval{\SecB + 25}}
\edef\IsbnCX{\fpeval{\IsbnG + 20}}
\edef\IsbnCY{\fpeval{\SecB + 12.5}}

% Cartouches de cotes, dans la marge technique. Ancrés vers l'INTÉRIEUR du
% papier : ancrés vers l'extérieur, ils débordaient et étaient rognés.
\edef\CoteBasY{\fpeval{1}}
\edef\CoteHautY{\fpeval{\PapierH - 1}}

% ---------------------------------------------------------------------
% 6. COULEURS ET POLICES — reprises de la page de titre (build_book.py)
% ---------------------------------------------------------------------

\definecolor{TitleBand}{cmyk}{.9,.55,.1,.35}   % bleu profond
\definecolor{TitleAccent}{cmyk}{.05,.8,1,0}    % rouge DARC (accent fin)
\definecolor{Repere}{cmyk}{0,1,1,0}            % magenta : repères de montage

\setmainfont{Libertinus Serif}[
  ItalicFont     = Libertinus Serif Italic,
  BoldFont       = Libertinus Serif Bold,
  BoldItalicFont = Libertinus Serif Bold Italic]

\pagestyle{empty}

% ---------------------------------------------------------------------
% 7. FILIGRANE DE CLASSE
%    Le texte est construit ici, et non dans le nœud : un \\ produit à
%    l'intérieur d'un \foreach referme le groupe d'alignement de TikZ.
%    Une lettre seule est grande ; plusieurs s'empilent et se réduisent
%    (build_book.py v0.20 : « NEA » à 105 pt débordait encore).
% ---------------------------------------------------------------------

\makeatletter
\def\FiligraneTexte{}
\newif\ifpremierelettre \premierelettretrue
\count255=0
\@for\@tempa:=\Lettres\do{%
  \advance\count255 by 1\relax
  \ifpremierelettre \premierelettrefalse
  \else \g@addto@macro\FiligraneTexte{\\}\fi
  \expandafter\g@addto@macro\expandafter\FiligraneTexte\expandafter{\@tempa}}%
\xdef\NbLettres{\the\count255}
\makeatother

\ifnum\NbLettres=1 \def\CorpsF{150}\else\def\CorpsF{86}\fi
\edef\CorpsFI{\fpeval{round(\CorpsF*1.02, 0)}}

% Texte du cartouche de cotes, lui aussi construit hors du nœud TikZ.
%
% NE PAS y employer \textsuperscript : dans un nœud « align=center », il
% casse le chemin TikZ (« Giving up on this path. Did you forget a
% semicolon? »), et l'erreur est signalée sur le nœud SUIVANT — elle
% accuse donc une ponctuation là où la faute est ici. Isolé par essais :
% $^2$ et le caractère ² passent, \textsuperscript non.
\def\CotesTexte{%
  \Classe\ --- \NbPages\ pages --- papier \Grammage\ g/m$^2$,
  bouffant \Bouffant\ --- fond perdu \FondPerdu\ mm\\
  papier \PapierL\ mm sur \PapierH\ mm --- fini \FiniL\ mm sur
  \PageHauteur\ mm\\
  EPREUVE, ne pas envoyer a l'impression%
}

% ---------------------------------------------------------------------
% 8. CONTRÔLE DE SUPERPOSITION
%
%  \MesureNoeud{nœud}{préfixe}{retrait} mesure la boîte d'un nœud par ses
%  quatre coins (un nœud pivoté a ses ancres pivotées : le min et le max des
%  quatre suffisent) et la ramène au repère du papier, en mm. Le point (0,0)
%  est lu par le même « let » que les ancres : les deux passent donc par la
%  même transformation, et la comparaison ne dépend pas de l'endroit où
%  TeX a posé l'image sur la page.
%
%  Le retrait sert au seul filigrane, pour mesurer les lettres et non leur
%  cadre : c'est le bord des lettres que le détourage rognerait.
%
%  PIÈGE MESURÉ le 16/09/2026 : la marge interne TikZ (.3333em) n'est PAS
%  évaluée dans la police du nœud, mais dans celle qui entoure le dessin.
%  Instrumenté : dans le nœud, em = 150 pt ; boîte brute du « N » :
%  41,7 x 37,5 mm, soit la lettre plus 2 x 1,3 mm — 1,3 mm étant .3333em
%  à 11 pt. Un retrait calculé à 150 pt (17,6 mm) ramenait la lettre à
%  6,6 mm de large, et le contrôle aurait dit OK sur une mesure absurde.
% ---------------------------------------------------------------------

\newcount\NbAlertes
\newdimen\FiligraneSep
\newsavebox\CreditsBoite

% Crédits de 4e : lignes coupées à la main. Leur largeur RÉELLE est celle
% de la plus longue ligne, mesurée dans une boîte à part (voir § 8, ISBN).
\def\CreditsTexte{%
    Réalisé à partir des contenus de 50ohm.de (en allemand)\\
    50ohm.de-Autorenteam, coordonné par le référat AJW du DARC e.\,V.\\
    \Auteur\\[3pt]
    Licence CC BY 4.0 --- \texttt{\Depot}}

\def\MesureNoeud#1#2#3{%
  \path let \p1=(#1.south west), \p2=(#1.north east),
            \p3=(#1.north west), \p4=(#1.south east), \p5=(0,0) in
    \pgfextra{%
      \expandafter\xdef\csname #2xmin\endcsname{\fpeval{round((min(\x1,\x2,\x3,\x4) - \x5 + #3)/1mm, 2)}}%
      \expandafter\xdef\csname #2xmax\endcsname{\fpeval{round((max(\x1,\x2,\x3,\x4) - \x5 - #3)/1mm, 2)}}%
      \expandafter\xdef\csname #2ymin\endcsname{\fpeval{round((min(\y1,\y2,\y3,\y4) - \y5 + #3)/1mm, 2)}}%
      \expandafter\xdef\csname #2ymax\endcsname{\fpeval{round((max(\y1,\y2,\y3,\y4) - \y5 - #3)/1mm, 2)}}%
    };}

% \ControleDans{libellé}{préfixe}{xmin}{xmax}{ymin}{ymax}
\def\ControleDans#1#2#3#4#5#6{%
  \ifnum\fpeval{\csname #2xmin\endcsname >= #3 && \csname #2xmax\endcsname <= #4
             && \csname #2ymin\endcsname >= #5 && \csname #2ymax\endcsname <= #6}=1
    \typeout{COUVERTURE-CONTROLE OK #1 : x \csname #2xmin\endcsname-\csname #2xmax\endcsname,
      y \csname #2ymin\endcsname-\csname #2ymax\endcsname\space dans x #3-#4, y #5-#6}%
  \else
    \global\advance\NbAlertes by 1
    \typeout{COUVERTURE-CONTROLE ALERTE #1 : x \csname #2xmin\endcsname-\csname #2xmax\endcsname,
      y \csname #2ymin\endcsname-\csname #2ymax\endcsname\space HORS de x #3-#4, y #5-#6}%
  \fi}

% \ControleDisjoints{libellé}{préfixe}{xmin}{xmax}{ymin}{ymax}
% Deux rectangles se chevauchent si et seulement s'ils se recouvrent sur
% les DEUX axes à la fois.
\def\ControleDisjoints#1#2#3#4#5#6{%
  \ifnum\fpeval{\csname #2xmin\endcsname < #4 && #3 < \csname #2xmax\endcsname
             && \csname #2ymin\endcsname < #6 && #5 < \csname #2ymax\endcsname}=1
    \global\advance\NbAlertes by 1
    \typeout{COUVERTURE-CONTROLE ALERTE #1 : chevauchement}%
  \else
    \typeout{COUVERTURE-CONTROLE OK #1 : disjoints}%
  \fi}

\begin{document}
% Marge interne du filigrane, évaluée dans la police qui ENTOURE le dessin :
% c'est là que TikZ évalue le .3333em de « inner sep » (cf. § 8).
\FiligraneSep=.3333em
% Largeur réelle des crédits : la plus longue ligne, dans leur police.
\sbox\CreditsBoite{\small\begin{tabular}{@{}l@{}}\CreditsTexte\end{tabular}}%
\begin{tikzpicture}[remember picture, overlay,
                    x=1mm, y=1mm, shift={(current page.south west)}]

% ---- Fonds ----------------------------------------------------------
% La 4e de couverture reste sur le blanc du papier.
% Le bandeau bleu couvre le dos, puis reprend au tiers extérieur de la
% 1re de couverture — même partition 2/3-1/3 que la page de titre.
% Couverture rigide : pas de filets au dos (ils tomberaient à côté des
% plis, vu la tolérance), et un dos bleu élargi de 2 mm de chaque côté.
\ifx\Reliure\MotRigide
  \fill[TitleBand] (\DosEtenduG, 0) rectangle (\DosEtenduD, \PapierH);
\else
  \fill[TitleBand] (\DosG, 0) rectangle (\DosD, \PapierH);
\fi
\fill[TitleBand] (\BandeauG, 0) rectangle (\PapierL, \PapierH);
% Filets d'accent : à la jonction sur la 1re, de part et d'autre du dos.
\fill[TitleAccent] (\BandeauAccent, 0) rectangle (\BandeauG, \PapierH);
\ifx\Reliure\MotRigide\else
  \fill[TitleAccent] (\DosAccentG, 0) rectangle (\DosG, \PapierH);
  \fill[TitleAccent] (\DosD, 0) rectangle (\DosAccentD, \PapierH);
\fi

% ---- 1re de couverture ----------------------------------------------
% Filigrane de classe, en réserve claire, détouré au bandeau.
\begin{scope}
\clip (\BandeauG, 0) rectangle (\PapierL, \PapierH);
\node[name=filigrane, anchor=north, align=center, text=white, opacity=0.16,
      font=\fontsize{\CorpsF}{\CorpsFI}\selectfont\bfseries]
  at (\FiligraneX, \FiligraneY) {\FiligraneTexte};
\end{scope}

\node[name=titre, anchor=west, align=left, text width=\TitreW mm]
  at (\TitreX, \TitreY) {%
    {\fontsize{54}{56}\selectfont\bfseries \TitreOuvrage}\\[10pt]
    {\Large\color{TitleBand}\bfseries \SousTitre}\\[16pt]
    {\large\color{black!75} \SousTitreClasse}};

\node[name=version, anchor=south west, align=left, font=\small, text=black!70]
  at (\TitreX, \VersionY) {Version \Version\\ Licence CC BY 4.0};

% ---- Dos -------------------------------------------------------------
% Un dos étroit ne reçoit pas de texte : en deçà de 6 mm, le débord au
% massicot ferait mordre le texte vertical sur les plats.
\ifnum\fpeval{\Dos >= 6}=1
  \node[name=dos, rotate=90, anchor=center, text=white, font=\bfseries,
        inner sep=0pt]
    at (\DosAxe, \MilieuH)
    {\TitreOuvrage\ \ --- \ \SousTitreClasse\ \ --- \ \Version};
\fi

% ---- 4e de couverture ------------------------------------------------
\node[name=quatre, anchor=north west, align=justify, text width=\QuatreW mm]
  at (\QuatreX, \QuatreY) {\QuatriemeTexte};

% Le texte vit dans \CreditsTexte (§ 8) pour être mesuré à part : la boîte
% du nœud a la largeur fixe \QuatreW et couvre l'emplacement ISBN sans
% qu'aucun caractère n'y soit. Retirer « text width » pour l'éviter change
% l'interlignage — essayé le 16/09/2026, 6 934 pixels différents.
\node[name=credits, anchor=south west, align=left, font=\small, text=black!70,
      text width=\QuatreW mm]
  at (\QuatreX, \QuatreBasY) {\CreditsTexte};

% ---- Repères de fabrication (mode épreuve uniquement) ----------------
\ifx\Reperes\ouimot
  % Traits de coupe, hors du fond perdu.
  \foreach \x in {\CoupeG, \CoupeD} {
    \draw[Repere, line width=0.25pt] (\x, \TechExt) -- (\x, \TechInt);
    \draw[Repere, line width=0.25pt] (\x, \HautInt) -- (\x, \HautExt);}
  \foreach \y in {\CoupeB, \CoupeH} {
    \draw[Repere, line width=0.25pt] (\TechExt, \y) -- (\TechInt, \y);
    \draw[Repere, line width=0.25pt] (\DroitInt, \y) -- (\DroitExt, \y);}
  % Limites du dos, en pointillés.
  \foreach \x in {\DosG, \DosD} {
    \draw[Repere, line width=0.25pt, dashed] (\x, \TechExt) -- (\x, \TechInt);
    \draw[Repere, line width=0.25pt, dashed] (\x, \HautInt) -- (\x, \HautExt);}
  % Zone de sécurité.
  \draw[Repere, line width=0.2pt, dotted]
    (\SecG, \SecB) rectangle (\SecD, \SecH);
  % Emplacement réservé au code-barres ISBN.
  \draw[Repere, line width=0.2pt] (\IsbnG, \SecB) rectangle (\IsbnD, \IsbnH);
  \node[font=\tiny\sffamily, text=Repere, align=center]
    at (\IsbnCX, \IsbnCY) {emplacement\\ISBN};
  % Cotes.
  \node[anchor=north, font=\tiny\sffamily, text=Repere]
    at (\DosAxe, \CoteHautY) {dos \Dos\ mm (\DosOrigine)};
  \node[anchor=south, font=\tiny\sffamily, text=Repere, align=center]
    at (\MilieuL, \CoteBasY) {\CotesTexte};
\fi

% ---- Contrôle de superposition (§ 8) ----------------------------------
% Rien n'est dessiné : des mesures et des lignes de journal.
\MesureNoeud{filigrane}{Fil}{\the\FiligraneSep}
\MesureNoeud{titre}{Tit}{0pt}
\MesureNoeud{version}{Ver}{0pt}
\MesureNoeud{quatre}{Qua}{0pt}
\MesureNoeud{credits}{Cre}{0pt}
% Bord droit des crédits = bord gauche + marge interne + plus longue ligne.
\xdef\Crexmax{\fpeval{round(\Crexmin + (\the\FiligraneSep + \the\wd\CreditsBoite)/1mm, 2)}}

% Le filigrane peut entrer dans la zone de sécurité (il est décoratif),
% pas franchir la coupe ni sortir du bandeau — le détourage le rognerait.
\ControleDans{filigrane dans le bandeau}{Fil}{\BandeauG}{\CoupeD}{\CoupeB}{\CoupeH}
\ControleDans{titre sur la 1re, hors bandeau}{Tit}{\PliPremiere}{\BandeauAccent}{\SecB}{\SecH}
\ControleDans{version sur la 1re, hors bandeau}{Ver}{\PliPremiere}{\BandeauAccent}{\SecB}{\SecH}
\ControleDans{texte de 4e dans la zone de securite}{Qua}{\SecG}{\PliQuatre}{\SecB}{\SecH}
\ControleDans{credits de 4e dans la zone de securite}{Cre}{\SecG}{\PliQuatre}{\SecB}{\SecH}
\ControleDisjoints{titre / version}{Tit}{\Verxmin}{\Verxmax}{\Verymin}{\Verymax}
\ControleDisjoints{texte de 4e / credits}{Qua}{\Crexmin}{\Crexmax}{\Creymin}{\Creymax}
\ControleDisjoints{texte de 4e / emplacement ISBN}{Qua}{\IsbnG}{\IsbnD}{\SecB}{\IsbnH}
\ControleDisjoints{credits / emplacement ISBN}{Cre}{\IsbnG}{\IsbnD}{\SecB}{\IsbnH}
\ifnum\fpeval{\Dos >= 6}=1
  \MesureNoeud{dos}{Dos}{0pt}
  \ControleDans{texte du dos entre les plis}{Dos}{\DosG}{\DosD}{\SecB}{\SecH}
\else
  \typeout{COUVERTURE-CONTROLE OK dos de \Dos\space mm : trop etroit, sans texte}%
\fi

\end{tikzpicture}
\typeout{COUVERTURE-BILAN alertes=\the\NbAlertes\space; classe=\Classe\space;
  format=\Format\space; reliure=\Reliure\space; pages=\NbPages\space;
  version=\Version\space;
  dos=\Dos\space mm (\DosOrigine); fini=\FiniL x \PageHauteur\space mm;
  papier=\PapierL x \PapierH\space mm}
\end{document}
"
%
%  Pour le fichier destiné à l'imprimeur (sans repères) :
%
%      lualatex "\def\Reperes{non}% =====================================================================
%  couverture.tex — couverture dos carré collé pour les livres 50ohm-fr
% =====================================================================
%
%  Produit une couverture d'un seul tenant : 4e de couverture | dos |
%  1re de couverture, aux dimensions exactes du livre relié, fond perdu
%  compris. L'épaisseur du dos est CALCULÉE à partir du nombre de pages,
%  du grammage et du bouffant du papier.
%
%  ORIGINE ET ÉTAT. Écrit le 20/08/2026 dans Fichierstravail/ (hors git),
%  versé au dépôt le 16/09/2026 pour que fenetre_compilation.py puisse
%  l'appeler. Ce qui a changé à cette occasion :
%
%    - le FORMAT est un paramètre (\Format = a4, 20x24 ou 20x24-marge,
%      mêmes noms que --format de build_book.py) ;
%    - la pagination et la version peuvent être IMPOSÉES de l'extérieur
%      (\PagesForcees, \VersionForcee) : la fenêtre lit la pagination dans
%      le journal du livre qu'elle vient de compiler, au lieu de se fier à
%      une table qui vieillit ;
%    - un dos imposé par l'imprimeur n'est appliqué QUE pour la pagination
%      et le format pour lesquels il a été donné ;
%    - un CONTRÔLE DE SUPERPOSITION est calculé par LaTeX lui-même, à partir
%      des boîtes réelles des nœuds, et écrit dans le journal ;
%    - le tome SWL est refusé explicitement : sa 4e de couverture n'est pas
%      rédigée (le texte commun parle des classes N, E et A).
%
%    Neutralité vérifiée : en A4, pour N (258 p., a.2) et NEA (816 p., a.3),
%    en épreuve comme en impression, le rendu est identique au pixel près à
%    celui du gabarit d'origine.
%
%  COMPILATION (LuaLaTeX, depuis la racine du dépôt) :
%
%      lualatex couverture.tex
%      lualatex couverture.tex        <-- DEUX FOIS, obligatoirement
%
%  DEUX PASSES SONT NÉCESSAIRES. Le dessin s'ancre sur « current page »,
%  dont la position n'est connue qu'après un premier passage : à la
%  première compilation, tout se tasse dans le coin supérieur et la
%  couverture paraît ratée alors qu'elle ne l'est pas. Aucune erreur n'est
%  émise — seul le rendu le montre. Le contrôle de superposition, lui
%  aussi, ne vaut qu'à la seconde passe.
%
%  Le plus simple est de passer par fenetre_compilation.py (case
%  « couverture pour l'imprimeur »), qui enchaîne les deux passes, produit
%  l'épreuve ET le fichier d'impression, et relit le contrôle.
%
%  À la main, pour une autre classe, un autre format, une pagination mesurée :
%
%      lualatex "\def\Classe{A}\def\Format{20x24}\def\PagesForcees{482}% =====================================================================
%  couverture.tex — couverture dos carré collé pour les livres 50ohm-fr
% =====================================================================
%
%  Produit une couverture d'un seul tenant : 4e de couverture | dos |
%  1re de couverture, aux dimensions exactes du livre relié, fond perdu
%  compris. L'épaisseur du dos est CALCULÉE à partir du nombre de pages,
%  du grammage et du bouffant du papier.
%
%  ORIGINE ET ÉTAT. Écrit le 20/08/2026 dans Fichierstravail/ (hors git),
%  versé au dépôt le 16/09/2026 pour que fenetre_compilation.py puisse
%  l'appeler. Ce qui a changé à cette occasion :
%
%    - le FORMAT est un paramètre (\Format = a4, 20x24 ou 20x24-marge,
%      mêmes noms que --format de build_book.py) ;
%    - la pagination et la version peuvent être IMPOSÉES de l'extérieur
%      (\PagesForcees, \VersionForcee) : la fenêtre lit la pagination dans
%      le journal du livre qu'elle vient de compiler, au lieu de se fier à
%      une table qui vieillit ;
%    - un dos imposé par l'imprimeur n'est appliqué QUE pour la pagination
%      et le format pour lesquels il a été donné ;
%    - un CONTRÔLE DE SUPERPOSITION est calculé par LaTeX lui-même, à partir
%      des boîtes réelles des nœuds, et écrit dans le journal ;
%    - le tome SWL est refusé explicitement : sa 4e de couverture n'est pas
%      rédigée (le texte commun parle des classes N, E et A).
%
%    Neutralité vérifiée : en A4, pour N (258 p., a.2) et NEA (816 p., a.3),
%    en épreuve comme en impression, le rendu est identique au pixel près à
%    celui du gabarit d'origine.
%
%  COMPILATION (LuaLaTeX, depuis la racine du dépôt) :
%
%      lualatex couverture.tex
%      lualatex couverture.tex        <-- DEUX FOIS, obligatoirement
%
%  DEUX PASSES SONT NÉCESSAIRES. Le dessin s'ancre sur « current page »,
%  dont la position n'est connue qu'après un premier passage : à la
%  première compilation, tout se tasse dans le coin supérieur et la
%  couverture paraît ratée alors qu'elle ne l'est pas. Aucune erreur n'est
%  émise — seul le rendu le montre. Le contrôle de superposition, lui
%  aussi, ne vaut qu'à la seconde passe.
%
%  Le plus simple est de passer par fenetre_compilation.py (case
%  « couverture pour l'imprimeur »), qui enchaîne les deux passes, produit
%  l'épreuve ET le fichier d'impression, et relit le contrôle.
%
%  À la main, pour une autre classe, un autre format, une pagination mesurée :
%
%      lualatex "\def\Classe{A}\def\Format{20x24}\def\PagesForcees{482}% =====================================================================
%  couverture.tex — couverture dos carré collé pour les livres 50ohm-fr
% =====================================================================
%
%  Produit une couverture d'un seul tenant : 4e de couverture | dos |
%  1re de couverture, aux dimensions exactes du livre relié, fond perdu
%  compris. L'épaisseur du dos est CALCULÉE à partir du nombre de pages,
%  du grammage et du bouffant du papier.
%
%  ORIGINE ET ÉTAT. Écrit le 20/08/2026 dans Fichierstravail/ (hors git),
%  versé au dépôt le 16/09/2026 pour que fenetre_compilation.py puisse
%  l'appeler. Ce qui a changé à cette occasion :
%
%    - le FORMAT est un paramètre (\Format = a4, 20x24 ou 20x24-marge,
%      mêmes noms que --format de build_book.py) ;
%    - la pagination et la version peuvent être IMPOSÉES de l'extérieur
%      (\PagesForcees, \VersionForcee) : la fenêtre lit la pagination dans
%      le journal du livre qu'elle vient de compiler, au lieu de se fier à
%      une table qui vieillit ;
%    - un dos imposé par l'imprimeur n'est appliqué QUE pour la pagination
%      et le format pour lesquels il a été donné ;
%    - un CONTRÔLE DE SUPERPOSITION est calculé par LaTeX lui-même, à partir
%      des boîtes réelles des nœuds, et écrit dans le journal ;
%    - le tome SWL est refusé explicitement : sa 4e de couverture n'est pas
%      rédigée (le texte commun parle des classes N, E et A).
%
%    Neutralité vérifiée : en A4, pour N (258 p., a.2) et NEA (816 p., a.3),
%    en épreuve comme en impression, le rendu est identique au pixel près à
%    celui du gabarit d'origine.
%
%  COMPILATION (LuaLaTeX, depuis la racine du dépôt) :
%
%      lualatex couverture.tex
%      lualatex couverture.tex        <-- DEUX FOIS, obligatoirement
%
%  DEUX PASSES SONT NÉCESSAIRES. Le dessin s'ancre sur « current page »,
%  dont la position n'est connue qu'après un premier passage : à la
%  première compilation, tout se tasse dans le coin supérieur et la
%  couverture paraît ratée alors qu'elle ne l'est pas. Aucune erreur n'est
%  émise — seul le rendu le montre. Le contrôle de superposition, lui
%  aussi, ne vaut qu'à la seconde passe.
%
%  Le plus simple est de passer par fenetre_compilation.py (case
%  « couverture pour l'imprimeur »), qui enchaîne les deux passes, produit
%  l'épreuve ET le fichier d'impression, et relit le contrôle.
%
%  À la main, pour une autre classe, un autre format, une pagination mesurée :
%
%      lualatex "\def\Classe{A}\def\Format{20x24}\def\PagesForcees{482}\input{couverture.tex}"
%
%  Pour le fichier destiné à l'imprimeur (sans repères) :
%
%      lualatex "\def\Reperes{non}\input{couverture.tex}"
%
%  ---------------------------------------------------------------------
%  DEUX MODES, un seul paramètre : \Reperes
%
%    oui  — mode ÉPREUVE. Ajoute une marge technique, les traits de coupe,
%           les limites du dos, la zone de sécurité et les cotes. C'est ce
%           qu'il faut pour vérifier à l'écran.
%    non  — mode PRODUCTION. Papier = format final + fond perdu, rien
%           d'autre. C'est ce fichier-là qu'on envoie à l'imprimeur.
%
%  ---------------------------------------------------------------------
%  CONTRÔLE DE SUPERPOSITION
%
%  À la fin du dessin, chaque bloc de texte est mesuré par ses ancres TikZ
%  (quatre coins, ce qui vaut aussi pour le texte pivoté du dos) et comparé
%  à la zone où il doit tenir, puis aux blocs voisins. Le journal porte une
%  ligne par vérification :
%
%      COUVERTURE-CONTROLE OK      <quoi> : ...
%      COUVERTURE-CONTROLE ALERTE  <quoi> : ...
%
%  et un bilan : COUVERTURE-BILAN alertes=<n> ; ...
%
%  Ce que le contrôle VOIT : un bloc qui franchit le pli du dos, la coupe,
%  la zone de sécurité ou le bord du bandeau ; deux blocs qui se chevauchent ;
%  un bloc qui empiète sur l'emplacement réservé au code-barres ISBN.
%  Ce qu'il NE VOIT PAS : une ligne trop longue à l'INTÉRIEUR d'un bloc à
%  largeur fixe (la boîte garde sa largeur, le texte déborde). Celle-là se
%  signale par un « Overfull \hbox » dans le même journal, que la fenêtre
%  compte aussi.
%
%  ---------------------------------------------------------------------
%  AVERTISSEMENT SUR L'ÉPAISSEUR DU DOS
%
%  Le dos se calcule ainsi :
%
%      épaisseur d'une feuille (mm) = grammage x bouffant / 1000
%      dos (mm)                     = (nombre de pages / 2) x épaisseur
%
%  On divise par deux parce qu'une feuille porte DEUX pages.
%
%  Le « bouffant » (main, volume, bulk) dépend du papier et n'est pas
%  devinable : un offset 80 g courant est autour de 1,30, un bouffant
%  d'édition monte à 1,75 et plus. UN ÉCART DE BOUFFANT DÉPLACE LE DOS DE
%  PLUSIEURS MILLIMÈTRES. Demander à l'imprimeur soit son bouffant, soit
%  directement l'épaisseur du dos, et reporter ci-dessous. Tant que cette
%  valeur n'est pas confirmée, la couverture n'est pas bonne à tirer.
%
%  Si l'imprimeur donne le dos en clair, forcer la valeur avec \DosImpose.
% =====================================================================

\documentclass[11pt]{article}

% ---------------------------------------------------------------------
% 1. PARAMÈTRES — c'est la seule zone à modifier
% ---------------------------------------------------------------------

% Classe de l'ouvrage : N, E, A, NE, EA ou NEA.
% (\providecommand : une valeur passée en ligne de commande l'emporte.)
\providecommand{\Classe}{NEA}

% Mode épreuve (oui) ou production (non).
\providecommand{\Reperes}{oui}

% Format du livre FINI : a4, 20x24 ou 20x24-marge — les noms de --format
% dans build_book.py. Les deux maquettes 20 x 24 ne diffèrent que par leur
% colonne de marge intérieure : le livre fini, donc la couverture, est le
% même (200 x 240 mm).
\providecommand{\Format}{a4}

% Reliure : « souple » (dos carré collé, fond perdu 3 mm) ou « rigide »
% (couverture rigide, dite hardcover). En rigide, consignes de l'imprimeur
% du 17/09/2026 : chaque plat déborde du livre de 5 mm en largeur et en
% hauteur (le carton dépasse du bloc), et le fichier porte 15 mm de rabat
% autour du carton au lieu de 3 mm de fond perdu. Les aplats vont déjà au
% bord du papier : ils remplissent le rabat d'eux-mêmes.
\providecommand{\Reliure}{souple}

% Pagination et version imposées de l'extérieur. Vides = tables du § 2.
% La fenêtre les renseigne toujours : la pagination vient du journal du
% livre réellement compilé, dans le format demandé.
\providecommand{\PagesForcees}{}
\providecommand{\VersionForcee}{}

% Papier de l'INTÉRIEUR du livre (pas de la couverture).
% CoolLibri, papier « Standard 90 g blanc ».
\def\Grammage{90}        % g/m²
\def\Bouffant{1.206}     % main DÉDUITE, voir ci-dessous

% Épaisseur du dos IMPOSÉE par l'imprimeur, en mm, par classe.
% Une valeur > 0 prime sur le calcul — MAIS SEULEMENT pour la pagination et
% le format pour lesquels l'imprimeur l'a donnée (\DosImposePages<classe>,
% \DosImposeFormat<classe>). Hors de ce cas, le dos est calculé et le
% journal le dit : un dos de 14 mm donné pour 258 pages A4 ne vaut rien pour
% 262 pages, ni pour le même livre recomposé en 20 x 24.
%
% ATTENTION : cette valeur est propre à CHAQUE ÉDITION. Un dos unique
% appliqué aux quatre livres serait faux pour trois d'entre eux.
%
%  N   — 14 mm, VALEUR DONNÉE PAR COOLIBRI pour 258 pages en 90 g, A4.
%        Fichier attendu : 440 x 303 mm, ce que ce gabarit produit en mode
%        production (2 x 210 + 14 + 2 x 3 = 440 ; 297 + 2 x 3 = 303).
%        Le N de la a.3 compte 262 pages : cette valeur ne s'y applique plus.
%  E, A, NEA — non demandés à l'imprimeur : ils sont CALCULÉS avec la main
%        déduite ci-dessous. À faire confirmer avant tout tirage.
%
% La main de 1,206 n'est pas publiée par CoolLibri : elle est déduite de
% leur propre chiffre, 14 mm / (0,090 x 258/2) = 1,206. Elle reproduit donc
% exactement le dos de N, et n'est qu'une ESTIMATION pour les autres.
% Écarts que cela recouvre, pour mémoire : une main de 1,0 donnerait
% 11,61 mm pour N, une main de 1,3 en donnerait 16,77. Il ne fallait pas
% deviner.
%  E   — 17 mm, VALEUR DONNÉE PAR L'IMPRIMEUR le 17/09/2026 pour 276 pages,
%        20 x 24, couverture RIGIDE, reliure collée (Klebebindung), papier
%        115 g. En reliure cousue (Fadenheftung), il annonce 18 mm.
%
% \DosImposeReliure<classe> : la reliure pour laquelle la valeur a été
% donnée. Le dos d'une couverture rigide n'est pas celui d'un dos carré
% collé souple, même à pagination égale.
\def\DosImposeN{14}   \def\DosImposePagesN{258}  \def\DosImposeFormatN{a4}
\def\DosImposeE{17}   \def\DosImposePagesE{276}  \def\DosImposeFormatE{20x24-marge}
\def\DosImposeA{0}    \def\DosImposePagesA{0}    \def\DosImposeFormatA{}
\def\DosImposeNE{0}   \def\DosImposePagesNE{0}   \def\DosImposeFormatNE{}
\def\DosImposeEA{0}   \def\DosImposePagesEA{0}   \def\DosImposeFormatEA{}
\def\DosImposeNEA{0}  \def\DosImposePagesNEA{0}  \def\DosImposeFormatNEA{}
\def\DosImposeReliureN{souple}  \def\DosImposeReliureE{rigide}
\def\DosImposeReliureA{souple}  \def\DosImposeReliureNE{souple}
\def\DosImposeReliureEA{souple} \def\DosImposeReliureNEA{souple}

% Fond perdu, en mm. 3 mm est la valeur courante ; certains imprimeurs
% demandent 5 mm.
\def\FondPerdu{3}

% Marge technique autour du fond perdu, en mm. Sert à loger les traits de
% coupe en mode épreuve. Ignorée en mode production.
\def\MargeTech{12}

% Zone de sécurité, en mm : aucun texte ne doit s'en approcher davantage
% du bord de coupe.
\def\Securite{8}

% ---------------------------------------------------------------------
% 2. DONNÉES PAR CLASSE
%    pages : nombre de pages du PDF intérieur, COUVERTURE NON COMPRISE.
%    Un dos carré collé impose un nombre de pages PAIR ; le fichier
%    refuse de compiler si le compte est impair.
%
%    Ces tables ne servent que si \PagesForcees et \VersionForcee sont
%    vides, c'est-à-dire à la main. Valeurs A4 mesurées sur les PDF de la
%    a.3 (CLAUDE.md § 1, 05/09/2026). En 20 x 24, la pagination change
%    entièrement : NEA passe de 816 à 1036 pages.
% ---------------------------------------------------------------------

\def\DonneesN{262}
\def\DonneesE{214}
\def\DonneesA{386}
\def\DonneesNE{0}      % édition jamais produite
\def\DonneesEA{0}      % édition jamais produite
\def\DonneesNEA{816}

% Version imprimée sur le dos et la 1re de couverture, PAR CLASSE.
% « Les versions sont suivies par classe, chacune évoluant à son rythme »
% (CHANGELOG.md).
\def\VersionN{a.3}
\def\VersionE{a.3}
\def\VersionA{a.3}
\def\VersionNE{a.3}
\def\VersionEA{a.3}
\def\VersionNEA{a.3}

% ---------------------------------------------------------------------
% 3. TEXTES
% ---------------------------------------------------------------------

\def\TitreOuvrage{50\,Ohm}
\def\SousTitre{Préparation à l'examen radioamateur}

% Sous-titre propre à chaque classe (repris de FR_TITLES, build_book.py).
\def\SousTitreN{Cours complet — classe N}
\def\SousTitreE{Cours complet — classe E}
\def\SousTitreA{Cours complet — classe A}
\def\SousTitreNE{Cours complet — classes N et E}
\def\SousTitreEA{Cours complet — classes E et A}
\def\SousTitreNEA{Cours complet — classes N, E et A}

% Lettres du filigrane, empilées de haut en bas (cf. build_book.py v0.20).
\def\LettresN{N}
\def\LettresE{E}
\def\LettresA{A}
\def\LettresNE{N,E}
\def\LettresEA{E,A}
\def\LettresNEA{N,E,A}

% Texte de 4e de couverture.
% NB : il parle des classes N, E et A quelle que soit l'édition. Juste pour
% N, E, A et NEA ; approximatif pour NE et EA ; faux pour SWL, d'où son refus.
\def\QuatriemeTexte{%
Ce guide d'étude prépare aux examens radioamateur allemands des classes
N, E et A. Il traduit en français les contenus du projet 50ohm.de, publiés
par le référat AJW du DARC e.\,V., et les complète d'encarts « En France »
qui signalent, sans rien retirer au texte allemand, ce que la
réglementation française prévoit sur le même sujet.\par\smallskip
Les questions officielles du catalogue allemand accompagnent chaque
notion, avec leurs réponses.%
}

\def\Auteur{Traduit avec l'aide d'une IA par Pierre F4JWI}
\def\Depot{github.com/DrPLG/50ohm-fr}

% =====================================================================
%  À PARTIR D'ICI, PLUS RIEN À MODIFIER
% =====================================================================

\usepackage{iftex}
\RequireLuaTeX
\usepackage{xfp}
\usepackage{tikz}
\usetikzlibrary{calc}
\usepackage{fontspec}
\usepackage[french]{babel}
\usepackage{geometry}

% ---------------------------------------------------------------------
% 4. RÉSOLUTION DES DONNÉES DE LA CLASSE
% ---------------------------------------------------------------------

\makeatletter
% Tome SWL : refus explicite, plutôt qu'un « classe inconnue » qui laisserait
% croire à un oubli. Le texte de 4e parle des classes N, E et A.
\def\@swlmot{SWL}
\ifx\Classe\@swlmot
  \PackageError{couverture}{Pas de couverture pour le tome SWL}%
    {Sa 4e de couverture n'est pas redigee : le texte commun parle des
     classes N, E et A. C'est une decision editoriale a prendre d'abord.}%
\fi

\def\@resoudre#1#2{%
  \expandafter\let\expandafter#1\csname #2\Classe\endcsname
  \ifx#1\relax
    \PackageError{couverture}{Classe \Classe\space inconnue : pas de #2\Classe}%
      {Valeurs admises : N, E, A, NE, EA, NEA.}%
  \fi}
\@resoudre\NbPages{Donnees}
\@resoudre\SousTitreClasse{SousTitre}
\@resoudre\Lettres{Lettres}
\@resoudre\DosImpose{DosImpose}
\@resoudre\DosImposePages{DosImposePages}
\@resoudre\DosImposeFormat{DosImposeFormat}
\@resoudre\DosImposeReliure{DosImposeReliure}
\@resoudre\Version{Version}

% Valeurs imposées de l'extérieur. Test de vacuité par \detokenize : il ne
% dépend pas de la façon dont la macro a été définie.
\if\relax\detokenize\expandafter{\PagesForcees}\relax\else
  \let\NbPages\PagesForcees
\fi
\if\relax\detokenize\expandafter{\VersionForcee}\relax\else
  \let\Version\VersionForcee
\fi

% Format du livre fini.
\def\@formatAquatre{a4}
\def\@formatVingt{20x24}
\def\@formatVingtMarge{20x24-marge}
\ifx\Format\@formatAquatre
  \def\PageLargeur{210}\def\PageHauteur{297}
\else\ifx\Format\@formatVingt
  \def\PageLargeur{200}\def\PageHauteur{240}
\else\ifx\Format\@formatVingtMarge
  \def\PageLargeur{200}\def\PageHauteur{240}
\else
  \PackageError{couverture}{Format \Format\space inconnu}%
    {Valeurs admises : a4, 20x24, 20x24-marge.}%
  \def\PageLargeur{210}\def\PageHauteur{297}
\fi\fi\fi
% Reliure. Hors de « souple » et « rigide » : erreur, pas de supposition.
\def\@reliuresouple{souple}
\def\@reliurerigide{rigide}
\ifx\Reliure\@reliurerigide
  \edef\PageLargeur{\fpeval{\PageLargeur + 5}}
  \edef\PageHauteur{\fpeval{\PageHauteur + 5}}
  \def\FondPerdu{15}
\else\ifx\Reliure\@reliuresouple\else
  \PackageError{couverture}{Reliure \Reliure\space inconnue}%
    {Valeurs admises : souple, rigide.}%
\fi\fi
\makeatother

% Garde-fou : une édition jamais produite n'a pas de couverture.
\ifnum\NbPages=0
  \PackageError{couverture}{L'edition \Classe\space n'a pas de pagination}%
    {Cette edition n'a jamais ete compilee. Renseigner \string\Donnees\Classe
     ou passer \string\PagesForcees.}%
\fi

% Avertissement : l'imprimeur demande un nombre de pages multiple de 4
% (17/09/2026). Pas une erreur : c'est son exigence, pas une loi du papier.
\ifnum\fpeval{\NbPages - 4*floor(\NbPages/4)}=0 \else
  \typeout{COUVERTURE-AVERTISSEMENT \NbPages\space pages : pas un multiple
    de 4, que l'imprimeur exige -- ajouter des pages a l'interieur}%
\fi

% Garde-fou : un dos carré collé impose un nombre de pages pair.
\ifodd\NbPages
  \PackageError{couverture}{Nombre de pages impair (\NbPages)}%
    {Un dos carre colle exige un nombre PAIR de pages : une feuille en
     porte deux. Ajouter une page blanche a l'interieur.}%
\fi

% ---------------------------------------------------------------------
% 5. GÉOMÉTRIE — c'est ici que se joue l'exactitude du dos
% ---------------------------------------------------------------------

% Épaisseur du dos : valeur imposée si elle est donnée ET qu'elle vaut pour
% cette pagination et ce format ; calcul sinon.
\edef\DosCalcule{\fpeval{round(\NbPages / 2 * \Grammage * \Bouffant / 1000, 2)}}
\def\DosOrigine{calculee}
\ifnum\fpeval{\DosImpose > 0}=1
  \ifnum\fpeval{\NbPages = \DosImposePages}=1
    \ifx\Format\DosImposeFormat
      \ifx\Reliure\DosImposeReliure
        \def\DosOrigine{imposee}
      \fi
    \fi
  \fi
  \def\MotCalculee{calculee}
  \ifx\DosOrigine\MotCalculee
    \typeout{COUVERTURE-AVERTISSEMENT dos impose de \DosImpose\space mm
      donne pour \DosImposePages\space pages en \DosImposeFormat\space
      (reliure \DosImposeReliure) -- inapplicable a \NbPages\space pages en
      \Format\space (reliure \Reliure), dos calcule}%
  \fi
\fi
% NB : pas de \@ dans ces noms. Cette section est hors \makeatletter, et
% « \def\@imposee{...} » y redéfinit \@ (suivi du texte « imposee ») : le
% test échouait toujours, en silence. Constaté le 16/09/2026 — le dos de N
% tombait juste par coïncidence, 14,0 mm calculés pour 14 imposés.
\def\MotRigide{rigide}
\def\MotImposee{imposee}
\ifx\DosOrigine\MotImposee
  \edef\Dos{\DosImpose}
\else
  \edef\Dos{\DosCalcule}
\fi

% Marge technique : uniquement en mode épreuve.
\def\ouimot{oui}
\ifx\Reperes\ouimot \edef\Tech{\MargeTech}\else \def\Tech{0}\fi

% Dimensions du papier.
%   largeur = marge + fond perdu + 4e + dos + 1re + fond perdu + marge
%   hauteur = marge + fond perdu + hauteur + fond perdu + marge
\edef\PapierL{\fpeval{2*\PageLargeur + \Dos + 2*\FondPerdu + 2*\Tech}}
\edef\PapierH{\fpeval{\PageHauteur + 2*\FondPerdu + 2*\Tech}}

\geometry{paperwidth=\PapierL mm, paperheight=\PapierH mm,
          margin=0pt, nohead, nofoot}

% ---------------------------------------------------------------------
%  TOUTES LES COTES SONT FIGÉES ICI, ET NULLE PART AILLEURS.
%
%  Raison, payée par trois compilations ratées : un \fpeval laissé dans
%  une coordonnée ou dans le texte d'un nœud TikZ n'est pas développé au
%  moment où TikZ analyse le chemin. On obtient « Giving up on this path.
%  Did you forget a semicolon? », qui accuse une ponctuation alors que la
%  faute est ailleurs. Le corps du document ci-dessous n'emploie donc que
%  des macros déjà calculées.
% ---------------------------------------------------------------------

% Repères horizontaux, mesurés depuis le bord GAUCHE du papier.
\edef\CoupeG{\fpeval{\Tech + \FondPerdu}}          % coupe gauche (4e de couv.)
\edef\DosG{\fpeval{\CoupeG + \PageLargeur}}         % dos, début
\edef\DosD{\fpeval{\DosG + \Dos}}                   % dos, fin
\edef\CoupeD{\fpeval{\DosD + \PageLargeur}}         % coupe droite (1re de couv.)
% Repères verticaux, depuis le bord BAS.
\edef\CoupeB{\fpeval{\Tech + \FondPerdu}}
\edef\CoupeH{\fpeval{\CoupeB + \PageHauteur}}
% Axes.
\edef\DosAxe{\fpeval{(\DosG + \DosD)/2}}
\edef\MilieuH{\fpeval{(\CoupeB + \CoupeH)/2}}
\edef\MilieuL{\fpeval{\PapierL/2}}

% Format fini, pour les cotes affichées.
\edef\FiniL{\fpeval{2*\PageLargeur + \Dos}}

% Bandeau de la 1re de couverture : même proportion que la page de titre.
\edef\BandeauL{\fpeval{round(0.34*\PageLargeur, 2)}}
\edef\BandeauG{\fpeval{\CoupeD - \BandeauL}}
\edef\BandeauAccent{\fpeval{\BandeauG - 0.7}}
\edef\DosAccentG{\fpeval{\DosG - 0.7}}
\edef\DosAccentD{\fpeval{\DosD + 0.7}}
% Couverture rigide : le bleu du dos déborde de 2 mm sur chaque plat, pour
% absorber la tolérance de l'imprimeur sur l'épaisseur du dos (±10 %).
\edef\DosEtenduG{\fpeval{\DosG - 2}}
\edef\DosEtenduD{\fpeval{\DosD + 2}}

% Filigrane de classe.
\edef\FiligraneX{\fpeval{(\BandeauG + \CoupeD)/2}}
\edef\FiligraneY{\fpeval{\CoupeH - 6}}

% Bloc-titre de la 1re de couverture.
\edef\TitreX{\fpeval{\DosD + 12}}
\edef\TitreY{\fpeval{\CoupeH - 72}}
\edef\TitreW{\fpeval{\BandeauG - \DosD - 16}}
\edef\VersionY{\fpeval{\CoupeB + 14}}

% 4e de couverture.
\edef\QuatreX{\fpeval{\CoupeG + \Securite + 6}}
\edef\QuatreY{\fpeval{\CoupeH - \Securite - 14}}
\edef\QuatreBasY{\fpeval{\CoupeB + \Securite + 22}}
\edef\QuatreW{\fpeval{\PageLargeur - 2*\Securite - 12}}

% Zone de sécurité.
\edef\SecG{\fpeval{\CoupeG + \Securite}}
\edef\SecD{\fpeval{\CoupeD - \Securite}}
\edef\SecB{\fpeval{\CoupeB + \Securite}}
\edef\SecH{\fpeval{\CoupeH - \Securite}}

% Zones de contrôle : ce qui doit rester de part et d'autre du pli du dos.
\edef\PliQuatre{\fpeval{\DosG - \Securite}}         % 4e : limite côté dos
\edef\PliPremiere{\fpeval{\DosD + \Securite}}       % 1re : limite côté dos

% Traits de coupe et cotes, dans la marge technique.
\edef\TechExt{\fpeval{\Tech - 8}}
\edef\TechInt{\fpeval{\Tech - 1}}
\edef\TechCote{\fpeval{\Tech - 9}}
\edef\HautInt{\fpeval{\PapierH - \Tech + 1}}
\edef\HautExt{\fpeval{\PapierH - \Tech + 8}}
\edef\HautCote{\fpeval{\PapierH - \Tech + 9}}
\edef\DroitInt{\fpeval{\PapierL - \Tech + 1}}
\edef\DroitExt{\fpeval{\PapierL - \Tech + 8}}

% Emplacement du code-barres ISBN : en bas à DROITE de la 4e de couverture,
% selon l'usage éditorial, et non à gauche où il chevaucherait les mentions
% de licence.
\edef\IsbnD{\fpeval{\DosG - \Securite}}
\edef\IsbnG{\fpeval{\IsbnD - 40}}
\edef\IsbnH{\fpeval{\SecB + 25}}
\edef\IsbnCX{\fpeval{\IsbnG + 20}}
\edef\IsbnCY{\fpeval{\SecB + 12.5}}

% Cartouches de cotes, dans la marge technique. Ancrés vers l'INTÉRIEUR du
% papier : ancrés vers l'extérieur, ils débordaient et étaient rognés.
\edef\CoteBasY{\fpeval{1}}
\edef\CoteHautY{\fpeval{\PapierH - 1}}

% ---------------------------------------------------------------------
% 6. COULEURS ET POLICES — reprises de la page de titre (build_book.py)
% ---------------------------------------------------------------------

\definecolor{TitleBand}{cmyk}{.9,.55,.1,.35}   % bleu profond
\definecolor{TitleAccent}{cmyk}{.05,.8,1,0}    % rouge DARC (accent fin)
\definecolor{Repere}{cmyk}{0,1,1,0}            % magenta : repères de montage

\setmainfont{Libertinus Serif}[
  ItalicFont     = Libertinus Serif Italic,
  BoldFont       = Libertinus Serif Bold,
  BoldItalicFont = Libertinus Serif Bold Italic]

\pagestyle{empty}

% ---------------------------------------------------------------------
% 7. FILIGRANE DE CLASSE
%    Le texte est construit ici, et non dans le nœud : un \\ produit à
%    l'intérieur d'un \foreach referme le groupe d'alignement de TikZ.
%    Une lettre seule est grande ; plusieurs s'empilent et se réduisent
%    (build_book.py v0.20 : « NEA » à 105 pt débordait encore).
% ---------------------------------------------------------------------

\makeatletter
\def\FiligraneTexte{}
\newif\ifpremierelettre \premierelettretrue
\count255=0
\@for\@tempa:=\Lettres\do{%
  \advance\count255 by 1\relax
  \ifpremierelettre \premierelettrefalse
  \else \g@addto@macro\FiligraneTexte{\\}\fi
  \expandafter\g@addto@macro\expandafter\FiligraneTexte\expandafter{\@tempa}}%
\xdef\NbLettres{\the\count255}
\makeatother

\ifnum\NbLettres=1 \def\CorpsF{150}\else\def\CorpsF{86}\fi
\edef\CorpsFI{\fpeval{round(\CorpsF*1.02, 0)}}

% Texte du cartouche de cotes, lui aussi construit hors du nœud TikZ.
%
% NE PAS y employer \textsuperscript : dans un nœud « align=center », il
% casse le chemin TikZ (« Giving up on this path. Did you forget a
% semicolon? »), et l'erreur est signalée sur le nœud SUIVANT — elle
% accuse donc une ponctuation là où la faute est ici. Isolé par essais :
% $^2$ et le caractère ² passent, \textsuperscript non.
\def\CotesTexte{%
  \Classe\ --- \NbPages\ pages --- papier \Grammage\ g/m$^2$,
  bouffant \Bouffant\ --- fond perdu \FondPerdu\ mm\\
  papier \PapierL\ mm sur \PapierH\ mm --- fini \FiniL\ mm sur
  \PageHauteur\ mm\\
  EPREUVE, ne pas envoyer a l'impression%
}

% ---------------------------------------------------------------------
% 8. CONTRÔLE DE SUPERPOSITION
%
%  \MesureNoeud{nœud}{préfixe}{retrait} mesure la boîte d'un nœud par ses
%  quatre coins (un nœud pivoté a ses ancres pivotées : le min et le max des
%  quatre suffisent) et la ramène au repère du papier, en mm. Le point (0,0)
%  est lu par le même « let » que les ancres : les deux passent donc par la
%  même transformation, et la comparaison ne dépend pas de l'endroit où
%  TeX a posé l'image sur la page.
%
%  Le retrait sert au seul filigrane, pour mesurer les lettres et non leur
%  cadre : c'est le bord des lettres que le détourage rognerait.
%
%  PIÈGE MESURÉ le 16/09/2026 : la marge interne TikZ (.3333em) n'est PAS
%  évaluée dans la police du nœud, mais dans celle qui entoure le dessin.
%  Instrumenté : dans le nœud, em = 150 pt ; boîte brute du « N » :
%  41,7 x 37,5 mm, soit la lettre plus 2 x 1,3 mm — 1,3 mm étant .3333em
%  à 11 pt. Un retrait calculé à 150 pt (17,6 mm) ramenait la lettre à
%  6,6 mm de large, et le contrôle aurait dit OK sur une mesure absurde.
% ---------------------------------------------------------------------

\newcount\NbAlertes
\newdimen\FiligraneSep
\newsavebox\CreditsBoite

% Crédits de 4e : lignes coupées à la main. Leur largeur RÉELLE est celle
% de la plus longue ligne, mesurée dans une boîte à part (voir § 8, ISBN).
\def\CreditsTexte{%
    Réalisé à partir des contenus de 50ohm.de (en allemand)\\
    50ohm.de-Autorenteam, coordonné par le référat AJW du DARC e.\,V.\\
    \Auteur\\[3pt]
    Licence CC BY 4.0 --- \texttt{\Depot}}

\def\MesureNoeud#1#2#3{%
  \path let \p1=(#1.south west), \p2=(#1.north east),
            \p3=(#1.north west), \p4=(#1.south east), \p5=(0,0) in
    \pgfextra{%
      \expandafter\xdef\csname #2xmin\endcsname{\fpeval{round((min(\x1,\x2,\x3,\x4) - \x5 + #3)/1mm, 2)}}%
      \expandafter\xdef\csname #2xmax\endcsname{\fpeval{round((max(\x1,\x2,\x3,\x4) - \x5 - #3)/1mm, 2)}}%
      \expandafter\xdef\csname #2ymin\endcsname{\fpeval{round((min(\y1,\y2,\y3,\y4) - \y5 + #3)/1mm, 2)}}%
      \expandafter\xdef\csname #2ymax\endcsname{\fpeval{round((max(\y1,\y2,\y3,\y4) - \y5 - #3)/1mm, 2)}}%
    };}

% \ControleDans{libellé}{préfixe}{xmin}{xmax}{ymin}{ymax}
\def\ControleDans#1#2#3#4#5#6{%
  \ifnum\fpeval{\csname #2xmin\endcsname >= #3 && \csname #2xmax\endcsname <= #4
             && \csname #2ymin\endcsname >= #5 && \csname #2ymax\endcsname <= #6}=1
    \typeout{COUVERTURE-CONTROLE OK #1 : x \csname #2xmin\endcsname-\csname #2xmax\endcsname,
      y \csname #2ymin\endcsname-\csname #2ymax\endcsname\space dans x #3-#4, y #5-#6}%
  \else
    \global\advance\NbAlertes by 1
    \typeout{COUVERTURE-CONTROLE ALERTE #1 : x \csname #2xmin\endcsname-\csname #2xmax\endcsname,
      y \csname #2ymin\endcsname-\csname #2ymax\endcsname\space HORS de x #3-#4, y #5-#6}%
  \fi}

% \ControleDisjoints{libellé}{préfixe}{xmin}{xmax}{ymin}{ymax}
% Deux rectangles se chevauchent si et seulement s'ils se recouvrent sur
% les DEUX axes à la fois.
\def\ControleDisjoints#1#2#3#4#5#6{%
  \ifnum\fpeval{\csname #2xmin\endcsname < #4 && #3 < \csname #2xmax\endcsname
             && \csname #2ymin\endcsname < #6 && #5 < \csname #2ymax\endcsname}=1
    \global\advance\NbAlertes by 1
    \typeout{COUVERTURE-CONTROLE ALERTE #1 : chevauchement}%
  \else
    \typeout{COUVERTURE-CONTROLE OK #1 : disjoints}%
  \fi}

\begin{document}
% Marge interne du filigrane, évaluée dans la police qui ENTOURE le dessin :
% c'est là que TikZ évalue le .3333em de « inner sep » (cf. § 8).
\FiligraneSep=.3333em
% Largeur réelle des crédits : la plus longue ligne, dans leur police.
\sbox\CreditsBoite{\small\begin{tabular}{@{}l@{}}\CreditsTexte\end{tabular}}%
\begin{tikzpicture}[remember picture, overlay,
                    x=1mm, y=1mm, shift={(current page.south west)}]

% ---- Fonds ----------------------------------------------------------
% La 4e de couverture reste sur le blanc du papier.
% Le bandeau bleu couvre le dos, puis reprend au tiers extérieur de la
% 1re de couverture — même partition 2/3-1/3 que la page de titre.
% Couverture rigide : pas de filets au dos (ils tomberaient à côté des
% plis, vu la tolérance), et un dos bleu élargi de 2 mm de chaque côté.
\ifx\Reliure\MotRigide
  \fill[TitleBand] (\DosEtenduG, 0) rectangle (\DosEtenduD, \PapierH);
\else
  \fill[TitleBand] (\DosG, 0) rectangle (\DosD, \PapierH);
\fi
\fill[TitleBand] (\BandeauG, 0) rectangle (\PapierL, \PapierH);
% Filets d'accent : à la jonction sur la 1re, de part et d'autre du dos.
\fill[TitleAccent] (\BandeauAccent, 0) rectangle (\BandeauG, \PapierH);
\ifx\Reliure\MotRigide\else
  \fill[TitleAccent] (\DosAccentG, 0) rectangle (\DosG, \PapierH);
  \fill[TitleAccent] (\DosD, 0) rectangle (\DosAccentD, \PapierH);
\fi

% ---- 1re de couverture ----------------------------------------------
% Filigrane de classe, en réserve claire, détouré au bandeau.
\begin{scope}
\clip (\BandeauG, 0) rectangle (\PapierL, \PapierH);
\node[name=filigrane, anchor=north, align=center, text=white, opacity=0.16,
      font=\fontsize{\CorpsF}{\CorpsFI}\selectfont\bfseries]
  at (\FiligraneX, \FiligraneY) {\FiligraneTexte};
\end{scope}

\node[name=titre, anchor=west, align=left, text width=\TitreW mm]
  at (\TitreX, \TitreY) {%
    {\fontsize{54}{56}\selectfont\bfseries \TitreOuvrage}\\[10pt]
    {\Large\color{TitleBand}\bfseries \SousTitre}\\[16pt]
    {\large\color{black!75} \SousTitreClasse}};

\node[name=version, anchor=south west, align=left, font=\small, text=black!70]
  at (\TitreX, \VersionY) {Version \Version\\ Licence CC BY 4.0};

% ---- Dos -------------------------------------------------------------
% Un dos étroit ne reçoit pas de texte : en deçà de 6 mm, le débord au
% massicot ferait mordre le texte vertical sur les plats.
\ifnum\fpeval{\Dos >= 6}=1
  \node[name=dos, rotate=90, anchor=center, text=white, font=\bfseries,
        inner sep=0pt]
    at (\DosAxe, \MilieuH)
    {\TitreOuvrage\ \ --- \ \SousTitreClasse\ \ --- \ \Version};
\fi

% ---- 4e de couverture ------------------------------------------------
\node[name=quatre, anchor=north west, align=justify, text width=\QuatreW mm]
  at (\QuatreX, \QuatreY) {\QuatriemeTexte};

% Le texte vit dans \CreditsTexte (§ 8) pour être mesuré à part : la boîte
% du nœud a la largeur fixe \QuatreW et couvre l'emplacement ISBN sans
% qu'aucun caractère n'y soit. Retirer « text width » pour l'éviter change
% l'interlignage — essayé le 16/09/2026, 6 934 pixels différents.
\node[name=credits, anchor=south west, align=left, font=\small, text=black!70,
      text width=\QuatreW mm]
  at (\QuatreX, \QuatreBasY) {\CreditsTexte};

% ---- Repères de fabrication (mode épreuve uniquement) ----------------
\ifx\Reperes\ouimot
  % Traits de coupe, hors du fond perdu.
  \foreach \x in {\CoupeG, \CoupeD} {
    \draw[Repere, line width=0.25pt] (\x, \TechExt) -- (\x, \TechInt);
    \draw[Repere, line width=0.25pt] (\x, \HautInt) -- (\x, \HautExt);}
  \foreach \y in {\CoupeB, \CoupeH} {
    \draw[Repere, line width=0.25pt] (\TechExt, \y) -- (\TechInt, \y);
    \draw[Repere, line width=0.25pt] (\DroitInt, \y) -- (\DroitExt, \y);}
  % Limites du dos, en pointillés.
  \foreach \x in {\DosG, \DosD} {
    \draw[Repere, line width=0.25pt, dashed] (\x, \TechExt) -- (\x, \TechInt);
    \draw[Repere, line width=0.25pt, dashed] (\x, \HautInt) -- (\x, \HautExt);}
  % Zone de sécurité.
  \draw[Repere, line width=0.2pt, dotted]
    (\SecG, \SecB) rectangle (\SecD, \SecH);
  % Emplacement réservé au code-barres ISBN.
  \draw[Repere, line width=0.2pt] (\IsbnG, \SecB) rectangle (\IsbnD, \IsbnH);
  \node[font=\tiny\sffamily, text=Repere, align=center]
    at (\IsbnCX, \IsbnCY) {emplacement\\ISBN};
  % Cotes.
  \node[anchor=north, font=\tiny\sffamily, text=Repere]
    at (\DosAxe, \CoteHautY) {dos \Dos\ mm (\DosOrigine)};
  \node[anchor=south, font=\tiny\sffamily, text=Repere, align=center]
    at (\MilieuL, \CoteBasY) {\CotesTexte};
\fi

% ---- Contrôle de superposition (§ 8) ----------------------------------
% Rien n'est dessiné : des mesures et des lignes de journal.
\MesureNoeud{filigrane}{Fil}{\the\FiligraneSep}
\MesureNoeud{titre}{Tit}{0pt}
\MesureNoeud{version}{Ver}{0pt}
\MesureNoeud{quatre}{Qua}{0pt}
\MesureNoeud{credits}{Cre}{0pt}
% Bord droit des crédits = bord gauche + marge interne + plus longue ligne.
\xdef\Crexmax{\fpeval{round(\Crexmin + (\the\FiligraneSep + \the\wd\CreditsBoite)/1mm, 2)}}

% Le filigrane peut entrer dans la zone de sécurité (il est décoratif),
% pas franchir la coupe ni sortir du bandeau — le détourage le rognerait.
\ControleDans{filigrane dans le bandeau}{Fil}{\BandeauG}{\CoupeD}{\CoupeB}{\CoupeH}
\ControleDans{titre sur la 1re, hors bandeau}{Tit}{\PliPremiere}{\BandeauAccent}{\SecB}{\SecH}
\ControleDans{version sur la 1re, hors bandeau}{Ver}{\PliPremiere}{\BandeauAccent}{\SecB}{\SecH}
\ControleDans{texte de 4e dans la zone de securite}{Qua}{\SecG}{\PliQuatre}{\SecB}{\SecH}
\ControleDans{credits de 4e dans la zone de securite}{Cre}{\SecG}{\PliQuatre}{\SecB}{\SecH}
\ControleDisjoints{titre / version}{Tit}{\Verxmin}{\Verxmax}{\Verymin}{\Verymax}
\ControleDisjoints{texte de 4e / credits}{Qua}{\Crexmin}{\Crexmax}{\Creymin}{\Creymax}
\ControleDisjoints{texte de 4e / emplacement ISBN}{Qua}{\IsbnG}{\IsbnD}{\SecB}{\IsbnH}
\ControleDisjoints{credits / emplacement ISBN}{Cre}{\IsbnG}{\IsbnD}{\SecB}{\IsbnH}
\ifnum\fpeval{\Dos >= 6}=1
  \MesureNoeud{dos}{Dos}{0pt}
  \ControleDans{texte du dos entre les plis}{Dos}{\DosG}{\DosD}{\SecB}{\SecH}
\else
  \typeout{COUVERTURE-CONTROLE OK dos de \Dos\space mm : trop etroit, sans texte}%
\fi

\end{tikzpicture}
\typeout{COUVERTURE-BILAN alertes=\the\NbAlertes\space; classe=\Classe\space;
  format=\Format\space; reliure=\Reliure\space; pages=\NbPages\space;
  version=\Version\space;
  dos=\Dos\space mm (\DosOrigine); fini=\FiniL x \PageHauteur\space mm;
  papier=\PapierL x \PapierH\space mm}
\end{document}
"
%
%  Pour le fichier destiné à l'imprimeur (sans repères) :
%
%      lualatex "\def\Reperes{non}% =====================================================================
%  couverture.tex — couverture dos carré collé pour les livres 50ohm-fr
% =====================================================================
%
%  Produit une couverture d'un seul tenant : 4e de couverture | dos |
%  1re de couverture, aux dimensions exactes du livre relié, fond perdu
%  compris. L'épaisseur du dos est CALCULÉE à partir du nombre de pages,
%  du grammage et du bouffant du papier.
%
%  ORIGINE ET ÉTAT. Écrit le 20/08/2026 dans Fichierstravail/ (hors git),
%  versé au dépôt le 16/09/2026 pour que fenetre_compilation.py puisse
%  l'appeler. Ce qui a changé à cette occasion :
%
%    - le FORMAT est un paramètre (\Format = a4, 20x24 ou 20x24-marge,
%      mêmes noms que --format de build_book.py) ;
%    - la pagination et la version peuvent être IMPOSÉES de l'extérieur
%      (\PagesForcees, \VersionForcee) : la fenêtre lit la pagination dans
%      le journal du livre qu'elle vient de compiler, au lieu de se fier à
%      une table qui vieillit ;
%    - un dos imposé par l'imprimeur n'est appliqué QUE pour la pagination
%      et le format pour lesquels il a été donné ;
%    - un CONTRÔLE DE SUPERPOSITION est calculé par LaTeX lui-même, à partir
%      des boîtes réelles des nœuds, et écrit dans le journal ;
%    - le tome SWL est refusé explicitement : sa 4e de couverture n'est pas
%      rédigée (le texte commun parle des classes N, E et A).
%
%    Neutralité vérifiée : en A4, pour N (258 p., a.2) et NEA (816 p., a.3),
%    en épreuve comme en impression, le rendu est identique au pixel près à
%    celui du gabarit d'origine.
%
%  COMPILATION (LuaLaTeX, depuis la racine du dépôt) :
%
%      lualatex couverture.tex
%      lualatex couverture.tex        <-- DEUX FOIS, obligatoirement
%
%  DEUX PASSES SONT NÉCESSAIRES. Le dessin s'ancre sur « current page »,
%  dont la position n'est connue qu'après un premier passage : à la
%  première compilation, tout se tasse dans le coin supérieur et la
%  couverture paraît ratée alors qu'elle ne l'est pas. Aucune erreur n'est
%  émise — seul le rendu le montre. Le contrôle de superposition, lui
%  aussi, ne vaut qu'à la seconde passe.
%
%  Le plus simple est de passer par fenetre_compilation.py (case
%  « couverture pour l'imprimeur »), qui enchaîne les deux passes, produit
%  l'épreuve ET le fichier d'impression, et relit le contrôle.
%
%  À la main, pour une autre classe, un autre format, une pagination mesurée :
%
%      lualatex "\def\Classe{A}\def\Format{20x24}\def\PagesForcees{482}\input{couverture.tex}"
%
%  Pour le fichier destiné à l'imprimeur (sans repères) :
%
%      lualatex "\def\Reperes{non}\input{couverture.tex}"
%
%  ---------------------------------------------------------------------
%  DEUX MODES, un seul paramètre : \Reperes
%
%    oui  — mode ÉPREUVE. Ajoute une marge technique, les traits de coupe,
%           les limites du dos, la zone de sécurité et les cotes. C'est ce
%           qu'il faut pour vérifier à l'écran.
%    non  — mode PRODUCTION. Papier = format final + fond perdu, rien
%           d'autre. C'est ce fichier-là qu'on envoie à l'imprimeur.
%
%  ---------------------------------------------------------------------
%  CONTRÔLE DE SUPERPOSITION
%
%  À la fin du dessin, chaque bloc de texte est mesuré par ses ancres TikZ
%  (quatre coins, ce qui vaut aussi pour le texte pivoté du dos) et comparé
%  à la zone où il doit tenir, puis aux blocs voisins. Le journal porte une
%  ligne par vérification :
%
%      COUVERTURE-CONTROLE OK      <quoi> : ...
%      COUVERTURE-CONTROLE ALERTE  <quoi> : ...
%
%  et un bilan : COUVERTURE-BILAN alertes=<n> ; ...
%
%  Ce que le contrôle VOIT : un bloc qui franchit le pli du dos, la coupe,
%  la zone de sécurité ou le bord du bandeau ; deux blocs qui se chevauchent ;
%  un bloc qui empiète sur l'emplacement réservé au code-barres ISBN.
%  Ce qu'il NE VOIT PAS : une ligne trop longue à l'INTÉRIEUR d'un bloc à
%  largeur fixe (la boîte garde sa largeur, le texte déborde). Celle-là se
%  signale par un « Overfull \hbox » dans le même journal, que la fenêtre
%  compte aussi.
%
%  ---------------------------------------------------------------------
%  AVERTISSEMENT SUR L'ÉPAISSEUR DU DOS
%
%  Le dos se calcule ainsi :
%
%      épaisseur d'une feuille (mm) = grammage x bouffant / 1000
%      dos (mm)                     = (nombre de pages / 2) x épaisseur
%
%  On divise par deux parce qu'une feuille porte DEUX pages.
%
%  Le « bouffant » (main, volume, bulk) dépend du papier et n'est pas
%  devinable : un offset 80 g courant est autour de 1,30, un bouffant
%  d'édition monte à 1,75 et plus. UN ÉCART DE BOUFFANT DÉPLACE LE DOS DE
%  PLUSIEURS MILLIMÈTRES. Demander à l'imprimeur soit son bouffant, soit
%  directement l'épaisseur du dos, et reporter ci-dessous. Tant que cette
%  valeur n'est pas confirmée, la couverture n'est pas bonne à tirer.
%
%  Si l'imprimeur donne le dos en clair, forcer la valeur avec \DosImpose.
% =====================================================================

\documentclass[11pt]{article}

% ---------------------------------------------------------------------
% 1. PARAMÈTRES — c'est la seule zone à modifier
% ---------------------------------------------------------------------

% Classe de l'ouvrage : N, E, A, NE, EA ou NEA.
% (\providecommand : une valeur passée en ligne de commande l'emporte.)
\providecommand{\Classe}{NEA}

% Mode épreuve (oui) ou production (non).
\providecommand{\Reperes}{oui}

% Format du livre FINI : a4, 20x24 ou 20x24-marge — les noms de --format
% dans build_book.py. Les deux maquettes 20 x 24 ne diffèrent que par leur
% colonne de marge intérieure : le livre fini, donc la couverture, est le
% même (200 x 240 mm).
\providecommand{\Format}{a4}

% Reliure : « souple » (dos carré collé, fond perdu 3 mm) ou « rigide »
% (couverture rigide, dite hardcover). En rigide, consignes de l'imprimeur
% du 17/09/2026 : chaque plat déborde du livre de 5 mm en largeur et en
% hauteur (le carton dépasse du bloc), et le fichier porte 15 mm de rabat
% autour du carton au lieu de 3 mm de fond perdu. Les aplats vont déjà au
% bord du papier : ils remplissent le rabat d'eux-mêmes.
\providecommand{\Reliure}{souple}

% Pagination et version imposées de l'extérieur. Vides = tables du § 2.
% La fenêtre les renseigne toujours : la pagination vient du journal du
% livre réellement compilé, dans le format demandé.
\providecommand{\PagesForcees}{}
\providecommand{\VersionForcee}{}

% Papier de l'INTÉRIEUR du livre (pas de la couverture).
% CoolLibri, papier « Standard 90 g blanc ».
\def\Grammage{90}        % g/m²
\def\Bouffant{1.206}     % main DÉDUITE, voir ci-dessous

% Épaisseur du dos IMPOSÉE par l'imprimeur, en mm, par classe.
% Une valeur > 0 prime sur le calcul — MAIS SEULEMENT pour la pagination et
% le format pour lesquels l'imprimeur l'a donnée (\DosImposePages<classe>,
% \DosImposeFormat<classe>). Hors de ce cas, le dos est calculé et le
% journal le dit : un dos de 14 mm donné pour 258 pages A4 ne vaut rien pour
% 262 pages, ni pour le même livre recomposé en 20 x 24.
%
% ATTENTION : cette valeur est propre à CHAQUE ÉDITION. Un dos unique
% appliqué aux quatre livres serait faux pour trois d'entre eux.
%
%  N   — 14 mm, VALEUR DONNÉE PAR COOLIBRI pour 258 pages en 90 g, A4.
%        Fichier attendu : 440 x 303 mm, ce que ce gabarit produit en mode
%        production (2 x 210 + 14 + 2 x 3 = 440 ; 297 + 2 x 3 = 303).
%        Le N de la a.3 compte 262 pages : cette valeur ne s'y applique plus.
%  E, A, NEA — non demandés à l'imprimeur : ils sont CALCULÉS avec la main
%        déduite ci-dessous. À faire confirmer avant tout tirage.
%
% La main de 1,206 n'est pas publiée par CoolLibri : elle est déduite de
% leur propre chiffre, 14 mm / (0,090 x 258/2) = 1,206. Elle reproduit donc
% exactement le dos de N, et n'est qu'une ESTIMATION pour les autres.
% Écarts que cela recouvre, pour mémoire : une main de 1,0 donnerait
% 11,61 mm pour N, une main de 1,3 en donnerait 16,77. Il ne fallait pas
% deviner.
%  E   — 17 mm, VALEUR DONNÉE PAR L'IMPRIMEUR le 17/09/2026 pour 276 pages,
%        20 x 24, couverture RIGIDE, reliure collée (Klebebindung), papier
%        115 g. En reliure cousue (Fadenheftung), il annonce 18 mm.
%
% \DosImposeReliure<classe> : la reliure pour laquelle la valeur a été
% donnée. Le dos d'une couverture rigide n'est pas celui d'un dos carré
% collé souple, même à pagination égale.
\def\DosImposeN{14}   \def\DosImposePagesN{258}  \def\DosImposeFormatN{a4}
\def\DosImposeE{17}   \def\DosImposePagesE{276}  \def\DosImposeFormatE{20x24-marge}
\def\DosImposeA{0}    \def\DosImposePagesA{0}    \def\DosImposeFormatA{}
\def\DosImposeNE{0}   \def\DosImposePagesNE{0}   \def\DosImposeFormatNE{}
\def\DosImposeEA{0}   \def\DosImposePagesEA{0}   \def\DosImposeFormatEA{}
\def\DosImposeNEA{0}  \def\DosImposePagesNEA{0}  \def\DosImposeFormatNEA{}
\def\DosImposeReliureN{souple}  \def\DosImposeReliureE{rigide}
\def\DosImposeReliureA{souple}  \def\DosImposeReliureNE{souple}
\def\DosImposeReliureEA{souple} \def\DosImposeReliureNEA{souple}

% Fond perdu, en mm. 3 mm est la valeur courante ; certains imprimeurs
% demandent 5 mm.
\def\FondPerdu{3}

% Marge technique autour du fond perdu, en mm. Sert à loger les traits de
% coupe en mode épreuve. Ignorée en mode production.
\def\MargeTech{12}

% Zone de sécurité, en mm : aucun texte ne doit s'en approcher davantage
% du bord de coupe.
\def\Securite{8}

% ---------------------------------------------------------------------
% 2. DONNÉES PAR CLASSE
%    pages : nombre de pages du PDF intérieur, COUVERTURE NON COMPRISE.
%    Un dos carré collé impose un nombre de pages PAIR ; le fichier
%    refuse de compiler si le compte est impair.
%
%    Ces tables ne servent que si \PagesForcees et \VersionForcee sont
%    vides, c'est-à-dire à la main. Valeurs A4 mesurées sur les PDF de la
%    a.3 (CLAUDE.md § 1, 05/09/2026). En 20 x 24, la pagination change
%    entièrement : NEA passe de 816 à 1036 pages.
% ---------------------------------------------------------------------

\def\DonneesN{262}
\def\DonneesE{214}
\def\DonneesA{386}
\def\DonneesNE{0}      % édition jamais produite
\def\DonneesEA{0}      % édition jamais produite
\def\DonneesNEA{816}

% Version imprimée sur le dos et la 1re de couverture, PAR CLASSE.
% « Les versions sont suivies par classe, chacune évoluant à son rythme »
% (CHANGELOG.md).
\def\VersionN{a.3}
\def\VersionE{a.3}
\def\VersionA{a.3}
\def\VersionNE{a.3}
\def\VersionEA{a.3}
\def\VersionNEA{a.3}

% ---------------------------------------------------------------------
% 3. TEXTES
% ---------------------------------------------------------------------

\def\TitreOuvrage{50\,Ohm}
\def\SousTitre{Préparation à l'examen radioamateur}

% Sous-titre propre à chaque classe (repris de FR_TITLES, build_book.py).
\def\SousTitreN{Cours complet — classe N}
\def\SousTitreE{Cours complet — classe E}
\def\SousTitreA{Cours complet — classe A}
\def\SousTitreNE{Cours complet — classes N et E}
\def\SousTitreEA{Cours complet — classes E et A}
\def\SousTitreNEA{Cours complet — classes N, E et A}

% Lettres du filigrane, empilées de haut en bas (cf. build_book.py v0.20).
\def\LettresN{N}
\def\LettresE{E}
\def\LettresA{A}
\def\LettresNE{N,E}
\def\LettresEA{E,A}
\def\LettresNEA{N,E,A}

% Texte de 4e de couverture.
% NB : il parle des classes N, E et A quelle que soit l'édition. Juste pour
% N, E, A et NEA ; approximatif pour NE et EA ; faux pour SWL, d'où son refus.
\def\QuatriemeTexte{%
Ce guide d'étude prépare aux examens radioamateur allemands des classes
N, E et A. Il traduit en français les contenus du projet 50ohm.de, publiés
par le référat AJW du DARC e.\,V., et les complète d'encarts « En France »
qui signalent, sans rien retirer au texte allemand, ce que la
réglementation française prévoit sur le même sujet.\par\smallskip
Les questions officielles du catalogue allemand accompagnent chaque
notion, avec leurs réponses.%
}

\def\Auteur{Traduit avec l'aide d'une IA par Pierre F4JWI}
\def\Depot{github.com/DrPLG/50ohm-fr}

% =====================================================================
%  À PARTIR D'ICI, PLUS RIEN À MODIFIER
% =====================================================================

\usepackage{iftex}
\RequireLuaTeX
\usepackage{xfp}
\usepackage{tikz}
\usetikzlibrary{calc}
\usepackage{fontspec}
\usepackage[french]{babel}
\usepackage{geometry}

% ---------------------------------------------------------------------
% 4. RÉSOLUTION DES DONNÉES DE LA CLASSE
% ---------------------------------------------------------------------

\makeatletter
% Tome SWL : refus explicite, plutôt qu'un « classe inconnue » qui laisserait
% croire à un oubli. Le texte de 4e parle des classes N, E et A.
\def\@swlmot{SWL}
\ifx\Classe\@swlmot
  \PackageError{couverture}{Pas de couverture pour le tome SWL}%
    {Sa 4e de couverture n'est pas redigee : le texte commun parle des
     classes N, E et A. C'est une decision editoriale a prendre d'abord.}%
\fi

\def\@resoudre#1#2{%
  \expandafter\let\expandafter#1\csname #2\Classe\endcsname
  \ifx#1\relax
    \PackageError{couverture}{Classe \Classe\space inconnue : pas de #2\Classe}%
      {Valeurs admises : N, E, A, NE, EA, NEA.}%
  \fi}
\@resoudre\NbPages{Donnees}
\@resoudre\SousTitreClasse{SousTitre}
\@resoudre\Lettres{Lettres}
\@resoudre\DosImpose{DosImpose}
\@resoudre\DosImposePages{DosImposePages}
\@resoudre\DosImposeFormat{DosImposeFormat}
\@resoudre\DosImposeReliure{DosImposeReliure}
\@resoudre\Version{Version}

% Valeurs imposées de l'extérieur. Test de vacuité par \detokenize : il ne
% dépend pas de la façon dont la macro a été définie.
\if\relax\detokenize\expandafter{\PagesForcees}\relax\else
  \let\NbPages\PagesForcees
\fi
\if\relax\detokenize\expandafter{\VersionForcee}\relax\else
  \let\Version\VersionForcee
\fi

% Format du livre fini.
\def\@formatAquatre{a4}
\def\@formatVingt{20x24}
\def\@formatVingtMarge{20x24-marge}
\ifx\Format\@formatAquatre
  \def\PageLargeur{210}\def\PageHauteur{297}
\else\ifx\Format\@formatVingt
  \def\PageLargeur{200}\def\PageHauteur{240}
\else\ifx\Format\@formatVingtMarge
  \def\PageLargeur{200}\def\PageHauteur{240}
\else
  \PackageError{couverture}{Format \Format\space inconnu}%
    {Valeurs admises : a4, 20x24, 20x24-marge.}%
  \def\PageLargeur{210}\def\PageHauteur{297}
\fi\fi\fi
% Reliure. Hors de « souple » et « rigide » : erreur, pas de supposition.
\def\@reliuresouple{souple}
\def\@reliurerigide{rigide}
\ifx\Reliure\@reliurerigide
  \edef\PageLargeur{\fpeval{\PageLargeur + 5}}
  \edef\PageHauteur{\fpeval{\PageHauteur + 5}}
  \def\FondPerdu{15}
\else\ifx\Reliure\@reliuresouple\else
  \PackageError{couverture}{Reliure \Reliure\space inconnue}%
    {Valeurs admises : souple, rigide.}%
\fi\fi
\makeatother

% Garde-fou : une édition jamais produite n'a pas de couverture.
\ifnum\NbPages=0
  \PackageError{couverture}{L'edition \Classe\space n'a pas de pagination}%
    {Cette edition n'a jamais ete compilee. Renseigner \string\Donnees\Classe
     ou passer \string\PagesForcees.}%
\fi

% Avertissement : l'imprimeur demande un nombre de pages multiple de 4
% (17/09/2026). Pas une erreur : c'est son exigence, pas une loi du papier.
\ifnum\fpeval{\NbPages - 4*floor(\NbPages/4)}=0 \else
  \typeout{COUVERTURE-AVERTISSEMENT \NbPages\space pages : pas un multiple
    de 4, que l'imprimeur exige -- ajouter des pages a l'interieur}%
\fi

% Garde-fou : un dos carré collé impose un nombre de pages pair.
\ifodd\NbPages
  \PackageError{couverture}{Nombre de pages impair (\NbPages)}%
    {Un dos carre colle exige un nombre PAIR de pages : une feuille en
     porte deux. Ajouter une page blanche a l'interieur.}%
\fi

% ---------------------------------------------------------------------
% 5. GÉOMÉTRIE — c'est ici que se joue l'exactitude du dos
% ---------------------------------------------------------------------

% Épaisseur du dos : valeur imposée si elle est donnée ET qu'elle vaut pour
% cette pagination et ce format ; calcul sinon.
\edef\DosCalcule{\fpeval{round(\NbPages / 2 * \Grammage * \Bouffant / 1000, 2)}}
\def\DosOrigine{calculee}
\ifnum\fpeval{\DosImpose > 0}=1
  \ifnum\fpeval{\NbPages = \DosImposePages}=1
    \ifx\Format\DosImposeFormat
      \ifx\Reliure\DosImposeReliure
        \def\DosOrigine{imposee}
      \fi
    \fi
  \fi
  \def\MotCalculee{calculee}
  \ifx\DosOrigine\MotCalculee
    \typeout{COUVERTURE-AVERTISSEMENT dos impose de \DosImpose\space mm
      donne pour \DosImposePages\space pages en \DosImposeFormat\space
      (reliure \DosImposeReliure) -- inapplicable a \NbPages\space pages en
      \Format\space (reliure \Reliure), dos calcule}%
  \fi
\fi
% NB : pas de \@ dans ces noms. Cette section est hors \makeatletter, et
% « \def\@imposee{...} » y redéfinit \@ (suivi du texte « imposee ») : le
% test échouait toujours, en silence. Constaté le 16/09/2026 — le dos de N
% tombait juste par coïncidence, 14,0 mm calculés pour 14 imposés.
\def\MotRigide{rigide}
\def\MotImposee{imposee}
\ifx\DosOrigine\MotImposee
  \edef\Dos{\DosImpose}
\else
  \edef\Dos{\DosCalcule}
\fi

% Marge technique : uniquement en mode épreuve.
\def\ouimot{oui}
\ifx\Reperes\ouimot \edef\Tech{\MargeTech}\else \def\Tech{0}\fi

% Dimensions du papier.
%   largeur = marge + fond perdu + 4e + dos + 1re + fond perdu + marge
%   hauteur = marge + fond perdu + hauteur + fond perdu + marge
\edef\PapierL{\fpeval{2*\PageLargeur + \Dos + 2*\FondPerdu + 2*\Tech}}
\edef\PapierH{\fpeval{\PageHauteur + 2*\FondPerdu + 2*\Tech}}

\geometry{paperwidth=\PapierL mm, paperheight=\PapierH mm,
          margin=0pt, nohead, nofoot}

% ---------------------------------------------------------------------
%  TOUTES LES COTES SONT FIGÉES ICI, ET NULLE PART AILLEURS.
%
%  Raison, payée par trois compilations ratées : un \fpeval laissé dans
%  une coordonnée ou dans le texte d'un nœud TikZ n'est pas développé au
%  moment où TikZ analyse le chemin. On obtient « Giving up on this path.
%  Did you forget a semicolon? », qui accuse une ponctuation alors que la
%  faute est ailleurs. Le corps du document ci-dessous n'emploie donc que
%  des macros déjà calculées.
% ---------------------------------------------------------------------

% Repères horizontaux, mesurés depuis le bord GAUCHE du papier.
\edef\CoupeG{\fpeval{\Tech + \FondPerdu}}          % coupe gauche (4e de couv.)
\edef\DosG{\fpeval{\CoupeG + \PageLargeur}}         % dos, début
\edef\DosD{\fpeval{\DosG + \Dos}}                   % dos, fin
\edef\CoupeD{\fpeval{\DosD + \PageLargeur}}         % coupe droite (1re de couv.)
% Repères verticaux, depuis le bord BAS.
\edef\CoupeB{\fpeval{\Tech + \FondPerdu}}
\edef\CoupeH{\fpeval{\CoupeB + \PageHauteur}}
% Axes.
\edef\DosAxe{\fpeval{(\DosG + \DosD)/2}}
\edef\MilieuH{\fpeval{(\CoupeB + \CoupeH)/2}}
\edef\MilieuL{\fpeval{\PapierL/2}}

% Format fini, pour les cotes affichées.
\edef\FiniL{\fpeval{2*\PageLargeur + \Dos}}

% Bandeau de la 1re de couverture : même proportion que la page de titre.
\edef\BandeauL{\fpeval{round(0.34*\PageLargeur, 2)}}
\edef\BandeauG{\fpeval{\CoupeD - \BandeauL}}
\edef\BandeauAccent{\fpeval{\BandeauG - 0.7}}
\edef\DosAccentG{\fpeval{\DosG - 0.7}}
\edef\DosAccentD{\fpeval{\DosD + 0.7}}
% Couverture rigide : le bleu du dos déborde de 2 mm sur chaque plat, pour
% absorber la tolérance de l'imprimeur sur l'épaisseur du dos (±10 %).
\edef\DosEtenduG{\fpeval{\DosG - 2}}
\edef\DosEtenduD{\fpeval{\DosD + 2}}

% Filigrane de classe.
\edef\FiligraneX{\fpeval{(\BandeauG + \CoupeD)/2}}
\edef\FiligraneY{\fpeval{\CoupeH - 6}}

% Bloc-titre de la 1re de couverture.
\edef\TitreX{\fpeval{\DosD + 12}}
\edef\TitreY{\fpeval{\CoupeH - 72}}
\edef\TitreW{\fpeval{\BandeauG - \DosD - 16}}
\edef\VersionY{\fpeval{\CoupeB + 14}}

% 4e de couverture.
\edef\QuatreX{\fpeval{\CoupeG + \Securite + 6}}
\edef\QuatreY{\fpeval{\CoupeH - \Securite - 14}}
\edef\QuatreBasY{\fpeval{\CoupeB + \Securite + 22}}
\edef\QuatreW{\fpeval{\PageLargeur - 2*\Securite - 12}}

% Zone de sécurité.
\edef\SecG{\fpeval{\CoupeG + \Securite}}
\edef\SecD{\fpeval{\CoupeD - \Securite}}
\edef\SecB{\fpeval{\CoupeB + \Securite}}
\edef\SecH{\fpeval{\CoupeH - \Securite}}

% Zones de contrôle : ce qui doit rester de part et d'autre du pli du dos.
\edef\PliQuatre{\fpeval{\DosG - \Securite}}         % 4e : limite côté dos
\edef\PliPremiere{\fpeval{\DosD + \Securite}}       % 1re : limite côté dos

% Traits de coupe et cotes, dans la marge technique.
\edef\TechExt{\fpeval{\Tech - 8}}
\edef\TechInt{\fpeval{\Tech - 1}}
\edef\TechCote{\fpeval{\Tech - 9}}
\edef\HautInt{\fpeval{\PapierH - \Tech + 1}}
\edef\HautExt{\fpeval{\PapierH - \Tech + 8}}
\edef\HautCote{\fpeval{\PapierH - \Tech + 9}}
\edef\DroitInt{\fpeval{\PapierL - \Tech + 1}}
\edef\DroitExt{\fpeval{\PapierL - \Tech + 8}}

% Emplacement du code-barres ISBN : en bas à DROITE de la 4e de couverture,
% selon l'usage éditorial, et non à gauche où il chevaucherait les mentions
% de licence.
\edef\IsbnD{\fpeval{\DosG - \Securite}}
\edef\IsbnG{\fpeval{\IsbnD - 40}}
\edef\IsbnH{\fpeval{\SecB + 25}}
\edef\IsbnCX{\fpeval{\IsbnG + 20}}
\edef\IsbnCY{\fpeval{\SecB + 12.5}}

% Cartouches de cotes, dans la marge technique. Ancrés vers l'INTÉRIEUR du
% papier : ancrés vers l'extérieur, ils débordaient et étaient rognés.
\edef\CoteBasY{\fpeval{1}}
\edef\CoteHautY{\fpeval{\PapierH - 1}}

% ---------------------------------------------------------------------
% 6. COULEURS ET POLICES — reprises de la page de titre (build_book.py)
% ---------------------------------------------------------------------

\definecolor{TitleBand}{cmyk}{.9,.55,.1,.35}   % bleu profond
\definecolor{TitleAccent}{cmyk}{.05,.8,1,0}    % rouge DARC (accent fin)
\definecolor{Repere}{cmyk}{0,1,1,0}            % magenta : repères de montage

\setmainfont{Libertinus Serif}[
  ItalicFont     = Libertinus Serif Italic,
  BoldFont       = Libertinus Serif Bold,
  BoldItalicFont = Libertinus Serif Bold Italic]

\pagestyle{empty}

% ---------------------------------------------------------------------
% 7. FILIGRANE DE CLASSE
%    Le texte est construit ici, et non dans le nœud : un \\ produit à
%    l'intérieur d'un \foreach referme le groupe d'alignement de TikZ.
%    Une lettre seule est grande ; plusieurs s'empilent et se réduisent
%    (build_book.py v0.20 : « NEA » à 105 pt débordait encore).
% ---------------------------------------------------------------------

\makeatletter
\def\FiligraneTexte{}
\newif\ifpremierelettre \premierelettretrue
\count255=0
\@for\@tempa:=\Lettres\do{%
  \advance\count255 by 1\relax
  \ifpremierelettre \premierelettrefalse
  \else \g@addto@macro\FiligraneTexte{\\}\fi
  \expandafter\g@addto@macro\expandafter\FiligraneTexte\expandafter{\@tempa}}%
\xdef\NbLettres{\the\count255}
\makeatother

\ifnum\NbLettres=1 \def\CorpsF{150}\else\def\CorpsF{86}\fi
\edef\CorpsFI{\fpeval{round(\CorpsF*1.02, 0)}}

% Texte du cartouche de cotes, lui aussi construit hors du nœud TikZ.
%
% NE PAS y employer \textsuperscript : dans un nœud « align=center », il
% casse le chemin TikZ (« Giving up on this path. Did you forget a
% semicolon? »), et l'erreur est signalée sur le nœud SUIVANT — elle
% accuse donc une ponctuation là où la faute est ici. Isolé par essais :
% $^2$ et le caractère ² passent, \textsuperscript non.
\def\CotesTexte{%
  \Classe\ --- \NbPages\ pages --- papier \Grammage\ g/m$^2$,
  bouffant \Bouffant\ --- fond perdu \FondPerdu\ mm\\
  papier \PapierL\ mm sur \PapierH\ mm --- fini \FiniL\ mm sur
  \PageHauteur\ mm\\
  EPREUVE, ne pas envoyer a l'impression%
}

% ---------------------------------------------------------------------
% 8. CONTRÔLE DE SUPERPOSITION
%
%  \MesureNoeud{nœud}{préfixe}{retrait} mesure la boîte d'un nœud par ses
%  quatre coins (un nœud pivoté a ses ancres pivotées : le min et le max des
%  quatre suffisent) et la ramène au repère du papier, en mm. Le point (0,0)
%  est lu par le même « let » que les ancres : les deux passent donc par la
%  même transformation, et la comparaison ne dépend pas de l'endroit où
%  TeX a posé l'image sur la page.
%
%  Le retrait sert au seul filigrane, pour mesurer les lettres et non leur
%  cadre : c'est le bord des lettres que le détourage rognerait.
%
%  PIÈGE MESURÉ le 16/09/2026 : la marge interne TikZ (.3333em) n'est PAS
%  évaluée dans la police du nœud, mais dans celle qui entoure le dessin.
%  Instrumenté : dans le nœud, em = 150 pt ; boîte brute du « N » :
%  41,7 x 37,5 mm, soit la lettre plus 2 x 1,3 mm — 1,3 mm étant .3333em
%  à 11 pt. Un retrait calculé à 150 pt (17,6 mm) ramenait la lettre à
%  6,6 mm de large, et le contrôle aurait dit OK sur une mesure absurde.
% ---------------------------------------------------------------------

\newcount\NbAlertes
\newdimen\FiligraneSep
\newsavebox\CreditsBoite

% Crédits de 4e : lignes coupées à la main. Leur largeur RÉELLE est celle
% de la plus longue ligne, mesurée dans une boîte à part (voir § 8, ISBN).
\def\CreditsTexte{%
    Réalisé à partir des contenus de 50ohm.de (en allemand)\\
    50ohm.de-Autorenteam, coordonné par le référat AJW du DARC e.\,V.\\
    \Auteur\\[3pt]
    Licence CC BY 4.0 --- \texttt{\Depot}}

\def\MesureNoeud#1#2#3{%
  \path let \p1=(#1.south west), \p2=(#1.north east),
            \p3=(#1.north west), \p4=(#1.south east), \p5=(0,0) in
    \pgfextra{%
      \expandafter\xdef\csname #2xmin\endcsname{\fpeval{round((min(\x1,\x2,\x3,\x4) - \x5 + #3)/1mm, 2)}}%
      \expandafter\xdef\csname #2xmax\endcsname{\fpeval{round((max(\x1,\x2,\x3,\x4) - \x5 - #3)/1mm, 2)}}%
      \expandafter\xdef\csname #2ymin\endcsname{\fpeval{round((min(\y1,\y2,\y3,\y4) - \y5 + #3)/1mm, 2)}}%
      \expandafter\xdef\csname #2ymax\endcsname{\fpeval{round((max(\y1,\y2,\y3,\y4) - \y5 - #3)/1mm, 2)}}%
    };}

% \ControleDans{libellé}{préfixe}{xmin}{xmax}{ymin}{ymax}
\def\ControleDans#1#2#3#4#5#6{%
  \ifnum\fpeval{\csname #2xmin\endcsname >= #3 && \csname #2xmax\endcsname <= #4
             && \csname #2ymin\endcsname >= #5 && \csname #2ymax\endcsname <= #6}=1
    \typeout{COUVERTURE-CONTROLE OK #1 : x \csname #2xmin\endcsname-\csname #2xmax\endcsname,
      y \csname #2ymin\endcsname-\csname #2ymax\endcsname\space dans x #3-#4, y #5-#6}%
  \else
    \global\advance\NbAlertes by 1
    \typeout{COUVERTURE-CONTROLE ALERTE #1 : x \csname #2xmin\endcsname-\csname #2xmax\endcsname,
      y \csname #2ymin\endcsname-\csname #2ymax\endcsname\space HORS de x #3-#4, y #5-#6}%
  \fi}

% \ControleDisjoints{libellé}{préfixe}{xmin}{xmax}{ymin}{ymax}
% Deux rectangles se chevauchent si et seulement s'ils se recouvrent sur
% les DEUX axes à la fois.
\def\ControleDisjoints#1#2#3#4#5#6{%
  \ifnum\fpeval{\csname #2xmin\endcsname < #4 && #3 < \csname #2xmax\endcsname
             && \csname #2ymin\endcsname < #6 && #5 < \csname #2ymax\endcsname}=1
    \global\advance\NbAlertes by 1
    \typeout{COUVERTURE-CONTROLE ALERTE #1 : chevauchement}%
  \else
    \typeout{COUVERTURE-CONTROLE OK #1 : disjoints}%
  \fi}

\begin{document}
% Marge interne du filigrane, évaluée dans la police qui ENTOURE le dessin :
% c'est là que TikZ évalue le .3333em de « inner sep » (cf. § 8).
\FiligraneSep=.3333em
% Largeur réelle des crédits : la plus longue ligne, dans leur police.
\sbox\CreditsBoite{\small\begin{tabular}{@{}l@{}}\CreditsTexte\end{tabular}}%
\begin{tikzpicture}[remember picture, overlay,
                    x=1mm, y=1mm, shift={(current page.south west)}]

% ---- Fonds ----------------------------------------------------------
% La 4e de couverture reste sur le blanc du papier.
% Le bandeau bleu couvre le dos, puis reprend au tiers extérieur de la
% 1re de couverture — même partition 2/3-1/3 que la page de titre.
% Couverture rigide : pas de filets au dos (ils tomberaient à côté des
% plis, vu la tolérance), et un dos bleu élargi de 2 mm de chaque côté.
\ifx\Reliure\MotRigide
  \fill[TitleBand] (\DosEtenduG, 0) rectangle (\DosEtenduD, \PapierH);
\else
  \fill[TitleBand] (\DosG, 0) rectangle (\DosD, \PapierH);
\fi
\fill[TitleBand] (\BandeauG, 0) rectangle (\PapierL, \PapierH);
% Filets d'accent : à la jonction sur la 1re, de part et d'autre du dos.
\fill[TitleAccent] (\BandeauAccent, 0) rectangle (\BandeauG, \PapierH);
\ifx\Reliure\MotRigide\else
  \fill[TitleAccent] (\DosAccentG, 0) rectangle (\DosG, \PapierH);
  \fill[TitleAccent] (\DosD, 0) rectangle (\DosAccentD, \PapierH);
\fi

% ---- 1re de couverture ----------------------------------------------
% Filigrane de classe, en réserve claire, détouré au bandeau.
\begin{scope}
\clip (\BandeauG, 0) rectangle (\PapierL, \PapierH);
\node[name=filigrane, anchor=north, align=center, text=white, opacity=0.16,
      font=\fontsize{\CorpsF}{\CorpsFI}\selectfont\bfseries]
  at (\FiligraneX, \FiligraneY) {\FiligraneTexte};
\end{scope}

\node[name=titre, anchor=west, align=left, text width=\TitreW mm]
  at (\TitreX, \TitreY) {%
    {\fontsize{54}{56}\selectfont\bfseries \TitreOuvrage}\\[10pt]
    {\Large\color{TitleBand}\bfseries \SousTitre}\\[16pt]
    {\large\color{black!75} \SousTitreClasse}};

\node[name=version, anchor=south west, align=left, font=\small, text=black!70]
  at (\TitreX, \VersionY) {Version \Version\\ Licence CC BY 4.0};

% ---- Dos -------------------------------------------------------------
% Un dos étroit ne reçoit pas de texte : en deçà de 6 mm, le débord au
% massicot ferait mordre le texte vertical sur les plats.
\ifnum\fpeval{\Dos >= 6}=1
  \node[name=dos, rotate=90, anchor=center, text=white, font=\bfseries,
        inner sep=0pt]
    at (\DosAxe, \MilieuH)
    {\TitreOuvrage\ \ --- \ \SousTitreClasse\ \ --- \ \Version};
\fi

% ---- 4e de couverture ------------------------------------------------
\node[name=quatre, anchor=north west, align=justify, text width=\QuatreW mm]
  at (\QuatreX, \QuatreY) {\QuatriemeTexte};

% Le texte vit dans \CreditsTexte (§ 8) pour être mesuré à part : la boîte
% du nœud a la largeur fixe \QuatreW et couvre l'emplacement ISBN sans
% qu'aucun caractère n'y soit. Retirer « text width » pour l'éviter change
% l'interlignage — essayé le 16/09/2026, 6 934 pixels différents.
\node[name=credits, anchor=south west, align=left, font=\small, text=black!70,
      text width=\QuatreW mm]
  at (\QuatreX, \QuatreBasY) {\CreditsTexte};

% ---- Repères de fabrication (mode épreuve uniquement) ----------------
\ifx\Reperes\ouimot
  % Traits de coupe, hors du fond perdu.
  \foreach \x in {\CoupeG, \CoupeD} {
    \draw[Repere, line width=0.25pt] (\x, \TechExt) -- (\x, \TechInt);
    \draw[Repere, line width=0.25pt] (\x, \HautInt) -- (\x, \HautExt);}
  \foreach \y in {\CoupeB, \CoupeH} {
    \draw[Repere, line width=0.25pt] (\TechExt, \y) -- (\TechInt, \y);
    \draw[Repere, line width=0.25pt] (\DroitInt, \y) -- (\DroitExt, \y);}
  % Limites du dos, en pointillés.
  \foreach \x in {\DosG, \DosD} {
    \draw[Repere, line width=0.25pt, dashed] (\x, \TechExt) -- (\x, \TechInt);
    \draw[Repere, line width=0.25pt, dashed] (\x, \HautInt) -- (\x, \HautExt);}
  % Zone de sécurité.
  \draw[Repere, line width=0.2pt, dotted]
    (\SecG, \SecB) rectangle (\SecD, \SecH);
  % Emplacement réservé au code-barres ISBN.
  \draw[Repere, line width=0.2pt] (\IsbnG, \SecB) rectangle (\IsbnD, \IsbnH);
  \node[font=\tiny\sffamily, text=Repere, align=center]
    at (\IsbnCX, \IsbnCY) {emplacement\\ISBN};
  % Cotes.
  \node[anchor=north, font=\tiny\sffamily, text=Repere]
    at (\DosAxe, \CoteHautY) {dos \Dos\ mm (\DosOrigine)};
  \node[anchor=south, font=\tiny\sffamily, text=Repere, align=center]
    at (\MilieuL, \CoteBasY) {\CotesTexte};
\fi

% ---- Contrôle de superposition (§ 8) ----------------------------------
% Rien n'est dessiné : des mesures et des lignes de journal.
\MesureNoeud{filigrane}{Fil}{\the\FiligraneSep}
\MesureNoeud{titre}{Tit}{0pt}
\MesureNoeud{version}{Ver}{0pt}
\MesureNoeud{quatre}{Qua}{0pt}
\MesureNoeud{credits}{Cre}{0pt}
% Bord droit des crédits = bord gauche + marge interne + plus longue ligne.
\xdef\Crexmax{\fpeval{round(\Crexmin + (\the\FiligraneSep + \the\wd\CreditsBoite)/1mm, 2)}}

% Le filigrane peut entrer dans la zone de sécurité (il est décoratif),
% pas franchir la coupe ni sortir du bandeau — le détourage le rognerait.
\ControleDans{filigrane dans le bandeau}{Fil}{\BandeauG}{\CoupeD}{\CoupeB}{\CoupeH}
\ControleDans{titre sur la 1re, hors bandeau}{Tit}{\PliPremiere}{\BandeauAccent}{\SecB}{\SecH}
\ControleDans{version sur la 1re, hors bandeau}{Ver}{\PliPremiere}{\BandeauAccent}{\SecB}{\SecH}
\ControleDans{texte de 4e dans la zone de securite}{Qua}{\SecG}{\PliQuatre}{\SecB}{\SecH}
\ControleDans{credits de 4e dans la zone de securite}{Cre}{\SecG}{\PliQuatre}{\SecB}{\SecH}
\ControleDisjoints{titre / version}{Tit}{\Verxmin}{\Verxmax}{\Verymin}{\Verymax}
\ControleDisjoints{texte de 4e / credits}{Qua}{\Crexmin}{\Crexmax}{\Creymin}{\Creymax}
\ControleDisjoints{texte de 4e / emplacement ISBN}{Qua}{\IsbnG}{\IsbnD}{\SecB}{\IsbnH}
\ControleDisjoints{credits / emplacement ISBN}{Cre}{\IsbnG}{\IsbnD}{\SecB}{\IsbnH}
\ifnum\fpeval{\Dos >= 6}=1
  \MesureNoeud{dos}{Dos}{0pt}
  \ControleDans{texte du dos entre les plis}{Dos}{\DosG}{\DosD}{\SecB}{\SecH}
\else
  \typeout{COUVERTURE-CONTROLE OK dos de \Dos\space mm : trop etroit, sans texte}%
\fi

\end{tikzpicture}
\typeout{COUVERTURE-BILAN alertes=\the\NbAlertes\space; classe=\Classe\space;
  format=\Format\space; reliure=\Reliure\space; pages=\NbPages\space;
  version=\Version\space;
  dos=\Dos\space mm (\DosOrigine); fini=\FiniL x \PageHauteur\space mm;
  papier=\PapierL x \PapierH\space mm}
\end{document}
"
%
%  ---------------------------------------------------------------------
%  DEUX MODES, un seul paramètre : \Reperes
%
%    oui  — mode ÉPREUVE. Ajoute une marge technique, les traits de coupe,
%           les limites du dos, la zone de sécurité et les cotes. C'est ce
%           qu'il faut pour vérifier à l'écran.
%    non  — mode PRODUCTION. Papier = format final + fond perdu, rien
%           d'autre. C'est ce fichier-là qu'on envoie à l'imprimeur.
%
%  ---------------------------------------------------------------------
%  CONTRÔLE DE SUPERPOSITION
%
%  À la fin du dessin, chaque bloc de texte est mesuré par ses ancres TikZ
%  (quatre coins, ce qui vaut aussi pour le texte pivoté du dos) et comparé
%  à la zone où il doit tenir, puis aux blocs voisins. Le journal porte une
%  ligne par vérification :
%
%      COUVERTURE-CONTROLE OK      <quoi> : ...
%      COUVERTURE-CONTROLE ALERTE  <quoi> : ...
%
%  et un bilan : COUVERTURE-BILAN alertes=<n> ; ...
%
%  Ce que le contrôle VOIT : un bloc qui franchit le pli du dos, la coupe,
%  la zone de sécurité ou le bord du bandeau ; deux blocs qui se chevauchent ;
%  un bloc qui empiète sur l'emplacement réservé au code-barres ISBN.
%  Ce qu'il NE VOIT PAS : une ligne trop longue à l'INTÉRIEUR d'un bloc à
%  largeur fixe (la boîte garde sa largeur, le texte déborde). Celle-là se
%  signale par un « Overfull \hbox » dans le même journal, que la fenêtre
%  compte aussi.
%
%  ---------------------------------------------------------------------
%  AVERTISSEMENT SUR L'ÉPAISSEUR DU DOS
%
%  Le dos se calcule ainsi :
%
%      épaisseur d'une feuille (mm) = grammage x bouffant / 1000
%      dos (mm)                     = (nombre de pages / 2) x épaisseur
%
%  On divise par deux parce qu'une feuille porte DEUX pages.
%
%  Le « bouffant » (main, volume, bulk) dépend du papier et n'est pas
%  devinable : un offset 80 g courant est autour de 1,30, un bouffant
%  d'édition monte à 1,75 et plus. UN ÉCART DE BOUFFANT DÉPLACE LE DOS DE
%  PLUSIEURS MILLIMÈTRES. Demander à l'imprimeur soit son bouffant, soit
%  directement l'épaisseur du dos, et reporter ci-dessous. Tant que cette
%  valeur n'est pas confirmée, la couverture n'est pas bonne à tirer.
%
%  Si l'imprimeur donne le dos en clair, forcer la valeur avec \DosImpose.
% =====================================================================

\documentclass[11pt]{article}

% ---------------------------------------------------------------------
% 1. PARAMÈTRES — c'est la seule zone à modifier
% ---------------------------------------------------------------------

% Classe de l'ouvrage : N, E, A, NE, EA ou NEA.
% (\providecommand : une valeur passée en ligne de commande l'emporte.)
\providecommand{\Classe}{NEA}

% Mode épreuve (oui) ou production (non).
\providecommand{\Reperes}{oui}

% Format du livre FINI : a4, 20x24 ou 20x24-marge — les noms de --format
% dans build_book.py. Les deux maquettes 20 x 24 ne diffèrent que par leur
% colonne de marge intérieure : le livre fini, donc la couverture, est le
% même (200 x 240 mm).
\providecommand{\Format}{a4}

% Reliure : « souple » (dos carré collé, fond perdu 3 mm) ou « rigide »
% (couverture rigide, dite hardcover). En rigide, consignes de l'imprimeur
% du 17/09/2026 : chaque plat déborde du livre de 5 mm en largeur et en
% hauteur (le carton dépasse du bloc), et le fichier porte 15 mm de rabat
% autour du carton au lieu de 3 mm de fond perdu. Les aplats vont déjà au
% bord du papier : ils remplissent le rabat d'eux-mêmes.
\providecommand{\Reliure}{souple}

% Pagination et version imposées de l'extérieur. Vides = tables du § 2.
% La fenêtre les renseigne toujours : la pagination vient du journal du
% livre réellement compilé, dans le format demandé.
\providecommand{\PagesForcees}{}
\providecommand{\VersionForcee}{}

% Papier de l'INTÉRIEUR du livre (pas de la couverture).
% CoolLibri, papier « Standard 90 g blanc ».
\def\Grammage{90}        % g/m²
\def\Bouffant{1.206}     % main DÉDUITE, voir ci-dessous

% Épaisseur du dos IMPOSÉE par l'imprimeur, en mm, par classe.
% Une valeur > 0 prime sur le calcul — MAIS SEULEMENT pour la pagination et
% le format pour lesquels l'imprimeur l'a donnée (\DosImposePages<classe>,
% \DosImposeFormat<classe>). Hors de ce cas, le dos est calculé et le
% journal le dit : un dos de 14 mm donné pour 258 pages A4 ne vaut rien pour
% 262 pages, ni pour le même livre recomposé en 20 x 24.
%
% ATTENTION : cette valeur est propre à CHAQUE ÉDITION. Un dos unique
% appliqué aux quatre livres serait faux pour trois d'entre eux.
%
%  N   — 14 mm, VALEUR DONNÉE PAR COOLIBRI pour 258 pages en 90 g, A4.
%        Fichier attendu : 440 x 303 mm, ce que ce gabarit produit en mode
%        production (2 x 210 + 14 + 2 x 3 = 440 ; 297 + 2 x 3 = 303).
%        Le N de la a.3 compte 262 pages : cette valeur ne s'y applique plus.
%  E, A, NEA — non demandés à l'imprimeur : ils sont CALCULÉS avec la main
%        déduite ci-dessous. À faire confirmer avant tout tirage.
%
% La main de 1,206 n'est pas publiée par CoolLibri : elle est déduite de
% leur propre chiffre, 14 mm / (0,090 x 258/2) = 1,206. Elle reproduit donc
% exactement le dos de N, et n'est qu'une ESTIMATION pour les autres.
% Écarts que cela recouvre, pour mémoire : une main de 1,0 donnerait
% 11,61 mm pour N, une main de 1,3 en donnerait 16,77. Il ne fallait pas
% deviner.
%  E   — 17 mm, VALEUR DONNÉE PAR L'IMPRIMEUR le 17/09/2026 pour 276 pages,
%        20 x 24, couverture RIGIDE, reliure collée (Klebebindung), papier
%        115 g. En reliure cousue (Fadenheftung), il annonce 18 mm.
%
% \DosImposeReliure<classe> : la reliure pour laquelle la valeur a été
% donnée. Le dos d'une couverture rigide n'est pas celui d'un dos carré
% collé souple, même à pagination égale.
\def\DosImposeN{14}   \def\DosImposePagesN{258}  \def\DosImposeFormatN{a4}
\def\DosImposeE{17}   \def\DosImposePagesE{276}  \def\DosImposeFormatE{20x24-marge}
\def\DosImposeA{0}    \def\DosImposePagesA{0}    \def\DosImposeFormatA{}
\def\DosImposeNE{0}   \def\DosImposePagesNE{0}   \def\DosImposeFormatNE{}
\def\DosImposeEA{0}   \def\DosImposePagesEA{0}   \def\DosImposeFormatEA{}
\def\DosImposeNEA{0}  \def\DosImposePagesNEA{0}  \def\DosImposeFormatNEA{}
\def\DosImposeReliureN{souple}  \def\DosImposeReliureE{rigide}
\def\DosImposeReliureA{souple}  \def\DosImposeReliureNE{souple}
\def\DosImposeReliureEA{souple} \def\DosImposeReliureNEA{souple}

% Fond perdu, en mm. 3 mm est la valeur courante ; certains imprimeurs
% demandent 5 mm.
\def\FondPerdu{3}

% Marge technique autour du fond perdu, en mm. Sert à loger les traits de
% coupe en mode épreuve. Ignorée en mode production.
\def\MargeTech{12}

% Zone de sécurité, en mm : aucun texte ne doit s'en approcher davantage
% du bord de coupe.
\def\Securite{8}

% ---------------------------------------------------------------------
% 2. DONNÉES PAR CLASSE
%    pages : nombre de pages du PDF intérieur, COUVERTURE NON COMPRISE.
%    Un dos carré collé impose un nombre de pages PAIR ; le fichier
%    refuse de compiler si le compte est impair.
%
%    Ces tables ne servent que si \PagesForcees et \VersionForcee sont
%    vides, c'est-à-dire à la main. Valeurs A4 mesurées sur les PDF de la
%    a.3 (CLAUDE.md § 1, 05/09/2026). En 20 x 24, la pagination change
%    entièrement : NEA passe de 816 à 1036 pages.
% ---------------------------------------------------------------------

\def\DonneesN{262}
\def\DonneesE{214}
\def\DonneesA{386}
\def\DonneesNE{0}      % édition jamais produite
\def\DonneesEA{0}      % édition jamais produite
\def\DonneesNEA{816}

% Version imprimée sur le dos et la 1re de couverture, PAR CLASSE.
% « Les versions sont suivies par classe, chacune évoluant à son rythme »
% (CHANGELOG.md).
\def\VersionN{a.3}
\def\VersionE{a.3}
\def\VersionA{a.3}
\def\VersionNE{a.3}
\def\VersionEA{a.3}
\def\VersionNEA{a.3}

% ---------------------------------------------------------------------
% 3. TEXTES
% ---------------------------------------------------------------------

\def\TitreOuvrage{50\,Ohm}
\def\SousTitre{Préparation à l'examen radioamateur}

% Sous-titre propre à chaque classe (repris de FR_TITLES, build_book.py).
\def\SousTitreN{Cours complet — classe N}
\def\SousTitreE{Cours complet — classe E}
\def\SousTitreA{Cours complet — classe A}
\def\SousTitreNE{Cours complet — classes N et E}
\def\SousTitreEA{Cours complet — classes E et A}
\def\SousTitreNEA{Cours complet — classes N, E et A}

% Lettres du filigrane, empilées de haut en bas (cf. build_book.py v0.20).
\def\LettresN{N}
\def\LettresE{E}
\def\LettresA{A}
\def\LettresNE{N,E}
\def\LettresEA{E,A}
\def\LettresNEA{N,E,A}

% Texte de 4e de couverture.
% NB : il parle des classes N, E et A quelle que soit l'édition. Juste pour
% N, E, A et NEA ; approximatif pour NE et EA ; faux pour SWL, d'où son refus.
\def\QuatriemeTexte{%
Ce guide d'étude prépare aux examens radioamateur allemands des classes
N, E et A. Il traduit en français les contenus du projet 50ohm.de, publiés
par le référat AJW du DARC e.\,V., et les complète d'encarts « En France »
qui signalent, sans rien retirer au texte allemand, ce que la
réglementation française prévoit sur le même sujet.\par\smallskip
Les questions officielles du catalogue allemand accompagnent chaque
notion, avec leurs réponses.%
}

\def\Auteur{Traduit avec l'aide d'une IA par Pierre F4JWI}
\def\Depot{github.com/DrPLG/50ohm-fr}

% =====================================================================
%  À PARTIR D'ICI, PLUS RIEN À MODIFIER
% =====================================================================

\usepackage{iftex}
\RequireLuaTeX
\usepackage{xfp}
\usepackage{tikz}
\usetikzlibrary{calc}
\usepackage{fontspec}
\usepackage[french]{babel}
\usepackage{geometry}

% ---------------------------------------------------------------------
% 4. RÉSOLUTION DES DONNÉES DE LA CLASSE
% ---------------------------------------------------------------------

\makeatletter
% Tome SWL : refus explicite, plutôt qu'un « classe inconnue » qui laisserait
% croire à un oubli. Le texte de 4e parle des classes N, E et A.
\def\@swlmot{SWL}
\ifx\Classe\@swlmot
  \PackageError{couverture}{Pas de couverture pour le tome SWL}%
    {Sa 4e de couverture n'est pas redigee : le texte commun parle des
     classes N, E et A. C'est une decision editoriale a prendre d'abord.}%
\fi

\def\@resoudre#1#2{%
  \expandafter\let\expandafter#1\csname #2\Classe\endcsname
  \ifx#1\relax
    \PackageError{couverture}{Classe \Classe\space inconnue : pas de #2\Classe}%
      {Valeurs admises : N, E, A, NE, EA, NEA.}%
  \fi}
\@resoudre\NbPages{Donnees}
\@resoudre\SousTitreClasse{SousTitre}
\@resoudre\Lettres{Lettres}
\@resoudre\DosImpose{DosImpose}
\@resoudre\DosImposePages{DosImposePages}
\@resoudre\DosImposeFormat{DosImposeFormat}
\@resoudre\DosImposeReliure{DosImposeReliure}
\@resoudre\Version{Version}

% Valeurs imposées de l'extérieur. Test de vacuité par \detokenize : il ne
% dépend pas de la façon dont la macro a été définie.
\if\relax\detokenize\expandafter{\PagesForcees}\relax\else
  \let\NbPages\PagesForcees
\fi
\if\relax\detokenize\expandafter{\VersionForcee}\relax\else
  \let\Version\VersionForcee
\fi

% Format du livre fini.
\def\@formatAquatre{a4}
\def\@formatVingt{20x24}
\def\@formatVingtMarge{20x24-marge}
\ifx\Format\@formatAquatre
  \def\PageLargeur{210}\def\PageHauteur{297}
\else\ifx\Format\@formatVingt
  \def\PageLargeur{200}\def\PageHauteur{240}
\else\ifx\Format\@formatVingtMarge
  \def\PageLargeur{200}\def\PageHauteur{240}
\else
  \PackageError{couverture}{Format \Format\space inconnu}%
    {Valeurs admises : a4, 20x24, 20x24-marge.}%
  \def\PageLargeur{210}\def\PageHauteur{297}
\fi\fi\fi
% Reliure. Hors de « souple » et « rigide » : erreur, pas de supposition.
\def\@reliuresouple{souple}
\def\@reliurerigide{rigide}
\ifx\Reliure\@reliurerigide
  \edef\PageLargeur{\fpeval{\PageLargeur + 5}}
  \edef\PageHauteur{\fpeval{\PageHauteur + 5}}
  \def\FondPerdu{15}
\else\ifx\Reliure\@reliuresouple\else
  \PackageError{couverture}{Reliure \Reliure\space inconnue}%
    {Valeurs admises : souple, rigide.}%
\fi\fi
\makeatother

% Garde-fou : une édition jamais produite n'a pas de couverture.
\ifnum\NbPages=0
  \PackageError{couverture}{L'edition \Classe\space n'a pas de pagination}%
    {Cette edition n'a jamais ete compilee. Renseigner \string\Donnees\Classe
     ou passer \string\PagesForcees.}%
\fi

% Avertissement : l'imprimeur demande un nombre de pages multiple de 4
% (17/09/2026). Pas une erreur : c'est son exigence, pas une loi du papier.
\ifnum\fpeval{\NbPages - 4*floor(\NbPages/4)}=0 \else
  \typeout{COUVERTURE-AVERTISSEMENT \NbPages\space pages : pas un multiple
    de 4, que l'imprimeur exige -- ajouter des pages a l'interieur}%
\fi

% Garde-fou : un dos carré collé impose un nombre de pages pair.
\ifodd\NbPages
  \PackageError{couverture}{Nombre de pages impair (\NbPages)}%
    {Un dos carre colle exige un nombre PAIR de pages : une feuille en
     porte deux. Ajouter une page blanche a l'interieur.}%
\fi

% ---------------------------------------------------------------------
% 5. GÉOMÉTRIE — c'est ici que se joue l'exactitude du dos
% ---------------------------------------------------------------------

% Épaisseur du dos : valeur imposée si elle est donnée ET qu'elle vaut pour
% cette pagination et ce format ; calcul sinon.
\edef\DosCalcule{\fpeval{round(\NbPages / 2 * \Grammage * \Bouffant / 1000, 2)}}
\def\DosOrigine{calculee}
\ifnum\fpeval{\DosImpose > 0}=1
  \ifnum\fpeval{\NbPages = \DosImposePages}=1
    \ifx\Format\DosImposeFormat
      \ifx\Reliure\DosImposeReliure
        \def\DosOrigine{imposee}
      \fi
    \fi
  \fi
  \def\MotCalculee{calculee}
  \ifx\DosOrigine\MotCalculee
    \typeout{COUVERTURE-AVERTISSEMENT dos impose de \DosImpose\space mm
      donne pour \DosImposePages\space pages en \DosImposeFormat\space
      (reliure \DosImposeReliure) -- inapplicable a \NbPages\space pages en
      \Format\space (reliure \Reliure), dos calcule}%
  \fi
\fi
% NB : pas de \@ dans ces noms. Cette section est hors \makeatletter, et
% « \def\@imposee{...} » y redéfinit \@ (suivi du texte « imposee ») : le
% test échouait toujours, en silence. Constaté le 16/09/2026 — le dos de N
% tombait juste par coïncidence, 14,0 mm calculés pour 14 imposés.
\def\MotRigide{rigide}
\def\MotImposee{imposee}
\ifx\DosOrigine\MotImposee
  \edef\Dos{\DosImpose}
\else
  \edef\Dos{\DosCalcule}
\fi

% Marge technique : uniquement en mode épreuve.
\def\ouimot{oui}
\ifx\Reperes\ouimot \edef\Tech{\MargeTech}\else \def\Tech{0}\fi

% Dimensions du papier.
%   largeur = marge + fond perdu + 4e + dos + 1re + fond perdu + marge
%   hauteur = marge + fond perdu + hauteur + fond perdu + marge
\edef\PapierL{\fpeval{2*\PageLargeur + \Dos + 2*\FondPerdu + 2*\Tech}}
\edef\PapierH{\fpeval{\PageHauteur + 2*\FondPerdu + 2*\Tech}}

\geometry{paperwidth=\PapierL mm, paperheight=\PapierH mm,
          margin=0pt, nohead, nofoot}

% ---------------------------------------------------------------------
%  TOUTES LES COTES SONT FIGÉES ICI, ET NULLE PART AILLEURS.
%
%  Raison, payée par trois compilations ratées : un \fpeval laissé dans
%  une coordonnée ou dans le texte d'un nœud TikZ n'est pas développé au
%  moment où TikZ analyse le chemin. On obtient « Giving up on this path.
%  Did you forget a semicolon? », qui accuse une ponctuation alors que la
%  faute est ailleurs. Le corps du document ci-dessous n'emploie donc que
%  des macros déjà calculées.
% ---------------------------------------------------------------------

% Repères horizontaux, mesurés depuis le bord GAUCHE du papier.
\edef\CoupeG{\fpeval{\Tech + \FondPerdu}}          % coupe gauche (4e de couv.)
\edef\DosG{\fpeval{\CoupeG + \PageLargeur}}         % dos, début
\edef\DosD{\fpeval{\DosG + \Dos}}                   % dos, fin
\edef\CoupeD{\fpeval{\DosD + \PageLargeur}}         % coupe droite (1re de couv.)
% Repères verticaux, depuis le bord BAS.
\edef\CoupeB{\fpeval{\Tech + \FondPerdu}}
\edef\CoupeH{\fpeval{\CoupeB + \PageHauteur}}
% Axes.
\edef\DosAxe{\fpeval{(\DosG + \DosD)/2}}
\edef\MilieuH{\fpeval{(\CoupeB + \CoupeH)/2}}
\edef\MilieuL{\fpeval{\PapierL/2}}

% Format fini, pour les cotes affichées.
\edef\FiniL{\fpeval{2*\PageLargeur + \Dos}}

% Bandeau de la 1re de couverture : même proportion que la page de titre.
\edef\BandeauL{\fpeval{round(0.34*\PageLargeur, 2)}}
\edef\BandeauG{\fpeval{\CoupeD - \BandeauL}}
\edef\BandeauAccent{\fpeval{\BandeauG - 0.7}}
\edef\DosAccentG{\fpeval{\DosG - 0.7}}
\edef\DosAccentD{\fpeval{\DosD + 0.7}}
% Couverture rigide : le bleu du dos déborde de 2 mm sur chaque plat, pour
% absorber la tolérance de l'imprimeur sur l'épaisseur du dos (±10 %).
\edef\DosEtenduG{\fpeval{\DosG - 2}}
\edef\DosEtenduD{\fpeval{\DosD + 2}}

% Filigrane de classe.
\edef\FiligraneX{\fpeval{(\BandeauG + \CoupeD)/2}}
\edef\FiligraneY{\fpeval{\CoupeH - 6}}

% Bloc-titre de la 1re de couverture.
\edef\TitreX{\fpeval{\DosD + 12}}
\edef\TitreY{\fpeval{\CoupeH - 72}}
\edef\TitreW{\fpeval{\BandeauG - \DosD - 16}}
\edef\VersionY{\fpeval{\CoupeB + 14}}

% 4e de couverture.
\edef\QuatreX{\fpeval{\CoupeG + \Securite + 6}}
\edef\QuatreY{\fpeval{\CoupeH - \Securite - 14}}
\edef\QuatreBasY{\fpeval{\CoupeB + \Securite + 22}}
\edef\QuatreW{\fpeval{\PageLargeur - 2*\Securite - 12}}

% Zone de sécurité.
\edef\SecG{\fpeval{\CoupeG + \Securite}}
\edef\SecD{\fpeval{\CoupeD - \Securite}}
\edef\SecB{\fpeval{\CoupeB + \Securite}}
\edef\SecH{\fpeval{\CoupeH - \Securite}}

% Zones de contrôle : ce qui doit rester de part et d'autre du pli du dos.
\edef\PliQuatre{\fpeval{\DosG - \Securite}}         % 4e : limite côté dos
\edef\PliPremiere{\fpeval{\DosD + \Securite}}       % 1re : limite côté dos

% Traits de coupe et cotes, dans la marge technique.
\edef\TechExt{\fpeval{\Tech - 8}}
\edef\TechInt{\fpeval{\Tech - 1}}
\edef\TechCote{\fpeval{\Tech - 9}}
\edef\HautInt{\fpeval{\PapierH - \Tech + 1}}
\edef\HautExt{\fpeval{\PapierH - \Tech + 8}}
\edef\HautCote{\fpeval{\PapierH - \Tech + 9}}
\edef\DroitInt{\fpeval{\PapierL - \Tech + 1}}
\edef\DroitExt{\fpeval{\PapierL - \Tech + 8}}

% Emplacement du code-barres ISBN : en bas à DROITE de la 4e de couverture,
% selon l'usage éditorial, et non à gauche où il chevaucherait les mentions
% de licence.
\edef\IsbnD{\fpeval{\DosG - \Securite}}
\edef\IsbnG{\fpeval{\IsbnD - 40}}
\edef\IsbnH{\fpeval{\SecB + 25}}
\edef\IsbnCX{\fpeval{\IsbnG + 20}}
\edef\IsbnCY{\fpeval{\SecB + 12.5}}

% Cartouches de cotes, dans la marge technique. Ancrés vers l'INTÉRIEUR du
% papier : ancrés vers l'extérieur, ils débordaient et étaient rognés.
\edef\CoteBasY{\fpeval{1}}
\edef\CoteHautY{\fpeval{\PapierH - 1}}

% ---------------------------------------------------------------------
% 6. COULEURS ET POLICES — reprises de la page de titre (build_book.py)
% ---------------------------------------------------------------------

\definecolor{TitleBand}{cmyk}{.9,.55,.1,.35}   % bleu profond
\definecolor{TitleAccent}{cmyk}{.05,.8,1,0}    % rouge DARC (accent fin)
\definecolor{Repere}{cmyk}{0,1,1,0}            % magenta : repères de montage

\setmainfont{Libertinus Serif}[
  ItalicFont     = Libertinus Serif Italic,
  BoldFont       = Libertinus Serif Bold,
  BoldItalicFont = Libertinus Serif Bold Italic]

\pagestyle{empty}

% ---------------------------------------------------------------------
% 7. FILIGRANE DE CLASSE
%    Le texte est construit ici, et non dans le nœud : un \\ produit à
%    l'intérieur d'un \foreach referme le groupe d'alignement de TikZ.
%    Une lettre seule est grande ; plusieurs s'empilent et se réduisent
%    (build_book.py v0.20 : « NEA » à 105 pt débordait encore).
% ---------------------------------------------------------------------

\makeatletter
\def\FiligraneTexte{}
\newif\ifpremierelettre \premierelettretrue
\count255=0
\@for\@tempa:=\Lettres\do{%
  \advance\count255 by 1\relax
  \ifpremierelettre \premierelettrefalse
  \else \g@addto@macro\FiligraneTexte{\\}\fi
  \expandafter\g@addto@macro\expandafter\FiligraneTexte\expandafter{\@tempa}}%
\xdef\NbLettres{\the\count255}
\makeatother

\ifnum\NbLettres=1 \def\CorpsF{150}\else\def\CorpsF{86}\fi
\edef\CorpsFI{\fpeval{round(\CorpsF*1.02, 0)}}

% Texte du cartouche de cotes, lui aussi construit hors du nœud TikZ.
%
% NE PAS y employer \textsuperscript : dans un nœud « align=center », il
% casse le chemin TikZ (« Giving up on this path. Did you forget a
% semicolon? »), et l'erreur est signalée sur le nœud SUIVANT — elle
% accuse donc une ponctuation là où la faute est ici. Isolé par essais :
% $^2$ et le caractère ² passent, \textsuperscript non.
\def\CotesTexte{%
  \Classe\ --- \NbPages\ pages --- papier \Grammage\ g/m$^2$,
  bouffant \Bouffant\ --- fond perdu \FondPerdu\ mm\\
  papier \PapierL\ mm sur \PapierH\ mm --- fini \FiniL\ mm sur
  \PageHauteur\ mm\\
  EPREUVE, ne pas envoyer a l'impression%
}

% ---------------------------------------------------------------------
% 8. CONTRÔLE DE SUPERPOSITION
%
%  \MesureNoeud{nœud}{préfixe}{retrait} mesure la boîte d'un nœud par ses
%  quatre coins (un nœud pivoté a ses ancres pivotées : le min et le max des
%  quatre suffisent) et la ramène au repère du papier, en mm. Le point (0,0)
%  est lu par le même « let » que les ancres : les deux passent donc par la
%  même transformation, et la comparaison ne dépend pas de l'endroit où
%  TeX a posé l'image sur la page.
%
%  Le retrait sert au seul filigrane, pour mesurer les lettres et non leur
%  cadre : c'est le bord des lettres que le détourage rognerait.
%
%  PIÈGE MESURÉ le 16/09/2026 : la marge interne TikZ (.3333em) n'est PAS
%  évaluée dans la police du nœud, mais dans celle qui entoure le dessin.
%  Instrumenté : dans le nœud, em = 150 pt ; boîte brute du « N » :
%  41,7 x 37,5 mm, soit la lettre plus 2 x 1,3 mm — 1,3 mm étant .3333em
%  à 11 pt. Un retrait calculé à 150 pt (17,6 mm) ramenait la lettre à
%  6,6 mm de large, et le contrôle aurait dit OK sur une mesure absurde.
% ---------------------------------------------------------------------

\newcount\NbAlertes
\newdimen\FiligraneSep
\newsavebox\CreditsBoite

% Crédits de 4e : lignes coupées à la main. Leur largeur RÉELLE est celle
% de la plus longue ligne, mesurée dans une boîte à part (voir § 8, ISBN).
\def\CreditsTexte{%
    Réalisé à partir des contenus de 50ohm.de (en allemand)\\
    50ohm.de-Autorenteam, coordonné par le référat AJW du DARC e.\,V.\\
    \Auteur\\[3pt]
    Licence CC BY 4.0 --- \texttt{\Depot}}

\def\MesureNoeud#1#2#3{%
  \path let \p1=(#1.south west), \p2=(#1.north east),
            \p3=(#1.north west), \p4=(#1.south east), \p5=(0,0) in
    \pgfextra{%
      \expandafter\xdef\csname #2xmin\endcsname{\fpeval{round((min(\x1,\x2,\x3,\x4) - \x5 + #3)/1mm, 2)}}%
      \expandafter\xdef\csname #2xmax\endcsname{\fpeval{round((max(\x1,\x2,\x3,\x4) - \x5 - #3)/1mm, 2)}}%
      \expandafter\xdef\csname #2ymin\endcsname{\fpeval{round((min(\y1,\y2,\y3,\y4) - \y5 + #3)/1mm, 2)}}%
      \expandafter\xdef\csname #2ymax\endcsname{\fpeval{round((max(\y1,\y2,\y3,\y4) - \y5 - #3)/1mm, 2)}}%
    };}

% \ControleDans{libellé}{préfixe}{xmin}{xmax}{ymin}{ymax}
\def\ControleDans#1#2#3#4#5#6{%
  \ifnum\fpeval{\csname #2xmin\endcsname >= #3 && \csname #2xmax\endcsname <= #4
             && \csname #2ymin\endcsname >= #5 && \csname #2ymax\endcsname <= #6}=1
    \typeout{COUVERTURE-CONTROLE OK #1 : x \csname #2xmin\endcsname-\csname #2xmax\endcsname,
      y \csname #2ymin\endcsname-\csname #2ymax\endcsname\space dans x #3-#4, y #5-#6}%
  \else
    \global\advance\NbAlertes by 1
    \typeout{COUVERTURE-CONTROLE ALERTE #1 : x \csname #2xmin\endcsname-\csname #2xmax\endcsname,
      y \csname #2ymin\endcsname-\csname #2ymax\endcsname\space HORS de x #3-#4, y #5-#6}%
  \fi}

% \ControleDisjoints{libellé}{préfixe}{xmin}{xmax}{ymin}{ymax}
% Deux rectangles se chevauchent si et seulement s'ils se recouvrent sur
% les DEUX axes à la fois.
\def\ControleDisjoints#1#2#3#4#5#6{%
  \ifnum\fpeval{\csname #2xmin\endcsname < #4 && #3 < \csname #2xmax\endcsname
             && \csname #2ymin\endcsname < #6 && #5 < \csname #2ymax\endcsname}=1
    \global\advance\NbAlertes by 1
    \typeout{COUVERTURE-CONTROLE ALERTE #1 : chevauchement}%
  \else
    \typeout{COUVERTURE-CONTROLE OK #1 : disjoints}%
  \fi}

\begin{document}
% Marge interne du filigrane, évaluée dans la police qui ENTOURE le dessin :
% c'est là que TikZ évalue le .3333em de « inner sep » (cf. § 8).
\FiligraneSep=.3333em
% Largeur réelle des crédits : la plus longue ligne, dans leur police.
\sbox\CreditsBoite{\small\begin{tabular}{@{}l@{}}\CreditsTexte\end{tabular}}%
\begin{tikzpicture}[remember picture, overlay,
                    x=1mm, y=1mm, shift={(current page.south west)}]

% ---- Fonds ----------------------------------------------------------
% La 4e de couverture reste sur le blanc du papier.
% Le bandeau bleu couvre le dos, puis reprend au tiers extérieur de la
% 1re de couverture — même partition 2/3-1/3 que la page de titre.
% Couverture rigide : pas de filets au dos (ils tomberaient à côté des
% plis, vu la tolérance), et un dos bleu élargi de 2 mm de chaque côté.
\ifx\Reliure\MotRigide
  \fill[TitleBand] (\DosEtenduG, 0) rectangle (\DosEtenduD, \PapierH);
\else
  \fill[TitleBand] (\DosG, 0) rectangle (\DosD, \PapierH);
\fi
\fill[TitleBand] (\BandeauG, 0) rectangle (\PapierL, \PapierH);
% Filets d'accent : à la jonction sur la 1re, de part et d'autre du dos.
\fill[TitleAccent] (\BandeauAccent, 0) rectangle (\BandeauG, \PapierH);
\ifx\Reliure\MotRigide\else
  \fill[TitleAccent] (\DosAccentG, 0) rectangle (\DosG, \PapierH);
  \fill[TitleAccent] (\DosD, 0) rectangle (\DosAccentD, \PapierH);
\fi

% ---- 1re de couverture ----------------------------------------------
% Filigrane de classe, en réserve claire, détouré au bandeau.
\begin{scope}
\clip (\BandeauG, 0) rectangle (\PapierL, \PapierH);
\node[name=filigrane, anchor=north, align=center, text=white, opacity=0.16,
      font=\fontsize{\CorpsF}{\CorpsFI}\selectfont\bfseries]
  at (\FiligraneX, \FiligraneY) {\FiligraneTexte};
\end{scope}

\node[name=titre, anchor=west, align=left, text width=\TitreW mm]
  at (\TitreX, \TitreY) {%
    {\fontsize{54}{56}\selectfont\bfseries \TitreOuvrage}\\[10pt]
    {\Large\color{TitleBand}\bfseries \SousTitre}\\[16pt]
    {\large\color{black!75} \SousTitreClasse}};

\node[name=version, anchor=south west, align=left, font=\small, text=black!70]
  at (\TitreX, \VersionY) {Version \Version\\ Licence CC BY 4.0};

% ---- Dos -------------------------------------------------------------
% Un dos étroit ne reçoit pas de texte : en deçà de 6 mm, le débord au
% massicot ferait mordre le texte vertical sur les plats.
\ifnum\fpeval{\Dos >= 6}=1
  \node[name=dos, rotate=90, anchor=center, text=white, font=\bfseries,
        inner sep=0pt]
    at (\DosAxe, \MilieuH)
    {\TitreOuvrage\ \ --- \ \SousTitreClasse\ \ --- \ \Version};
\fi

% ---- 4e de couverture ------------------------------------------------
\node[name=quatre, anchor=north west, align=justify, text width=\QuatreW mm]
  at (\QuatreX, \QuatreY) {\QuatriemeTexte};

% Le texte vit dans \CreditsTexte (§ 8) pour être mesuré à part : la boîte
% du nœud a la largeur fixe \QuatreW et couvre l'emplacement ISBN sans
% qu'aucun caractère n'y soit. Retirer « text width » pour l'éviter change
% l'interlignage — essayé le 16/09/2026, 6 934 pixels différents.
\node[name=credits, anchor=south west, align=left, font=\small, text=black!70,
      text width=\QuatreW mm]
  at (\QuatreX, \QuatreBasY) {\CreditsTexte};

% ---- Repères de fabrication (mode épreuve uniquement) ----------------
\ifx\Reperes\ouimot
  % Traits de coupe, hors du fond perdu.
  \foreach \x in {\CoupeG, \CoupeD} {
    \draw[Repere, line width=0.25pt] (\x, \TechExt) -- (\x, \TechInt);
    \draw[Repere, line width=0.25pt] (\x, \HautInt) -- (\x, \HautExt);}
  \foreach \y in {\CoupeB, \CoupeH} {
    \draw[Repere, line width=0.25pt] (\TechExt, \y) -- (\TechInt, \y);
    \draw[Repere, line width=0.25pt] (\DroitInt, \y) -- (\DroitExt, \y);}
  % Limites du dos, en pointillés.
  \foreach \x in {\DosG, \DosD} {
    \draw[Repere, line width=0.25pt, dashed] (\x, \TechExt) -- (\x, \TechInt);
    \draw[Repere, line width=0.25pt, dashed] (\x, \HautInt) -- (\x, \HautExt);}
  % Zone de sécurité.
  \draw[Repere, line width=0.2pt, dotted]
    (\SecG, \SecB) rectangle (\SecD, \SecH);
  % Emplacement réservé au code-barres ISBN.
  \draw[Repere, line width=0.2pt] (\IsbnG, \SecB) rectangle (\IsbnD, \IsbnH);
  \node[font=\tiny\sffamily, text=Repere, align=center]
    at (\IsbnCX, \IsbnCY) {emplacement\\ISBN};
  % Cotes.
  \node[anchor=north, font=\tiny\sffamily, text=Repere]
    at (\DosAxe, \CoteHautY) {dos \Dos\ mm (\DosOrigine)};
  \node[anchor=south, font=\tiny\sffamily, text=Repere, align=center]
    at (\MilieuL, \CoteBasY) {\CotesTexte};
\fi

% ---- Contrôle de superposition (§ 8) ----------------------------------
% Rien n'est dessiné : des mesures et des lignes de journal.
\MesureNoeud{filigrane}{Fil}{\the\FiligraneSep}
\MesureNoeud{titre}{Tit}{0pt}
\MesureNoeud{version}{Ver}{0pt}
\MesureNoeud{quatre}{Qua}{0pt}
\MesureNoeud{credits}{Cre}{0pt}
% Bord droit des crédits = bord gauche + marge interne + plus longue ligne.
\xdef\Crexmax{\fpeval{round(\Crexmin + (\the\FiligraneSep + \the\wd\CreditsBoite)/1mm, 2)}}

% Le filigrane peut entrer dans la zone de sécurité (il est décoratif),
% pas franchir la coupe ni sortir du bandeau — le détourage le rognerait.
\ControleDans{filigrane dans le bandeau}{Fil}{\BandeauG}{\CoupeD}{\CoupeB}{\CoupeH}
\ControleDans{titre sur la 1re, hors bandeau}{Tit}{\PliPremiere}{\BandeauAccent}{\SecB}{\SecH}
\ControleDans{version sur la 1re, hors bandeau}{Ver}{\PliPremiere}{\BandeauAccent}{\SecB}{\SecH}
\ControleDans{texte de 4e dans la zone de securite}{Qua}{\SecG}{\PliQuatre}{\SecB}{\SecH}
\ControleDans{credits de 4e dans la zone de securite}{Cre}{\SecG}{\PliQuatre}{\SecB}{\SecH}
\ControleDisjoints{titre / version}{Tit}{\Verxmin}{\Verxmax}{\Verymin}{\Verymax}
\ControleDisjoints{texte de 4e / credits}{Qua}{\Crexmin}{\Crexmax}{\Creymin}{\Creymax}
\ControleDisjoints{texte de 4e / emplacement ISBN}{Qua}{\IsbnG}{\IsbnD}{\SecB}{\IsbnH}
\ControleDisjoints{credits / emplacement ISBN}{Cre}{\IsbnG}{\IsbnD}{\SecB}{\IsbnH}
\ifnum\fpeval{\Dos >= 6}=1
  \MesureNoeud{dos}{Dos}{0pt}
  \ControleDans{texte du dos entre les plis}{Dos}{\DosG}{\DosD}{\SecB}{\SecH}
\else
  \typeout{COUVERTURE-CONTROLE OK dos de \Dos\space mm : trop etroit, sans texte}%
\fi

\end{tikzpicture}
\typeout{COUVERTURE-BILAN alertes=\the\NbAlertes\space; classe=\Classe\space;
  format=\Format\space; reliure=\Reliure\space; pages=\NbPages\space;
  version=\Version\space;
  dos=\Dos\space mm (\DosOrigine); fini=\FiniL x \PageHauteur\space mm;
  papier=\PapierL x \PapierH\space mm}
\end{document}
"
%
%  Pour le fichier destiné à l'imprimeur (sans repères) :
%
%      lualatex "\def\Reperes{non}% =====================================================================
%  couverture.tex — couverture dos carré collé pour les livres 50ohm-fr
% =====================================================================
%
%  Produit une couverture d'un seul tenant : 4e de couverture | dos |
%  1re de couverture, aux dimensions exactes du livre relié, fond perdu
%  compris. L'épaisseur du dos est CALCULÉE à partir du nombre de pages,
%  du grammage et du bouffant du papier.
%
%  ORIGINE ET ÉTAT. Écrit le 20/08/2026 dans Fichierstravail/ (hors git),
%  versé au dépôt le 16/09/2026 pour que fenetre_compilation.py puisse
%  l'appeler. Ce qui a changé à cette occasion :
%
%    - le FORMAT est un paramètre (\Format = a4, 20x24 ou 20x24-marge,
%      mêmes noms que --format de build_book.py) ;
%    - la pagination et la version peuvent être IMPOSÉES de l'extérieur
%      (\PagesForcees, \VersionForcee) : la fenêtre lit la pagination dans
%      le journal du livre qu'elle vient de compiler, au lieu de se fier à
%      une table qui vieillit ;
%    - un dos imposé par l'imprimeur n'est appliqué QUE pour la pagination
%      et le format pour lesquels il a été donné ;
%    - un CONTRÔLE DE SUPERPOSITION est calculé par LaTeX lui-même, à partir
%      des boîtes réelles des nœuds, et écrit dans le journal ;
%    - le tome SWL est refusé explicitement : sa 4e de couverture n'est pas
%      rédigée (le texte commun parle des classes N, E et A).
%
%    Neutralité vérifiée : en A4, pour N (258 p., a.2) et NEA (816 p., a.3),
%    en épreuve comme en impression, le rendu est identique au pixel près à
%    celui du gabarit d'origine.
%
%  COMPILATION (LuaLaTeX, depuis la racine du dépôt) :
%
%      lualatex couverture.tex
%      lualatex couverture.tex        <-- DEUX FOIS, obligatoirement
%
%  DEUX PASSES SONT NÉCESSAIRES. Le dessin s'ancre sur « current page »,
%  dont la position n'est connue qu'après un premier passage : à la
%  première compilation, tout se tasse dans le coin supérieur et la
%  couverture paraît ratée alors qu'elle ne l'est pas. Aucune erreur n'est
%  émise — seul le rendu le montre. Le contrôle de superposition, lui
%  aussi, ne vaut qu'à la seconde passe.
%
%  Le plus simple est de passer par fenetre_compilation.py (case
%  « couverture pour l'imprimeur »), qui enchaîne les deux passes, produit
%  l'épreuve ET le fichier d'impression, et relit le contrôle.
%
%  À la main, pour une autre classe, un autre format, une pagination mesurée :
%
%      lualatex "\def\Classe{A}\def\Format{20x24}\def\PagesForcees{482}% =====================================================================
%  couverture.tex — couverture dos carré collé pour les livres 50ohm-fr
% =====================================================================
%
%  Produit une couverture d'un seul tenant : 4e de couverture | dos |
%  1re de couverture, aux dimensions exactes du livre relié, fond perdu
%  compris. L'épaisseur du dos est CALCULÉE à partir du nombre de pages,
%  du grammage et du bouffant du papier.
%
%  ORIGINE ET ÉTAT. Écrit le 20/08/2026 dans Fichierstravail/ (hors git),
%  versé au dépôt le 16/09/2026 pour que fenetre_compilation.py puisse
%  l'appeler. Ce qui a changé à cette occasion :
%
%    - le FORMAT est un paramètre (\Format = a4, 20x24 ou 20x24-marge,
%      mêmes noms que --format de build_book.py) ;
%    - la pagination et la version peuvent être IMPOSÉES de l'extérieur
%      (\PagesForcees, \VersionForcee) : la fenêtre lit la pagination dans
%      le journal du livre qu'elle vient de compiler, au lieu de se fier à
%      une table qui vieillit ;
%    - un dos imposé par l'imprimeur n'est appliqué QUE pour la pagination
%      et le format pour lesquels il a été donné ;
%    - un CONTRÔLE DE SUPERPOSITION est calculé par LaTeX lui-même, à partir
%      des boîtes réelles des nœuds, et écrit dans le journal ;
%    - le tome SWL est refusé explicitement : sa 4e de couverture n'est pas
%      rédigée (le texte commun parle des classes N, E et A).
%
%    Neutralité vérifiée : en A4, pour N (258 p., a.2) et NEA (816 p., a.3),
%    en épreuve comme en impression, le rendu est identique au pixel près à
%    celui du gabarit d'origine.
%
%  COMPILATION (LuaLaTeX, depuis la racine du dépôt) :
%
%      lualatex couverture.tex
%      lualatex couverture.tex        <-- DEUX FOIS, obligatoirement
%
%  DEUX PASSES SONT NÉCESSAIRES. Le dessin s'ancre sur « current page »,
%  dont la position n'est connue qu'après un premier passage : à la
%  première compilation, tout se tasse dans le coin supérieur et la
%  couverture paraît ratée alors qu'elle ne l'est pas. Aucune erreur n'est
%  émise — seul le rendu le montre. Le contrôle de superposition, lui
%  aussi, ne vaut qu'à la seconde passe.
%
%  Le plus simple est de passer par fenetre_compilation.py (case
%  « couverture pour l'imprimeur »), qui enchaîne les deux passes, produit
%  l'épreuve ET le fichier d'impression, et relit le contrôle.
%
%  À la main, pour une autre classe, un autre format, une pagination mesurée :
%
%      lualatex "\def\Classe{A}\def\Format{20x24}\def\PagesForcees{482}\input{couverture.tex}"
%
%  Pour le fichier destiné à l'imprimeur (sans repères) :
%
%      lualatex "\def\Reperes{non}\input{couverture.tex}"
%
%  ---------------------------------------------------------------------
%  DEUX MODES, un seul paramètre : \Reperes
%
%    oui  — mode ÉPREUVE. Ajoute une marge technique, les traits de coupe,
%           les limites du dos, la zone de sécurité et les cotes. C'est ce
%           qu'il faut pour vérifier à l'écran.
%    non  — mode PRODUCTION. Papier = format final + fond perdu, rien
%           d'autre. C'est ce fichier-là qu'on envoie à l'imprimeur.
%
%  ---------------------------------------------------------------------
%  CONTRÔLE DE SUPERPOSITION
%
%  À la fin du dessin, chaque bloc de texte est mesuré par ses ancres TikZ
%  (quatre coins, ce qui vaut aussi pour le texte pivoté du dos) et comparé
%  à la zone où il doit tenir, puis aux blocs voisins. Le journal porte une
%  ligne par vérification :
%
%      COUVERTURE-CONTROLE OK      <quoi> : ...
%      COUVERTURE-CONTROLE ALERTE  <quoi> : ...
%
%  et un bilan : COUVERTURE-BILAN alertes=<n> ; ...
%
%  Ce que le contrôle VOIT : un bloc qui franchit le pli du dos, la coupe,
%  la zone de sécurité ou le bord du bandeau ; deux blocs qui se chevauchent ;
%  un bloc qui empiète sur l'emplacement réservé au code-barres ISBN.
%  Ce qu'il NE VOIT PAS : une ligne trop longue à l'INTÉRIEUR d'un bloc à
%  largeur fixe (la boîte garde sa largeur, le texte déborde). Celle-là se
%  signale par un « Overfull \hbox » dans le même journal, que la fenêtre
%  compte aussi.
%
%  ---------------------------------------------------------------------
%  AVERTISSEMENT SUR L'ÉPAISSEUR DU DOS
%
%  Le dos se calcule ainsi :
%
%      épaisseur d'une feuille (mm) = grammage x bouffant / 1000
%      dos (mm)                     = (nombre de pages / 2) x épaisseur
%
%  On divise par deux parce qu'une feuille porte DEUX pages.
%
%  Le « bouffant » (main, volume, bulk) dépend du papier et n'est pas
%  devinable : un offset 80 g courant est autour de 1,30, un bouffant
%  d'édition monte à 1,75 et plus. UN ÉCART DE BOUFFANT DÉPLACE LE DOS DE
%  PLUSIEURS MILLIMÈTRES. Demander à l'imprimeur soit son bouffant, soit
%  directement l'épaisseur du dos, et reporter ci-dessous. Tant que cette
%  valeur n'est pas confirmée, la couverture n'est pas bonne à tirer.
%
%  Si l'imprimeur donne le dos en clair, forcer la valeur avec \DosImpose.
% =====================================================================

\documentclass[11pt]{article}

% ---------------------------------------------------------------------
% 1. PARAMÈTRES — c'est la seule zone à modifier
% ---------------------------------------------------------------------

% Classe de l'ouvrage : N, E, A, NE, EA ou NEA.
% (\providecommand : une valeur passée en ligne de commande l'emporte.)
\providecommand{\Classe}{NEA}

% Mode épreuve (oui) ou production (non).
\providecommand{\Reperes}{oui}

% Format du livre FINI : a4, 20x24 ou 20x24-marge — les noms de --format
% dans build_book.py. Les deux maquettes 20 x 24 ne diffèrent que par leur
% colonne de marge intérieure : le livre fini, donc la couverture, est le
% même (200 x 240 mm).
\providecommand{\Format}{a4}

% Reliure : « souple » (dos carré collé, fond perdu 3 mm) ou « rigide »
% (couverture rigide, dite hardcover). En rigide, consignes de l'imprimeur
% du 17/09/2026 : chaque plat déborde du livre de 5 mm en largeur et en
% hauteur (le carton dépasse du bloc), et le fichier porte 15 mm de rabat
% autour du carton au lieu de 3 mm de fond perdu. Les aplats vont déjà au
% bord du papier : ils remplissent le rabat d'eux-mêmes.
\providecommand{\Reliure}{souple}

% Pagination et version imposées de l'extérieur. Vides = tables du § 2.
% La fenêtre les renseigne toujours : la pagination vient du journal du
% livre réellement compilé, dans le format demandé.
\providecommand{\PagesForcees}{}
\providecommand{\VersionForcee}{}

% Papier de l'INTÉRIEUR du livre (pas de la couverture).
% CoolLibri, papier « Standard 90 g blanc ».
\def\Grammage{90}        % g/m²
\def\Bouffant{1.206}     % main DÉDUITE, voir ci-dessous

% Épaisseur du dos IMPOSÉE par l'imprimeur, en mm, par classe.
% Une valeur > 0 prime sur le calcul — MAIS SEULEMENT pour la pagination et
% le format pour lesquels l'imprimeur l'a donnée (\DosImposePages<classe>,
% \DosImposeFormat<classe>). Hors de ce cas, le dos est calculé et le
% journal le dit : un dos de 14 mm donné pour 258 pages A4 ne vaut rien pour
% 262 pages, ni pour le même livre recomposé en 20 x 24.
%
% ATTENTION : cette valeur est propre à CHAQUE ÉDITION. Un dos unique
% appliqué aux quatre livres serait faux pour trois d'entre eux.
%
%  N   — 14 mm, VALEUR DONNÉE PAR COOLIBRI pour 258 pages en 90 g, A4.
%        Fichier attendu : 440 x 303 mm, ce que ce gabarit produit en mode
%        production (2 x 210 + 14 + 2 x 3 = 440 ; 297 + 2 x 3 = 303).
%        Le N de la a.3 compte 262 pages : cette valeur ne s'y applique plus.
%  E, A, NEA — non demandés à l'imprimeur : ils sont CALCULÉS avec la main
%        déduite ci-dessous. À faire confirmer avant tout tirage.
%
% La main de 1,206 n'est pas publiée par CoolLibri : elle est déduite de
% leur propre chiffre, 14 mm / (0,090 x 258/2) = 1,206. Elle reproduit donc
% exactement le dos de N, et n'est qu'une ESTIMATION pour les autres.
% Écarts que cela recouvre, pour mémoire : une main de 1,0 donnerait
% 11,61 mm pour N, une main de 1,3 en donnerait 16,77. Il ne fallait pas
% deviner.
%  E   — 17 mm, VALEUR DONNÉE PAR L'IMPRIMEUR le 17/09/2026 pour 276 pages,
%        20 x 24, couverture RIGIDE, reliure collée (Klebebindung), papier
%        115 g. En reliure cousue (Fadenheftung), il annonce 18 mm.
%
% \DosImposeReliure<classe> : la reliure pour laquelle la valeur a été
% donnée. Le dos d'une couverture rigide n'est pas celui d'un dos carré
% collé souple, même à pagination égale.
\def\DosImposeN{14}   \def\DosImposePagesN{258}  \def\DosImposeFormatN{a4}
\def\DosImposeE{17}   \def\DosImposePagesE{276}  \def\DosImposeFormatE{20x24-marge}
\def\DosImposeA{0}    \def\DosImposePagesA{0}    \def\DosImposeFormatA{}
\def\DosImposeNE{0}   \def\DosImposePagesNE{0}   \def\DosImposeFormatNE{}
\def\DosImposeEA{0}   \def\DosImposePagesEA{0}   \def\DosImposeFormatEA{}
\def\DosImposeNEA{0}  \def\DosImposePagesNEA{0}  \def\DosImposeFormatNEA{}
\def\DosImposeReliureN{souple}  \def\DosImposeReliureE{rigide}
\def\DosImposeReliureA{souple}  \def\DosImposeReliureNE{souple}
\def\DosImposeReliureEA{souple} \def\DosImposeReliureNEA{souple}

% Fond perdu, en mm. 3 mm est la valeur courante ; certains imprimeurs
% demandent 5 mm.
\def\FondPerdu{3}

% Marge technique autour du fond perdu, en mm. Sert à loger les traits de
% coupe en mode épreuve. Ignorée en mode production.
\def\MargeTech{12}

% Zone de sécurité, en mm : aucun texte ne doit s'en approcher davantage
% du bord de coupe.
\def\Securite{8}

% ---------------------------------------------------------------------
% 2. DONNÉES PAR CLASSE
%    pages : nombre de pages du PDF intérieur, COUVERTURE NON COMPRISE.
%    Un dos carré collé impose un nombre de pages PAIR ; le fichier
%    refuse de compiler si le compte est impair.
%
%    Ces tables ne servent que si \PagesForcees et \VersionForcee sont
%    vides, c'est-à-dire à la main. Valeurs A4 mesurées sur les PDF de la
%    a.3 (CLAUDE.md § 1, 05/09/2026). En 20 x 24, la pagination change
%    entièrement : NEA passe de 816 à 1036 pages.
% ---------------------------------------------------------------------

\def\DonneesN{262}
\def\DonneesE{214}
\def\DonneesA{386}
\def\DonneesNE{0}      % édition jamais produite
\def\DonneesEA{0}      % édition jamais produite
\def\DonneesNEA{816}

% Version imprimée sur le dos et la 1re de couverture, PAR CLASSE.
% « Les versions sont suivies par classe, chacune évoluant à son rythme »
% (CHANGELOG.md).
\def\VersionN{a.3}
\def\VersionE{a.3}
\def\VersionA{a.3}
\def\VersionNE{a.3}
\def\VersionEA{a.3}
\def\VersionNEA{a.3}

% ---------------------------------------------------------------------
% 3. TEXTES
% ---------------------------------------------------------------------

\def\TitreOuvrage{50\,Ohm}
\def\SousTitre{Préparation à l'examen radioamateur}

% Sous-titre propre à chaque classe (repris de FR_TITLES, build_book.py).
\def\SousTitreN{Cours complet — classe N}
\def\SousTitreE{Cours complet — classe E}
\def\SousTitreA{Cours complet — classe A}
\def\SousTitreNE{Cours complet — classes N et E}
\def\SousTitreEA{Cours complet — classes E et A}
\def\SousTitreNEA{Cours complet — classes N, E et A}

% Lettres du filigrane, empilées de haut en bas (cf. build_book.py v0.20).
\def\LettresN{N}
\def\LettresE{E}
\def\LettresA{A}
\def\LettresNE{N,E}
\def\LettresEA{E,A}
\def\LettresNEA{N,E,A}

% Texte de 4e de couverture.
% NB : il parle des classes N, E et A quelle que soit l'édition. Juste pour
% N, E, A et NEA ; approximatif pour NE et EA ; faux pour SWL, d'où son refus.
\def\QuatriemeTexte{%
Ce guide d'étude prépare aux examens radioamateur allemands des classes
N, E et A. Il traduit en français les contenus du projet 50ohm.de, publiés
par le référat AJW du DARC e.\,V., et les complète d'encarts « En France »
qui signalent, sans rien retirer au texte allemand, ce que la
réglementation française prévoit sur le même sujet.\par\smallskip
Les questions officielles du catalogue allemand accompagnent chaque
notion, avec leurs réponses.%
}

\def\Auteur{Traduit avec l'aide d'une IA par Pierre F4JWI}
\def\Depot{github.com/DrPLG/50ohm-fr}

% =====================================================================
%  À PARTIR D'ICI, PLUS RIEN À MODIFIER
% =====================================================================

\usepackage{iftex}
\RequireLuaTeX
\usepackage{xfp}
\usepackage{tikz}
\usetikzlibrary{calc}
\usepackage{fontspec}
\usepackage[french]{babel}
\usepackage{geometry}

% ---------------------------------------------------------------------
% 4. RÉSOLUTION DES DONNÉES DE LA CLASSE
% ---------------------------------------------------------------------

\makeatletter
% Tome SWL : refus explicite, plutôt qu'un « classe inconnue » qui laisserait
% croire à un oubli. Le texte de 4e parle des classes N, E et A.
\def\@swlmot{SWL}
\ifx\Classe\@swlmot
  \PackageError{couverture}{Pas de couverture pour le tome SWL}%
    {Sa 4e de couverture n'est pas redigee : le texte commun parle des
     classes N, E et A. C'est une decision editoriale a prendre d'abord.}%
\fi

\def\@resoudre#1#2{%
  \expandafter\let\expandafter#1\csname #2\Classe\endcsname
  \ifx#1\relax
    \PackageError{couverture}{Classe \Classe\space inconnue : pas de #2\Classe}%
      {Valeurs admises : N, E, A, NE, EA, NEA.}%
  \fi}
\@resoudre\NbPages{Donnees}
\@resoudre\SousTitreClasse{SousTitre}
\@resoudre\Lettres{Lettres}
\@resoudre\DosImpose{DosImpose}
\@resoudre\DosImposePages{DosImposePages}
\@resoudre\DosImposeFormat{DosImposeFormat}
\@resoudre\DosImposeReliure{DosImposeReliure}
\@resoudre\Version{Version}

% Valeurs imposées de l'extérieur. Test de vacuité par \detokenize : il ne
% dépend pas de la façon dont la macro a été définie.
\if\relax\detokenize\expandafter{\PagesForcees}\relax\else
  \let\NbPages\PagesForcees
\fi
\if\relax\detokenize\expandafter{\VersionForcee}\relax\else
  \let\Version\VersionForcee
\fi

% Format du livre fini.
\def\@formatAquatre{a4}
\def\@formatVingt{20x24}
\def\@formatVingtMarge{20x24-marge}
\ifx\Format\@formatAquatre
  \def\PageLargeur{210}\def\PageHauteur{297}
\else\ifx\Format\@formatVingt
  \def\PageLargeur{200}\def\PageHauteur{240}
\else\ifx\Format\@formatVingtMarge
  \def\PageLargeur{200}\def\PageHauteur{240}
\else
  \PackageError{couverture}{Format \Format\space inconnu}%
    {Valeurs admises : a4, 20x24, 20x24-marge.}%
  \def\PageLargeur{210}\def\PageHauteur{297}
\fi\fi\fi
% Reliure. Hors de « souple » et « rigide » : erreur, pas de supposition.
\def\@reliuresouple{souple}
\def\@reliurerigide{rigide}
\ifx\Reliure\@reliurerigide
  \edef\PageLargeur{\fpeval{\PageLargeur + 5}}
  \edef\PageHauteur{\fpeval{\PageHauteur + 5}}
  \def\FondPerdu{15}
\else\ifx\Reliure\@reliuresouple\else
  \PackageError{couverture}{Reliure \Reliure\space inconnue}%
    {Valeurs admises : souple, rigide.}%
\fi\fi
\makeatother

% Garde-fou : une édition jamais produite n'a pas de couverture.
\ifnum\NbPages=0
  \PackageError{couverture}{L'edition \Classe\space n'a pas de pagination}%
    {Cette edition n'a jamais ete compilee. Renseigner \string\Donnees\Classe
     ou passer \string\PagesForcees.}%
\fi

% Avertissement : l'imprimeur demande un nombre de pages multiple de 4
% (17/09/2026). Pas une erreur : c'est son exigence, pas une loi du papier.
\ifnum\fpeval{\NbPages - 4*floor(\NbPages/4)}=0 \else
  \typeout{COUVERTURE-AVERTISSEMENT \NbPages\space pages : pas un multiple
    de 4, que l'imprimeur exige -- ajouter des pages a l'interieur}%
\fi

% Garde-fou : un dos carré collé impose un nombre de pages pair.
\ifodd\NbPages
  \PackageError{couverture}{Nombre de pages impair (\NbPages)}%
    {Un dos carre colle exige un nombre PAIR de pages : une feuille en
     porte deux. Ajouter une page blanche a l'interieur.}%
\fi

% ---------------------------------------------------------------------
% 5. GÉOMÉTRIE — c'est ici que se joue l'exactitude du dos
% ---------------------------------------------------------------------

% Épaisseur du dos : valeur imposée si elle est donnée ET qu'elle vaut pour
% cette pagination et ce format ; calcul sinon.
\edef\DosCalcule{\fpeval{round(\NbPages / 2 * \Grammage * \Bouffant / 1000, 2)}}
\def\DosOrigine{calculee}
\ifnum\fpeval{\DosImpose > 0}=1
  \ifnum\fpeval{\NbPages = \DosImposePages}=1
    \ifx\Format\DosImposeFormat
      \ifx\Reliure\DosImposeReliure
        \def\DosOrigine{imposee}
      \fi
    \fi
  \fi
  \def\MotCalculee{calculee}
  \ifx\DosOrigine\MotCalculee
    \typeout{COUVERTURE-AVERTISSEMENT dos impose de \DosImpose\space mm
      donne pour \DosImposePages\space pages en \DosImposeFormat\space
      (reliure \DosImposeReliure) -- inapplicable a \NbPages\space pages en
      \Format\space (reliure \Reliure), dos calcule}%
  \fi
\fi
% NB : pas de \@ dans ces noms. Cette section est hors \makeatletter, et
% « \def\@imposee{...} » y redéfinit \@ (suivi du texte « imposee ») : le
% test échouait toujours, en silence. Constaté le 16/09/2026 — le dos de N
% tombait juste par coïncidence, 14,0 mm calculés pour 14 imposés.
\def\MotRigide{rigide}
\def\MotImposee{imposee}
\ifx\DosOrigine\MotImposee
  \edef\Dos{\DosImpose}
\else
  \edef\Dos{\DosCalcule}
\fi

% Marge technique : uniquement en mode épreuve.
\def\ouimot{oui}
\ifx\Reperes\ouimot \edef\Tech{\MargeTech}\else \def\Tech{0}\fi

% Dimensions du papier.
%   largeur = marge + fond perdu + 4e + dos + 1re + fond perdu + marge
%   hauteur = marge + fond perdu + hauteur + fond perdu + marge
\edef\PapierL{\fpeval{2*\PageLargeur + \Dos + 2*\FondPerdu + 2*\Tech}}
\edef\PapierH{\fpeval{\PageHauteur + 2*\FondPerdu + 2*\Tech}}

\geometry{paperwidth=\PapierL mm, paperheight=\PapierH mm,
          margin=0pt, nohead, nofoot}

% ---------------------------------------------------------------------
%  TOUTES LES COTES SONT FIGÉES ICI, ET NULLE PART AILLEURS.
%
%  Raison, payée par trois compilations ratées : un \fpeval laissé dans
%  une coordonnée ou dans le texte d'un nœud TikZ n'est pas développé au
%  moment où TikZ analyse le chemin. On obtient « Giving up on this path.
%  Did you forget a semicolon? », qui accuse une ponctuation alors que la
%  faute est ailleurs. Le corps du document ci-dessous n'emploie donc que
%  des macros déjà calculées.
% ---------------------------------------------------------------------

% Repères horizontaux, mesurés depuis le bord GAUCHE du papier.
\edef\CoupeG{\fpeval{\Tech + \FondPerdu}}          % coupe gauche (4e de couv.)
\edef\DosG{\fpeval{\CoupeG + \PageLargeur}}         % dos, début
\edef\DosD{\fpeval{\DosG + \Dos}}                   % dos, fin
\edef\CoupeD{\fpeval{\DosD + \PageLargeur}}         % coupe droite (1re de couv.)
% Repères verticaux, depuis le bord BAS.
\edef\CoupeB{\fpeval{\Tech + \FondPerdu}}
\edef\CoupeH{\fpeval{\CoupeB + \PageHauteur}}
% Axes.
\edef\DosAxe{\fpeval{(\DosG + \DosD)/2}}
\edef\MilieuH{\fpeval{(\CoupeB + \CoupeH)/2}}
\edef\MilieuL{\fpeval{\PapierL/2}}

% Format fini, pour les cotes affichées.
\edef\FiniL{\fpeval{2*\PageLargeur + \Dos}}

% Bandeau de la 1re de couverture : même proportion que la page de titre.
\edef\BandeauL{\fpeval{round(0.34*\PageLargeur, 2)}}
\edef\BandeauG{\fpeval{\CoupeD - \BandeauL}}
\edef\BandeauAccent{\fpeval{\BandeauG - 0.7}}
\edef\DosAccentG{\fpeval{\DosG - 0.7}}
\edef\DosAccentD{\fpeval{\DosD + 0.7}}
% Couverture rigide : le bleu du dos déborde de 2 mm sur chaque plat, pour
% absorber la tolérance de l'imprimeur sur l'épaisseur du dos (±10 %).
\edef\DosEtenduG{\fpeval{\DosG - 2}}
\edef\DosEtenduD{\fpeval{\DosD + 2}}

% Filigrane de classe.
\edef\FiligraneX{\fpeval{(\BandeauG + \CoupeD)/2}}
\edef\FiligraneY{\fpeval{\CoupeH - 6}}

% Bloc-titre de la 1re de couverture.
\edef\TitreX{\fpeval{\DosD + 12}}
\edef\TitreY{\fpeval{\CoupeH - 72}}
\edef\TitreW{\fpeval{\BandeauG - \DosD - 16}}
\edef\VersionY{\fpeval{\CoupeB + 14}}

% 4e de couverture.
\edef\QuatreX{\fpeval{\CoupeG + \Securite + 6}}
\edef\QuatreY{\fpeval{\CoupeH - \Securite - 14}}
\edef\QuatreBasY{\fpeval{\CoupeB + \Securite + 22}}
\edef\QuatreW{\fpeval{\PageLargeur - 2*\Securite - 12}}

% Zone de sécurité.
\edef\SecG{\fpeval{\CoupeG + \Securite}}
\edef\SecD{\fpeval{\CoupeD - \Securite}}
\edef\SecB{\fpeval{\CoupeB + \Securite}}
\edef\SecH{\fpeval{\CoupeH - \Securite}}

% Zones de contrôle : ce qui doit rester de part et d'autre du pli du dos.
\edef\PliQuatre{\fpeval{\DosG - \Securite}}         % 4e : limite côté dos
\edef\PliPremiere{\fpeval{\DosD + \Securite}}       % 1re : limite côté dos

% Traits de coupe et cotes, dans la marge technique.
\edef\TechExt{\fpeval{\Tech - 8}}
\edef\TechInt{\fpeval{\Tech - 1}}
\edef\TechCote{\fpeval{\Tech - 9}}
\edef\HautInt{\fpeval{\PapierH - \Tech + 1}}
\edef\HautExt{\fpeval{\PapierH - \Tech + 8}}
\edef\HautCote{\fpeval{\PapierH - \Tech + 9}}
\edef\DroitInt{\fpeval{\PapierL - \Tech + 1}}
\edef\DroitExt{\fpeval{\PapierL - \Tech + 8}}

% Emplacement du code-barres ISBN : en bas à DROITE de la 4e de couverture,
% selon l'usage éditorial, et non à gauche où il chevaucherait les mentions
% de licence.
\edef\IsbnD{\fpeval{\DosG - \Securite}}
\edef\IsbnG{\fpeval{\IsbnD - 40}}
\edef\IsbnH{\fpeval{\SecB + 25}}
\edef\IsbnCX{\fpeval{\IsbnG + 20}}
\edef\IsbnCY{\fpeval{\SecB + 12.5}}

% Cartouches de cotes, dans la marge technique. Ancrés vers l'INTÉRIEUR du
% papier : ancrés vers l'extérieur, ils débordaient et étaient rognés.
\edef\CoteBasY{\fpeval{1}}
\edef\CoteHautY{\fpeval{\PapierH - 1}}

% ---------------------------------------------------------------------
% 6. COULEURS ET POLICES — reprises de la page de titre (build_book.py)
% ---------------------------------------------------------------------

\definecolor{TitleBand}{cmyk}{.9,.55,.1,.35}   % bleu profond
\definecolor{TitleAccent}{cmyk}{.05,.8,1,0}    % rouge DARC (accent fin)
\definecolor{Repere}{cmyk}{0,1,1,0}            % magenta : repères de montage

\setmainfont{Libertinus Serif}[
  ItalicFont     = Libertinus Serif Italic,
  BoldFont       = Libertinus Serif Bold,
  BoldItalicFont = Libertinus Serif Bold Italic]

\pagestyle{empty}

% ---------------------------------------------------------------------
% 7. FILIGRANE DE CLASSE
%    Le texte est construit ici, et non dans le nœud : un \\ produit à
%    l'intérieur d'un \foreach referme le groupe d'alignement de TikZ.
%    Une lettre seule est grande ; plusieurs s'empilent et se réduisent
%    (build_book.py v0.20 : « NEA » à 105 pt débordait encore).
% ---------------------------------------------------------------------

\makeatletter
\def\FiligraneTexte{}
\newif\ifpremierelettre \premierelettretrue
\count255=0
\@for\@tempa:=\Lettres\do{%
  \advance\count255 by 1\relax
  \ifpremierelettre \premierelettrefalse
  \else \g@addto@macro\FiligraneTexte{\\}\fi
  \expandafter\g@addto@macro\expandafter\FiligraneTexte\expandafter{\@tempa}}%
\xdef\NbLettres{\the\count255}
\makeatother

\ifnum\NbLettres=1 \def\CorpsF{150}\else\def\CorpsF{86}\fi
\edef\CorpsFI{\fpeval{round(\CorpsF*1.02, 0)}}

% Texte du cartouche de cotes, lui aussi construit hors du nœud TikZ.
%
% NE PAS y employer \textsuperscript : dans un nœud « align=center », il
% casse le chemin TikZ (« Giving up on this path. Did you forget a
% semicolon? »), et l'erreur est signalée sur le nœud SUIVANT — elle
% accuse donc une ponctuation là où la faute est ici. Isolé par essais :
% $^2$ et le caractère ² passent, \textsuperscript non.
\def\CotesTexte{%
  \Classe\ --- \NbPages\ pages --- papier \Grammage\ g/m$^2$,
  bouffant \Bouffant\ --- fond perdu \FondPerdu\ mm\\
  papier \PapierL\ mm sur \PapierH\ mm --- fini \FiniL\ mm sur
  \PageHauteur\ mm\\
  EPREUVE, ne pas envoyer a l'impression%
}

% ---------------------------------------------------------------------
% 8. CONTRÔLE DE SUPERPOSITION
%
%  \MesureNoeud{nœud}{préfixe}{retrait} mesure la boîte d'un nœud par ses
%  quatre coins (un nœud pivoté a ses ancres pivotées : le min et le max des
%  quatre suffisent) et la ramène au repère du papier, en mm. Le point (0,0)
%  est lu par le même « let » que les ancres : les deux passent donc par la
%  même transformation, et la comparaison ne dépend pas de l'endroit où
%  TeX a posé l'image sur la page.
%
%  Le retrait sert au seul filigrane, pour mesurer les lettres et non leur
%  cadre : c'est le bord des lettres que le détourage rognerait.
%
%  PIÈGE MESURÉ le 16/09/2026 : la marge interne TikZ (.3333em) n'est PAS
%  évaluée dans la police du nœud, mais dans celle qui entoure le dessin.
%  Instrumenté : dans le nœud, em = 150 pt ; boîte brute du « N » :
%  41,7 x 37,5 mm, soit la lettre plus 2 x 1,3 mm — 1,3 mm étant .3333em
%  à 11 pt. Un retrait calculé à 150 pt (17,6 mm) ramenait la lettre à
%  6,6 mm de large, et le contrôle aurait dit OK sur une mesure absurde.
% ---------------------------------------------------------------------

\newcount\NbAlertes
\newdimen\FiligraneSep
\newsavebox\CreditsBoite

% Crédits de 4e : lignes coupées à la main. Leur largeur RÉELLE est celle
% de la plus longue ligne, mesurée dans une boîte à part (voir § 8, ISBN).
\def\CreditsTexte{%
    Réalisé à partir des contenus de 50ohm.de (en allemand)\\
    50ohm.de-Autorenteam, coordonné par le référat AJW du DARC e.\,V.\\
    \Auteur\\[3pt]
    Licence CC BY 4.0 --- \texttt{\Depot}}

\def\MesureNoeud#1#2#3{%
  \path let \p1=(#1.south west), \p2=(#1.north east),
            \p3=(#1.north west), \p4=(#1.south east), \p5=(0,0) in
    \pgfextra{%
      \expandafter\xdef\csname #2xmin\endcsname{\fpeval{round((min(\x1,\x2,\x3,\x4) - \x5 + #3)/1mm, 2)}}%
      \expandafter\xdef\csname #2xmax\endcsname{\fpeval{round((max(\x1,\x2,\x3,\x4) - \x5 - #3)/1mm, 2)}}%
      \expandafter\xdef\csname #2ymin\endcsname{\fpeval{round((min(\y1,\y2,\y3,\y4) - \y5 + #3)/1mm, 2)}}%
      \expandafter\xdef\csname #2ymax\endcsname{\fpeval{round((max(\y1,\y2,\y3,\y4) - \y5 - #3)/1mm, 2)}}%
    };}

% \ControleDans{libellé}{préfixe}{xmin}{xmax}{ymin}{ymax}
\def\ControleDans#1#2#3#4#5#6{%
  \ifnum\fpeval{\csname #2xmin\endcsname >= #3 && \csname #2xmax\endcsname <= #4
             && \csname #2ymin\endcsname >= #5 && \csname #2ymax\endcsname <= #6}=1
    \typeout{COUVERTURE-CONTROLE OK #1 : x \csname #2xmin\endcsname-\csname #2xmax\endcsname,
      y \csname #2ymin\endcsname-\csname #2ymax\endcsname\space dans x #3-#4, y #5-#6}%
  \else
    \global\advance\NbAlertes by 1
    \typeout{COUVERTURE-CONTROLE ALERTE #1 : x \csname #2xmin\endcsname-\csname #2xmax\endcsname,
      y \csname #2ymin\endcsname-\csname #2ymax\endcsname\space HORS de x #3-#4, y #5-#6}%
  \fi}

% \ControleDisjoints{libellé}{préfixe}{xmin}{xmax}{ymin}{ymax}
% Deux rectangles se chevauchent si et seulement s'ils se recouvrent sur
% les DEUX axes à la fois.
\def\ControleDisjoints#1#2#3#4#5#6{%
  \ifnum\fpeval{\csname #2xmin\endcsname < #4 && #3 < \csname #2xmax\endcsname
             && \csname #2ymin\endcsname < #6 && #5 < \csname #2ymax\endcsname}=1
    \global\advance\NbAlertes by 1
    \typeout{COUVERTURE-CONTROLE ALERTE #1 : chevauchement}%
  \else
    \typeout{COUVERTURE-CONTROLE OK #1 : disjoints}%
  \fi}

\begin{document}
% Marge interne du filigrane, évaluée dans la police qui ENTOURE le dessin :
% c'est là que TikZ évalue le .3333em de « inner sep » (cf. § 8).
\FiligraneSep=.3333em
% Largeur réelle des crédits : la plus longue ligne, dans leur police.
\sbox\CreditsBoite{\small\begin{tabular}{@{}l@{}}\CreditsTexte\end{tabular}}%
\begin{tikzpicture}[remember picture, overlay,
                    x=1mm, y=1mm, shift={(current page.south west)}]

% ---- Fonds ----------------------------------------------------------
% La 4e de couverture reste sur le blanc du papier.
% Le bandeau bleu couvre le dos, puis reprend au tiers extérieur de la
% 1re de couverture — même partition 2/3-1/3 que la page de titre.
% Couverture rigide : pas de filets au dos (ils tomberaient à côté des
% plis, vu la tolérance), et un dos bleu élargi de 2 mm de chaque côté.
\ifx\Reliure\MotRigide
  \fill[TitleBand] (\DosEtenduG, 0) rectangle (\DosEtenduD, \PapierH);
\else
  \fill[TitleBand] (\DosG, 0) rectangle (\DosD, \PapierH);
\fi
\fill[TitleBand] (\BandeauG, 0) rectangle (\PapierL, \PapierH);
% Filets d'accent : à la jonction sur la 1re, de part et d'autre du dos.
\fill[TitleAccent] (\BandeauAccent, 0) rectangle (\BandeauG, \PapierH);
\ifx\Reliure\MotRigide\else
  \fill[TitleAccent] (\DosAccentG, 0) rectangle (\DosG, \PapierH);
  \fill[TitleAccent] (\DosD, 0) rectangle (\DosAccentD, \PapierH);
\fi

% ---- 1re de couverture ----------------------------------------------
% Filigrane de classe, en réserve claire, détouré au bandeau.
\begin{scope}
\clip (\BandeauG, 0) rectangle (\PapierL, \PapierH);
\node[name=filigrane, anchor=north, align=center, text=white, opacity=0.16,
      font=\fontsize{\CorpsF}{\CorpsFI}\selectfont\bfseries]
  at (\FiligraneX, \FiligraneY) {\FiligraneTexte};
\end{scope}

\node[name=titre, anchor=west, align=left, text width=\TitreW mm]
  at (\TitreX, \TitreY) {%
    {\fontsize{54}{56}\selectfont\bfseries \TitreOuvrage}\\[10pt]
    {\Large\color{TitleBand}\bfseries \SousTitre}\\[16pt]
    {\large\color{black!75} \SousTitreClasse}};

\node[name=version, anchor=south west, align=left, font=\small, text=black!70]
  at (\TitreX, \VersionY) {Version \Version\\ Licence CC BY 4.0};

% ---- Dos -------------------------------------------------------------
% Un dos étroit ne reçoit pas de texte : en deçà de 6 mm, le débord au
% massicot ferait mordre le texte vertical sur les plats.
\ifnum\fpeval{\Dos >= 6}=1
  \node[name=dos, rotate=90, anchor=center, text=white, font=\bfseries,
        inner sep=0pt]
    at (\DosAxe, \MilieuH)
    {\TitreOuvrage\ \ --- \ \SousTitreClasse\ \ --- \ \Version};
\fi

% ---- 4e de couverture ------------------------------------------------
\node[name=quatre, anchor=north west, align=justify, text width=\QuatreW mm]
  at (\QuatreX, \QuatreY) {\QuatriemeTexte};

% Le texte vit dans \CreditsTexte (§ 8) pour être mesuré à part : la boîte
% du nœud a la largeur fixe \QuatreW et couvre l'emplacement ISBN sans
% qu'aucun caractère n'y soit. Retirer « text width » pour l'éviter change
% l'interlignage — essayé le 16/09/2026, 6 934 pixels différents.
\node[name=credits, anchor=south west, align=left, font=\small, text=black!70,
      text width=\QuatreW mm]
  at (\QuatreX, \QuatreBasY) {\CreditsTexte};

% ---- Repères de fabrication (mode épreuve uniquement) ----------------
\ifx\Reperes\ouimot
  % Traits de coupe, hors du fond perdu.
  \foreach \x in {\CoupeG, \CoupeD} {
    \draw[Repere, line width=0.25pt] (\x, \TechExt) -- (\x, \TechInt);
    \draw[Repere, line width=0.25pt] (\x, \HautInt) -- (\x, \HautExt);}
  \foreach \y in {\CoupeB, \CoupeH} {
    \draw[Repere, line width=0.25pt] (\TechExt, \y) -- (\TechInt, \y);
    \draw[Repere, line width=0.25pt] (\DroitInt, \y) -- (\DroitExt, \y);}
  % Limites du dos, en pointillés.
  \foreach \x in {\DosG, \DosD} {
    \draw[Repere, line width=0.25pt, dashed] (\x, \TechExt) -- (\x, \TechInt);
    \draw[Repere, line width=0.25pt, dashed] (\x, \HautInt) -- (\x, \HautExt);}
  % Zone de sécurité.
  \draw[Repere, line width=0.2pt, dotted]
    (\SecG, \SecB) rectangle (\SecD, \SecH);
  % Emplacement réservé au code-barres ISBN.
  \draw[Repere, line width=0.2pt] (\IsbnG, \SecB) rectangle (\IsbnD, \IsbnH);
  \node[font=\tiny\sffamily, text=Repere, align=center]
    at (\IsbnCX, \IsbnCY) {emplacement\\ISBN};
  % Cotes.
  \node[anchor=north, font=\tiny\sffamily, text=Repere]
    at (\DosAxe, \CoteHautY) {dos \Dos\ mm (\DosOrigine)};
  \node[anchor=south, font=\tiny\sffamily, text=Repere, align=center]
    at (\MilieuL, \CoteBasY) {\CotesTexte};
\fi

% ---- Contrôle de superposition (§ 8) ----------------------------------
% Rien n'est dessiné : des mesures et des lignes de journal.
\MesureNoeud{filigrane}{Fil}{\the\FiligraneSep}
\MesureNoeud{titre}{Tit}{0pt}
\MesureNoeud{version}{Ver}{0pt}
\MesureNoeud{quatre}{Qua}{0pt}
\MesureNoeud{credits}{Cre}{0pt}
% Bord droit des crédits = bord gauche + marge interne + plus longue ligne.
\xdef\Crexmax{\fpeval{round(\Crexmin + (\the\FiligraneSep + \the\wd\CreditsBoite)/1mm, 2)}}

% Le filigrane peut entrer dans la zone de sécurité (il est décoratif),
% pas franchir la coupe ni sortir du bandeau — le détourage le rognerait.
\ControleDans{filigrane dans le bandeau}{Fil}{\BandeauG}{\CoupeD}{\CoupeB}{\CoupeH}
\ControleDans{titre sur la 1re, hors bandeau}{Tit}{\PliPremiere}{\BandeauAccent}{\SecB}{\SecH}
\ControleDans{version sur la 1re, hors bandeau}{Ver}{\PliPremiere}{\BandeauAccent}{\SecB}{\SecH}
\ControleDans{texte de 4e dans la zone de securite}{Qua}{\SecG}{\PliQuatre}{\SecB}{\SecH}
\ControleDans{credits de 4e dans la zone de securite}{Cre}{\SecG}{\PliQuatre}{\SecB}{\SecH}
\ControleDisjoints{titre / version}{Tit}{\Verxmin}{\Verxmax}{\Verymin}{\Verymax}
\ControleDisjoints{texte de 4e / credits}{Qua}{\Crexmin}{\Crexmax}{\Creymin}{\Creymax}
\ControleDisjoints{texte de 4e / emplacement ISBN}{Qua}{\IsbnG}{\IsbnD}{\SecB}{\IsbnH}
\ControleDisjoints{credits / emplacement ISBN}{Cre}{\IsbnG}{\IsbnD}{\SecB}{\IsbnH}
\ifnum\fpeval{\Dos >= 6}=1
  \MesureNoeud{dos}{Dos}{0pt}
  \ControleDans{texte du dos entre les plis}{Dos}{\DosG}{\DosD}{\SecB}{\SecH}
\else
  \typeout{COUVERTURE-CONTROLE OK dos de \Dos\space mm : trop etroit, sans texte}%
\fi

\end{tikzpicture}
\typeout{COUVERTURE-BILAN alertes=\the\NbAlertes\space; classe=\Classe\space;
  format=\Format\space; reliure=\Reliure\space; pages=\NbPages\space;
  version=\Version\space;
  dos=\Dos\space mm (\DosOrigine); fini=\FiniL x \PageHauteur\space mm;
  papier=\PapierL x \PapierH\space mm}
\end{document}
"
%
%  Pour le fichier destiné à l'imprimeur (sans repères) :
%
%      lualatex "\def\Reperes{non}% =====================================================================
%  couverture.tex — couverture dos carré collé pour les livres 50ohm-fr
% =====================================================================
%
%  Produit une couverture d'un seul tenant : 4e de couverture | dos |
%  1re de couverture, aux dimensions exactes du livre relié, fond perdu
%  compris. L'épaisseur du dos est CALCULÉE à partir du nombre de pages,
%  du grammage et du bouffant du papier.
%
%  ORIGINE ET ÉTAT. Écrit le 20/08/2026 dans Fichierstravail/ (hors git),
%  versé au dépôt le 16/09/2026 pour que fenetre_compilation.py puisse
%  l'appeler. Ce qui a changé à cette occasion :
%
%    - le FORMAT est un paramètre (\Format = a4, 20x24 ou 20x24-marge,
%      mêmes noms que --format de build_book.py) ;
%    - la pagination et la version peuvent être IMPOSÉES de l'extérieur
%      (\PagesForcees, \VersionForcee) : la fenêtre lit la pagination dans
%      le journal du livre qu'elle vient de compiler, au lieu de se fier à
%      une table qui vieillit ;
%    - un dos imposé par l'imprimeur n'est appliqué QUE pour la pagination
%      et le format pour lesquels il a été donné ;
%    - un CONTRÔLE DE SUPERPOSITION est calculé par LaTeX lui-même, à partir
%      des boîtes réelles des nœuds, et écrit dans le journal ;
%    - le tome SWL est refusé explicitement : sa 4e de couverture n'est pas
%      rédigée (le texte commun parle des classes N, E et A).
%
%    Neutralité vérifiée : en A4, pour N (258 p., a.2) et NEA (816 p., a.3),
%    en épreuve comme en impression, le rendu est identique au pixel près à
%    celui du gabarit d'origine.
%
%  COMPILATION (LuaLaTeX, depuis la racine du dépôt) :
%
%      lualatex couverture.tex
%      lualatex couverture.tex        <-- DEUX FOIS, obligatoirement
%
%  DEUX PASSES SONT NÉCESSAIRES. Le dessin s'ancre sur « current page »,
%  dont la position n'est connue qu'après un premier passage : à la
%  première compilation, tout se tasse dans le coin supérieur et la
%  couverture paraît ratée alors qu'elle ne l'est pas. Aucune erreur n'est
%  émise — seul le rendu le montre. Le contrôle de superposition, lui
%  aussi, ne vaut qu'à la seconde passe.
%
%  Le plus simple est de passer par fenetre_compilation.py (case
%  « couverture pour l'imprimeur »), qui enchaîne les deux passes, produit
%  l'épreuve ET le fichier d'impression, et relit le contrôle.
%
%  À la main, pour une autre classe, un autre format, une pagination mesurée :
%
%      lualatex "\def\Classe{A}\def\Format{20x24}\def\PagesForcees{482}\input{couverture.tex}"
%
%  Pour le fichier destiné à l'imprimeur (sans repères) :
%
%      lualatex "\def\Reperes{non}\input{couverture.tex}"
%
%  ---------------------------------------------------------------------
%  DEUX MODES, un seul paramètre : \Reperes
%
%    oui  — mode ÉPREUVE. Ajoute une marge technique, les traits de coupe,
%           les limites du dos, la zone de sécurité et les cotes. C'est ce
%           qu'il faut pour vérifier à l'écran.
%    non  — mode PRODUCTION. Papier = format final + fond perdu, rien
%           d'autre. C'est ce fichier-là qu'on envoie à l'imprimeur.
%
%  ---------------------------------------------------------------------
%  CONTRÔLE DE SUPERPOSITION
%
%  À la fin du dessin, chaque bloc de texte est mesuré par ses ancres TikZ
%  (quatre coins, ce qui vaut aussi pour le texte pivoté du dos) et comparé
%  à la zone où il doit tenir, puis aux blocs voisins. Le journal porte une
%  ligne par vérification :
%
%      COUVERTURE-CONTROLE OK      <quoi> : ...
%      COUVERTURE-CONTROLE ALERTE  <quoi> : ...
%
%  et un bilan : COUVERTURE-BILAN alertes=<n> ; ...
%
%  Ce que le contrôle VOIT : un bloc qui franchit le pli du dos, la coupe,
%  la zone de sécurité ou le bord du bandeau ; deux blocs qui se chevauchent ;
%  un bloc qui empiète sur l'emplacement réservé au code-barres ISBN.
%  Ce qu'il NE VOIT PAS : une ligne trop longue à l'INTÉRIEUR d'un bloc à
%  largeur fixe (la boîte garde sa largeur, le texte déborde). Celle-là se
%  signale par un « Overfull \hbox » dans le même journal, que la fenêtre
%  compte aussi.
%
%  ---------------------------------------------------------------------
%  AVERTISSEMENT SUR L'ÉPAISSEUR DU DOS
%
%  Le dos se calcule ainsi :
%
%      épaisseur d'une feuille (mm) = grammage x bouffant / 1000
%      dos (mm)                     = (nombre de pages / 2) x épaisseur
%
%  On divise par deux parce qu'une feuille porte DEUX pages.
%
%  Le « bouffant » (main, volume, bulk) dépend du papier et n'est pas
%  devinable : un offset 80 g courant est autour de 1,30, un bouffant
%  d'édition monte à 1,75 et plus. UN ÉCART DE BOUFFANT DÉPLACE LE DOS DE
%  PLUSIEURS MILLIMÈTRES. Demander à l'imprimeur soit son bouffant, soit
%  directement l'épaisseur du dos, et reporter ci-dessous. Tant que cette
%  valeur n'est pas confirmée, la couverture n'est pas bonne à tirer.
%
%  Si l'imprimeur donne le dos en clair, forcer la valeur avec \DosImpose.
% =====================================================================

\documentclass[11pt]{article}

% ---------------------------------------------------------------------
% 1. PARAMÈTRES — c'est la seule zone à modifier
% ---------------------------------------------------------------------

% Classe de l'ouvrage : N, E, A, NE, EA ou NEA.
% (\providecommand : une valeur passée en ligne de commande l'emporte.)
\providecommand{\Classe}{NEA}

% Mode épreuve (oui) ou production (non).
\providecommand{\Reperes}{oui}

% Format du livre FINI : a4, 20x24 ou 20x24-marge — les noms de --format
% dans build_book.py. Les deux maquettes 20 x 24 ne diffèrent que par leur
% colonne de marge intérieure : le livre fini, donc la couverture, est le
% même (200 x 240 mm).
\providecommand{\Format}{a4}

% Reliure : « souple » (dos carré collé, fond perdu 3 mm) ou « rigide »
% (couverture rigide, dite hardcover). En rigide, consignes de l'imprimeur
% du 17/09/2026 : chaque plat déborde du livre de 5 mm en largeur et en
% hauteur (le carton dépasse du bloc), et le fichier porte 15 mm de rabat
% autour du carton au lieu de 3 mm de fond perdu. Les aplats vont déjà au
% bord du papier : ils remplissent le rabat d'eux-mêmes.
\providecommand{\Reliure}{souple}

% Pagination et version imposées de l'extérieur. Vides = tables du § 2.
% La fenêtre les renseigne toujours : la pagination vient du journal du
% livre réellement compilé, dans le format demandé.
\providecommand{\PagesForcees}{}
\providecommand{\VersionForcee}{}

% Papier de l'INTÉRIEUR du livre (pas de la couverture).
% CoolLibri, papier « Standard 90 g blanc ».
\def\Grammage{90}        % g/m²
\def\Bouffant{1.206}     % main DÉDUITE, voir ci-dessous

% Épaisseur du dos IMPOSÉE par l'imprimeur, en mm, par classe.
% Une valeur > 0 prime sur le calcul — MAIS SEULEMENT pour la pagination et
% le format pour lesquels l'imprimeur l'a donnée (\DosImposePages<classe>,
% \DosImposeFormat<classe>). Hors de ce cas, le dos est calculé et le
% journal le dit : un dos de 14 mm donné pour 258 pages A4 ne vaut rien pour
% 262 pages, ni pour le même livre recomposé en 20 x 24.
%
% ATTENTION : cette valeur est propre à CHAQUE ÉDITION. Un dos unique
% appliqué aux quatre livres serait faux pour trois d'entre eux.
%
%  N   — 14 mm, VALEUR DONNÉE PAR COOLIBRI pour 258 pages en 90 g, A4.
%        Fichier attendu : 440 x 303 mm, ce que ce gabarit produit en mode
%        production (2 x 210 + 14 + 2 x 3 = 440 ; 297 + 2 x 3 = 303).
%        Le N de la a.3 compte 262 pages : cette valeur ne s'y applique plus.
%  E, A, NEA — non demandés à l'imprimeur : ils sont CALCULÉS avec la main
%        déduite ci-dessous. À faire confirmer avant tout tirage.
%
% La main de 1,206 n'est pas publiée par CoolLibri : elle est déduite de
% leur propre chiffre, 14 mm / (0,090 x 258/2) = 1,206. Elle reproduit donc
% exactement le dos de N, et n'est qu'une ESTIMATION pour les autres.
% Écarts que cela recouvre, pour mémoire : une main de 1,0 donnerait
% 11,61 mm pour N, une main de 1,3 en donnerait 16,77. Il ne fallait pas
% deviner.
%  E   — 17 mm, VALEUR DONNÉE PAR L'IMPRIMEUR le 17/09/2026 pour 276 pages,
%        20 x 24, couverture RIGIDE, reliure collée (Klebebindung), papier
%        115 g. En reliure cousue (Fadenheftung), il annonce 18 mm.
%
% \DosImposeReliure<classe> : la reliure pour laquelle la valeur a été
% donnée. Le dos d'une couverture rigide n'est pas celui d'un dos carré
% collé souple, même à pagination égale.
\def\DosImposeN{14}   \def\DosImposePagesN{258}  \def\DosImposeFormatN{a4}
\def\DosImposeE{17}   \def\DosImposePagesE{276}  \def\DosImposeFormatE{20x24-marge}
\def\DosImposeA{0}    \def\DosImposePagesA{0}    \def\DosImposeFormatA{}
\def\DosImposeNE{0}   \def\DosImposePagesNE{0}   \def\DosImposeFormatNE{}
\def\DosImposeEA{0}   \def\DosImposePagesEA{0}   \def\DosImposeFormatEA{}
\def\DosImposeNEA{0}  \def\DosImposePagesNEA{0}  \def\DosImposeFormatNEA{}
\def\DosImposeReliureN{souple}  \def\DosImposeReliureE{rigide}
\def\DosImposeReliureA{souple}  \def\DosImposeReliureNE{souple}
\def\DosImposeReliureEA{souple} \def\DosImposeReliureNEA{souple}

% Fond perdu, en mm. 3 mm est la valeur courante ; certains imprimeurs
% demandent 5 mm.
\def\FondPerdu{3}

% Marge technique autour du fond perdu, en mm. Sert à loger les traits de
% coupe en mode épreuve. Ignorée en mode production.
\def\MargeTech{12}

% Zone de sécurité, en mm : aucun texte ne doit s'en approcher davantage
% du bord de coupe.
\def\Securite{8}

% ---------------------------------------------------------------------
% 2. DONNÉES PAR CLASSE
%    pages : nombre de pages du PDF intérieur, COUVERTURE NON COMPRISE.
%    Un dos carré collé impose un nombre de pages PAIR ; le fichier
%    refuse de compiler si le compte est impair.
%
%    Ces tables ne servent que si \PagesForcees et \VersionForcee sont
%    vides, c'est-à-dire à la main. Valeurs A4 mesurées sur les PDF de la
%    a.3 (CLAUDE.md § 1, 05/09/2026). En 20 x 24, la pagination change
%    entièrement : NEA passe de 816 à 1036 pages.
% ---------------------------------------------------------------------

\def\DonneesN{262}
\def\DonneesE{214}
\def\DonneesA{386}
\def\DonneesNE{0}      % édition jamais produite
\def\DonneesEA{0}      % édition jamais produite
\def\DonneesNEA{816}

% Version imprimée sur le dos et la 1re de couverture, PAR CLASSE.
% « Les versions sont suivies par classe, chacune évoluant à son rythme »
% (CHANGELOG.md).
\def\VersionN{a.3}
\def\VersionE{a.3}
\def\VersionA{a.3}
\def\VersionNE{a.3}
\def\VersionEA{a.3}
\def\VersionNEA{a.3}

% ---------------------------------------------------------------------
% 3. TEXTES
% ---------------------------------------------------------------------

\def\TitreOuvrage{50\,Ohm}
\def\SousTitre{Préparation à l'examen radioamateur}

% Sous-titre propre à chaque classe (repris de FR_TITLES, build_book.py).
\def\SousTitreN{Cours complet — classe N}
\def\SousTitreE{Cours complet — classe E}
\def\SousTitreA{Cours complet — classe A}
\def\SousTitreNE{Cours complet — classes N et E}
\def\SousTitreEA{Cours complet — classes E et A}
\def\SousTitreNEA{Cours complet — classes N, E et A}

% Lettres du filigrane, empilées de haut en bas (cf. build_book.py v0.20).
\def\LettresN{N}
\def\LettresE{E}
\def\LettresA{A}
\def\LettresNE{N,E}
\def\LettresEA{E,A}
\def\LettresNEA{N,E,A}

% Texte de 4e de couverture.
% NB : il parle des classes N, E et A quelle que soit l'édition. Juste pour
% N, E, A et NEA ; approximatif pour NE et EA ; faux pour SWL, d'où son refus.
\def\QuatriemeTexte{%
Ce guide d'étude prépare aux examens radioamateur allemands des classes
N, E et A. Il traduit en français les contenus du projet 50ohm.de, publiés
par le référat AJW du DARC e.\,V., et les complète d'encarts « En France »
qui signalent, sans rien retirer au texte allemand, ce que la
réglementation française prévoit sur le même sujet.\par\smallskip
Les questions officielles du catalogue allemand accompagnent chaque
notion, avec leurs réponses.%
}

\def\Auteur{Traduit avec l'aide d'une IA par Pierre F4JWI}
\def\Depot{github.com/DrPLG/50ohm-fr}

% =====================================================================
%  À PARTIR D'ICI, PLUS RIEN À MODIFIER
% =====================================================================

\usepackage{iftex}
\RequireLuaTeX
\usepackage{xfp}
\usepackage{tikz}
\usetikzlibrary{calc}
\usepackage{fontspec}
\usepackage[french]{babel}
\usepackage{geometry}

% ---------------------------------------------------------------------
% 4. RÉSOLUTION DES DONNÉES DE LA CLASSE
% ---------------------------------------------------------------------

\makeatletter
% Tome SWL : refus explicite, plutôt qu'un « classe inconnue » qui laisserait
% croire à un oubli. Le texte de 4e parle des classes N, E et A.
\def\@swlmot{SWL}
\ifx\Classe\@swlmot
  \PackageError{couverture}{Pas de couverture pour le tome SWL}%
    {Sa 4e de couverture n'est pas redigee : le texte commun parle des
     classes N, E et A. C'est une decision editoriale a prendre d'abord.}%
\fi

\def\@resoudre#1#2{%
  \expandafter\let\expandafter#1\csname #2\Classe\endcsname
  \ifx#1\relax
    \PackageError{couverture}{Classe \Classe\space inconnue : pas de #2\Classe}%
      {Valeurs admises : N, E, A, NE, EA, NEA.}%
  \fi}
\@resoudre\NbPages{Donnees}
\@resoudre\SousTitreClasse{SousTitre}
\@resoudre\Lettres{Lettres}
\@resoudre\DosImpose{DosImpose}
\@resoudre\DosImposePages{DosImposePages}
\@resoudre\DosImposeFormat{DosImposeFormat}
\@resoudre\DosImposeReliure{DosImposeReliure}
\@resoudre\Version{Version}

% Valeurs imposées de l'extérieur. Test de vacuité par \detokenize : il ne
% dépend pas de la façon dont la macro a été définie.
\if\relax\detokenize\expandafter{\PagesForcees}\relax\else
  \let\NbPages\PagesForcees
\fi
\if\relax\detokenize\expandafter{\VersionForcee}\relax\else
  \let\Version\VersionForcee
\fi

% Format du livre fini.
\def\@formatAquatre{a4}
\def\@formatVingt{20x24}
\def\@formatVingtMarge{20x24-marge}
\ifx\Format\@formatAquatre
  \def\PageLargeur{210}\def\PageHauteur{297}
\else\ifx\Format\@formatVingt
  \def\PageLargeur{200}\def\PageHauteur{240}
\else\ifx\Format\@formatVingtMarge
  \def\PageLargeur{200}\def\PageHauteur{240}
\else
  \PackageError{couverture}{Format \Format\space inconnu}%
    {Valeurs admises : a4, 20x24, 20x24-marge.}%
  \def\PageLargeur{210}\def\PageHauteur{297}
\fi\fi\fi
% Reliure. Hors de « souple » et « rigide » : erreur, pas de supposition.
\def\@reliuresouple{souple}
\def\@reliurerigide{rigide}
\ifx\Reliure\@reliurerigide
  \edef\PageLargeur{\fpeval{\PageLargeur + 5}}
  \edef\PageHauteur{\fpeval{\PageHauteur + 5}}
  \def\FondPerdu{15}
\else\ifx\Reliure\@reliuresouple\else
  \PackageError{couverture}{Reliure \Reliure\space inconnue}%
    {Valeurs admises : souple, rigide.}%
\fi\fi
\makeatother

% Garde-fou : une édition jamais produite n'a pas de couverture.
\ifnum\NbPages=0
  \PackageError{couverture}{L'edition \Classe\space n'a pas de pagination}%
    {Cette edition n'a jamais ete compilee. Renseigner \string\Donnees\Classe
     ou passer \string\PagesForcees.}%
\fi

% Avertissement : l'imprimeur demande un nombre de pages multiple de 4
% (17/09/2026). Pas une erreur : c'est son exigence, pas une loi du papier.
\ifnum\fpeval{\NbPages - 4*floor(\NbPages/4)}=0 \else
  \typeout{COUVERTURE-AVERTISSEMENT \NbPages\space pages : pas un multiple
    de 4, que l'imprimeur exige -- ajouter des pages a l'interieur}%
\fi

% Garde-fou : un dos carré collé impose un nombre de pages pair.
\ifodd\NbPages
  \PackageError{couverture}{Nombre de pages impair (\NbPages)}%
    {Un dos carre colle exige un nombre PAIR de pages : une feuille en
     porte deux. Ajouter une page blanche a l'interieur.}%
\fi

% ---------------------------------------------------------------------
% 5. GÉOMÉTRIE — c'est ici que se joue l'exactitude du dos
% ---------------------------------------------------------------------

% Épaisseur du dos : valeur imposée si elle est donnée ET qu'elle vaut pour
% cette pagination et ce format ; calcul sinon.
\edef\DosCalcule{\fpeval{round(\NbPages / 2 * \Grammage * \Bouffant / 1000, 2)}}
\def\DosOrigine{calculee}
\ifnum\fpeval{\DosImpose > 0}=1
  \ifnum\fpeval{\NbPages = \DosImposePages}=1
    \ifx\Format\DosImposeFormat
      \ifx\Reliure\DosImposeReliure
        \def\DosOrigine{imposee}
      \fi
    \fi
  \fi
  \def\MotCalculee{calculee}
  \ifx\DosOrigine\MotCalculee
    \typeout{COUVERTURE-AVERTISSEMENT dos impose de \DosImpose\space mm
      donne pour \DosImposePages\space pages en \DosImposeFormat\space
      (reliure \DosImposeReliure) -- inapplicable a \NbPages\space pages en
      \Format\space (reliure \Reliure), dos calcule}%
  \fi
\fi
% NB : pas de \@ dans ces noms. Cette section est hors \makeatletter, et
% « \def\@imposee{...} » y redéfinit \@ (suivi du texte « imposee ») : le
% test échouait toujours, en silence. Constaté le 16/09/2026 — le dos de N
% tombait juste par coïncidence, 14,0 mm calculés pour 14 imposés.
\def\MotRigide{rigide}
\def\MotImposee{imposee}
\ifx\DosOrigine\MotImposee
  \edef\Dos{\DosImpose}
\else
  \edef\Dos{\DosCalcule}
\fi

% Marge technique : uniquement en mode épreuve.
\def\ouimot{oui}
\ifx\Reperes\ouimot \edef\Tech{\MargeTech}\else \def\Tech{0}\fi

% Dimensions du papier.
%   largeur = marge + fond perdu + 4e + dos + 1re + fond perdu + marge
%   hauteur = marge + fond perdu + hauteur + fond perdu + marge
\edef\PapierL{\fpeval{2*\PageLargeur + \Dos + 2*\FondPerdu + 2*\Tech}}
\edef\PapierH{\fpeval{\PageHauteur + 2*\FondPerdu + 2*\Tech}}

\geometry{paperwidth=\PapierL mm, paperheight=\PapierH mm,
          margin=0pt, nohead, nofoot}

% ---------------------------------------------------------------------
%  TOUTES LES COTES SONT FIGÉES ICI, ET NULLE PART AILLEURS.
%
%  Raison, payée par trois compilations ratées : un \fpeval laissé dans
%  une coordonnée ou dans le texte d'un nœud TikZ n'est pas développé au
%  moment où TikZ analyse le chemin. On obtient « Giving up on this path.
%  Did you forget a semicolon? », qui accuse une ponctuation alors que la
%  faute est ailleurs. Le corps du document ci-dessous n'emploie donc que
%  des macros déjà calculées.
% ---------------------------------------------------------------------

% Repères horizontaux, mesurés depuis le bord GAUCHE du papier.
\edef\CoupeG{\fpeval{\Tech + \FondPerdu}}          % coupe gauche (4e de couv.)
\edef\DosG{\fpeval{\CoupeG + \PageLargeur}}         % dos, début
\edef\DosD{\fpeval{\DosG + \Dos}}                   % dos, fin
\edef\CoupeD{\fpeval{\DosD + \PageLargeur}}         % coupe droite (1re de couv.)
% Repères verticaux, depuis le bord BAS.
\edef\CoupeB{\fpeval{\Tech + \FondPerdu}}
\edef\CoupeH{\fpeval{\CoupeB + \PageHauteur}}
% Axes.
\edef\DosAxe{\fpeval{(\DosG + \DosD)/2}}
\edef\MilieuH{\fpeval{(\CoupeB + \CoupeH)/2}}
\edef\MilieuL{\fpeval{\PapierL/2}}

% Format fini, pour les cotes affichées.
\edef\FiniL{\fpeval{2*\PageLargeur + \Dos}}

% Bandeau de la 1re de couverture : même proportion que la page de titre.
\edef\BandeauL{\fpeval{round(0.34*\PageLargeur, 2)}}
\edef\BandeauG{\fpeval{\CoupeD - \BandeauL}}
\edef\BandeauAccent{\fpeval{\BandeauG - 0.7}}
\edef\DosAccentG{\fpeval{\DosG - 0.7}}
\edef\DosAccentD{\fpeval{\DosD + 0.7}}
% Couverture rigide : le bleu du dos déborde de 2 mm sur chaque plat, pour
% absorber la tolérance de l'imprimeur sur l'épaisseur du dos (±10 %).
\edef\DosEtenduG{\fpeval{\DosG - 2}}
\edef\DosEtenduD{\fpeval{\DosD + 2}}

% Filigrane de classe.
\edef\FiligraneX{\fpeval{(\BandeauG + \CoupeD)/2}}
\edef\FiligraneY{\fpeval{\CoupeH - 6}}

% Bloc-titre de la 1re de couverture.
\edef\TitreX{\fpeval{\DosD + 12}}
\edef\TitreY{\fpeval{\CoupeH - 72}}
\edef\TitreW{\fpeval{\BandeauG - \DosD - 16}}
\edef\VersionY{\fpeval{\CoupeB + 14}}

% 4e de couverture.
\edef\QuatreX{\fpeval{\CoupeG + \Securite + 6}}
\edef\QuatreY{\fpeval{\CoupeH - \Securite - 14}}
\edef\QuatreBasY{\fpeval{\CoupeB + \Securite + 22}}
\edef\QuatreW{\fpeval{\PageLargeur - 2*\Securite - 12}}

% Zone de sécurité.
\edef\SecG{\fpeval{\CoupeG + \Securite}}
\edef\SecD{\fpeval{\CoupeD - \Securite}}
\edef\SecB{\fpeval{\CoupeB + \Securite}}
\edef\SecH{\fpeval{\CoupeH - \Securite}}

% Zones de contrôle : ce qui doit rester de part et d'autre du pli du dos.
\edef\PliQuatre{\fpeval{\DosG - \Securite}}         % 4e : limite côté dos
\edef\PliPremiere{\fpeval{\DosD + \Securite}}       % 1re : limite côté dos

% Traits de coupe et cotes, dans la marge technique.
\edef\TechExt{\fpeval{\Tech - 8}}
\edef\TechInt{\fpeval{\Tech - 1}}
\edef\TechCote{\fpeval{\Tech - 9}}
\edef\HautInt{\fpeval{\PapierH - \Tech + 1}}
\edef\HautExt{\fpeval{\PapierH - \Tech + 8}}
\edef\HautCote{\fpeval{\PapierH - \Tech + 9}}
\edef\DroitInt{\fpeval{\PapierL - \Tech + 1}}
\edef\DroitExt{\fpeval{\PapierL - \Tech + 8}}

% Emplacement du code-barres ISBN : en bas à DROITE de la 4e de couverture,
% selon l'usage éditorial, et non à gauche où il chevaucherait les mentions
% de licence.
\edef\IsbnD{\fpeval{\DosG - \Securite}}
\edef\IsbnG{\fpeval{\IsbnD - 40}}
\edef\IsbnH{\fpeval{\SecB + 25}}
\edef\IsbnCX{\fpeval{\IsbnG + 20}}
\edef\IsbnCY{\fpeval{\SecB + 12.5}}

% Cartouches de cotes, dans la marge technique. Ancrés vers l'INTÉRIEUR du
% papier : ancrés vers l'extérieur, ils débordaient et étaient rognés.
\edef\CoteBasY{\fpeval{1}}
\edef\CoteHautY{\fpeval{\PapierH - 1}}

% ---------------------------------------------------------------------
% 6. COULEURS ET POLICES — reprises de la page de titre (build_book.py)
% ---------------------------------------------------------------------

\definecolor{TitleBand}{cmyk}{.9,.55,.1,.35}   % bleu profond
\definecolor{TitleAccent}{cmyk}{.05,.8,1,0}    % rouge DARC (accent fin)
\definecolor{Repere}{cmyk}{0,1,1,0}            % magenta : repères de montage

\setmainfont{Libertinus Serif}[
  ItalicFont     = Libertinus Serif Italic,
  BoldFont       = Libertinus Serif Bold,
  BoldItalicFont = Libertinus Serif Bold Italic]

\pagestyle{empty}

% ---------------------------------------------------------------------
% 7. FILIGRANE DE CLASSE
%    Le texte est construit ici, et non dans le nœud : un \\ produit à
%    l'intérieur d'un \foreach referme le groupe d'alignement de TikZ.
%    Une lettre seule est grande ; plusieurs s'empilent et se réduisent
%    (build_book.py v0.20 : « NEA » à 105 pt débordait encore).
% ---------------------------------------------------------------------

\makeatletter
\def\FiligraneTexte{}
\newif\ifpremierelettre \premierelettretrue
\count255=0
\@for\@tempa:=\Lettres\do{%
  \advance\count255 by 1\relax
  \ifpremierelettre \premierelettrefalse
  \else \g@addto@macro\FiligraneTexte{\\}\fi
  \expandafter\g@addto@macro\expandafter\FiligraneTexte\expandafter{\@tempa}}%
\xdef\NbLettres{\the\count255}
\makeatother

\ifnum\NbLettres=1 \def\CorpsF{150}\else\def\CorpsF{86}\fi
\edef\CorpsFI{\fpeval{round(\CorpsF*1.02, 0)}}

% Texte du cartouche de cotes, lui aussi construit hors du nœud TikZ.
%
% NE PAS y employer \textsuperscript : dans un nœud « align=center », il
% casse le chemin TikZ (« Giving up on this path. Did you forget a
% semicolon? »), et l'erreur est signalée sur le nœud SUIVANT — elle
% accuse donc une ponctuation là où la faute est ici. Isolé par essais :
% $^2$ et le caractère ² passent, \textsuperscript non.
\def\CotesTexte{%
  \Classe\ --- \NbPages\ pages --- papier \Grammage\ g/m$^2$,
  bouffant \Bouffant\ --- fond perdu \FondPerdu\ mm\\
  papier \PapierL\ mm sur \PapierH\ mm --- fini \FiniL\ mm sur
  \PageHauteur\ mm\\
  EPREUVE, ne pas envoyer a l'impression%
}

% ---------------------------------------------------------------------
% 8. CONTRÔLE DE SUPERPOSITION
%
%  \MesureNoeud{nœud}{préfixe}{retrait} mesure la boîte d'un nœud par ses
%  quatre coins (un nœud pivoté a ses ancres pivotées : le min et le max des
%  quatre suffisent) et la ramène au repère du papier, en mm. Le point (0,0)
%  est lu par le même « let » que les ancres : les deux passent donc par la
%  même transformation, et la comparaison ne dépend pas de l'endroit où
%  TeX a posé l'image sur la page.
%
%  Le retrait sert au seul filigrane, pour mesurer les lettres et non leur
%  cadre : c'est le bord des lettres que le détourage rognerait.
%
%  PIÈGE MESURÉ le 16/09/2026 : la marge interne TikZ (.3333em) n'est PAS
%  évaluée dans la police du nœud, mais dans celle qui entoure le dessin.
%  Instrumenté : dans le nœud, em = 150 pt ; boîte brute du « N » :
%  41,7 x 37,5 mm, soit la lettre plus 2 x 1,3 mm — 1,3 mm étant .3333em
%  à 11 pt. Un retrait calculé à 150 pt (17,6 mm) ramenait la lettre à
%  6,6 mm de large, et le contrôle aurait dit OK sur une mesure absurde.
% ---------------------------------------------------------------------

\newcount\NbAlertes
\newdimen\FiligraneSep
\newsavebox\CreditsBoite

% Crédits de 4e : lignes coupées à la main. Leur largeur RÉELLE est celle
% de la plus longue ligne, mesurée dans une boîte à part (voir § 8, ISBN).
\def\CreditsTexte{%
    Réalisé à partir des contenus de 50ohm.de (en allemand)\\
    50ohm.de-Autorenteam, coordonné par le référat AJW du DARC e.\,V.\\
    \Auteur\\[3pt]
    Licence CC BY 4.0 --- \texttt{\Depot}}

\def\MesureNoeud#1#2#3{%
  \path let \p1=(#1.south west), \p2=(#1.north east),
            \p3=(#1.north west), \p4=(#1.south east), \p5=(0,0) in
    \pgfextra{%
      \expandafter\xdef\csname #2xmin\endcsname{\fpeval{round((min(\x1,\x2,\x3,\x4) - \x5 + #3)/1mm, 2)}}%
      \expandafter\xdef\csname #2xmax\endcsname{\fpeval{round((max(\x1,\x2,\x3,\x4) - \x5 - #3)/1mm, 2)}}%
      \expandafter\xdef\csname #2ymin\endcsname{\fpeval{round((min(\y1,\y2,\y3,\y4) - \y5 + #3)/1mm, 2)}}%
      \expandafter\xdef\csname #2ymax\endcsname{\fpeval{round((max(\y1,\y2,\y3,\y4) - \y5 - #3)/1mm, 2)}}%
    };}

% \ControleDans{libellé}{préfixe}{xmin}{xmax}{ymin}{ymax}
\def\ControleDans#1#2#3#4#5#6{%
  \ifnum\fpeval{\csname #2xmin\endcsname >= #3 && \csname #2xmax\endcsname <= #4
             && \csname #2ymin\endcsname >= #5 && \csname #2ymax\endcsname <= #6}=1
    \typeout{COUVERTURE-CONTROLE OK #1 : x \csname #2xmin\endcsname-\csname #2xmax\endcsname,
      y \csname #2ymin\endcsname-\csname #2ymax\endcsname\space dans x #3-#4, y #5-#6}%
  \else
    \global\advance\NbAlertes by 1
    \typeout{COUVERTURE-CONTROLE ALERTE #1 : x \csname #2xmin\endcsname-\csname #2xmax\endcsname,
      y \csname #2ymin\endcsname-\csname #2ymax\endcsname\space HORS de x #3-#4, y #5-#6}%
  \fi}

% \ControleDisjoints{libellé}{préfixe}{xmin}{xmax}{ymin}{ymax}
% Deux rectangles se chevauchent si et seulement s'ils se recouvrent sur
% les DEUX axes à la fois.
\def\ControleDisjoints#1#2#3#4#5#6{%
  \ifnum\fpeval{\csname #2xmin\endcsname < #4 && #3 < \csname #2xmax\endcsname
             && \csname #2ymin\endcsname < #6 && #5 < \csname #2ymax\endcsname}=1
    \global\advance\NbAlertes by 1
    \typeout{COUVERTURE-CONTROLE ALERTE #1 : chevauchement}%
  \else
    \typeout{COUVERTURE-CONTROLE OK #1 : disjoints}%
  \fi}

\begin{document}
% Marge interne du filigrane, évaluée dans la police qui ENTOURE le dessin :
% c'est là que TikZ évalue le .3333em de « inner sep » (cf. § 8).
\FiligraneSep=.3333em
% Largeur réelle des crédits : la plus longue ligne, dans leur police.
\sbox\CreditsBoite{\small\begin{tabular}{@{}l@{}}\CreditsTexte\end{tabular}}%
\begin{tikzpicture}[remember picture, overlay,
                    x=1mm, y=1mm, shift={(current page.south west)}]

% ---- Fonds ----------------------------------------------------------
% La 4e de couverture reste sur le blanc du papier.
% Le bandeau bleu couvre le dos, puis reprend au tiers extérieur de la
% 1re de couverture — même partition 2/3-1/3 que la page de titre.
% Couverture rigide : pas de filets au dos (ils tomberaient à côté des
% plis, vu la tolérance), et un dos bleu élargi de 2 mm de chaque côté.
\ifx\Reliure\MotRigide
  \fill[TitleBand] (\DosEtenduG, 0) rectangle (\DosEtenduD, \PapierH);
\else
  \fill[TitleBand] (\DosG, 0) rectangle (\DosD, \PapierH);
\fi
\fill[TitleBand] (\BandeauG, 0) rectangle (\PapierL, \PapierH);
% Filets d'accent : à la jonction sur la 1re, de part et d'autre du dos.
\fill[TitleAccent] (\BandeauAccent, 0) rectangle (\BandeauG, \PapierH);
\ifx\Reliure\MotRigide\else
  \fill[TitleAccent] (\DosAccentG, 0) rectangle (\DosG, \PapierH);
  \fill[TitleAccent] (\DosD, 0) rectangle (\DosAccentD, \PapierH);
\fi

% ---- 1re de couverture ----------------------------------------------
% Filigrane de classe, en réserve claire, détouré au bandeau.
\begin{scope}
\clip (\BandeauG, 0) rectangle (\PapierL, \PapierH);
\node[name=filigrane, anchor=north, align=center, text=white, opacity=0.16,
      font=\fontsize{\CorpsF}{\CorpsFI}\selectfont\bfseries]
  at (\FiligraneX, \FiligraneY) {\FiligraneTexte};
\end{scope}

\node[name=titre, anchor=west, align=left, text width=\TitreW mm]
  at (\TitreX, \TitreY) {%
    {\fontsize{54}{56}\selectfont\bfseries \TitreOuvrage}\\[10pt]
    {\Large\color{TitleBand}\bfseries \SousTitre}\\[16pt]
    {\large\color{black!75} \SousTitreClasse}};

\node[name=version, anchor=south west, align=left, font=\small, text=black!70]
  at (\TitreX, \VersionY) {Version \Version\\ Licence CC BY 4.0};

% ---- Dos -------------------------------------------------------------
% Un dos étroit ne reçoit pas de texte : en deçà de 6 mm, le débord au
% massicot ferait mordre le texte vertical sur les plats.
\ifnum\fpeval{\Dos >= 6}=1
  \node[name=dos, rotate=90, anchor=center, text=white, font=\bfseries,
        inner sep=0pt]
    at (\DosAxe, \MilieuH)
    {\TitreOuvrage\ \ --- \ \SousTitreClasse\ \ --- \ \Version};
\fi

% ---- 4e de couverture ------------------------------------------------
\node[name=quatre, anchor=north west, align=justify, text width=\QuatreW mm]
  at (\QuatreX, \QuatreY) {\QuatriemeTexte};

% Le texte vit dans \CreditsTexte (§ 8) pour être mesuré à part : la boîte
% du nœud a la largeur fixe \QuatreW et couvre l'emplacement ISBN sans
% qu'aucun caractère n'y soit. Retirer « text width » pour l'éviter change
% l'interlignage — essayé le 16/09/2026, 6 934 pixels différents.
\node[name=credits, anchor=south west, align=left, font=\small, text=black!70,
      text width=\QuatreW mm]
  at (\QuatreX, \QuatreBasY) {\CreditsTexte};

% ---- Repères de fabrication (mode épreuve uniquement) ----------------
\ifx\Reperes\ouimot
  % Traits de coupe, hors du fond perdu.
  \foreach \x in {\CoupeG, \CoupeD} {
    \draw[Repere, line width=0.25pt] (\x, \TechExt) -- (\x, \TechInt);
    \draw[Repere, line width=0.25pt] (\x, \HautInt) -- (\x, \HautExt);}
  \foreach \y in {\CoupeB, \CoupeH} {
    \draw[Repere, line width=0.25pt] (\TechExt, \y) -- (\TechInt, \y);
    \draw[Repere, line width=0.25pt] (\DroitInt, \y) -- (\DroitExt, \y);}
  % Limites du dos, en pointillés.
  \foreach \x in {\DosG, \DosD} {
    \draw[Repere, line width=0.25pt, dashed] (\x, \TechExt) -- (\x, \TechInt);
    \draw[Repere, line width=0.25pt, dashed] (\x, \HautInt) -- (\x, \HautExt);}
  % Zone de sécurité.
  \draw[Repere, line width=0.2pt, dotted]
    (\SecG, \SecB) rectangle (\SecD, \SecH);
  % Emplacement réservé au code-barres ISBN.
  \draw[Repere, line width=0.2pt] (\IsbnG, \SecB) rectangle (\IsbnD, \IsbnH);
  \node[font=\tiny\sffamily, text=Repere, align=center]
    at (\IsbnCX, \IsbnCY) {emplacement\\ISBN};
  % Cotes.
  \node[anchor=north, font=\tiny\sffamily, text=Repere]
    at (\DosAxe, \CoteHautY) {dos \Dos\ mm (\DosOrigine)};
  \node[anchor=south, font=\tiny\sffamily, text=Repere, align=center]
    at (\MilieuL, \CoteBasY) {\CotesTexte};
\fi

% ---- Contrôle de superposition (§ 8) ----------------------------------
% Rien n'est dessiné : des mesures et des lignes de journal.
\MesureNoeud{filigrane}{Fil}{\the\FiligraneSep}
\MesureNoeud{titre}{Tit}{0pt}
\MesureNoeud{version}{Ver}{0pt}
\MesureNoeud{quatre}{Qua}{0pt}
\MesureNoeud{credits}{Cre}{0pt}
% Bord droit des crédits = bord gauche + marge interne + plus longue ligne.
\xdef\Crexmax{\fpeval{round(\Crexmin + (\the\FiligraneSep + \the\wd\CreditsBoite)/1mm, 2)}}

% Le filigrane peut entrer dans la zone de sécurité (il est décoratif),
% pas franchir la coupe ni sortir du bandeau — le détourage le rognerait.
\ControleDans{filigrane dans le bandeau}{Fil}{\BandeauG}{\CoupeD}{\CoupeB}{\CoupeH}
\ControleDans{titre sur la 1re, hors bandeau}{Tit}{\PliPremiere}{\BandeauAccent}{\SecB}{\SecH}
\ControleDans{version sur la 1re, hors bandeau}{Ver}{\PliPremiere}{\BandeauAccent}{\SecB}{\SecH}
\ControleDans{texte de 4e dans la zone de securite}{Qua}{\SecG}{\PliQuatre}{\SecB}{\SecH}
\ControleDans{credits de 4e dans la zone de securite}{Cre}{\SecG}{\PliQuatre}{\SecB}{\SecH}
\ControleDisjoints{titre / version}{Tit}{\Verxmin}{\Verxmax}{\Verymin}{\Verymax}
\ControleDisjoints{texte de 4e / credits}{Qua}{\Crexmin}{\Crexmax}{\Creymin}{\Creymax}
\ControleDisjoints{texte de 4e / emplacement ISBN}{Qua}{\IsbnG}{\IsbnD}{\SecB}{\IsbnH}
\ControleDisjoints{credits / emplacement ISBN}{Cre}{\IsbnG}{\IsbnD}{\SecB}{\IsbnH}
\ifnum\fpeval{\Dos >= 6}=1
  \MesureNoeud{dos}{Dos}{0pt}
  \ControleDans{texte du dos entre les plis}{Dos}{\DosG}{\DosD}{\SecB}{\SecH}
\else
  \typeout{COUVERTURE-CONTROLE OK dos de \Dos\space mm : trop etroit, sans texte}%
\fi

\end{tikzpicture}
\typeout{COUVERTURE-BILAN alertes=\the\NbAlertes\space; classe=\Classe\space;
  format=\Format\space; reliure=\Reliure\space; pages=\NbPages\space;
  version=\Version\space;
  dos=\Dos\space mm (\DosOrigine); fini=\FiniL x \PageHauteur\space mm;
  papier=\PapierL x \PapierH\space mm}
\end{document}
"
%
%  ---------------------------------------------------------------------
%  DEUX MODES, un seul paramètre : \Reperes
%
%    oui  — mode ÉPREUVE. Ajoute une marge technique, les traits de coupe,
%           les limites du dos, la zone de sécurité et les cotes. C'est ce
%           qu'il faut pour vérifier à l'écran.
%    non  — mode PRODUCTION. Papier = format final + fond perdu, rien
%           d'autre. C'est ce fichier-là qu'on envoie à l'imprimeur.
%
%  ---------------------------------------------------------------------
%  CONTRÔLE DE SUPERPOSITION
%
%  À la fin du dessin, chaque bloc de texte est mesuré par ses ancres TikZ
%  (quatre coins, ce qui vaut aussi pour le texte pivoté du dos) et comparé
%  à la zone où il doit tenir, puis aux blocs voisins. Le journal porte une
%  ligne par vérification :
%
%      COUVERTURE-CONTROLE OK      <quoi> : ...
%      COUVERTURE-CONTROLE ALERTE  <quoi> : ...
%
%  et un bilan : COUVERTURE-BILAN alertes=<n> ; ...
%
%  Ce que le contrôle VOIT : un bloc qui franchit le pli du dos, la coupe,
%  la zone de sécurité ou le bord du bandeau ; deux blocs qui se chevauchent ;
%  un bloc qui empiète sur l'emplacement réservé au code-barres ISBN.
%  Ce qu'il NE VOIT PAS : une ligne trop longue à l'INTÉRIEUR d'un bloc à
%  largeur fixe (la boîte garde sa largeur, le texte déborde). Celle-là se
%  signale par un « Overfull \hbox » dans le même journal, que la fenêtre
%  compte aussi.
%
%  ---------------------------------------------------------------------
%  AVERTISSEMENT SUR L'ÉPAISSEUR DU DOS
%
%  Le dos se calcule ainsi :
%
%      épaisseur d'une feuille (mm) = grammage x bouffant / 1000
%      dos (mm)                     = (nombre de pages / 2) x épaisseur
%
%  On divise par deux parce qu'une feuille porte DEUX pages.
%
%  Le « bouffant » (main, volume, bulk) dépend du papier et n'est pas
%  devinable : un offset 80 g courant est autour de 1,30, un bouffant
%  d'édition monte à 1,75 et plus. UN ÉCART DE BOUFFANT DÉPLACE LE DOS DE
%  PLUSIEURS MILLIMÈTRES. Demander à l'imprimeur soit son bouffant, soit
%  directement l'épaisseur du dos, et reporter ci-dessous. Tant que cette
%  valeur n'est pas confirmée, la couverture n'est pas bonne à tirer.
%
%  Si l'imprimeur donne le dos en clair, forcer la valeur avec \DosImpose.
% =====================================================================

\documentclass[11pt]{article}

% ---------------------------------------------------------------------
% 1. PARAMÈTRES — c'est la seule zone à modifier
% ---------------------------------------------------------------------

% Classe de l'ouvrage : N, E, A, NE, EA ou NEA.
% (\providecommand : une valeur passée en ligne de commande l'emporte.)
\providecommand{\Classe}{NEA}

% Mode épreuve (oui) ou production (non).
\providecommand{\Reperes}{oui}

% Format du livre FINI : a4, 20x24 ou 20x24-marge — les noms de --format
% dans build_book.py. Les deux maquettes 20 x 24 ne diffèrent que par leur
% colonne de marge intérieure : le livre fini, donc la couverture, est le
% même (200 x 240 mm).
\providecommand{\Format}{a4}

% Reliure : « souple » (dos carré collé, fond perdu 3 mm) ou « rigide »
% (couverture rigide, dite hardcover). En rigide, consignes de l'imprimeur
% du 17/09/2026 : chaque plat déborde du livre de 5 mm en largeur et en
% hauteur (le carton dépasse du bloc), et le fichier porte 15 mm de rabat
% autour du carton au lieu de 3 mm de fond perdu. Les aplats vont déjà au
% bord du papier : ils remplissent le rabat d'eux-mêmes.
\providecommand{\Reliure}{souple}

% Pagination et version imposées de l'extérieur. Vides = tables du § 2.
% La fenêtre les renseigne toujours : la pagination vient du journal du
% livre réellement compilé, dans le format demandé.
\providecommand{\PagesForcees}{}
\providecommand{\VersionForcee}{}

% Papier de l'INTÉRIEUR du livre (pas de la couverture).
% CoolLibri, papier « Standard 90 g blanc ».
\def\Grammage{90}        % g/m²
\def\Bouffant{1.206}     % main DÉDUITE, voir ci-dessous

% Épaisseur du dos IMPOSÉE par l'imprimeur, en mm, par classe.
% Une valeur > 0 prime sur le calcul — MAIS SEULEMENT pour la pagination et
% le format pour lesquels l'imprimeur l'a donnée (\DosImposePages<classe>,
% \DosImposeFormat<classe>). Hors de ce cas, le dos est calculé et le
% journal le dit : un dos de 14 mm donné pour 258 pages A4 ne vaut rien pour
% 262 pages, ni pour le même livre recomposé en 20 x 24.
%
% ATTENTION : cette valeur est propre à CHAQUE ÉDITION. Un dos unique
% appliqué aux quatre livres serait faux pour trois d'entre eux.
%
%  N   — 14 mm, VALEUR DONNÉE PAR COOLIBRI pour 258 pages en 90 g, A4.
%        Fichier attendu : 440 x 303 mm, ce que ce gabarit produit en mode
%        production (2 x 210 + 14 + 2 x 3 = 440 ; 297 + 2 x 3 = 303).
%        Le N de la a.3 compte 262 pages : cette valeur ne s'y applique plus.
%  E, A, NEA — non demandés à l'imprimeur : ils sont CALCULÉS avec la main
%        déduite ci-dessous. À faire confirmer avant tout tirage.
%
% La main de 1,206 n'est pas publiée par CoolLibri : elle est déduite de
% leur propre chiffre, 14 mm / (0,090 x 258/2) = 1,206. Elle reproduit donc
% exactement le dos de N, et n'est qu'une ESTIMATION pour les autres.
% Écarts que cela recouvre, pour mémoire : une main de 1,0 donnerait
% 11,61 mm pour N, une main de 1,3 en donnerait 16,77. Il ne fallait pas
% deviner.
%  E   — 17 mm, VALEUR DONNÉE PAR L'IMPRIMEUR le 17/09/2026 pour 276 pages,
%        20 x 24, couverture RIGIDE, reliure collée (Klebebindung), papier
%        115 g. En reliure cousue (Fadenheftung), il annonce 18 mm.
%
% \DosImposeReliure<classe> : la reliure pour laquelle la valeur a été
% donnée. Le dos d'une couverture rigide n'est pas celui d'un dos carré
% collé souple, même à pagination égale.
\def\DosImposeN{14}   \def\DosImposePagesN{258}  \def\DosImposeFormatN{a4}
\def\DosImposeE{17}   \def\DosImposePagesE{276}  \def\DosImposeFormatE{20x24-marge}
\def\DosImposeA{0}    \def\DosImposePagesA{0}    \def\DosImposeFormatA{}
\def\DosImposeNE{0}   \def\DosImposePagesNE{0}   \def\DosImposeFormatNE{}
\def\DosImposeEA{0}   \def\DosImposePagesEA{0}   \def\DosImposeFormatEA{}
\def\DosImposeNEA{0}  \def\DosImposePagesNEA{0}  \def\DosImposeFormatNEA{}
\def\DosImposeReliureN{souple}  \def\DosImposeReliureE{rigide}
\def\DosImposeReliureA{souple}  \def\DosImposeReliureNE{souple}
\def\DosImposeReliureEA{souple} \def\DosImposeReliureNEA{souple}

% Fond perdu, en mm. 3 mm est la valeur courante ; certains imprimeurs
% demandent 5 mm.
\def\FondPerdu{3}

% Marge technique autour du fond perdu, en mm. Sert à loger les traits de
% coupe en mode épreuve. Ignorée en mode production.
\def\MargeTech{12}

% Zone de sécurité, en mm : aucun texte ne doit s'en approcher davantage
% du bord de coupe.
\def\Securite{8}

% ---------------------------------------------------------------------
% 2. DONNÉES PAR CLASSE
%    pages : nombre de pages du PDF intérieur, COUVERTURE NON COMPRISE.
%    Un dos carré collé impose un nombre de pages PAIR ; le fichier
%    refuse de compiler si le compte est impair.
%
%    Ces tables ne servent que si \PagesForcees et \VersionForcee sont
%    vides, c'est-à-dire à la main. Valeurs A4 mesurées sur les PDF de la
%    a.3 (CLAUDE.md § 1, 05/09/2026). En 20 x 24, la pagination change
%    entièrement : NEA passe de 816 à 1036 pages.
% ---------------------------------------------------------------------

\def\DonneesN{262}
\def\DonneesE{214}
\def\DonneesA{386}
\def\DonneesNE{0}      % édition jamais produite
\def\DonneesEA{0}      % édition jamais produite
\def\DonneesNEA{816}

% Version imprimée sur le dos et la 1re de couverture, PAR CLASSE.
% « Les versions sont suivies par classe, chacune évoluant à son rythme »
% (CHANGELOG.md).
\def\VersionN{a.3}
\def\VersionE{a.3}
\def\VersionA{a.3}
\def\VersionNE{a.3}
\def\VersionEA{a.3}
\def\VersionNEA{a.3}

% ---------------------------------------------------------------------
% 3. TEXTES
% ---------------------------------------------------------------------

\def\TitreOuvrage{50\,Ohm}
\def\SousTitre{Préparation à l'examen radioamateur}

% Sous-titre propre à chaque classe (repris de FR_TITLES, build_book.py).
\def\SousTitreN{Cours complet — classe N}
\def\SousTitreE{Cours complet — classe E}
\def\SousTitreA{Cours complet — classe A}
\def\SousTitreNE{Cours complet — classes N et E}
\def\SousTitreEA{Cours complet — classes E et A}
\def\SousTitreNEA{Cours complet — classes N, E et A}

% Lettres du filigrane, empilées de haut en bas (cf. build_book.py v0.20).
\def\LettresN{N}
\def\LettresE{E}
\def\LettresA{A}
\def\LettresNE{N,E}
\def\LettresEA{E,A}
\def\LettresNEA{N,E,A}

% Texte de 4e de couverture.
% NB : il parle des classes N, E et A quelle que soit l'édition. Juste pour
% N, E, A et NEA ; approximatif pour NE et EA ; faux pour SWL, d'où son refus.
\def\QuatriemeTexte{%
Ce guide d'étude prépare aux examens radioamateur allemands des classes
N, E et A. Il traduit en français les contenus du projet 50ohm.de, publiés
par le référat AJW du DARC e.\,V., et les complète d'encarts « En France »
qui signalent, sans rien retirer au texte allemand, ce que la
réglementation française prévoit sur le même sujet.\par\smallskip
Les questions officielles du catalogue allemand accompagnent chaque
notion, avec leurs réponses.%
}

\def\Auteur{Traduit avec l'aide d'une IA par Pierre F4JWI}
\def\Depot{github.com/DrPLG/50ohm-fr}

% =====================================================================
%  À PARTIR D'ICI, PLUS RIEN À MODIFIER
% =====================================================================

\usepackage{iftex}
\RequireLuaTeX
\usepackage{xfp}
\usepackage{tikz}
\usetikzlibrary{calc}
\usepackage{fontspec}
\usepackage[french]{babel}
\usepackage{geometry}

% ---------------------------------------------------------------------
% 4. RÉSOLUTION DES DONNÉES DE LA CLASSE
% ---------------------------------------------------------------------

\makeatletter
% Tome SWL : refus explicite, plutôt qu'un « classe inconnue » qui laisserait
% croire à un oubli. Le texte de 4e parle des classes N, E et A.
\def\@swlmot{SWL}
\ifx\Classe\@swlmot
  \PackageError{couverture}{Pas de couverture pour le tome SWL}%
    {Sa 4e de couverture n'est pas redigee : le texte commun parle des
     classes N, E et A. C'est une decision editoriale a prendre d'abord.}%
\fi

\def\@resoudre#1#2{%
  \expandafter\let\expandafter#1\csname #2\Classe\endcsname
  \ifx#1\relax
    \PackageError{couverture}{Classe \Classe\space inconnue : pas de #2\Classe}%
      {Valeurs admises : N, E, A, NE, EA, NEA.}%
  \fi}
\@resoudre\NbPages{Donnees}
\@resoudre\SousTitreClasse{SousTitre}
\@resoudre\Lettres{Lettres}
\@resoudre\DosImpose{DosImpose}
\@resoudre\DosImposePages{DosImposePages}
\@resoudre\DosImposeFormat{DosImposeFormat}
\@resoudre\DosImposeReliure{DosImposeReliure}
\@resoudre\Version{Version}

% Valeurs imposées de l'extérieur. Test de vacuité par \detokenize : il ne
% dépend pas de la façon dont la macro a été définie.
\if\relax\detokenize\expandafter{\PagesForcees}\relax\else
  \let\NbPages\PagesForcees
\fi
\if\relax\detokenize\expandafter{\VersionForcee}\relax\else
  \let\Version\VersionForcee
\fi

% Format du livre fini.
\def\@formatAquatre{a4}
\def\@formatVingt{20x24}
\def\@formatVingtMarge{20x24-marge}
\ifx\Format\@formatAquatre
  \def\PageLargeur{210}\def\PageHauteur{297}
\else\ifx\Format\@formatVingt
  \def\PageLargeur{200}\def\PageHauteur{240}
\else\ifx\Format\@formatVingtMarge
  \def\PageLargeur{200}\def\PageHauteur{240}
\else
  \PackageError{couverture}{Format \Format\space inconnu}%
    {Valeurs admises : a4, 20x24, 20x24-marge.}%
  \def\PageLargeur{210}\def\PageHauteur{297}
\fi\fi\fi
% Reliure. Hors de « souple » et « rigide » : erreur, pas de supposition.
\def\@reliuresouple{souple}
\def\@reliurerigide{rigide}
\ifx\Reliure\@reliurerigide
  \edef\PageLargeur{\fpeval{\PageLargeur + 5}}
  \edef\PageHauteur{\fpeval{\PageHauteur + 5}}
  \def\FondPerdu{15}
\else\ifx\Reliure\@reliuresouple\else
  \PackageError{couverture}{Reliure \Reliure\space inconnue}%
    {Valeurs admises : souple, rigide.}%
\fi\fi
\makeatother

% Garde-fou : une édition jamais produite n'a pas de couverture.
\ifnum\NbPages=0
  \PackageError{couverture}{L'edition \Classe\space n'a pas de pagination}%
    {Cette edition n'a jamais ete compilee. Renseigner \string\Donnees\Classe
     ou passer \string\PagesForcees.}%
\fi

% Avertissement : l'imprimeur demande un nombre de pages multiple de 4
% (17/09/2026). Pas une erreur : c'est son exigence, pas une loi du papier.
\ifnum\fpeval{\NbPages - 4*floor(\NbPages/4)}=0 \else
  \typeout{COUVERTURE-AVERTISSEMENT \NbPages\space pages : pas un multiple
    de 4, que l'imprimeur exige -- ajouter des pages a l'interieur}%
\fi

% Garde-fou : un dos carré collé impose un nombre de pages pair.
\ifodd\NbPages
  \PackageError{couverture}{Nombre de pages impair (\NbPages)}%
    {Un dos carre colle exige un nombre PAIR de pages : une feuille en
     porte deux. Ajouter une page blanche a l'interieur.}%
\fi

% ---------------------------------------------------------------------
% 5. GÉOMÉTRIE — c'est ici que se joue l'exactitude du dos
% ---------------------------------------------------------------------

% Épaisseur du dos : valeur imposée si elle est donnée ET qu'elle vaut pour
% cette pagination et ce format ; calcul sinon.
\edef\DosCalcule{\fpeval{round(\NbPages / 2 * \Grammage * \Bouffant / 1000, 2)}}
\def\DosOrigine{calculee}
\ifnum\fpeval{\DosImpose > 0}=1
  \ifnum\fpeval{\NbPages = \DosImposePages}=1
    \ifx\Format\DosImposeFormat
      \ifx\Reliure\DosImposeReliure
        \def\DosOrigine{imposee}
      \fi
    \fi
  \fi
  \def\MotCalculee{calculee}
  \ifx\DosOrigine\MotCalculee
    \typeout{COUVERTURE-AVERTISSEMENT dos impose de \DosImpose\space mm
      donne pour \DosImposePages\space pages en \DosImposeFormat\space
      (reliure \DosImposeReliure) -- inapplicable a \NbPages\space pages en
      \Format\space (reliure \Reliure), dos calcule}%
  \fi
\fi
% NB : pas de \@ dans ces noms. Cette section est hors \makeatletter, et
% « \def\@imposee{...} » y redéfinit \@ (suivi du texte « imposee ») : le
% test échouait toujours, en silence. Constaté le 16/09/2026 — le dos de N
% tombait juste par coïncidence, 14,0 mm calculés pour 14 imposés.
\def\MotRigide{rigide}
\def\MotImposee{imposee}
\ifx\DosOrigine\MotImposee
  \edef\Dos{\DosImpose}
\else
  \edef\Dos{\DosCalcule}
\fi

% Marge technique : uniquement en mode épreuve.
\def\ouimot{oui}
\ifx\Reperes\ouimot \edef\Tech{\MargeTech}\else \def\Tech{0}\fi

% Dimensions du papier.
%   largeur = marge + fond perdu + 4e + dos + 1re + fond perdu + marge
%   hauteur = marge + fond perdu + hauteur + fond perdu + marge
\edef\PapierL{\fpeval{2*\PageLargeur + \Dos + 2*\FondPerdu + 2*\Tech}}
\edef\PapierH{\fpeval{\PageHauteur + 2*\FondPerdu + 2*\Tech}}

\geometry{paperwidth=\PapierL mm, paperheight=\PapierH mm,
          margin=0pt, nohead, nofoot}

% ---------------------------------------------------------------------
%  TOUTES LES COTES SONT FIGÉES ICI, ET NULLE PART AILLEURS.
%
%  Raison, payée par trois compilations ratées : un \fpeval laissé dans
%  une coordonnée ou dans le texte d'un nœud TikZ n'est pas développé au
%  moment où TikZ analyse le chemin. On obtient « Giving up on this path.
%  Did you forget a semicolon? », qui accuse une ponctuation alors que la
%  faute est ailleurs. Le corps du document ci-dessous n'emploie donc que
%  des macros déjà calculées.
% ---------------------------------------------------------------------

% Repères horizontaux, mesurés depuis le bord GAUCHE du papier.
\edef\CoupeG{\fpeval{\Tech + \FondPerdu}}          % coupe gauche (4e de couv.)
\edef\DosG{\fpeval{\CoupeG + \PageLargeur}}         % dos, début
\edef\DosD{\fpeval{\DosG + \Dos}}                   % dos, fin
\edef\CoupeD{\fpeval{\DosD + \PageLargeur}}         % coupe droite (1re de couv.)
% Repères verticaux, depuis le bord BAS.
\edef\CoupeB{\fpeval{\Tech + \FondPerdu}}
\edef\CoupeH{\fpeval{\CoupeB + \PageHauteur}}
% Axes.
\edef\DosAxe{\fpeval{(\DosG + \DosD)/2}}
\edef\MilieuH{\fpeval{(\CoupeB + \CoupeH)/2}}
\edef\MilieuL{\fpeval{\PapierL/2}}

% Format fini, pour les cotes affichées.
\edef\FiniL{\fpeval{2*\PageLargeur + \Dos}}

% Bandeau de la 1re de couverture : même proportion que la page de titre.
\edef\BandeauL{\fpeval{round(0.34*\PageLargeur, 2)}}
\edef\BandeauG{\fpeval{\CoupeD - \BandeauL}}
\edef\BandeauAccent{\fpeval{\BandeauG - 0.7}}
\edef\DosAccentG{\fpeval{\DosG - 0.7}}
\edef\DosAccentD{\fpeval{\DosD + 0.7}}
% Couverture rigide : le bleu du dos déborde de 2 mm sur chaque plat, pour
% absorber la tolérance de l'imprimeur sur l'épaisseur du dos (±10 %).
\edef\DosEtenduG{\fpeval{\DosG - 2}}
\edef\DosEtenduD{\fpeval{\DosD + 2}}

% Filigrane de classe.
\edef\FiligraneX{\fpeval{(\BandeauG + \CoupeD)/2}}
\edef\FiligraneY{\fpeval{\CoupeH - 6}}

% Bloc-titre de la 1re de couverture.
\edef\TitreX{\fpeval{\DosD + 12}}
\edef\TitreY{\fpeval{\CoupeH - 72}}
\edef\TitreW{\fpeval{\BandeauG - \DosD - 16}}
\edef\VersionY{\fpeval{\CoupeB + 14}}

% 4e de couverture.
\edef\QuatreX{\fpeval{\CoupeG + \Securite + 6}}
\edef\QuatreY{\fpeval{\CoupeH - \Securite - 14}}
\edef\QuatreBasY{\fpeval{\CoupeB + \Securite + 22}}
\edef\QuatreW{\fpeval{\PageLargeur - 2*\Securite - 12}}

% Zone de sécurité.
\edef\SecG{\fpeval{\CoupeG + \Securite}}
\edef\SecD{\fpeval{\CoupeD - \Securite}}
\edef\SecB{\fpeval{\CoupeB + \Securite}}
\edef\SecH{\fpeval{\CoupeH - \Securite}}

% Zones de contrôle : ce qui doit rester de part et d'autre du pli du dos.
\edef\PliQuatre{\fpeval{\DosG - \Securite}}         % 4e : limite côté dos
\edef\PliPremiere{\fpeval{\DosD + \Securite}}       % 1re : limite côté dos

% Traits de coupe et cotes, dans la marge technique.
\edef\TechExt{\fpeval{\Tech - 8}}
\edef\TechInt{\fpeval{\Tech - 1}}
\edef\TechCote{\fpeval{\Tech - 9}}
\edef\HautInt{\fpeval{\PapierH - \Tech + 1}}
\edef\HautExt{\fpeval{\PapierH - \Tech + 8}}
\edef\HautCote{\fpeval{\PapierH - \Tech + 9}}
\edef\DroitInt{\fpeval{\PapierL - \Tech + 1}}
\edef\DroitExt{\fpeval{\PapierL - \Tech + 8}}

% Emplacement du code-barres ISBN : en bas à DROITE de la 4e de couverture,
% selon l'usage éditorial, et non à gauche où il chevaucherait les mentions
% de licence.
\edef\IsbnD{\fpeval{\DosG - \Securite}}
\edef\IsbnG{\fpeval{\IsbnD - 40}}
\edef\IsbnH{\fpeval{\SecB + 25}}
\edef\IsbnCX{\fpeval{\IsbnG + 20}}
\edef\IsbnCY{\fpeval{\SecB + 12.5}}

% Cartouches de cotes, dans la marge technique. Ancrés vers l'INTÉRIEUR du
% papier : ancrés vers l'extérieur, ils débordaient et étaient rognés.
\edef\CoteBasY{\fpeval{1}}
\edef\CoteHautY{\fpeval{\PapierH - 1}}

% ---------------------------------------------------------------------
% 6. COULEURS ET POLICES — reprises de la page de titre (build_book.py)
% ---------------------------------------------------------------------

\definecolor{TitleBand}{cmyk}{.9,.55,.1,.35}   % bleu profond
\definecolor{TitleAccent}{cmyk}{.05,.8,1,0}    % rouge DARC (accent fin)
\definecolor{Repere}{cmyk}{0,1,1,0}            % magenta : repères de montage

\setmainfont{Libertinus Serif}[
  ItalicFont     = Libertinus Serif Italic,
  BoldFont       = Libertinus Serif Bold,
  BoldItalicFont = Libertinus Serif Bold Italic]

\pagestyle{empty}

% ---------------------------------------------------------------------
% 7. FILIGRANE DE CLASSE
%    Le texte est construit ici, et non dans le nœud : un \\ produit à
%    l'intérieur d'un \foreach referme le groupe d'alignement de TikZ.
%    Une lettre seule est grande ; plusieurs s'empilent et se réduisent
%    (build_book.py v0.20 : « NEA » à 105 pt débordait encore).
% ---------------------------------------------------------------------

\makeatletter
\def\FiligraneTexte{}
\newif\ifpremierelettre \premierelettretrue
\count255=0
\@for\@tempa:=\Lettres\do{%
  \advance\count255 by 1\relax
  \ifpremierelettre \premierelettrefalse
  \else \g@addto@macro\FiligraneTexte{\\}\fi
  \expandafter\g@addto@macro\expandafter\FiligraneTexte\expandafter{\@tempa}}%
\xdef\NbLettres{\the\count255}
\makeatother

\ifnum\NbLettres=1 \def\CorpsF{150}\else\def\CorpsF{86}\fi
\edef\CorpsFI{\fpeval{round(\CorpsF*1.02, 0)}}

% Texte du cartouche de cotes, lui aussi construit hors du nœud TikZ.
%
% NE PAS y employer \textsuperscript : dans un nœud « align=center », il
% casse le chemin TikZ (« Giving up on this path. Did you forget a
% semicolon? »), et l'erreur est signalée sur le nœud SUIVANT — elle
% accuse donc une ponctuation là où la faute est ici. Isolé par essais :
% $^2$ et le caractère ² passent, \textsuperscript non.
\def\CotesTexte{%
  \Classe\ --- \NbPages\ pages --- papier \Grammage\ g/m$^2$,
  bouffant \Bouffant\ --- fond perdu \FondPerdu\ mm\\
  papier \PapierL\ mm sur \PapierH\ mm --- fini \FiniL\ mm sur
  \PageHauteur\ mm\\
  EPREUVE, ne pas envoyer a l'impression%
}

% ---------------------------------------------------------------------
% 8. CONTRÔLE DE SUPERPOSITION
%
%  \MesureNoeud{nœud}{préfixe}{retrait} mesure la boîte d'un nœud par ses
%  quatre coins (un nœud pivoté a ses ancres pivotées : le min et le max des
%  quatre suffisent) et la ramène au repère du papier, en mm. Le point (0,0)
%  est lu par le même « let » que les ancres : les deux passent donc par la
%  même transformation, et la comparaison ne dépend pas de l'endroit où
%  TeX a posé l'image sur la page.
%
%  Le retrait sert au seul filigrane, pour mesurer les lettres et non leur
%  cadre : c'est le bord des lettres que le détourage rognerait.
%
%  PIÈGE MESURÉ le 16/09/2026 : la marge interne TikZ (.3333em) n'est PAS
%  évaluée dans la police du nœud, mais dans celle qui entoure le dessin.
%  Instrumenté : dans le nœud, em = 150 pt ; boîte brute du « N » :
%  41,7 x 37,5 mm, soit la lettre plus 2 x 1,3 mm — 1,3 mm étant .3333em
%  à 11 pt. Un retrait calculé à 150 pt (17,6 mm) ramenait la lettre à
%  6,6 mm de large, et le contrôle aurait dit OK sur une mesure absurde.
% ---------------------------------------------------------------------

\newcount\NbAlertes
\newdimen\FiligraneSep
\newsavebox\CreditsBoite

% Crédits de 4e : lignes coupées à la main. Leur largeur RÉELLE est celle
% de la plus longue ligne, mesurée dans une boîte à part (voir § 8, ISBN).
\def\CreditsTexte{%
    Réalisé à partir des contenus de 50ohm.de (en allemand)\\
    50ohm.de-Autorenteam, coordonné par le référat AJW du DARC e.\,V.\\
    \Auteur\\[3pt]
    Licence CC BY 4.0 --- \texttt{\Depot}}

\def\MesureNoeud#1#2#3{%
  \path let \p1=(#1.south west), \p2=(#1.north east),
            \p3=(#1.north west), \p4=(#1.south east), \p5=(0,0) in
    \pgfextra{%
      \expandafter\xdef\csname #2xmin\endcsname{\fpeval{round((min(\x1,\x2,\x3,\x4) - \x5 + #3)/1mm, 2)}}%
      \expandafter\xdef\csname #2xmax\endcsname{\fpeval{round((max(\x1,\x2,\x3,\x4) - \x5 - #3)/1mm, 2)}}%
      \expandafter\xdef\csname #2ymin\endcsname{\fpeval{round((min(\y1,\y2,\y3,\y4) - \y5 + #3)/1mm, 2)}}%
      \expandafter\xdef\csname #2ymax\endcsname{\fpeval{round((max(\y1,\y2,\y3,\y4) - \y5 - #3)/1mm, 2)}}%
    };}

% \ControleDans{libellé}{préfixe}{xmin}{xmax}{ymin}{ymax}
\def\ControleDans#1#2#3#4#5#6{%
  \ifnum\fpeval{\csname #2xmin\endcsname >= #3 && \csname #2xmax\endcsname <= #4
             && \csname #2ymin\endcsname >= #5 && \csname #2ymax\endcsname <= #6}=1
    \typeout{COUVERTURE-CONTROLE OK #1 : x \csname #2xmin\endcsname-\csname #2xmax\endcsname,
      y \csname #2ymin\endcsname-\csname #2ymax\endcsname\space dans x #3-#4, y #5-#6}%
  \else
    \global\advance\NbAlertes by 1
    \typeout{COUVERTURE-CONTROLE ALERTE #1 : x \csname #2xmin\endcsname-\csname #2xmax\endcsname,
      y \csname #2ymin\endcsname-\csname #2ymax\endcsname\space HORS de x #3-#4, y #5-#6}%
  \fi}

% \ControleDisjoints{libellé}{préfixe}{xmin}{xmax}{ymin}{ymax}
% Deux rectangles se chevauchent si et seulement s'ils se recouvrent sur
% les DEUX axes à la fois.
\def\ControleDisjoints#1#2#3#4#5#6{%
  \ifnum\fpeval{\csname #2xmin\endcsname < #4 && #3 < \csname #2xmax\endcsname
             && \csname #2ymin\endcsname < #6 && #5 < \csname #2ymax\endcsname}=1
    \global\advance\NbAlertes by 1
    \typeout{COUVERTURE-CONTROLE ALERTE #1 : chevauchement}%
  \else
    \typeout{COUVERTURE-CONTROLE OK #1 : disjoints}%
  \fi}

\begin{document}
% Marge interne du filigrane, évaluée dans la police qui ENTOURE le dessin :
% c'est là que TikZ évalue le .3333em de « inner sep » (cf. § 8).
\FiligraneSep=.3333em
% Largeur réelle des crédits : la plus longue ligne, dans leur police.
\sbox\CreditsBoite{\small\begin{tabular}{@{}l@{}}\CreditsTexte\end{tabular}}%
\begin{tikzpicture}[remember picture, overlay,
                    x=1mm, y=1mm, shift={(current page.south west)}]

% ---- Fonds ----------------------------------------------------------
% La 4e de couverture reste sur le blanc du papier.
% Le bandeau bleu couvre le dos, puis reprend au tiers extérieur de la
% 1re de couverture — même partition 2/3-1/3 que la page de titre.
% Couverture rigide : pas de filets au dos (ils tomberaient à côté des
% plis, vu la tolérance), et un dos bleu élargi de 2 mm de chaque côté.
\ifx\Reliure\MotRigide
  \fill[TitleBand] (\DosEtenduG, 0) rectangle (\DosEtenduD, \PapierH);
\else
  \fill[TitleBand] (\DosG, 0) rectangle (\DosD, \PapierH);
\fi
\fill[TitleBand] (\BandeauG, 0) rectangle (\PapierL, \PapierH);
% Filets d'accent : à la jonction sur la 1re, de part et d'autre du dos.
\fill[TitleAccent] (\BandeauAccent, 0) rectangle (\BandeauG, \PapierH);
\ifx\Reliure\MotRigide\else
  \fill[TitleAccent] (\DosAccentG, 0) rectangle (\DosG, \PapierH);
  \fill[TitleAccent] (\DosD, 0) rectangle (\DosAccentD, \PapierH);
\fi

% ---- 1re de couverture ----------------------------------------------
% Filigrane de classe, en réserve claire, détouré au bandeau.
\begin{scope}
\clip (\BandeauG, 0) rectangle (\PapierL, \PapierH);
\node[name=filigrane, anchor=north, align=center, text=white, opacity=0.16,
      font=\fontsize{\CorpsF}{\CorpsFI}\selectfont\bfseries]
  at (\FiligraneX, \FiligraneY) {\FiligraneTexte};
\end{scope}

\node[name=titre, anchor=west, align=left, text width=\TitreW mm]
  at (\TitreX, \TitreY) {%
    {\fontsize{54}{56}\selectfont\bfseries \TitreOuvrage}\\[10pt]
    {\Large\color{TitleBand}\bfseries \SousTitre}\\[16pt]
    {\large\color{black!75} \SousTitreClasse}};

\node[name=version, anchor=south west, align=left, font=\small, text=black!70]
  at (\TitreX, \VersionY) {Version \Version\\ Licence CC BY 4.0};

% ---- Dos -------------------------------------------------------------
% Un dos étroit ne reçoit pas de texte : en deçà de 6 mm, le débord au
% massicot ferait mordre le texte vertical sur les plats.
\ifnum\fpeval{\Dos >= 6}=1
  \node[name=dos, rotate=90, anchor=center, text=white, font=\bfseries,
        inner sep=0pt]
    at (\DosAxe, \MilieuH)
    {\TitreOuvrage\ \ --- \ \SousTitreClasse\ \ --- \ \Version};
\fi

% ---- 4e de couverture ------------------------------------------------
\node[name=quatre, anchor=north west, align=justify, text width=\QuatreW mm]
  at (\QuatreX, \QuatreY) {\QuatriemeTexte};

% Le texte vit dans \CreditsTexte (§ 8) pour être mesuré à part : la boîte
% du nœud a la largeur fixe \QuatreW et couvre l'emplacement ISBN sans
% qu'aucun caractère n'y soit. Retirer « text width » pour l'éviter change
% l'interlignage — essayé le 16/09/2026, 6 934 pixels différents.
\node[name=credits, anchor=south west, align=left, font=\small, text=black!70,
      text width=\QuatreW mm]
  at (\QuatreX, \QuatreBasY) {\CreditsTexte};

% ---- Repères de fabrication (mode épreuve uniquement) ----------------
\ifx\Reperes\ouimot
  % Traits de coupe, hors du fond perdu.
  \foreach \x in {\CoupeG, \CoupeD} {
    \draw[Repere, line width=0.25pt] (\x, \TechExt) -- (\x, \TechInt);
    \draw[Repere, line width=0.25pt] (\x, \HautInt) -- (\x, \HautExt);}
  \foreach \y in {\CoupeB, \CoupeH} {
    \draw[Repere, line width=0.25pt] (\TechExt, \y) -- (\TechInt, \y);
    \draw[Repere, line width=0.25pt] (\DroitInt, \y) -- (\DroitExt, \y);}
  % Limites du dos, en pointillés.
  \foreach \x in {\DosG, \DosD} {
    \draw[Repere, line width=0.25pt, dashed] (\x, \TechExt) -- (\x, \TechInt);
    \draw[Repere, line width=0.25pt, dashed] (\x, \HautInt) -- (\x, \HautExt);}
  % Zone de sécurité.
  \draw[Repere, line width=0.2pt, dotted]
    (\SecG, \SecB) rectangle (\SecD, \SecH);
  % Emplacement réservé au code-barres ISBN.
  \draw[Repere, line width=0.2pt] (\IsbnG, \SecB) rectangle (\IsbnD, \IsbnH);
  \node[font=\tiny\sffamily, text=Repere, align=center]
    at (\IsbnCX, \IsbnCY) {emplacement\\ISBN};
  % Cotes.
  \node[anchor=north, font=\tiny\sffamily, text=Repere]
    at (\DosAxe, \CoteHautY) {dos \Dos\ mm (\DosOrigine)};
  \node[anchor=south, font=\tiny\sffamily, text=Repere, align=center]
    at (\MilieuL, \CoteBasY) {\CotesTexte};
\fi

% ---- Contrôle de superposition (§ 8) ----------------------------------
% Rien n'est dessiné : des mesures et des lignes de journal.
\MesureNoeud{filigrane}{Fil}{\the\FiligraneSep}
\MesureNoeud{titre}{Tit}{0pt}
\MesureNoeud{version}{Ver}{0pt}
\MesureNoeud{quatre}{Qua}{0pt}
\MesureNoeud{credits}{Cre}{0pt}
% Bord droit des crédits = bord gauche + marge interne + plus longue ligne.
\xdef\Crexmax{\fpeval{round(\Crexmin + (\the\FiligraneSep + \the\wd\CreditsBoite)/1mm, 2)}}

% Le filigrane peut entrer dans la zone de sécurité (il est décoratif),
% pas franchir la coupe ni sortir du bandeau — le détourage le rognerait.
\ControleDans{filigrane dans le bandeau}{Fil}{\BandeauG}{\CoupeD}{\CoupeB}{\CoupeH}
\ControleDans{titre sur la 1re, hors bandeau}{Tit}{\PliPremiere}{\BandeauAccent}{\SecB}{\SecH}
\ControleDans{version sur la 1re, hors bandeau}{Ver}{\PliPremiere}{\BandeauAccent}{\SecB}{\SecH}
\ControleDans{texte de 4e dans la zone de securite}{Qua}{\SecG}{\PliQuatre}{\SecB}{\SecH}
\ControleDans{credits de 4e dans la zone de securite}{Cre}{\SecG}{\PliQuatre}{\SecB}{\SecH}
\ControleDisjoints{titre / version}{Tit}{\Verxmin}{\Verxmax}{\Verymin}{\Verymax}
\ControleDisjoints{texte de 4e / credits}{Qua}{\Crexmin}{\Crexmax}{\Creymin}{\Creymax}
\ControleDisjoints{texte de 4e / emplacement ISBN}{Qua}{\IsbnG}{\IsbnD}{\SecB}{\IsbnH}
\ControleDisjoints{credits / emplacement ISBN}{Cre}{\IsbnG}{\IsbnD}{\SecB}{\IsbnH}
\ifnum\fpeval{\Dos >= 6}=1
  \MesureNoeud{dos}{Dos}{0pt}
  \ControleDans{texte du dos entre les plis}{Dos}{\DosG}{\DosD}{\SecB}{\SecH}
\else
  \typeout{COUVERTURE-CONTROLE OK dos de \Dos\space mm : trop etroit, sans texte}%
\fi

\end{tikzpicture}
\typeout{COUVERTURE-BILAN alertes=\the\NbAlertes\space; classe=\Classe\space;
  format=\Format\space; reliure=\Reliure\space; pages=\NbPages\space;
  version=\Version\space;
  dos=\Dos\space mm (\DosOrigine); fini=\FiniL x \PageHauteur\space mm;
  papier=\PapierL x \PapierH\space mm}
\end{document}
"
%
%  ---------------------------------------------------------------------
%  DEUX MODES, un seul paramètre : \Reperes
%
%    oui  — mode ÉPREUVE. Ajoute une marge technique, les traits de coupe,
%           les limites du dos, la zone de sécurité et les cotes. C'est ce
%           qu'il faut pour vérifier à l'écran.
%    non  — mode PRODUCTION. Papier = format final + fond perdu, rien
%           d'autre. C'est ce fichier-là qu'on envoie à l'imprimeur.
%
%  ---------------------------------------------------------------------
%  CONTRÔLE DE SUPERPOSITION
%
%  À la fin du dessin, chaque bloc de texte est mesuré par ses ancres TikZ
%  (quatre coins, ce qui vaut aussi pour le texte pivoté du dos) et comparé
%  à la zone où il doit tenir, puis aux blocs voisins. Le journal porte une
%  ligne par vérification :
%
%      COUVERTURE-CONTROLE OK      <quoi> : ...
%      COUVERTURE-CONTROLE ALERTE  <quoi> : ...
%
%  et un bilan : COUVERTURE-BILAN alertes=<n> ; ...
%
%  Ce que le contrôle VOIT : un bloc qui franchit le pli du dos, la coupe,
%  la zone de sécurité ou le bord du bandeau ; deux blocs qui se chevauchent ;
%  un bloc qui empiète sur l'emplacement réservé au code-barres ISBN.
%  Ce qu'il NE VOIT PAS : une ligne trop longue à l'INTÉRIEUR d'un bloc à
%  largeur fixe (la boîte garde sa largeur, le texte déborde). Celle-là se
%  signale par un « Overfull \hbox » dans le même journal, que la fenêtre
%  compte aussi.
%
%  ---------------------------------------------------------------------
%  AVERTISSEMENT SUR L'ÉPAISSEUR DU DOS
%
%  Le dos se calcule ainsi :
%
%      épaisseur d'une feuille (mm) = grammage x bouffant / 1000
%      dos (mm)                     = (nombre de pages / 2) x épaisseur
%
%  On divise par deux parce qu'une feuille porte DEUX pages.
%
%  Le « bouffant » (main, volume, bulk) dépend du papier et n'est pas
%  devinable : un offset 80 g courant est autour de 1,30, un bouffant
%  d'édition monte à 1,75 et plus. UN ÉCART DE BOUFFANT DÉPLACE LE DOS DE
%  PLUSIEURS MILLIMÈTRES. Demander à l'imprimeur soit son bouffant, soit
%  directement l'épaisseur du dos, et reporter ci-dessous. Tant que cette
%  valeur n'est pas confirmée, la couverture n'est pas bonne à tirer.
%
%  Si l'imprimeur donne le dos en clair, forcer la valeur avec \DosImpose.
% =====================================================================

\documentclass[11pt]{article}

% ---------------------------------------------------------------------
% 1. PARAMÈTRES — c'est la seule zone à modifier
% ---------------------------------------------------------------------

% Classe de l'ouvrage : N, E, A, NE, EA ou NEA.
% (\providecommand : une valeur passée en ligne de commande l'emporte.)
\providecommand{\Classe}{NEA}

% Mode épreuve (oui) ou production (non).
\providecommand{\Reperes}{oui}

% Format du livre FINI : a4, 20x24 ou 20x24-marge — les noms de --format
% dans build_book.py. Les deux maquettes 20 x 24 ne diffèrent que par leur
% colonne de marge intérieure : le livre fini, donc la couverture, est le
% même (200 x 240 mm).
\providecommand{\Format}{a4}

% Reliure : « souple » (dos carré collé, fond perdu 3 mm) ou « rigide »
% (couverture rigide, dite hardcover). En rigide, consignes de l'imprimeur
% du 17/09/2026 : chaque plat déborde du livre de 5 mm en largeur et en
% hauteur (le carton dépasse du bloc), et le fichier porte 15 mm de rabat
% autour du carton au lieu de 3 mm de fond perdu. Les aplats vont déjà au
% bord du papier : ils remplissent le rabat d'eux-mêmes.
\providecommand{\Reliure}{souple}

% Pagination et version imposées de l'extérieur. Vides = tables du § 2.
% La fenêtre les renseigne toujours : la pagination vient du journal du
% livre réellement compilé, dans le format demandé.
\providecommand{\PagesForcees}{}
\providecommand{\VersionForcee}{}

% Papier de l'INTÉRIEUR du livre (pas de la couverture).
% CoolLibri, papier « Standard 90 g blanc ».
\def\Grammage{90}        % g/m²
\def\Bouffant{1.206}     % main DÉDUITE, voir ci-dessous

% Épaisseur du dos IMPOSÉE par l'imprimeur, en mm, par classe.
% Une valeur > 0 prime sur le calcul — MAIS SEULEMENT pour la pagination et
% le format pour lesquels l'imprimeur l'a donnée (\DosImposePages<classe>,
% \DosImposeFormat<classe>). Hors de ce cas, le dos est calculé et le
% journal le dit : un dos de 14 mm donné pour 258 pages A4 ne vaut rien pour
% 262 pages, ni pour le même livre recomposé en 20 x 24.
%
% ATTENTION : cette valeur est propre à CHAQUE ÉDITION. Un dos unique
% appliqué aux quatre livres serait faux pour trois d'entre eux.
%
%  N   — 14 mm, VALEUR DONNÉE PAR COOLIBRI pour 258 pages en 90 g, A4.
%        Fichier attendu : 440 x 303 mm, ce que ce gabarit produit en mode
%        production (2 x 210 + 14 + 2 x 3 = 440 ; 297 + 2 x 3 = 303).
%        Le N de la a.3 compte 262 pages : cette valeur ne s'y applique plus.
%  E, A, NEA — non demandés à l'imprimeur : ils sont CALCULÉS avec la main
%        déduite ci-dessous. À faire confirmer avant tout tirage.
%
% La main de 1,206 n'est pas publiée par CoolLibri : elle est déduite de
% leur propre chiffre, 14 mm / (0,090 x 258/2) = 1,206. Elle reproduit donc
% exactement le dos de N, et n'est qu'une ESTIMATION pour les autres.
% Écarts que cela recouvre, pour mémoire : une main de 1,0 donnerait
% 11,61 mm pour N, une main de 1,3 en donnerait 16,77. Il ne fallait pas
% deviner.
%  E   — 17 mm, VALEUR DONNÉE PAR L'IMPRIMEUR le 17/09/2026 pour 276 pages,
%        20 x 24, couverture RIGIDE, reliure collée (Klebebindung), papier
%        115 g. En reliure cousue (Fadenheftung), il annonce 18 mm.
%
% \DosImposeReliure<classe> : la reliure pour laquelle la valeur a été
% donnée. Le dos d'une couverture rigide n'est pas celui d'un dos carré
% collé souple, même à pagination égale.
\def\DosImposeN{14}   \def\DosImposePagesN{258}  \def\DosImposeFormatN{a4}
\def\DosImposeE{17}   \def\DosImposePagesE{276}  \def\DosImposeFormatE{20x24-marge}
\def\DosImposeA{0}    \def\DosImposePagesA{0}    \def\DosImposeFormatA{}
\def\DosImposeNE{0}   \def\DosImposePagesNE{0}   \def\DosImposeFormatNE{}
\def\DosImposeEA{0}   \def\DosImposePagesEA{0}   \def\DosImposeFormatEA{}
\def\DosImposeNEA{0}  \def\DosImposePagesNEA{0}  \def\DosImposeFormatNEA{}
\def\DosImposeReliureN{souple}  \def\DosImposeReliureE{rigide}
\def\DosImposeReliureA{souple}  \def\DosImposeReliureNE{souple}
\def\DosImposeReliureEA{souple} \def\DosImposeReliureNEA{souple}

% Fond perdu, en mm. 3 mm est la valeur courante ; certains imprimeurs
% demandent 5 mm.
\def\FondPerdu{3}

% Marge technique autour du fond perdu, en mm. Sert à loger les traits de
% coupe en mode épreuve. Ignorée en mode production.
\def\MargeTech{12}

% Zone de sécurité, en mm : aucun texte ne doit s'en approcher davantage
% du bord de coupe.
\def\Securite{8}

% ---------------------------------------------------------------------
% 2. DONNÉES PAR CLASSE
%    pages : nombre de pages du PDF intérieur, COUVERTURE NON COMPRISE.
%    Un dos carré collé impose un nombre de pages PAIR ; le fichier
%    refuse de compiler si le compte est impair.
%
%    Ces tables ne servent que si \PagesForcees et \VersionForcee sont
%    vides, c'est-à-dire à la main. Valeurs A4 mesurées sur les PDF de la
%    a.3 (CLAUDE.md § 1, 05/09/2026). En 20 x 24, la pagination change
%    entièrement : NEA passe de 816 à 1036 pages.
% ---------------------------------------------------------------------

\def\DonneesN{262}
\def\DonneesE{214}
\def\DonneesA{386}
\def\DonneesNE{0}      % édition jamais produite
\def\DonneesEA{0}      % édition jamais produite
\def\DonneesNEA{816}

% Version imprimée sur le dos et la 1re de couverture, PAR CLASSE.
% « Les versions sont suivies par classe, chacune évoluant à son rythme »
% (CHANGELOG.md).
\def\VersionN{a.3}
\def\VersionE{a.3}
\def\VersionA{a.3}
\def\VersionNE{a.3}
\def\VersionEA{a.3}
\def\VersionNEA{a.3}

% ---------------------------------------------------------------------
% 3. TEXTES
% ---------------------------------------------------------------------

\def\TitreOuvrage{50\,Ohm}
\def\SousTitre{Préparation à l'examen radioamateur}

% Sous-titre propre à chaque classe (repris de FR_TITLES, build_book.py).
\def\SousTitreN{Cours complet — classe N}
\def\SousTitreE{Cours complet — classe E}
\def\SousTitreA{Cours complet — classe A}
\def\SousTitreNE{Cours complet — classes N et E}
\def\SousTitreEA{Cours complet — classes E et A}
\def\SousTitreNEA{Cours complet — classes N, E et A}

% Lettres du filigrane, empilées de haut en bas (cf. build_book.py v0.20).
\def\LettresN{N}
\def\LettresE{E}
\def\LettresA{A}
\def\LettresNE{N,E}
\def\LettresEA{E,A}
\def\LettresNEA{N,E,A}

% Texte de 4e de couverture.
% NB : il parle des classes N, E et A quelle que soit l'édition. Juste pour
% N, E, A et NEA ; approximatif pour NE et EA ; faux pour SWL, d'où son refus.
\def\QuatriemeTexte{%
Ce guide d'étude prépare aux examens radioamateur allemands des classes
N, E et A. Il traduit en français les contenus du projet 50ohm.de, publiés
par le référat AJW du DARC e.\,V., et les complète d'encarts « En France »
qui signalent, sans rien retirer au texte allemand, ce que la
réglementation française prévoit sur le même sujet.\par\smallskip
Les questions officielles du catalogue allemand accompagnent chaque
notion, avec leurs réponses.%
}

\def\Auteur{Traduit avec l'aide d'une IA par Pierre F4JWI}
\def\Depot{github.com/DrPLG/50ohm-fr}

% =====================================================================
%  À PARTIR D'ICI, PLUS RIEN À MODIFIER
% =====================================================================

\usepackage{iftex}
\RequireLuaTeX
\usepackage{xfp}
\usepackage{tikz}
\usetikzlibrary{calc}
\usepackage{fontspec}
\usepackage[french]{babel}
\usepackage{geometry}

% ---------------------------------------------------------------------
% 4. RÉSOLUTION DES DONNÉES DE LA CLASSE
% ---------------------------------------------------------------------

\makeatletter
% Tome SWL : refus explicite, plutôt qu'un « classe inconnue » qui laisserait
% croire à un oubli. Le texte de 4e parle des classes N, E et A.
\def\@swlmot{SWL}
\ifx\Classe\@swlmot
  \PackageError{couverture}{Pas de couverture pour le tome SWL}%
    {Sa 4e de couverture n'est pas redigee : le texte commun parle des
     classes N, E et A. C'est une decision editoriale a prendre d'abord.}%
\fi

\def\@resoudre#1#2{%
  \expandafter\let\expandafter#1\csname #2\Classe\endcsname
  \ifx#1\relax
    \PackageError{couverture}{Classe \Classe\space inconnue : pas de #2\Classe}%
      {Valeurs admises : N, E, A, NE, EA, NEA.}%
  \fi}
\@resoudre\NbPages{Donnees}
\@resoudre\SousTitreClasse{SousTitre}
\@resoudre\Lettres{Lettres}
\@resoudre\DosImpose{DosImpose}
\@resoudre\DosImposePages{DosImposePages}
\@resoudre\DosImposeFormat{DosImposeFormat}
\@resoudre\DosImposeReliure{DosImposeReliure}
\@resoudre\Version{Version}

% Valeurs imposées de l'extérieur. Test de vacuité par \detokenize : il ne
% dépend pas de la façon dont la macro a été définie.
\if\relax\detokenize\expandafter{\PagesForcees}\relax\else
  \let\NbPages\PagesForcees
\fi
\if\relax\detokenize\expandafter{\VersionForcee}\relax\else
  \let\Version\VersionForcee
\fi

% Format du livre fini.
\def\@formatAquatre{a4}
\def\@formatVingt{20x24}
\def\@formatVingtMarge{20x24-marge}
\ifx\Format\@formatAquatre
  \def\PageLargeur{210}\def\PageHauteur{297}
\else\ifx\Format\@formatVingt
  \def\PageLargeur{200}\def\PageHauteur{240}
\else\ifx\Format\@formatVingtMarge
  \def\PageLargeur{200}\def\PageHauteur{240}
\else
  \PackageError{couverture}{Format \Format\space inconnu}%
    {Valeurs admises : a4, 20x24, 20x24-marge.}%
  \def\PageLargeur{210}\def\PageHauteur{297}
\fi\fi\fi
% Reliure. Hors de « souple » et « rigide » : erreur, pas de supposition.
\def\@reliuresouple{souple}
\def\@reliurerigide{rigide}
\ifx\Reliure\@reliurerigide
  \edef\PageLargeur{\fpeval{\PageLargeur + 5}}
  \edef\PageHauteur{\fpeval{\PageHauteur + 5}}
  \def\FondPerdu{15}
\else\ifx\Reliure\@reliuresouple\else
  \PackageError{couverture}{Reliure \Reliure\space inconnue}%
    {Valeurs admises : souple, rigide.}%
\fi\fi
\makeatother

% Garde-fou : une édition jamais produite n'a pas de couverture.
\ifnum\NbPages=0
  \PackageError{couverture}{L'edition \Classe\space n'a pas de pagination}%
    {Cette edition n'a jamais ete compilee. Renseigner \string\Donnees\Classe
     ou passer \string\PagesForcees.}%
\fi

% Avertissement : l'imprimeur demande un nombre de pages multiple de 4
% (17/09/2026). Pas une erreur : c'est son exigence, pas une loi du papier.
\ifnum\fpeval{\NbPages - 4*floor(\NbPages/4)}=0 \else
  \typeout{COUVERTURE-AVERTISSEMENT \NbPages\space pages : pas un multiple
    de 4, que l'imprimeur exige -- ajouter des pages a l'interieur}%
\fi

% Garde-fou : un dos carré collé impose un nombre de pages pair.
\ifodd\NbPages
  \PackageError{couverture}{Nombre de pages impair (\NbPages)}%
    {Un dos carre colle exige un nombre PAIR de pages : une feuille en
     porte deux. Ajouter une page blanche a l'interieur.}%
\fi

% ---------------------------------------------------------------------
% 5. GÉOMÉTRIE — c'est ici que se joue l'exactitude du dos
% ---------------------------------------------------------------------

% Épaisseur du dos : valeur imposée si elle est donnée ET qu'elle vaut pour
% cette pagination et ce format ; calcul sinon.
\edef\DosCalcule{\fpeval{round(\NbPages / 2 * \Grammage * \Bouffant / 1000, 2)}}
\def\DosOrigine{calculee}
\ifnum\fpeval{\DosImpose > 0}=1
  \ifnum\fpeval{\NbPages = \DosImposePages}=1
    \ifx\Format\DosImposeFormat
      \ifx\Reliure\DosImposeReliure
        \def\DosOrigine{imposee}
      \fi
    \fi
  \fi
  \def\MotCalculee{calculee}
  \ifx\DosOrigine\MotCalculee
    \typeout{COUVERTURE-AVERTISSEMENT dos impose de \DosImpose\space mm
      donne pour \DosImposePages\space pages en \DosImposeFormat\space
      (reliure \DosImposeReliure) -- inapplicable a \NbPages\space pages en
      \Format\space (reliure \Reliure), dos calcule}%
  \fi
\fi
% NB : pas de \@ dans ces noms. Cette section est hors \makeatletter, et
% « \def\@imposee{...} » y redéfinit \@ (suivi du texte « imposee ») : le
% test échouait toujours, en silence. Constaté le 16/09/2026 — le dos de N
% tombait juste par coïncidence, 14,0 mm calculés pour 14 imposés.
\def\MotRigide{rigide}
\def\MotImposee{imposee}
\ifx\DosOrigine\MotImposee
  \edef\Dos{\DosImpose}
\else
  \edef\Dos{\DosCalcule}
\fi

% Marge technique : uniquement en mode épreuve.
\def\ouimot{oui}
\ifx\Reperes\ouimot \edef\Tech{\MargeTech}\else \def\Tech{0}\fi

% Dimensions du papier.
%   largeur = marge + fond perdu + 4e + dos + 1re + fond perdu + marge
%   hauteur = marge + fond perdu + hauteur + fond perdu + marge
\edef\PapierL{\fpeval{2*\PageLargeur + \Dos + 2*\FondPerdu + 2*\Tech}}
\edef\PapierH{\fpeval{\PageHauteur + 2*\FondPerdu + 2*\Tech}}

\geometry{paperwidth=\PapierL mm, paperheight=\PapierH mm,
          margin=0pt, nohead, nofoot}

% ---------------------------------------------------------------------
%  TOUTES LES COTES SONT FIGÉES ICI, ET NULLE PART AILLEURS.
%
%  Raison, payée par trois compilations ratées : un \fpeval laissé dans
%  une coordonnée ou dans le texte d'un nœud TikZ n'est pas développé au
%  moment où TikZ analyse le chemin. On obtient « Giving up on this path.
%  Did you forget a semicolon? », qui accuse une ponctuation alors que la
%  faute est ailleurs. Le corps du document ci-dessous n'emploie donc que
%  des macros déjà calculées.
% ---------------------------------------------------------------------

% Repères horizontaux, mesurés depuis le bord GAUCHE du papier.
\edef\CoupeG{\fpeval{\Tech + \FondPerdu}}          % coupe gauche (4e de couv.)
\edef\DosG{\fpeval{\CoupeG + \PageLargeur}}         % dos, début
\edef\DosD{\fpeval{\DosG + \Dos}}                   % dos, fin
\edef\CoupeD{\fpeval{\DosD + \PageLargeur}}         % coupe droite (1re de couv.)
% Repères verticaux, depuis le bord BAS.
\edef\CoupeB{\fpeval{\Tech + \FondPerdu}}
\edef\CoupeH{\fpeval{\CoupeB + \PageHauteur}}
% Axes.
\edef\DosAxe{\fpeval{(\DosG + \DosD)/2}}
\edef\MilieuH{\fpeval{(\CoupeB + \CoupeH)/2}}
\edef\MilieuL{\fpeval{\PapierL/2}}

% Format fini, pour les cotes affichées.
\edef\FiniL{\fpeval{2*\PageLargeur + \Dos}}

% Bandeau de la 1re de couverture : même proportion que la page de titre.
\edef\BandeauL{\fpeval{round(0.34*\PageLargeur, 2)}}
\edef\BandeauG{\fpeval{\CoupeD - \BandeauL}}
\edef\BandeauAccent{\fpeval{\BandeauG - 0.7}}
\edef\DosAccentG{\fpeval{\DosG - 0.7}}
\edef\DosAccentD{\fpeval{\DosD + 0.7}}
% Couverture rigide : le bleu du dos déborde de 2 mm sur chaque plat, pour
% absorber la tolérance de l'imprimeur sur l'épaisseur du dos (±10 %).
\edef\DosEtenduG{\fpeval{\DosG - 2}}
\edef\DosEtenduD{\fpeval{\DosD + 2}}

% Filigrane de classe.
\edef\FiligraneX{\fpeval{(\BandeauG + \CoupeD)/2}}
\edef\FiligraneY{\fpeval{\CoupeH - 6}}

% Bloc-titre de la 1re de couverture.
\edef\TitreX{\fpeval{\DosD + 12}}
\edef\TitreY{\fpeval{\CoupeH - 72}}
\edef\TitreW{\fpeval{\BandeauG - \DosD - 16}}
\edef\VersionY{\fpeval{\CoupeB + 14}}

% 4e de couverture.
\edef\QuatreX{\fpeval{\CoupeG + \Securite + 6}}
\edef\QuatreY{\fpeval{\CoupeH - \Securite - 14}}
\edef\QuatreBasY{\fpeval{\CoupeB + \Securite + 22}}
\edef\QuatreW{\fpeval{\PageLargeur - 2*\Securite - 12}}

% Zone de sécurité.
\edef\SecG{\fpeval{\CoupeG + \Securite}}
\edef\SecD{\fpeval{\CoupeD - \Securite}}
\edef\SecB{\fpeval{\CoupeB + \Securite}}
\edef\SecH{\fpeval{\CoupeH - \Securite}}

% Zones de contrôle : ce qui doit rester de part et d'autre du pli du dos.
\edef\PliQuatre{\fpeval{\DosG - \Securite}}         % 4e : limite côté dos
\edef\PliPremiere{\fpeval{\DosD + \Securite}}       % 1re : limite côté dos

% Traits de coupe et cotes, dans la marge technique.
\edef\TechExt{\fpeval{\Tech - 8}}
\edef\TechInt{\fpeval{\Tech - 1}}
\edef\TechCote{\fpeval{\Tech - 9}}
\edef\HautInt{\fpeval{\PapierH - \Tech + 1}}
\edef\HautExt{\fpeval{\PapierH - \Tech + 8}}
\edef\HautCote{\fpeval{\PapierH - \Tech + 9}}
\edef\DroitInt{\fpeval{\PapierL - \Tech + 1}}
\edef\DroitExt{\fpeval{\PapierL - \Tech + 8}}

% Emplacement du code-barres ISBN : en bas à DROITE de la 4e de couverture,
% selon l'usage éditorial, et non à gauche où il chevaucherait les mentions
% de licence.
\edef\IsbnD{\fpeval{\DosG - \Securite}}
\edef\IsbnG{\fpeval{\IsbnD - 40}}
\edef\IsbnH{\fpeval{\SecB + 25}}
\edef\IsbnCX{\fpeval{\IsbnG + 20}}
\edef\IsbnCY{\fpeval{\SecB + 12.5}}

% Cartouches de cotes, dans la marge technique. Ancrés vers l'INTÉRIEUR du
% papier : ancrés vers l'extérieur, ils débordaient et étaient rognés.
\edef\CoteBasY{\fpeval{1}}
\edef\CoteHautY{\fpeval{\PapierH - 1}}

% ---------------------------------------------------------------------
% 6. COULEURS ET POLICES — reprises de la page de titre (build_book.py)
% ---------------------------------------------------------------------

\definecolor{TitleBand}{cmyk}{.9,.55,.1,.35}   % bleu profond
\definecolor{TitleAccent}{cmyk}{.05,.8,1,0}    % rouge DARC (accent fin)
\definecolor{Repere}{cmyk}{0,1,1,0}            % magenta : repères de montage

\setmainfont{Libertinus Serif}[
  ItalicFont     = Libertinus Serif Italic,
  BoldFont       = Libertinus Serif Bold,
  BoldItalicFont = Libertinus Serif Bold Italic]

\pagestyle{empty}

% ---------------------------------------------------------------------
% 7. FILIGRANE DE CLASSE
%    Le texte est construit ici, et non dans le nœud : un \\ produit à
%    l'intérieur d'un \foreach referme le groupe d'alignement de TikZ.
%    Une lettre seule est grande ; plusieurs s'empilent et se réduisent
%    (build_book.py v0.20 : « NEA » à 105 pt débordait encore).
% ---------------------------------------------------------------------

\makeatletter
\def\FiligraneTexte{}
\newif\ifpremierelettre \premierelettretrue
\count255=0
\@for\@tempa:=\Lettres\do{%
  \advance\count255 by 1\relax
  \ifpremierelettre \premierelettrefalse
  \else \g@addto@macro\FiligraneTexte{\\}\fi
  \expandafter\g@addto@macro\expandafter\FiligraneTexte\expandafter{\@tempa}}%
\xdef\NbLettres{\the\count255}
\makeatother

\ifnum\NbLettres=1 \def\CorpsF{150}\else\def\CorpsF{86}\fi
\edef\CorpsFI{\fpeval{round(\CorpsF*1.02, 0)}}

% Texte du cartouche de cotes, lui aussi construit hors du nœud TikZ.
%
% NE PAS y employer \textsuperscript : dans un nœud « align=center », il
% casse le chemin TikZ (« Giving up on this path. Did you forget a
% semicolon? »), et l'erreur est signalée sur le nœud SUIVANT — elle
% accuse donc une ponctuation là où la faute est ici. Isolé par essais :
% $^2$ et le caractère ² passent, \textsuperscript non.
\def\CotesTexte{%
  \Classe\ --- \NbPages\ pages --- papier \Grammage\ g/m$^2$,
  bouffant \Bouffant\ --- fond perdu \FondPerdu\ mm\\
  papier \PapierL\ mm sur \PapierH\ mm --- fini \FiniL\ mm sur
  \PageHauteur\ mm\\
  EPREUVE, ne pas envoyer a l'impression%
}

% ---------------------------------------------------------------------
% 8. CONTRÔLE DE SUPERPOSITION
%
%  \MesureNoeud{nœud}{préfixe}{retrait} mesure la boîte d'un nœud par ses
%  quatre coins (un nœud pivoté a ses ancres pivotées : le min et le max des
%  quatre suffisent) et la ramène au repère du papier, en mm. Le point (0,0)
%  est lu par le même « let » que les ancres : les deux passent donc par la
%  même transformation, et la comparaison ne dépend pas de l'endroit où
%  TeX a posé l'image sur la page.
%
%  Le retrait sert au seul filigrane, pour mesurer les lettres et non leur
%  cadre : c'est le bord des lettres que le détourage rognerait.
%
%  PIÈGE MESURÉ le 16/09/2026 : la marge interne TikZ (.3333em) n'est PAS
%  évaluée dans la police du nœud, mais dans celle qui entoure le dessin.
%  Instrumenté : dans le nœud, em = 150 pt ; boîte brute du « N » :
%  41,7 x 37,5 mm, soit la lettre plus 2 x 1,3 mm — 1,3 mm étant .3333em
%  à 11 pt. Un retrait calculé à 150 pt (17,6 mm) ramenait la lettre à
%  6,6 mm de large, et le contrôle aurait dit OK sur une mesure absurde.
% ---------------------------------------------------------------------

\newcount\NbAlertes
\newdimen\FiligraneSep
\newsavebox\CreditsBoite

% Crédits de 4e : lignes coupées à la main. Leur largeur RÉELLE est celle
% de la plus longue ligne, mesurée dans une boîte à part (voir § 8, ISBN).
\def\CreditsTexte{%
    Réalisé à partir des contenus de 50ohm.de (en allemand)\\
    50ohm.de-Autorenteam, coordonné par le référat AJW du DARC e.\,V.\\
    \Auteur\\[3pt]
    Licence CC BY 4.0 --- \texttt{\Depot}}

\def\MesureNoeud#1#2#3{%
  \path let \p1=(#1.south west), \p2=(#1.north east),
            \p3=(#1.north west), \p4=(#1.south east), \p5=(0,0) in
    \pgfextra{%
      \expandafter\xdef\csname #2xmin\endcsname{\fpeval{round((min(\x1,\x2,\x3,\x4) - \x5 + #3)/1mm, 2)}}%
      \expandafter\xdef\csname #2xmax\endcsname{\fpeval{round((max(\x1,\x2,\x3,\x4) - \x5 - #3)/1mm, 2)}}%
      \expandafter\xdef\csname #2ymin\endcsname{\fpeval{round((min(\y1,\y2,\y3,\y4) - \y5 + #3)/1mm, 2)}}%
      \expandafter\xdef\csname #2ymax\endcsname{\fpeval{round((max(\y1,\y2,\y3,\y4) - \y5 - #3)/1mm, 2)}}%
    };}

% \ControleDans{libellé}{préfixe}{xmin}{xmax}{ymin}{ymax}
\def\ControleDans#1#2#3#4#5#6{%
  \ifnum\fpeval{\csname #2xmin\endcsname >= #3 && \csname #2xmax\endcsname <= #4
             && \csname #2ymin\endcsname >= #5 && \csname #2ymax\endcsname <= #6}=1
    \typeout{COUVERTURE-CONTROLE OK #1 : x \csname #2xmin\endcsname-\csname #2xmax\endcsname,
      y \csname #2ymin\endcsname-\csname #2ymax\endcsname\space dans x #3-#4, y #5-#6}%
  \else
    \global\advance\NbAlertes by 1
    \typeout{COUVERTURE-CONTROLE ALERTE #1 : x \csname #2xmin\endcsname-\csname #2xmax\endcsname,
      y \csname #2ymin\endcsname-\csname #2ymax\endcsname\space HORS de x #3-#4, y #5-#6}%
  \fi}

% \ControleDisjoints{libellé}{préfixe}{xmin}{xmax}{ymin}{ymax}
% Deux rectangles se chevauchent si et seulement s'ils se recouvrent sur
% les DEUX axes à la fois.
\def\ControleDisjoints#1#2#3#4#5#6{%
  \ifnum\fpeval{\csname #2xmin\endcsname < #4 && #3 < \csname #2xmax\endcsname
             && \csname #2ymin\endcsname < #6 && #5 < \csname #2ymax\endcsname}=1
    \global\advance\NbAlertes by 1
    \typeout{COUVERTURE-CONTROLE ALERTE #1 : chevauchement}%
  \else
    \typeout{COUVERTURE-CONTROLE OK #1 : disjoints}%
  \fi}

\begin{document}
% Marge interne du filigrane, évaluée dans la police qui ENTOURE le dessin :
% c'est là que TikZ évalue le .3333em de « inner sep » (cf. § 8).
\FiligraneSep=.3333em
% Largeur réelle des crédits : la plus longue ligne, dans leur police.
\sbox\CreditsBoite{\small\begin{tabular}{@{}l@{}}\CreditsTexte\end{tabular}}%
\begin{tikzpicture}[remember picture, overlay,
                    x=1mm, y=1mm, shift={(current page.south west)}]

% ---- Fonds ----------------------------------------------------------
% La 4e de couverture reste sur le blanc du papier.
% Le bandeau bleu couvre le dos, puis reprend au tiers extérieur de la
% 1re de couverture — même partition 2/3-1/3 que la page de titre.
% Couverture rigide : pas de filets au dos (ils tomberaient à côté des
% plis, vu la tolérance), et un dos bleu élargi de 2 mm de chaque côté.
\ifx\Reliure\MotRigide
  \fill[TitleBand] (\DosEtenduG, 0) rectangle (\DosEtenduD, \PapierH);
\else
  \fill[TitleBand] (\DosG, 0) rectangle (\DosD, \PapierH);
\fi
\fill[TitleBand] (\BandeauG, 0) rectangle (\PapierL, \PapierH);
% Filets d'accent : à la jonction sur la 1re, de part et d'autre du dos.
\fill[TitleAccent] (\BandeauAccent, 0) rectangle (\BandeauG, \PapierH);
\ifx\Reliure\MotRigide\else
  \fill[TitleAccent] (\DosAccentG, 0) rectangle (\DosG, \PapierH);
  \fill[TitleAccent] (\DosD, 0) rectangle (\DosAccentD, \PapierH);
\fi

% ---- 1re de couverture ----------------------------------------------
% Filigrane de classe, en réserve claire, détouré au bandeau.
\begin{scope}
\clip (\BandeauG, 0) rectangle (\PapierL, \PapierH);
\node[name=filigrane, anchor=north, align=center, text=white, opacity=0.16,
      font=\fontsize{\CorpsF}{\CorpsFI}\selectfont\bfseries]
  at (\FiligraneX, \FiligraneY) {\FiligraneTexte};
\end{scope}

\node[name=titre, anchor=west, align=left, text width=\TitreW mm]
  at (\TitreX, \TitreY) {%
    {\fontsize{54}{56}\selectfont\bfseries \TitreOuvrage}\\[10pt]
    {\Large\color{TitleBand}\bfseries \SousTitre}\\[16pt]
    {\large\color{black!75} \SousTitreClasse}};

\node[name=version, anchor=south west, align=left, font=\small, text=black!70]
  at (\TitreX, \VersionY) {Version \Version\\ Licence CC BY 4.0};

% ---- Dos -------------------------------------------------------------
% Un dos étroit ne reçoit pas de texte : en deçà de 6 mm, le débord au
% massicot ferait mordre le texte vertical sur les plats.
\ifnum\fpeval{\Dos >= 6}=1
  \node[name=dos, rotate=90, anchor=center, text=white, font=\bfseries,
        inner sep=0pt]
    at (\DosAxe, \MilieuH)
    {\TitreOuvrage\ \ --- \ \SousTitreClasse\ \ --- \ \Version};
\fi

% ---- 4e de couverture ------------------------------------------------
\node[name=quatre, anchor=north west, align=justify, text width=\QuatreW mm]
  at (\QuatreX, \QuatreY) {\QuatriemeTexte};

% Le texte vit dans \CreditsTexte (§ 8) pour être mesuré à part : la boîte
% du nœud a la largeur fixe \QuatreW et couvre l'emplacement ISBN sans
% qu'aucun caractère n'y soit. Retirer « text width » pour l'éviter change
% l'interlignage — essayé le 16/09/2026, 6 934 pixels différents.
\node[name=credits, anchor=south west, align=left, font=\small, text=black!70,
      text width=\QuatreW mm]
  at (\QuatreX, \QuatreBasY) {\CreditsTexte};

% ---- Repères de fabrication (mode épreuve uniquement) ----------------
\ifx\Reperes\ouimot
  % Traits de coupe, hors du fond perdu.
  \foreach \x in {\CoupeG, \CoupeD} {
    \draw[Repere, line width=0.25pt] (\x, \TechExt) -- (\x, \TechInt);
    \draw[Repere, line width=0.25pt] (\x, \HautInt) -- (\x, \HautExt);}
  \foreach \y in {\CoupeB, \CoupeH} {
    \draw[Repere, line width=0.25pt] (\TechExt, \y) -- (\TechInt, \y);
    \draw[Repere, line width=0.25pt] (\DroitInt, \y) -- (\DroitExt, \y);}
  % Limites du dos, en pointillés.
  \foreach \x in {\DosG, \DosD} {
    \draw[Repere, line width=0.25pt, dashed] (\x, \TechExt) -- (\x, \TechInt);
    \draw[Repere, line width=0.25pt, dashed] (\x, \HautInt) -- (\x, \HautExt);}
  % Zone de sécurité.
  \draw[Repere, line width=0.2pt, dotted]
    (\SecG, \SecB) rectangle (\SecD, \SecH);
  % Emplacement réservé au code-barres ISBN.
  \draw[Repere, line width=0.2pt] (\IsbnG, \SecB) rectangle (\IsbnD, \IsbnH);
  \node[font=\tiny\sffamily, text=Repere, align=center]
    at (\IsbnCX, \IsbnCY) {emplacement\\ISBN};
  % Cotes.
  \node[anchor=north, font=\tiny\sffamily, text=Repere]
    at (\DosAxe, \CoteHautY) {dos \Dos\ mm (\DosOrigine)};
  \node[anchor=south, font=\tiny\sffamily, text=Repere, align=center]
    at (\MilieuL, \CoteBasY) {\CotesTexte};
\fi

% ---- Contrôle de superposition (§ 8) ----------------------------------
% Rien n'est dessiné : des mesures et des lignes de journal.
\MesureNoeud{filigrane}{Fil}{\the\FiligraneSep}
\MesureNoeud{titre}{Tit}{0pt}
\MesureNoeud{version}{Ver}{0pt}
\MesureNoeud{quatre}{Qua}{0pt}
\MesureNoeud{credits}{Cre}{0pt}
% Bord droit des crédits = bord gauche + marge interne + plus longue ligne.
\xdef\Crexmax{\fpeval{round(\Crexmin + (\the\FiligraneSep + \the\wd\CreditsBoite)/1mm, 2)}}

% Le filigrane peut entrer dans la zone de sécurité (il est décoratif),
% pas franchir la coupe ni sortir du bandeau — le détourage le rognerait.
\ControleDans{filigrane dans le bandeau}{Fil}{\BandeauG}{\CoupeD}{\CoupeB}{\CoupeH}
\ControleDans{titre sur la 1re, hors bandeau}{Tit}{\PliPremiere}{\BandeauAccent}{\SecB}{\SecH}
\ControleDans{version sur la 1re, hors bandeau}{Ver}{\PliPremiere}{\BandeauAccent}{\SecB}{\SecH}
\ControleDans{texte de 4e dans la zone de securite}{Qua}{\SecG}{\PliQuatre}{\SecB}{\SecH}
\ControleDans{credits de 4e dans la zone de securite}{Cre}{\SecG}{\PliQuatre}{\SecB}{\SecH}
\ControleDisjoints{titre / version}{Tit}{\Verxmin}{\Verxmax}{\Verymin}{\Verymax}
\ControleDisjoints{texte de 4e / credits}{Qua}{\Crexmin}{\Crexmax}{\Creymin}{\Creymax}
\ControleDisjoints{texte de 4e / emplacement ISBN}{Qua}{\IsbnG}{\IsbnD}{\SecB}{\IsbnH}
\ControleDisjoints{credits / emplacement ISBN}{Cre}{\IsbnG}{\IsbnD}{\SecB}{\IsbnH}
\ifnum\fpeval{\Dos >= 6}=1
  \MesureNoeud{dos}{Dos}{0pt}
  \ControleDans{texte du dos entre les plis}{Dos}{\DosG}{\DosD}{\SecB}{\SecH}
\else
  \typeout{COUVERTURE-CONTROLE OK dos de \Dos\space mm : trop etroit, sans texte}%
\fi

\end{tikzpicture}
\typeout{COUVERTURE-BILAN alertes=\the\NbAlertes\space; classe=\Classe\space;
  format=\Format\space; reliure=\Reliure\space; pages=\NbPages\space;
  version=\Version\space;
  dos=\Dos\space mm (\DosOrigine); fini=\FiniL x \PageHauteur\space mm;
  papier=\PapierL x \PapierH\space mm}
\end{document}
"
%
%  ---------------------------------------------------------------------
%  DEUX MODES, un seul paramètre : \Reperes
%
%    oui  — mode ÉPREUVE. Ajoute une marge technique, les traits de coupe,
%           les limites du dos, la zone de sécurité et les cotes. C'est ce
%           qu'il faut pour vérifier à l'écran.
%    non  — mode PRODUCTION. Papier = format final + fond perdu, rien
%           d'autre. C'est ce fichier-là qu'on envoie à l'imprimeur.
%
%  ---------------------------------------------------------------------
%  CONTRÔLE DE SUPERPOSITION
%
%  À la fin du dessin, chaque bloc de texte est mesuré par ses ancres TikZ
%  (quatre coins, ce qui vaut aussi pour le texte pivoté du dos) et comparé
%  à la zone où il doit tenir, puis aux blocs voisins. Le journal porte une
%  ligne par vérification :
%
%      COUVERTURE-CONTROLE OK      <quoi> : ...
%      COUVERTURE-CONTROLE ALERTE  <quoi> : ...
%
%  et un bilan : COUVERTURE-BILAN alertes=<n> ; ...
%
%  Ce que le contrôle VOIT : un bloc qui franchit le pli du dos, la coupe,
%  la zone de sécurité ou le bord du bandeau ; deux blocs qui se chevauchent ;
%  un bloc qui empiète sur l'emplacement réservé au code-barres ISBN.
%  Ce qu'il NE VOIT PAS : une ligne trop longue à l'INTÉRIEUR d'un bloc à
%  largeur fixe (la boîte garde sa largeur, le texte déborde). Celle-là se
%  signale par un « Overfull \hbox » dans le même journal, que la fenêtre
%  compte aussi.
%
%  ---------------------------------------------------------------------
%  AVERTISSEMENT SUR L'ÉPAISSEUR DU DOS
%
%  Le dos se calcule ainsi :
%
%      épaisseur d'une feuille (mm) = grammage x bouffant / 1000
%      dos (mm)                     = (nombre de pages / 2) x épaisseur
%
%  On divise par deux parce qu'une feuille porte DEUX pages.
%
%  Le « bouffant » (main, volume, bulk) dépend du papier et n'est pas
%  devinable : un offset 80 g courant est autour de 1,30, un bouffant
%  d'édition monte à 1,75 et plus. UN ÉCART DE BOUFFANT DÉPLACE LE DOS DE
%  PLUSIEURS MILLIMÈTRES. Demander à l'imprimeur soit son bouffant, soit
%  directement l'épaisseur du dos, et reporter ci-dessous. Tant que cette
%  valeur n'est pas confirmée, la couverture n'est pas bonne à tirer.
%
%  Si l'imprimeur donne le dos en clair, forcer la valeur avec \DosImpose.
% =====================================================================

\documentclass[11pt]{article}

% ---------------------------------------------------------------------
% 1. PARAMÈTRES — c'est la seule zone à modifier
% ---------------------------------------------------------------------

% Classe de l'ouvrage : N, E, A, NE, EA ou NEA.
% (\providecommand : une valeur passée en ligne de commande l'emporte.)
\providecommand{\Classe}{NEA}

% Mode épreuve (oui) ou production (non).
\providecommand{\Reperes}{oui}

% Format du livre FINI : a4, 20x24 ou 20x24-marge — les noms de --format
% dans build_book.py. Les deux maquettes 20 x 24 ne diffèrent que par leur
% colonne de marge intérieure : le livre fini, donc la couverture, est le
% même (200 x 240 mm).
\providecommand{\Format}{a4}

% Reliure : « souple » (dos carré collé, fond perdu 3 mm) ou « rigide »
% (couverture rigide, dite hardcover). En rigide, consignes de l'imprimeur
% du 17/09/2026 : chaque plat déborde du livre de 5 mm en largeur et en
% hauteur (le carton dépasse du bloc), et le fichier porte 15 mm de rabat
% autour du carton au lieu de 3 mm de fond perdu. Les aplats vont déjà au
% bord du papier : ils remplissent le rabat d'eux-mêmes.
\providecommand{\Reliure}{souple}

% Pagination et version imposées de l'extérieur. Vides = tables du § 2.
% La fenêtre les renseigne toujours : la pagination vient du journal du
% livre réellement compilé, dans le format demandé.
\providecommand{\PagesForcees}{}
\providecommand{\VersionForcee}{}

% Papier de l'INTÉRIEUR du livre (pas de la couverture).
% CoolLibri, papier « Standard 90 g blanc ».
\def\Grammage{90}        % g/m²
\def\Bouffant{1.206}     % main DÉDUITE, voir ci-dessous

% Épaisseur du dos IMPOSÉE par l'imprimeur, en mm, par classe.
% Une valeur > 0 prime sur le calcul — MAIS SEULEMENT pour la pagination et
% le format pour lesquels l'imprimeur l'a donnée (\DosImposePages<classe>,
% \DosImposeFormat<classe>). Hors de ce cas, le dos est calculé et le
% journal le dit : un dos de 14 mm donné pour 258 pages A4 ne vaut rien pour
% 262 pages, ni pour le même livre recomposé en 20 x 24.
%
% ATTENTION : cette valeur est propre à CHAQUE ÉDITION. Un dos unique
% appliqué aux quatre livres serait faux pour trois d'entre eux.
%
%  N   — 14 mm, VALEUR DONNÉE PAR COOLIBRI pour 258 pages en 90 g, A4.
%        Fichier attendu : 440 x 303 mm, ce que ce gabarit produit en mode
%        production (2 x 210 + 14 + 2 x 3 = 440 ; 297 + 2 x 3 = 303).
%        Le N de la a.3 compte 262 pages : cette valeur ne s'y applique plus.
%  E, A, NEA — non demandés à l'imprimeur : ils sont CALCULÉS avec la main
%        déduite ci-dessous. À faire confirmer avant tout tirage.
%
% La main de 1,206 n'est pas publiée par CoolLibri : elle est déduite de
% leur propre chiffre, 14 mm / (0,090 x 258/2) = 1,206. Elle reproduit donc
% exactement le dos de N, et n'est qu'une ESTIMATION pour les autres.
% Écarts que cela recouvre, pour mémoire : une main de 1,0 donnerait
% 11,61 mm pour N, une main de 1,3 en donnerait 16,77. Il ne fallait pas
% deviner.
%  E   — 17 mm, VALEUR DONNÉE PAR L'IMPRIMEUR le 17/09/2026 pour 276 pages,
%        20 x 24, couverture RIGIDE, reliure collée (Klebebindung), papier
%        115 g. En reliure cousue (Fadenheftung), il annonce 18 mm.
%
% \DosImposeReliure<classe> : la reliure pour laquelle la valeur a été
% donnée. Le dos d'une couverture rigide n'est pas celui d'un dos carré
% collé souple, même à pagination égale.
\def\DosImposeN{14}   \def\DosImposePagesN{258}  \def\DosImposeFormatN{a4}
\def\DosImposeE{17}   \def\DosImposePagesE{276}  \def\DosImposeFormatE{20x24-marge}
\def\DosImposeA{0}    \def\DosImposePagesA{0}    \def\DosImposeFormatA{}
\def\DosImposeNE{0}   \def\DosImposePagesNE{0}   \def\DosImposeFormatNE{}
\def\DosImposeEA{0}   \def\DosImposePagesEA{0}   \def\DosImposeFormatEA{}
\def\DosImposeNEA{0}  \def\DosImposePagesNEA{0}  \def\DosImposeFormatNEA{}
\def\DosImposeReliureN{souple}  \def\DosImposeReliureE{rigide}
\def\DosImposeReliureA{souple}  \def\DosImposeReliureNE{souple}
\def\DosImposeReliureEA{souple} \def\DosImposeReliureNEA{souple}

% Fond perdu, en mm. 3 mm est la valeur courante ; certains imprimeurs
% demandent 5 mm.
\def\FondPerdu{3}

% Marge technique autour du fond perdu, en mm. Sert à loger les traits de
% coupe en mode épreuve. Ignorée en mode production.
\def\MargeTech{12}

% Zone de sécurité, en mm : aucun texte ne doit s'en approcher davantage
% du bord de coupe.
\def\Securite{8}

% ---------------------------------------------------------------------
% 2. DONNÉES PAR CLASSE
%    pages : nombre de pages du PDF intérieur, COUVERTURE NON COMPRISE.
%    Un dos carré collé impose un nombre de pages PAIR ; le fichier
%    refuse de compiler si le compte est impair.
%
%    Ces tables ne servent que si \PagesForcees et \VersionForcee sont
%    vides, c'est-à-dire à la main. Valeurs A4 mesurées sur les PDF de la
%    a.3 (CLAUDE.md § 1, 05/09/2026). En 20 x 24, la pagination change
%    entièrement : NEA passe de 816 à 1036 pages.
% ---------------------------------------------------------------------

\def\DonneesN{262}
\def\DonneesE{214}
\def\DonneesA{386}
\def\DonneesNE{0}      % édition jamais produite
\def\DonneesEA{0}      % édition jamais produite
\def\DonneesNEA{816}

% Version imprimée sur le dos et la 1re de couverture, PAR CLASSE.
% « Les versions sont suivies par classe, chacune évoluant à son rythme »
% (CHANGELOG.md).
\def\VersionN{a.3}
\def\VersionE{a.3}
\def\VersionA{a.3}
\def\VersionNE{a.3}
\def\VersionEA{a.3}
\def\VersionNEA{a.3}

% ---------------------------------------------------------------------
% 3. TEXTES
% ---------------------------------------------------------------------

\def\TitreOuvrage{50\,Ohm}
\def\SousTitre{Préparation à l'examen radioamateur}

% Sous-titre propre à chaque classe (repris de FR_TITLES, build_book.py).
\def\SousTitreN{Cours complet — classe N}
\def\SousTitreE{Cours complet — classe E}
\def\SousTitreA{Cours complet — classe A}
\def\SousTitreNE{Cours complet — classes N et E}
\def\SousTitreEA{Cours complet — classes E et A}
\def\SousTitreNEA{Cours complet — classes N, E et A}

% Lettres du filigrane, empilées de haut en bas (cf. build_book.py v0.20).
\def\LettresN{N}
\def\LettresE{E}
\def\LettresA{A}
\def\LettresNE{N,E}
\def\LettresEA{E,A}
\def\LettresNEA{N,E,A}

% Texte de 4e de couverture.
% NB : il parle des classes N, E et A quelle que soit l'édition. Juste pour
% N, E, A et NEA ; approximatif pour NE et EA ; faux pour SWL, d'où son refus.
\def\QuatriemeTexte{%
Ce guide d'étude prépare aux examens radioamateur allemands des classes
N, E et A. Il traduit en français les contenus du projet 50ohm.de, publiés
par le référat AJW du DARC e.\,V., et les complète d'encarts « En France »
qui signalent, sans rien retirer au texte allemand, ce que la
réglementation française prévoit sur le même sujet.\par\smallskip
Les questions officielles du catalogue allemand accompagnent chaque
notion, avec leurs réponses.%
}

\def\Auteur{Traduit avec l'aide d'une IA par Pierre F4JWI}
\def\Depot{github.com/DrPLG/50ohm-fr}

% =====================================================================
%  À PARTIR D'ICI, PLUS RIEN À MODIFIER
% =====================================================================

\usepackage{iftex}
\RequireLuaTeX
\usepackage{xfp}
\usepackage{tikz}
\usetikzlibrary{calc}
\usepackage{fontspec}
\usepackage[french]{babel}
\usepackage{geometry}

% ---------------------------------------------------------------------
% 4. RÉSOLUTION DES DONNÉES DE LA CLASSE
% ---------------------------------------------------------------------

\makeatletter
% Tome SWL : refus explicite, plutôt qu'un « classe inconnue » qui laisserait
% croire à un oubli. Le texte de 4e parle des classes N, E et A.
\def\@swlmot{SWL}
\ifx\Classe\@swlmot
  \PackageError{couverture}{Pas de couverture pour le tome SWL}%
    {Sa 4e de couverture n'est pas redigee : le texte commun parle des
     classes N, E et A. C'est une decision editoriale a prendre d'abord.}%
\fi

\def\@resoudre#1#2{%
  \expandafter\let\expandafter#1\csname #2\Classe\endcsname
  \ifx#1\relax
    \PackageError{couverture}{Classe \Classe\space inconnue : pas de #2\Classe}%
      {Valeurs admises : N, E, A, NE, EA, NEA.}%
  \fi}
\@resoudre\NbPages{Donnees}
\@resoudre\SousTitreClasse{SousTitre}
\@resoudre\Lettres{Lettres}
\@resoudre\DosImpose{DosImpose}
\@resoudre\DosImposePages{DosImposePages}
\@resoudre\DosImposeFormat{DosImposeFormat}
\@resoudre\DosImposeReliure{DosImposeReliure}
\@resoudre\Version{Version}

% Valeurs imposées de l'extérieur. Test de vacuité par \detokenize : il ne
% dépend pas de la façon dont la macro a été définie.
\if\relax\detokenize\expandafter{\PagesForcees}\relax\else
  \let\NbPages\PagesForcees
\fi
\if\relax\detokenize\expandafter{\VersionForcee}\relax\else
  \let\Version\VersionForcee
\fi

% Format du livre fini.
\def\@formatAquatre{a4}
\def\@formatVingt{20x24}
\def\@formatVingtMarge{20x24-marge}
\ifx\Format\@formatAquatre
  \def\PageLargeur{210}\def\PageHauteur{297}
\else\ifx\Format\@formatVingt
  \def\PageLargeur{200}\def\PageHauteur{240}
\else\ifx\Format\@formatVingtMarge
  \def\PageLargeur{200}\def\PageHauteur{240}
\else
  \PackageError{couverture}{Format \Format\space inconnu}%
    {Valeurs admises : a4, 20x24, 20x24-marge.}%
  \def\PageLargeur{210}\def\PageHauteur{297}
\fi\fi\fi
% Reliure. Hors de « souple » et « rigide » : erreur, pas de supposition.
\def\@reliuresouple{souple}
\def\@reliurerigide{rigide}
\ifx\Reliure\@reliurerigide
  \edef\PageLargeur{\fpeval{\PageLargeur + 5}}
  \edef\PageHauteur{\fpeval{\PageHauteur + 5}}
  \def\FondPerdu{15}
\else\ifx\Reliure\@reliuresouple\else
  \PackageError{couverture}{Reliure \Reliure\space inconnue}%
    {Valeurs admises : souple, rigide.}%
\fi\fi
\makeatother

% Garde-fou : une édition jamais produite n'a pas de couverture.
\ifnum\NbPages=0
  \PackageError{couverture}{L'edition \Classe\space n'a pas de pagination}%
    {Cette edition n'a jamais ete compilee. Renseigner \string\Donnees\Classe
     ou passer \string\PagesForcees.}%
\fi

% Avertissement : l'imprimeur demande un nombre de pages multiple de 4
% (17/09/2026). Pas une erreur : c'est son exigence, pas une loi du papier.
\ifnum\fpeval{\NbPages - 4*floor(\NbPages/4)}=0 \else
  \typeout{COUVERTURE-AVERTISSEMENT \NbPages\space pages : pas un multiple
    de 4, que l'imprimeur exige -- ajouter des pages a l'interieur}%
\fi

% Garde-fou : un dos carré collé impose un nombre de pages pair.
\ifodd\NbPages
  \PackageError{couverture}{Nombre de pages impair (\NbPages)}%
    {Un dos carre colle exige un nombre PAIR de pages : une feuille en
     porte deux. Ajouter une page blanche a l'interieur.}%
\fi

% ---------------------------------------------------------------------
% 5. GÉOMÉTRIE — c'est ici que se joue l'exactitude du dos
% ---------------------------------------------------------------------

% Épaisseur du dos : valeur imposée si elle est donnée ET qu'elle vaut pour
% cette pagination et ce format ; calcul sinon.
\edef\DosCalcule{\fpeval{round(\NbPages / 2 * \Grammage * \Bouffant / 1000, 2)}}
\def\DosOrigine{calculee}
\ifnum\fpeval{\DosImpose > 0}=1
  \ifnum\fpeval{\NbPages = \DosImposePages}=1
    \ifx\Format\DosImposeFormat
      \ifx\Reliure\DosImposeReliure
        \def\DosOrigine{imposee}
      \fi
    \fi
  \fi
  \def\MotCalculee{calculee}
  \ifx\DosOrigine\MotCalculee
    \typeout{COUVERTURE-AVERTISSEMENT dos impose de \DosImpose\space mm
      donne pour \DosImposePages\space pages en \DosImposeFormat\space
      (reliure \DosImposeReliure) -- inapplicable a \NbPages\space pages en
      \Format\space (reliure \Reliure), dos calcule}%
  \fi
\fi
% NB : pas de \@ dans ces noms. Cette section est hors \makeatletter, et
% « \def\@imposee{...} » y redéfinit \@ (suivi du texte « imposee ») : le
% test échouait toujours, en silence. Constaté le 16/09/2026 — le dos de N
% tombait juste par coïncidence, 14,0 mm calculés pour 14 imposés.
\def\MotRigide{rigide}
\def\MotImposee{imposee}
\ifx\DosOrigine\MotImposee
  \edef\Dos{\DosImpose}
\else
  \edef\Dos{\DosCalcule}
\fi

% Marge technique : uniquement en mode épreuve.
\def\ouimot{oui}
\ifx\Reperes\ouimot \edef\Tech{\MargeTech}\else \def\Tech{0}\fi

% Dimensions du papier.
%   largeur = marge + fond perdu + 4e + dos + 1re + fond perdu + marge
%   hauteur = marge + fond perdu + hauteur + fond perdu + marge
\edef\PapierL{\fpeval{2*\PageLargeur + \Dos + 2*\FondPerdu + 2*\Tech}}
\edef\PapierH{\fpeval{\PageHauteur + 2*\FondPerdu + 2*\Tech}}

\geometry{paperwidth=\PapierL mm, paperheight=\PapierH mm,
          margin=0pt, nohead, nofoot}

% ---------------------------------------------------------------------
%  TOUTES LES COTES SONT FIGÉES ICI, ET NULLE PART AILLEURS.
%
%  Raison, payée par trois compilations ratées : un \fpeval laissé dans
%  une coordonnée ou dans le texte d'un nœud TikZ n'est pas développé au
%  moment où TikZ analyse le chemin. On obtient « Giving up on this path.
%  Did you forget a semicolon? », qui accuse une ponctuation alors que la
%  faute est ailleurs. Le corps du document ci-dessous n'emploie donc que
%  des macros déjà calculées.
% ---------------------------------------------------------------------

% Repères horizontaux, mesurés depuis le bord GAUCHE du papier.
\edef\CoupeG{\fpeval{\Tech + \FondPerdu}}          % coupe gauche (4e de couv.)
\edef\DosG{\fpeval{\CoupeG + \PageLargeur}}         % dos, début
\edef\DosD{\fpeval{\DosG + \Dos}}                   % dos, fin
\edef\CoupeD{\fpeval{\DosD + \PageLargeur}}         % coupe droite (1re de couv.)
% Repères verticaux, depuis le bord BAS.
\edef\CoupeB{\fpeval{\Tech + \FondPerdu}}
\edef\CoupeH{\fpeval{\CoupeB + \PageHauteur}}
% Axes.
\edef\DosAxe{\fpeval{(\DosG + \DosD)/2}}
\edef\MilieuH{\fpeval{(\CoupeB + \CoupeH)/2}}
\edef\MilieuL{\fpeval{\PapierL/2}}

% Format fini, pour les cotes affichées.
\edef\FiniL{\fpeval{2*\PageLargeur + \Dos}}

% Bandeau de la 1re de couverture : même proportion que la page de titre.
\edef\BandeauL{\fpeval{round(0.34*\PageLargeur, 2)}}
\edef\BandeauG{\fpeval{\CoupeD - \BandeauL}}
\edef\BandeauAccent{\fpeval{\BandeauG - 0.7}}
\edef\DosAccentG{\fpeval{\DosG - 0.7}}
\edef\DosAccentD{\fpeval{\DosD + 0.7}}
% Couverture rigide : le bleu du dos déborde de 2 mm sur chaque plat, pour
% absorber la tolérance de l'imprimeur sur l'épaisseur du dos (±10 %).
\edef\DosEtenduG{\fpeval{\DosG - 2}}
\edef\DosEtenduD{\fpeval{\DosD + 2}}

% Filigrane de classe.
\edef\FiligraneX{\fpeval{(\BandeauG + \CoupeD)/2}}
\edef\FiligraneY{\fpeval{\CoupeH - 6}}

% Bloc-titre de la 1re de couverture.
\edef\TitreX{\fpeval{\DosD + 12}}
\edef\TitreY{\fpeval{\CoupeH - 72}}
\edef\TitreW{\fpeval{\BandeauG - \DosD - 16}}
\edef\VersionY{\fpeval{\CoupeB + 14}}

% 4e de couverture.
\edef\QuatreX{\fpeval{\CoupeG + \Securite + 6}}
\edef\QuatreY{\fpeval{\CoupeH - \Securite - 14}}
\edef\QuatreBasY{\fpeval{\CoupeB + \Securite + 22}}
\edef\QuatreW{\fpeval{\PageLargeur - 2*\Securite - 12}}

% Zone de sécurité.
\edef\SecG{\fpeval{\CoupeG + \Securite}}
\edef\SecD{\fpeval{\CoupeD - \Securite}}
\edef\SecB{\fpeval{\CoupeB + \Securite}}
\edef\SecH{\fpeval{\CoupeH - \Securite}}

% Zones de contrôle : ce qui doit rester de part et d'autre du pli du dos.
\edef\PliQuatre{\fpeval{\DosG - \Securite}}         % 4e : limite côté dos
\edef\PliPremiere{\fpeval{\DosD + \Securite}}       % 1re : limite côté dos

% Traits de coupe et cotes, dans la marge technique.
\edef\TechExt{\fpeval{\Tech - 8}}
\edef\TechInt{\fpeval{\Tech - 1}}
\edef\TechCote{\fpeval{\Tech - 9}}
\edef\HautInt{\fpeval{\PapierH - \Tech + 1}}
\edef\HautExt{\fpeval{\PapierH - \Tech + 8}}
\edef\HautCote{\fpeval{\PapierH - \Tech + 9}}
\edef\DroitInt{\fpeval{\PapierL - \Tech + 1}}
\edef\DroitExt{\fpeval{\PapierL - \Tech + 8}}

% Emplacement du code-barres ISBN : en bas à DROITE de la 4e de couverture,
% selon l'usage éditorial, et non à gauche où il chevaucherait les mentions
% de licence.
\edef\IsbnD{\fpeval{\DosG - \Securite}}
\edef\IsbnG{\fpeval{\IsbnD - 40}}
\edef\IsbnH{\fpeval{\SecB + 25}}
\edef\IsbnCX{\fpeval{\IsbnG + 20}}
\edef\IsbnCY{\fpeval{\SecB + 12.5}}

% Cartouches de cotes, dans la marge technique. Ancrés vers l'INTÉRIEUR du
% papier : ancrés vers l'extérieur, ils débordaient et étaient rognés.
\edef\CoteBasY{\fpeval{1}}
\edef\CoteHautY{\fpeval{\PapierH - 1}}

% ---------------------------------------------------------------------
% 6. COULEURS ET POLICES — reprises de la page de titre (build_book.py)
% ---------------------------------------------------------------------

\definecolor{TitleBand}{cmyk}{.9,.55,.1,.35}   % bleu profond
\definecolor{TitleAccent}{cmyk}{.05,.8,1,0}    % rouge DARC (accent fin)
\definecolor{Repere}{cmyk}{0,1,1,0}            % magenta : repères de montage

\setmainfont{Libertinus Serif}[
  ItalicFont     = Libertinus Serif Italic,
  BoldFont       = Libertinus Serif Bold,
  BoldItalicFont = Libertinus Serif Bold Italic]

\pagestyle{empty}

% ---------------------------------------------------------------------
% 7. FILIGRANE DE CLASSE
%    Le texte est construit ici, et non dans le nœud : un \\ produit à
%    l'intérieur d'un \foreach referme le groupe d'alignement de TikZ.
%    Une lettre seule est grande ; plusieurs s'empilent et se réduisent
%    (build_book.py v0.20 : « NEA » à 105 pt débordait encore).
% ---------------------------------------------------------------------

\makeatletter
\def\FiligraneTexte{}
\newif\ifpremierelettre \premierelettretrue
\count255=0
\@for\@tempa:=\Lettres\do{%
  \advance\count255 by 1\relax
  \ifpremierelettre \premierelettrefalse
  \else \g@addto@macro\FiligraneTexte{\\}\fi
  \expandafter\g@addto@macro\expandafter\FiligraneTexte\expandafter{\@tempa}}%
\xdef\NbLettres{\the\count255}
\makeatother

\ifnum\NbLettres=1 \def\CorpsF{150}\else\def\CorpsF{86}\fi
\edef\CorpsFI{\fpeval{round(\CorpsF*1.02, 0)}}

% Texte du cartouche de cotes, lui aussi construit hors du nœud TikZ.
%
% NE PAS y employer \textsuperscript : dans un nœud « align=center », il
% casse le chemin TikZ (« Giving up on this path. Did you forget a
% semicolon? »), et l'erreur est signalée sur le nœud SUIVANT — elle
% accuse donc une ponctuation là où la faute est ici. Isolé par essais :
% $^2$ et le caractère ² passent, \textsuperscript non.
\def\CotesTexte{%
  \Classe\ --- \NbPages\ pages --- papier \Grammage\ g/m$^2$,
  bouffant \Bouffant\ --- fond perdu \FondPerdu\ mm\\
  papier \PapierL\ mm sur \PapierH\ mm --- fini \FiniL\ mm sur
  \PageHauteur\ mm\\
  EPREUVE, ne pas envoyer a l'impression%
}

% ---------------------------------------------------------------------
% 8. CONTRÔLE DE SUPERPOSITION
%
%  \MesureNoeud{nœud}{préfixe}{retrait} mesure la boîte d'un nœud par ses
%  quatre coins (un nœud pivoté a ses ancres pivotées : le min et le max des
%  quatre suffisent) et la ramène au repère du papier, en mm. Le point (0,0)
%  est lu par le même « let » que les ancres : les deux passent donc par la
%  même transformation, et la comparaison ne dépend pas de l'endroit où
%  TeX a posé l'image sur la page.
%
%  Le retrait sert au seul filigrane, pour mesurer les lettres et non leur
%  cadre : c'est le bord des lettres que le détourage rognerait.
%
%  PIÈGE MESURÉ le 16/09/2026 : la marge interne TikZ (.3333em) n'est PAS
%  évaluée dans la police du nœud, mais dans celle qui entoure le dessin.
%  Instrumenté : dans le nœud, em = 150 pt ; boîte brute du « N » :
%  41,7 x 37,5 mm, soit la lettre plus 2 x 1,3 mm — 1,3 mm étant .3333em
%  à 11 pt. Un retrait calculé à 150 pt (17,6 mm) ramenait la lettre à
%  6,6 mm de large, et le contrôle aurait dit OK sur une mesure absurde.
% ---------------------------------------------------------------------

\newcount\NbAlertes
\newdimen\FiligraneSep
\newsavebox\CreditsBoite

% Crédits de 4e : lignes coupées à la main. Leur largeur RÉELLE est celle
% de la plus longue ligne, mesurée dans une boîte à part (voir § 8, ISBN).
\def\CreditsTexte{%
    Réalisé à partir des contenus de 50ohm.de (en allemand)\\
    50ohm.de-Autorenteam, coordonné par le référat AJW du DARC e.\,V.\\
    \Auteur\\[3pt]
    Licence CC BY 4.0 --- \texttt{\Depot}}

\def\MesureNoeud#1#2#3{%
  \path let \p1=(#1.south west), \p2=(#1.north east),
            \p3=(#1.north west), \p4=(#1.south east), \p5=(0,0) in
    \pgfextra{%
      \expandafter\xdef\csname #2xmin\endcsname{\fpeval{round((min(\x1,\x2,\x3,\x4) - \x5 + #3)/1mm, 2)}}%
      \expandafter\xdef\csname #2xmax\endcsname{\fpeval{round((max(\x1,\x2,\x3,\x4) - \x5 - #3)/1mm, 2)}}%
      \expandafter\xdef\csname #2ymin\endcsname{\fpeval{round((min(\y1,\y2,\y3,\y4) - \y5 + #3)/1mm, 2)}}%
      \expandafter\xdef\csname #2ymax\endcsname{\fpeval{round((max(\y1,\y2,\y3,\y4) - \y5 - #3)/1mm, 2)}}%
    };}

% \ControleDans{libellé}{préfixe}{xmin}{xmax}{ymin}{ymax}
\def\ControleDans#1#2#3#4#5#6{%
  \ifnum\fpeval{\csname #2xmin\endcsname >= #3 && \csname #2xmax\endcsname <= #4
             && \csname #2ymin\endcsname >= #5 && \csname #2ymax\endcsname <= #6}=1
    \typeout{COUVERTURE-CONTROLE OK #1 : x \csname #2xmin\endcsname-\csname #2xmax\endcsname,
      y \csname #2ymin\endcsname-\csname #2ymax\endcsname\space dans x #3-#4, y #5-#6}%
  \else
    \global\advance\NbAlertes by 1
    \typeout{COUVERTURE-CONTROLE ALERTE #1 : x \csname #2xmin\endcsname-\csname #2xmax\endcsname,
      y \csname #2ymin\endcsname-\csname #2ymax\endcsname\space HORS de x #3-#4, y #5-#6}%
  \fi}

% \ControleDisjoints{libellé}{préfixe}{xmin}{xmax}{ymin}{ymax}
% Deux rectangles se chevauchent si et seulement s'ils se recouvrent sur
% les DEUX axes à la fois.
\def\ControleDisjoints#1#2#3#4#5#6{%
  \ifnum\fpeval{\csname #2xmin\endcsname < #4 && #3 < \csname #2xmax\endcsname
             && \csname #2ymin\endcsname < #6 && #5 < \csname #2ymax\endcsname}=1
    \global\advance\NbAlertes by 1
    \typeout{COUVERTURE-CONTROLE ALERTE #1 : chevauchement}%
  \else
    \typeout{COUVERTURE-CONTROLE OK #1 : disjoints}%
  \fi}

\begin{document}
% Marge interne du filigrane, évaluée dans la police qui ENTOURE le dessin :
% c'est là que TikZ évalue le .3333em de « inner sep » (cf. § 8).
\FiligraneSep=.3333em
% Largeur réelle des crédits : la plus longue ligne, dans leur police.
\sbox\CreditsBoite{\small\begin{tabular}{@{}l@{}}\CreditsTexte\end{tabular}}%
\begin{tikzpicture}[remember picture, overlay,
                    x=1mm, y=1mm, shift={(current page.south west)}]

% ---- Fonds ----------------------------------------------------------
% La 4e de couverture reste sur le blanc du papier.
% Le bandeau bleu couvre le dos, puis reprend au tiers extérieur de la
% 1re de couverture — même partition 2/3-1/3 que la page de titre.
% Couverture rigide : pas de filets au dos (ils tomberaient à côté des
% plis, vu la tolérance), et un dos bleu élargi de 2 mm de chaque côté.
\ifx\Reliure\MotRigide
  \fill[TitleBand] (\DosEtenduG, 0) rectangle (\DosEtenduD, \PapierH);
\else
  \fill[TitleBand] (\DosG, 0) rectangle (\DosD, \PapierH);
\fi
\fill[TitleBand] (\BandeauG, 0) rectangle (\PapierL, \PapierH);
% Filets d'accent : à la jonction sur la 1re, de part et d'autre du dos.
\fill[TitleAccent] (\BandeauAccent, 0) rectangle (\BandeauG, \PapierH);
\ifx\Reliure\MotRigide\else
  \fill[TitleAccent] (\DosAccentG, 0) rectangle (\DosG, \PapierH);
  \fill[TitleAccent] (\DosD, 0) rectangle (\DosAccentD, \PapierH);
\fi

% ---- 1re de couverture ----------------------------------------------
% Filigrane de classe, en réserve claire, détouré au bandeau.
\begin{scope}
\clip (\BandeauG, 0) rectangle (\PapierL, \PapierH);
\node[name=filigrane, anchor=north, align=center, text=white, opacity=0.16,
      font=\fontsize{\CorpsF}{\CorpsFI}\selectfont\bfseries]
  at (\FiligraneX, \FiligraneY) {\FiligraneTexte};
\end{scope}

\node[name=titre, anchor=west, align=left, text width=\TitreW mm]
  at (\TitreX, \TitreY) {%
    {\fontsize{54}{56}\selectfont\bfseries \TitreOuvrage}\\[10pt]
    {\Large\color{TitleBand}\bfseries \SousTitre}\\[16pt]
    {\large\color{black!75} \SousTitreClasse}};

\node[name=version, anchor=south west, align=left, font=\small, text=black!70]
  at (\TitreX, \VersionY) {Version \Version\\ Licence CC BY 4.0};

% ---- Dos -------------------------------------------------------------
% Un dos étroit ne reçoit pas de texte : en deçà de 6 mm, le débord au
% massicot ferait mordre le texte vertical sur les plats.
\ifnum\fpeval{\Dos >= 6}=1
  \node[name=dos, rotate=90, anchor=center, text=white, font=\bfseries,
        inner sep=0pt]
    at (\DosAxe, \MilieuH)
    {\TitreOuvrage\ \ --- \ \SousTitreClasse\ \ --- \ \Version};
\fi

% ---- 4e de couverture ------------------------------------------------
\node[name=quatre, anchor=north west, align=justify, text width=\QuatreW mm]
  at (\QuatreX, \QuatreY) {\QuatriemeTexte};

% Le texte vit dans \CreditsTexte (§ 8) pour être mesuré à part : la boîte
% du nœud a la largeur fixe \QuatreW et couvre l'emplacement ISBN sans
% qu'aucun caractère n'y soit. Retirer « text width » pour l'éviter change
% l'interlignage — essayé le 16/09/2026, 6 934 pixels différents.
\node[name=credits, anchor=south west, align=left, font=\small, text=black!70,
      text width=\QuatreW mm]
  at (\QuatreX, \QuatreBasY) {\CreditsTexte};

% ---- Repères de fabrication (mode épreuve uniquement) ----------------
\ifx\Reperes\ouimot
  % Traits de coupe, hors du fond perdu.
  \foreach \x in {\CoupeG, \CoupeD} {
    \draw[Repere, line width=0.25pt] (\x, \TechExt) -- (\x, \TechInt);
    \draw[Repere, line width=0.25pt] (\x, \HautInt) -- (\x, \HautExt);}
  \foreach \y in {\CoupeB, \CoupeH} {
    \draw[Repere, line width=0.25pt] (\TechExt, \y) -- (\TechInt, \y);
    \draw[Repere, line width=0.25pt] (\DroitInt, \y) -- (\DroitExt, \y);}
  % Limites du dos, en pointillés.
  \foreach \x in {\DosG, \DosD} {
    \draw[Repere, line width=0.25pt, dashed] (\x, \TechExt) -- (\x, \TechInt);
    \draw[Repere, line width=0.25pt, dashed] (\x, \HautInt) -- (\x, \HautExt);}
  % Zone de sécurité.
  \draw[Repere, line width=0.2pt, dotted]
    (\SecG, \SecB) rectangle (\SecD, \SecH);
  % Emplacement réservé au code-barres ISBN.
  \draw[Repere, line width=0.2pt] (\IsbnG, \SecB) rectangle (\IsbnD, \IsbnH);
  \node[font=\tiny\sffamily, text=Repere, align=center]
    at (\IsbnCX, \IsbnCY) {emplacement\\ISBN};
  % Cotes.
  \node[anchor=north, font=\tiny\sffamily, text=Repere]
    at (\DosAxe, \CoteHautY) {dos \Dos\ mm (\DosOrigine)};
  \node[anchor=south, font=\tiny\sffamily, text=Repere, align=center]
    at (\MilieuL, \CoteBasY) {\CotesTexte};
\fi

% ---- Contrôle de superposition (§ 8) ----------------------------------
% Rien n'est dessiné : des mesures et des lignes de journal.
\MesureNoeud{filigrane}{Fil}{\the\FiligraneSep}
\MesureNoeud{titre}{Tit}{0pt}
\MesureNoeud{version}{Ver}{0pt}
\MesureNoeud{quatre}{Qua}{0pt}
\MesureNoeud{credits}{Cre}{0pt}
% Bord droit des crédits = bord gauche + marge interne + plus longue ligne.
\xdef\Crexmax{\fpeval{round(\Crexmin + (\the\FiligraneSep + \the\wd\CreditsBoite)/1mm, 2)}}

% Le filigrane peut entrer dans la zone de sécurité (il est décoratif),
% pas franchir la coupe ni sortir du bandeau — le détourage le rognerait.
\ControleDans{filigrane dans le bandeau}{Fil}{\BandeauG}{\CoupeD}{\CoupeB}{\CoupeH}
\ControleDans{titre sur la 1re, hors bandeau}{Tit}{\PliPremiere}{\BandeauAccent}{\SecB}{\SecH}
\ControleDans{version sur la 1re, hors bandeau}{Ver}{\PliPremiere}{\BandeauAccent}{\SecB}{\SecH}
\ControleDans{texte de 4e dans la zone de securite}{Qua}{\SecG}{\PliQuatre}{\SecB}{\SecH}
\ControleDans{credits de 4e dans la zone de securite}{Cre}{\SecG}{\PliQuatre}{\SecB}{\SecH}
\ControleDisjoints{titre / version}{Tit}{\Verxmin}{\Verxmax}{\Verymin}{\Verymax}
\ControleDisjoints{texte de 4e / credits}{Qua}{\Crexmin}{\Crexmax}{\Creymin}{\Creymax}
\ControleDisjoints{texte de 4e / emplacement ISBN}{Qua}{\IsbnG}{\IsbnD}{\SecB}{\IsbnH}
\ControleDisjoints{credits / emplacement ISBN}{Cre}{\IsbnG}{\IsbnD}{\SecB}{\IsbnH}
\ifnum\fpeval{\Dos >= 6}=1
  \MesureNoeud{dos}{Dos}{0pt}
  \ControleDans{texte du dos entre les plis}{Dos}{\DosG}{\DosD}{\SecB}{\SecH}
\else
  \typeout{COUVERTURE-CONTROLE OK dos de \Dos\space mm : trop etroit, sans texte}%
\fi

\end{tikzpicture}
\typeout{COUVERTURE-BILAN alertes=\the\NbAlertes\space; classe=\Classe\space;
  format=\Format\space; reliure=\Reliure\space; pages=\NbPages\space;
  version=\Version\space;
  dos=\Dos\space mm (\DosOrigine); fini=\FiniL x \PageHauteur\space mm;
  papier=\PapierL x \PapierH\space mm}
\end{document}
